%% filename: Quantum-Boolean-memoirs.tex
%% Adapted from doc/latex/amscls/amsbook-template.tex for Memoirs/AMS monograph submission.
\documentclass{amsbook}

% Additional packages used by the manuscript.
\usepackage{mathtools}
\usepackage{amssymb}
\usepackage{mathrsfs}
\usepackage{dsfont}
\usepackage{upref}
\usepackage{graphicx}
\usepackage{enumerate}
\usepackage{color}
\usepackage[hidelinks]{hyperref}

\newtheorem{theorem}{Theorem}[chapter]
\newtheorem{proposition}[theorem]{Proposition}
\newtheorem{lemma}[theorem]{Lemma}
\newtheorem{corollary}[theorem]{Corollary}
\newtheorem{conjecture}[theorem]{Conjecture}
\newtheorem{problem}[theorem]{Problem}

\theoremstyle{definition}
\newtheorem{definition}[theorem]{Definition}
\newtheorem{example}[theorem]{Example}
\newtheorem{question}[theorem]{Question}
\newtheorem{assumption}[theorem]{Assumption}

\theoremstyle{remark}
\newtheorem{remark}[theorem]{Remark}

\numberwithin{section}{chapter}
\numberwithin{equation}{chapter}

% Notation
\newcommand{\Ex}{\mathcal{E}}
\newcommand{\M}{\mathbb{M}}
\newcommand{\1}{\mathbf{1}}
\newcommand{\s}{\mathbf{s}}
\newcommand{\Ra}{\mathcal{R}}
\newcommand{\Tr}{\mathcal{T}^{\mathcal{R}}}
\newcommand{\T}{\mathbb{T}}
\newcommand{\E}{\mathbb{E}}
\newcommand{\La}{\mathcal{L}}
\newcommand{\R}{\mathbb{R}}
\newcommand{\var}{\mathrm{Var}}
\newcommand{\tr}{\mathrm{tr}}
\newcommand{\dc}{\mathbf{d}}
\newcommand{\Na}{\mathbb{N}}
\newcommand{\Z}{\mathbb{Z}}
\newcommand{\Dd}{\mathscr{D}}
\newcommand{\supp}{{\mathrm supp}\,}

\begin{document}

\frontmatter

\title[Quantum Functional Inequalities]
{Functional Inequalities on the Quantum Boolean Cube with Applications to Quantum Learning Theory}

%\author[Jiao]{Yong Jiao}
%\address{School of Mathematics and Statistics, HNP-LAMA, Central South University, Changsha 410083, China}
%\email{jiaoyong@csu.edu.cn}
%
%\author[Luo]{Sijie Luo}
%\address{School of Mathematics and Statistics, HNP-LAMA, Central South University, Changsha 410083, China}
%\email{sijieluo@csu.edu.cn}
%
%\author[Zhao]{Tiantian Zhao}
%\address{Institute for Advanced Study in Mathematics, Harbin Institute of Technology, Harbin 150001, China}
%\email{zhaotiantian@hit.edu.cn}
%
%\author[Zhou]{Dejian Zhou}
%\address{School of Mathematics and Statistics, HNP-LAMA, Central South University, Changsha 410083, China}
%\email{zhoudejian@csu.edu.cn}


\subjclass[2020]{Primary 46L52; Secondary 42B25, 46L53, 81P45.}
\keywords{Quantum Boolean cube, noncommutative $L_p$-spaces, quantum semigroup, functional inequalities, Boolean analysis, learning quantum observables.}
\date{}

\begin{abstract}
This memoir develops a systematic theory of quantum Boolean analysis on the $n$-qubit matrix algebra $\mathbb{M}_{2^{n}}$, emphasizing the interplay among the Pauli expansion, Fourier degree, and the quantum Ornstein--Uhlenbeck semigroup. We first prove that this semigroup is holomorphic on noncommutative $L_p$-spaces for every $1<p<\infty$, providing the analytic tool for degree-dependent estimates.
For observables supported on high Fourier degrees, we establish heat-smoothing estimates and reverse Bernstein--Markov inequalities. We also obtain a strong contraction estimate for the first tail space. 
For  low-degree observables, we obtain lower bounds for the decay of the heat semigroup and Bernstein--Markov inequalities for both the quantum Laplacian and gradient. Combined with hypercontractivity and semigroup representations of covariance, these results yield refined Talagrand-type and covariance inequalities in quantum tail spaces.



We next introduce a notion of monotone quantum observables and prove quantum versions of the Harris--Kleitman inequality, Talagrand's correlation inequality, and the Keller--Mossel--Sen correlation inequality. 
In the Talagrand correlation inequality, an additional similarity condition on the first-order Pauli coefficients is needed because each quantum coordinate has three Pauli directions. We also establish quantitative BKS and inverse BKS inequalities, relating quantum influences to noise sensitivity. The BKS estimate is proved by  semigroup and noncommutative martingale methods, while the inverse inequality gives a converse estimate for monotone quantum observables.


Finally, we improve the  quantum Bohnenblust--Hille inequality  by showing that its constant is at most $3^{d/2}$ times the corresponding classical Boolean constant. A reduction to classical Boolean functions also yields quantum extensions of approximation results for Walsh tail spaces. Applying these estimates to quantum learning theory,  we prove, for bounded self-adjoint observables of fixed degree, a lower bound matching the upper bound in its dependence on the number of qubits.  Hence, this determines the sharp logarithmic order of the random-query complexity as $n\to\infty$. We also obtain results on exact learning, sparse approximation, and metric entropy.
\end{abstract}

\maketitle

\setcounter{page}{4}
\tableofcontents

\mainmatter
\chapter{Introduction}
The analysis of Boolean functions lies at the intersection of functional and harmonic analysis, which is built on specific analytic structures: the Walsh expansion, the discrete gradient, the heat semigroup, and the decomposition of a function according to Fourier degree. These fundamental structures underlie many of the central inequalities on the discrete cube, including Poincar\'e and logarithmic Sobolev inequalities, hypercontractivity, Talagrand's influence inequality, KKL-type estimates, correlation inequalities, noise-sensitivity theorems, Bohnenblust--Hille inequalities, and learning theorems for low-degree functions. A common feature of these results is that analytic estimates for the heat semigroup convert information about Fourier degree into quantitative information about influences, covariance, correlation, approximation, and learnability.

In recent years, motivated by developments in quantum computing and quantum information theory, it has become natural to extend these classical results to the noncommutative setting. Analysis on the classical Boolean cube can be naturally viewed as the commutative algebra generated by $n$ independent self-adjoint unitaries. In this interpretation, the coordinate functions are elementary two-valued observables, the Walsh characters are their products, and the Fourier degree measures how many coordinates an observable depends on. In the quantum Boolean cube, such
commuting observables are replaced by observables on $n$-qubit systems, namely self-adjoint elements of $\mathbb{M}_{2}(\mathbb{C})^{\otimes n}$. Several central problems in quantum information---such as the design of quantum error-correcting codes, and the analysis of quantum learning algorithms---rely on quantitative estimates in quantum Boolean cubes, which promotes the development of \emph{quantum} Boolean analysis.

The historical roots of non-commutative functional inequalities trace back to the fermionic setting, which serves as a fundamental tool in quantum field theory. Notably, Carlen and Lieb \cite{CL1993} established optimal hypercontractivity for Fermi fields by using the optimal uniform convexity estimate of noncommutative $L_{p}$ spaces, while Lust-Piquard established Riesz transform estimates and the $L_{p}$ Poincar\'e inequality in the fermionic setting in  \cite{ELP2008, Lu1998}. Building on the fundamental developments of noncommutative functional inequalities, Montanaro and Osborne \cite{MO2010} systematically formulated the modern framework for quantum Boolean functions. In their framework, a quantum Boolean function is a self-adjoint unitary $Q\in \mathbb{M}_{2}(\mathbb{C})^{\otimes n}$, and the Pauli expansion gives a Fourier analysis adapted to $n$-qubit observables. The associated quantum Ornstein--Uhlenbeck semigroup is the product depolarizing semigroup; on Pauli strings it acts by $\sigma_{\s}\mapsto e^{-t|\s|}\sigma_{\s}$ exactly as the classical heat semigroup acts on Walsh characters (we refer to Section \ref{sec-2} for the definitions). This structural correspondence ensures that concepts such as hypercontractivity, influence, noise sensitivity, and Fourier degree are mathematically well-defined in the quantum setting \cite{MO2010,Ki2014}. Applying the noncommutative semigroup theory and the Fourier analysis techniques, Montanaro and Osborne extended the Talagrand influence inequality and the Friedgut-Kalai-Naor theorem to the quantum Boolean cube setting.

It is worthwhile to mention here that extending classical Boolean analysis to the quantum setting invokes significant difficulties. The noncommutativity of $\mathbb{M}_{2}(\mathbb{C})^{\otimes n}$ causes certain fundamental classical identities (such as the exact equivalence between the $L_1$ and $L_2$ norms of gradients for Boolean functions) to fail. To overcome these essential obstacles, one must rely on the extensive developments in noncommutative harmonic analysis and noncommutative martingale theory achieved over the past few decades. Specifically, the noncommutative semigroup theory, the structures of noncommutative $L_{p}$ spaces, completely bounded maps, and noncommutative martingale inequalities (such as the Burkholder-Gundy inequality and the Doob maximal-type inequalities) provide fundamental tools to address the difficulties arising from the noncommutative phenomena.

Recently, motivated by the study of the quantum KKL conjecture posed by Montanaro and Osborne \cite{MO2010}, significant progress has been made in establishing noncommutative functional inequalities. In particular,  Talagrand-type influence inequalities, quantum KKL-type inequalities for geometric influences, quantum Friedgut-type theorems has been extensively studied in \cite{BGX2025,JLLZ2025,RWZ2024}, providing new insights into the complexity of quantum Boolean functions. 

Despite these remarkable advances, a systematic and self-contained analytical framework capable of governing the finer, higher-order structural properties of the quantum Boolean cube remains visibly absent. Therefore, the main purpose of this monograph is to provide a systematic investigation of the theory of functional inequalities in quantum Boolean cubes. We focus on estimates which lie at the interaction
between the quantum Ornstein--Uhlenbeck semigroup and Fourier degree. These include heat smoothing for quantum tail spaces, lower decay bounds and Bernstein--Markov inequalities for low degree observables, Talagrand-type inequalities, correlation inequalities for quantum monotone observables, BKS-type inequalities, Bohnenblust--Hille estimates, and applications to learning quantum observables. To highlight the obstacles imposed by the noncommutativity, it is natural to first trace the historical and conceptual evolution of these inequalities. In Sections \ref{boolean-cube}--\ref{quantum-boolean-cube}, we review the relevant scalar and Banach space-valued theories on the classical Boolean cubes, alongside existing functional inequalities on the quantum Boolean cube. A detailed outline of the memoir is provided in Section \ref{roadmap}.

%The main objective of this paper is to provide a systematic study of functional inequalities on quantum Boolean cube. 
\section{Functional inequalities on the Boolean cube}\label{boolean-cube}
 Throughout this monograph, let $(\{-1,1\}^{n},\mu)$ be the Boolean cube equipped with the uniform probability measure, which forms a basic model in Boolean analysis. Its Fourier--Walsh expansion, discrete gradient, and heat semigroup provide the standard analytic framework for Boolean function analysis and for many of its applications in probability, theoretical computer science, and social choice theory. We briefly recall the parts of this framework that will serve as the commutative model for the quantum results
developed below.

For each $A\subseteq [n]\coloneqq \{1,\dots, n\}$, the \emph{Walsh function} is given by
\[
w_A(\varepsilon)=\prod_{j\in A}\varepsilon_j,\quad \varepsilon=(\varepsilon_1,\varepsilon_2,\cdots,\varepsilon_n)\in \{-1,1\}^n.
\]
The collection of Walsh functions $\{w_A : A \subset [n]\}$ is called the \emph{Walsh system} in the sequel. Each function $f:\{-1,1\}^n\to \mathbb{C}$ has the Walsh-Fourier expansion
\begin{equation}
f(\varepsilon)=\sum_{A\subseteq [n]}\widehat{f}(A)w_A(\varepsilon), \quad \varepsilon=(\varepsilon_1,\cdots,\varepsilon_n)\in \{-1,1\}^{n},
\end{equation}
where $\widehat{f}(A)=\mathbb{E}(fw_A)$ is the Fourier coefficient. For each $\varepsilon\in \{-1,1\}^n$ and $j\in [n]$, denote by $\varepsilon^{\oplus j}$ the flipping of the $j$-th variable of $\varepsilon$. For each $j\in [n]$, the discrete partial derivatives  and the Laplacian operator on the Boolean cube as 
 \begin{equation}\label{lap-1}
 D_j(f)(\varepsilon)=\frac{f(\varepsilon)-f(\varepsilon^{\oplus j})}{2}, \quad \Delta =\sum_{j=1}^n D_j.
 \end{equation}
The heat semigroup generated by $-\Delta$ is diagonal in the Walsh system:
\begin{equation}\label{heat-semigroup}
e^{-t\Delta}(f)=\sum_{A\subseteq [n]}e^{-t|A|}\widehat{f}(A)w_{A},\quad t\geq 0,
\end{equation}
where $f:\{-1,1\}^{n}\to \mathbb{C}$ and $|A|=\#\{j\in [n]:j\in A\}$ stands for the cardinality of the subset $A$. Thus, $e^{-t\Delta}$ is indeed a \emph{Fourier multiplier} with symbol $A\mapsto e^{-t|A|}$ for each $A\subseteq [n]$. This spectral relation is fundamental for many of the inequalities recalled below. By the spectral gap of the operator $-\Delta$, it is clear that the following $L_{2}$-Poincar\'e inequality holds:
\[
\|f-\mathbb{E}(f)\|_{L_{2}(\{-1,1\}^{n})}\leq \left(\sum_{j=1}^{n}\|D_{j}(f)\|^{2}_{L_{2}(\{-1,1\}^{n})}\right)^{1/2},
\]
for every $f:\{-1,1\}^{n}\to \mathbb{R}$.

A first nontrivial and essential inequality is the $L_{p}$-Poincar\'e inequality, which we now recall as follows. For $1\leq p<\infty$, there exists a constant $K_{p}>0$ depending only on $p$ such that the following holds for each $f:\{-1,1\}^n\to\mathbb{R}$,
 \begin{equation}\label{ppi}
\|f-\mathbb{E}(f)\|_{L_{p}(\{-1,1\}^{n})}\leq K_{p}\left\|\left(\sum_{j=1}^{n}|D_jf|^{2}\right)^{1/2}\right\|_{L_{p}(\{-1,1\}^{n})}.
\end{equation}
This inequality can be found in \cite{Pi1986,Ta1993}, and the best $K_p=O(\sqrt{p})$ (as $p\to \infty$) was proved in \cite{EL2008} via a noncommutative method. Extensions and refinements of the Poincar\'e inequality \eqref{ppi} have been a central object of study in the probability and analysis
literature for decades.

In the seminal paper \cite{Gr1975}, Gross obtained the (gaussian) logarithmic Sobolev inequality, which can be viewed as an improvement of the $L_2$-Poincar\'e inequality. More precisely, for every $f:\{-1,1\}^n\to\mathbb{R}$ the following \emph{logarithmic Sobolev inequalities} 
\begin{equation}\label{lsi}
\mathbb{E}\left(|f|^2\log(|f|)\right)-\|f\|_{L_{2}(\{-1,1\}^{n})}^2\log\|f\|_{L_{2}(\{-1,1\}^{n})}\leq\sum_{j=1}^{n}\|D_{j}(f)\|^{2}_{L_{2}(\{-1,1\}^{n})}.
\end{equation}
Furthermore, Gross showed that \eqref{lsi} is equivalent to the (optimal) hypercontractivity inequality of the corresponding semigroup $\{e^{-t\Delta}\}_{t\geq 0}$: for $1<p\leq q <\infty$, 
\begin{equation}\label{bonami-1975}
\|e^{-t\Delta }\|_{L_{p}(\{-1,1\}^{n})\to L_{q}(\{-1,1\}^{n})}=1 \quad \mbox{if and only if} \quad e^{-2t}\leq \frac{p-1}{q-1}.
\end{equation}
The estimate \eqref{bonami-1975} was first proved by Bonami \cite{Bo1970}; Beckner \cite{Bec1975} subsequently extended it to complex-valued functions and related it to the sharp Hausdorff-Young inequality. We refer to \cite{BKS1999,FKN2002} for the applications of hypercontractivity to Boolean analysis, and to \cite{IN2022} for a recent advance in this direction. A similar hypercontractive phenomenon
also holds for other semigroups such as the Ornstein-Uhlenbeck semigroup on the Gaussian space $(\mathbb{R}^{n},\gamma_{n})$, which plays an essential role in Gaussian analysis; see e.g. \cite{Gr1975}. 

In the remarkable paper, Talagrand \cite{Ta1993} extended \eqref{lsi} to general $1\leq p<\infty$ in the following manner. For $1\leq p<\infty$, there exists $K_{p}>0$ such that
 \begin{equation}\label{plsi}
\|f-\mathbb{E}(f)\|_{L_p(\log L)^{\frac{p}{2}}(\{-1,1\}^{n})}\leq K_p\Big\|\Big(\sum_{j=1}^{n}|D_jf|^2\Big)^{\frac{1}{2}}\Big\|_{L_p(\{-1,1\}^{n})},
\end{equation}
for every $f:\{-1,1\}^{n}\to \mathbb{R}$. The $p$-logarithmic Sobolev inequality is a refinement of the $L_{p}$-Poincar\'e inequality \eqref{ppi}, providing an \emph{optimal} Sobolev embedding result. Gaussian analogues were obtained in \cite{Ad1979,Fe1975} for $1<p<\infty$ and in \cite{Le1988} for $p=1$. 

Another major improvement of the $L_2$-Poincar\'e inequality is known as the \emph{Talagrand influence inequality}, which states that  for every $f:\{-1,1\}^{n}\to \R$, the following inequality holds (\cite{Ta1994})
\begin{equation}\label{talagrand-inequality}
\var{f}\coloneqq \|f-\mathbb{E}(f)\|_{L_2(\{-1,1\}^{n})}^2\leq K\sum\limits_{j=1}^{n}\frac{\left\|D_{j}(f)\right\|_{L_2(\{-1,1\}^{n})}^{2}}{1+\log\left(\|D_{j}(f)\|_{L_2(\{-1,1\}^{n})}/\|D_{j}(f)\|_{L_1(\{-1,1\}^{n})}\right)}.
\end{equation}
For Boolean function $f:\{-1,1\}^{n}\to \{0,1\}$, \eqref{talagrand-inequality} yields the following Kahn-Kalai-Linial (KKL) inequality (\cite{KKL1988})(see also \cite{KKKMS2021, Ta1994}):
\begin{equation}\label{functional-kkl}
\var(f)\leq K\frac{\sum_{j=1}^{n}\|D_{j}(f)\|^{2}_{L_{2}(\{-1,1\}^{n})}}{\log\left(1/\max_{j\in [n]}\|D_{j}(f)\|^{2}_{L_{2}(\{-1,1\}^{n})}\right)},
\end{equation}
where $K>0$ is a universal constant. Consequently, we have
\begin{equation}\label{kkl}
\max_{j\in[n]}\|D_j(f)\|_{L_2(\{-1,1\}^{n})}^2\geq K\frac{\var(f)\log(n)}{n},
\end{equation}
for every Boolean function  $f:\{-1,1\}^{n}\to \{0,1\}$.

Interpreting the Boolean function $f$ as a voting rule, \eqref{kkl} says that if a voting
rule has all influences $\|D_j(f)\|_{L_2}^2$ are small, then $f$ is necessarily an unfair system in the sense that its expectation is very close to either $0$ or $1$. Talagrand's inequality has since become a standard tool in the study of threshold phenomena, communication complexity, superconcentration, and metric nonembeddability; see for example, \cite{Cha2014, CEE2020, EG2020, KMS2014, KR2009, OW2013}.


More recently, motivated by a conjecture of Talagrand \cite{Ta1997}, Eldan and Gross \cite{EG2020} (see also \cite{EG2022}) applied stochastic analysis techniques to prove the following result which is known as the \emph{Eldan-Gross inequality}
\begin{equation}\label{eg-2022}
\var(f)\sqrt{\log\left(1+\frac{e}{\sum_{j=1}^{n}\|D_{j}(f)\|^{2}_{L_1(\{-1,1\}^{n})}}\right)}\leq C\left\|\left(\sum_{j=1}^{n}|D_{j}(f)|^{2}\right)^{1/2}\right\|_{L_1(\{-1,1\}^{n})},
\end{equation}
for every Boolean function $f:\{-1,1\}^{n}\to\{0,1\}$. Notably, the Eldan-Gross inequality \eqref{eg-2022} unifies the KKL inequality  and Talagrand's isoperimetric inequality \cite[Theorem 1.1]{Ta1993}, making \eqref{eg-2022} into one of the most efficient tools in Boolean analysis. An alternative proof invoking the Bakry--\'Emery theory is given by Ivanisvili and Zhang \cite{IZ2026}. Other recent proofs and related developments can be found in \cite{BIM2023,EKLM2022,Ros2020}.

We conclude this section by recalling the \emph{Bohnenblust--Hille (BH) inequality}, which was established by Bohnenblust and Hille in \cite{B-H1931} to resolve Bohr's absolute convergence problem for Dirichlet series.  In the Boolean cube setting, the Bohnenblust--Hille inequality was established by Defant, Masty\l o, and P\'erez \cite{D-M-P2019} to explore the structures of low-degree polynomials: it asserts that for any polynomial $f: \{-1,1\}^n \to \mathbb{R}$ of degree at most $d$, its Fourier coefficients satisfy
\[
\left( \sum_{A \subseteq [n], |A| \le d} |\widehat{f}(A)|^{\frac{2d}{d+1}} \right)^{\frac{d+1}{2d}} \le C_d \|f\|_{L_\infty(\{-1,1\}^{n})},
\]
where $C_{d}$ is a constant independent of the dimension $n$. Such estimates connect the analytic property of a polynomial with the summability properties of its Fourier coefficients. In modern theoretical computer science and learning theory, the BH inequality in Boolean cubes serves as the mathematical foundation for analyzing and reconstructing low-degree Boolean functions. Most notably, the BH inequality in Boolean cubes provides optimal and dimension-free quantitative bounds that are intimately connected to the study of the Aaronson--Ambainis conjecture concerning the maximum influence of bounded polynomials \cite{A-A2014}. Furthermore, within the Probably Approximately Correct (PAC) learning framework, the BH inequality provides a rigorous structural justification for the celebrated Linial--Mansour--Nisan (LMN) algorithm \cite{L-M-N1993}. By guaranteeing that the tail of the Fourier spectrum of a bounded polynomial decays rapidly and concentrates in a highly summable regime, the BH inequality enables efficient reconstruction from low-frequency data. Consequently, it represents the fundamental link between analytic smoothness and sparsity in classical Boolean analysis. We refer interested readers to \cite{A-S-E-P2025, B-S2026, EI2022, EZ2024, S-V-Z2024,V-Z2024} for more detail on the BH inequality and related topics.
\section{Banach space-valued functional inequalities}\label{vector-valued}

The classical inequalities for scalar-valued functions recalled in the preceding section have several vector-valued counterparts. This direction is relevant here for two reasons. Firstly, it clarifies the connections between Boolean analysis and the non-linear geometry of Banach spaces which forms an essential part of the \emph{Ribe program}. Secondly, the methods developed in the vector-valued theory, in particular Pisier's use of $K$-convexity and holomorphic semigroups, provide one of the main sources for the noncommutative semigroup estimates proved later in this paper.

We begin with Pisier's vector-valued Poincar\'e inequality. For a fixed Banach space $(X,\|\cdot\|)$, let $1\leq p<\infty$ and $\varepsilon,\delta$ be independent random vectors that are uniformly distributed on the discrete cube $\{-1,1\}^{n}$. In the remarkable work \cite{Pi1986}, Pisier  discovered the following  inequalities for vector-valued functions on the discrete cube:
\begin{equation}\label{eq-pisier-ineq}
\|f-\mathbb{E}(f)\|^{p}_{L_{p}(\{-1,1\}^{n};X)} \le C_n^p \mathbb{E}\left\|\sum_{j=1}^n \delta_j D_j f(\varepsilon)\right\|^{p},
\end{equation}
where $f:\{-1,1\}^{n}\to X$. It was proved by Pisier that $C_n \sim \log n$ for \emph{arbitrary} Banach space $X$ (\cite[Lemma~7.3]{Pi1986}). Subsequently, it was shown by Talagrand \cite{Ta1993} that this order of growth is optimal for Banach $X$ with trivial cotype. 

Determining which Banach spaces admit a dimension-free Pisier inequality \eqref{eq-pisier-ineq} is one of the fundamental problems in Banach space theory. This inequality serves as a key ingredient in establishing the coincidence of Rademacher type and Enflo type. Although the dimension-free form of \eqref{eq-pisier-ineq} is false in general, Ivanisvili, van Handel and Volberg \cite{IvHV2020} resolved Enflo's problem by establishing the following new dimension-free \emph{Poincar\'e-type inequality} invoking heat semigroup theory in the Boolean. More precisely, they proved that for each $1\leq p<\infty$ there exists a constant $K_{p}>0$ depending only on $p$ such that for \emph{arbitrary} Banach space $X$, $n\in\mathbb{N}$, and every $f:\{-1,1\}^{n}\to X$ we have
\begin{equation}
\left(\mathbb{E}\|f- \mathbb{E}(f)\|^p\right)^{\frac{1}{p}} \le \int_0^\infty \left(\mathbb{E}\left\|\sum_{j=1}^n \delta_j(t)D_{j}(f)\right\|^p\right)^{\frac{1}{p}} dt,
\end{equation}
where
\[
\delta_j(t) := \frac{\xi_j(t) - e^{-t}}{\sqrt{1 - e^{-2t}}}
\]
and $\{\xi_j(t)\}_{j=1}^n$ are i.i.d.\ biased Rademacher random variables, independent of $\varepsilon$, with
\[
\mathbb{P}(\xi_j(t) = 1) = \frac{1 + e^{-t}}{2} \quad \text{and} \quad \mathbb{P}(\xi_j(t) = -1) = \frac{1 - e^{-t}}{2}.
\]
Inspired by the remarkable paper \cite{IvHV2020}, specific functional inequalities which refined the Poincar\'e inequalities have been extended to the Banach space-valued setting. For instance, in \cite{CE2023, CE2024}, Cordero-Erausquin and Eskenazis extended Talagrand's influence inequality \eqref{talagrand-inequality} and the logarithmic Sobolev inequality to Banach space-valued functions under natural type or cotype assumptions on the target space. We also refer the reader to \cite{IS2024} for the connection between Banach space-valued KKL inequality and the Rademacher type of the target space. More recently, Gordin \cite{Go2026} extended the main results of \cite{IvHV2020} to biased Boolean cubes. These Banach space-valued inequalities were then applied to obtain lower bounds for the average distortion, with respect to biased product measures, of embeddings of the Boolean cube into Banach spaces.

Motivated by the main purpose of the present paper, one of the main focuses is the ``heat smoothing estimate'' and Bernstein-Markov type inequalities. For each $d\in [n]$, let
\[
\mathcal{T}^{\geq d}(\{-1,1\}^n)=\left\{T\in L_{2}(\{-1,1\}^n):f=\sum_{|A|\geq d}\widehat{f}(A)w_{A}\right\}
\]
and 
\[
\mathcal{P}^{\leq d}(\{-1,1\}^n)=\left\{T\in L_{2}(\{-1,1\}^n):f=\sum_{|A|\leq d}\widehat{f}(A)w_{A}\right\}.
\]
The space $\mathcal{T}^{\geq d}(\{-1,1\}^n)$ is called the $d$-tail space. 
In \cite[Remark 5.5]{MN2014}, Mendel and Naor proposed the following problem which is known as the \emph{heat smoothing conjecture}:
\begin{conjecture}\label{conj-1}
Let $1<p<\infty$ and $d\in [n]$. Then there exist $c_p,C_p \in (0,\infty)$ such that the following inequality holds for every $f\in \mathcal{T}^{\geq d}(\{-1,1\}^n)$, 
\begin{equation}\label{con-1}
\|e^{-t\Delta} f\|_{L_p(\{-1,1\}^n)} \leq  C_p e^{-c_pdt} \|f\|_{L_p(\{-1,1\}^n)}.
\end{equation}
\end{conjecture}
\begin{conjecture}\label{conj-1-2}
Let $1<p<\infty$ and $d\in [n]$. For each $f\in \mathcal{T}^{\geq d}(\{-1,1\}^n)$, there exists $c_p\in  (0,\infty)$ such that 
\begin{equation}\label{con-2}
\|\Delta f\|_{L_p(\{-1,1\}^n)} \geq c_p d\|f\|_{L_p(\{-1,1\}^n)}.
\end{equation}
\end{conjecture}

The original question of Mendel and Naor  \cite{MN2014} was vector-valued, but the scalar form above is already open in this generality. Conjecture \ref{conj-1-2} is the corresponding reverse Bernstein--Markov inequality and would follow from Conjecture \ref{conj-1} by integrating the heat semigroup estimate. Related estimates of \eqref{con-1} date back to Meyer \cite{Me1984} in the Gaussian setting (see also \cite[Lemma 5.4]{MN2014} for the Boolean cube).

Heilman, Mossel and Oleszkiewicz \cite{HMO2017} resolved Conjecture \ref{conj-1} for $d=1$. While Eskenazis and Ivanisvili \cite{EI2023} considered a Gaussian analog of this conjecture and proved sharp heat smoothing estimates for holomorphic functions in the Gaussian setting.

There is also a complementary low-degree problem. Motivated by the heat smoothing conjecture, Eskenazis and Ivanisvili \cite{EI2020} proposed the following dual estimates for functions of degree
at most $d$.

\begin{conjecture}\label{conj-2}
Let $1<p<\infty$ and $d\in [n]$. Then there exist constants $c_p,C_p \in (0,\infty)$ such that the following inequality holds for $f\in \mathcal{P}^{\leq d}(\{-1,1\}^n)$,
\begin{equation}\label{con-3}
\|e^{-t\Delta} f\|_{L_p(\{-1,1\}^n)} \geq  C_p e^{-c_pdt} \|f\|_{L_p(\{-1,1\}^n)}.
\end{equation}
\end{conjecture}
\begin{conjecture}\label{conj-2-2}
Let $1<p<\infty$ and $d\in [n]$. For each $f\in \mathcal{P}^{\leq d}(\{-1,1\}^n)$, there exists $c_p\in  (0,\infty)$ such that 
\begin{equation}\label{con-4}
\|\Delta f\|_{L_p} \leq c_p d\|f\|_{L_p}.
\end{equation}
\end{conjecture}

The relevance of these conjectures to the present paper is twofold.  On the one hand, the heat smoothing estimates describe the decay of high Fourier levels under the semigroup; on the other hand, their dual low-degree counterparts control the stability of low-degree observables under the same semigroup.  These are exactly the degree-sensitive results that will be developed on the quantum Boolean cube in Sections \ref{sec-qhse} and \ref{sec-bm}.

The best general results in this direction for Banach space-valued functions are due to Mendel and Naor \cite{MN2014}. For a $K$-convex Banach space $X$\footnote{A Banach space $(X,\|\cdot\|_X)$ is said $K$-convex if each Rademacher projection $\mathrm{Rad}^{=k}(f)=\sum_{|A|=k}\widehat{f}(A)w_{A}$ is bounded on $L_p(\{-1,1\}^n; X)$, $1<p<\infty$.} and $1<p<\infty$, there exist $c(p,X)$ and $C(p,X)\in (0,\infty)$, and $A(p,X)\geq 1$ such that the following holds for every $f\in \{-1,1\}^{n}\to X$ lying in the $d$-tail space,

\begin{equation}\label{eq-mn}
\|e^{-t\Delta} f\|_{L_p(\{-1,1\}^{n}; X)} \le C(p,X) e^{-c(p,X)d\min\{t,t^{A(p,X)}\}} \|f\|_{L_p(\{-1,1\}^n;X)},
\end{equation}
and there exist $c(p,X) \in (0,\infty)$ and $\alpha\in (0,1]$ such that
\begin{equation}\label{eq-mn-2}
\|\Delta f\|_{L_p(\{-1,1\}^n;X)}\geq c(p,X) d^{\alpha} \|f\|_{L_p(\{-1,1\}^n;X)}. 
\end{equation}
Eskenazis and Ivanisvili later provided a short proof of \eqref{eq-mn} based on a duality argument developed in \cite{EI2020}, and Volberg proved \eqref{eq-mn-2} by analytic paraproduct methods \cite{Vo2024}.  Partial results for the low-degree conjectures, including the equivalence between Conjectures \ref{conj-2} and \ref{conj-2-2}, were obtained in \cite{DDFKOSW2025}; see also \cite{FHKL2016,EI2020,Vo2024}.  Pisier's holomorphic semigroup theorem for $K$-convex targets \cite{Pi1982} is a fundamental result behind these developments, which is one of the prototypes for the quantum holomorphic semigroup theory derived in Section \ref{sec-hqs}.

\section{Functional inequalities on the Quantum Boolean cube}\label{quantum-boolean-cube}

We next recall the noncommutative background most directly related to the present paper. This section is not meant to provide a comprehensive survey of noncommutative functional inequalities; its purpose is to isolate several foundational results that motivate the quantum Boolean analysis developed in the sequel.

Functional inequalities on the canonical anticommutation relations algebra, or CAR algebra for short, are fundamental tools in noncommutative analysis and quantum field theory.  Motivated by Gross' work on logarithmic Sobolev inequalities and hypercontractivity \cite{Gr1975,Gr1975-Duke}, Carlen and Lieb proved optimal hypercontractivity for Fermi fields \cite[Theorem~4]{CL1993}.  Lust-Piquard established $L_{p}$ estimates for the Riesz transforms associated with the fermionic number operator \cite{Lu1998}.  Efraim and Lust-Piquard subsequently obtained Poincar\'e type inequalities on the discrete cube and in the CAR algebra by a one-parameter group method \cite[Theorem 6.1]{ELP2008}; this approach also gives the correct order of the Poincar\'e constants.  More recently, Jiao, Luo and Zhou proved Eldan--Gross and KKL-type inequalities in the CAR setting \cite{JLZ2025}.  Heat smoothing estimates have also been studied in other noncommutative settings, for example, in holomorphic subalgebras of $q$-Gaussian algebras \cite{Zh2023}.

The quantum Boolean cube shares structural features with the
CAR algebra but differs from it in an
essential way: it  provides a finite-dimensional noncommutative model of the Boolean cube:
\[
\mathbb M_{2^n}(\mathbb C)=\mathbb M_2(\mathbb C)^{\otimes n}
\]
with normalized trace, equipped with the Pauli basis.  The Pauli strings play the role of Walsh characters, and their support size gives the Fourier degree.  The corresponding quantum Ornstein--Uhlenbeck semigroup is the product depolarizing semigroup and acts diagonally by
\[
        \sigma_{\mathbf s}\mapsto e^{-t|\mathbf s|}\sigma_{\mathbf s}.
\]
This is the structural reason why the semigroup and Fourier-degree questions recalled in Section \ref{vector-valued} have natural quantum analogues.

Montanaro and Osborne introduced a systematic Fourier-analytic framework for quantum Boolean functions \cite{MO2010}.  In their setting, a quantum Boolean function is a self-adjoint unitary in $\mathbb M_{2^n}(\mathbb C)$.  They proved, among other results, hypercontractivity for the quantum Ornstein--Uhlenbeck semigroup and quantum analogues of influence inequalities and Friedgut--Kalai--Naor type results; see also King's hypercontractivity theorem for unital qubit semigroups \cite{Ki2014}.  Montanaro and Osborne also formulated a quantum KKL problem.  A direct transfer of the classical Talagrand-to-KKL argument is obstructed by a basic noncommutative phenomenon: for quantum Boolean functions one no longer has the classical identity relating the geometric and $L_{2}$ influences, namely $\|d_{j}(A)\|_{L_{1}(\mathbb M_{2^{n}})}=\|d_{j}(A)\|_{L_{2}(\mathbb M_{2^{n}})}^{2}$. This failure is one of the simplest indications that the quantum theory cannot be obtained by a formal replacement of Walsh characters with Pauli counterparts.

Subsequent work has clarified several aspects of the quantum influence theory.  Rouz\'e, Wirth and Zhang  \cite{RWZ2024} proved quantum Talagrand, KKL and Friedgut-type theorems using geometric influence.  Blecher, Gao and Xu \cite{BGX2025}, and independently Jiao, Luo, Lin and Zhou \cite{JLLZ2025}, gave further approaches to quantum geometric Talagrand inequalities; the latter work also obtained a quantum Eldan--Gross inequality.  These results show that many qualitative features of Boolean analysis persist on the quantum Boolean cube, but they also show that the appropriate formulations often involve theory from the geometry of noncommutative $L_{p}$ spaces, noncommutative martingales, and noncommutative gradients.



The present paper continues this line of research in a degree-sensitive direction on the quantum Boolean cube. While the existing theory has focused mainly on quantum analogues of influence, KKL-type, Friedgut-type, and Eldan--Gross-type inequalities, the results developed below focus on the interaction between the quantum Ornstein--Uhlenbeck semigroup and the Pauli--Fourier degree. This perspective leads to heat smoothing estimates, Bernstein--Markov inequalities, tail Talagrand and covariance inequalities, correlation and BKS estimates, Bohnenblust--Hille inequalities, and applications to learning quantum observables. A recurring technical feature is that the classical pointwise and commutative semigroup arguments must be replaced by noncommutative $L_{p}$-methods, holomorphic functional calculus, noncommutative Khintchine inequalities, and martingale techniques.

 Existing results in this direction have focused mainly on quantum analogues of influence, KKL-type, Friedgut-type, and Eldan--Gross-type phenomena. Here we take a degree-sensitive viewpoint, organizing the analysis around the interaction between the quantum Ornstein--Uhlenbeck semigroup and the Pauli--Fourier degree. This leads to heat smoothing estimates, Bernstein--Markov inequalities, tail Talagrand and covariance inequalities, correlation and BKS estimates, Bohnenblust--Hille inequalities, and applications to learning quantum observables. Several classical arguments must consequently be replaced by tools from noncommutative $L_{p}$ space theory, functional calculus, noncommutative Khintchine inequalities, and noncommutative martingale theory.

\newpage
\section{Outline of the paper}\label{roadmap}
 The monograph begins with an analytic treatment of the quantum Boolean cube and then turns to degree-sensitive estimates. It concludes with applications to approximation and learning theory.

In Chapter \ref{sec-2}, we collect basic notations and elementary analytic facts used throughout the monograph.  We recall the Pauli-Walsh expansion of elements of $\mathbb{M}_{2^{n}}$, the quantum Ornstein--Uhlenbeck semigroup, the partial derivatives $d_{j}$, and the low-degree  subspace $\mathcal P^{\leq d}(\mathbb{M}_{2^{n}})$ and the tail subspace $\mathcal{T}^{\geq d}(\mathbb{M}_{2^{n}})$, which are fundamental objects that will be  investigated in detail through the monograph. Moreover, the (optimal) hypercontractivity of the quantum Ornstein--Uhlenbeck semigroup, a new anti-isomorphism of $\mathbb{M}_{2^{n}}$ (i.e., Lemma \ref{var-prop}) used to handle sign multipliers are introduced.  These preliminary results provide a basic bridge between the classical heat semigroup on the Boolean and the quantum Ornstein--Uhlenbeck semigroup on quantum Boolean cube.

Chapter \ref{sec-hqs} is devoted to the holomorphicity of the semigroup, which serves as a basic tool in subsequent chapters. We prove that the quantum Ornstein--Uhlenbeck semigroup is holomorphic on $L_{p}(\mathbb{M}_{2^{n}})$, uniformly in $n$, for every $1<p<\infty$.  This is the noncommutative counterpart of Pisier's holomorphic semigroup theorem for $K$-convex Banach spaces. Notably, the classical  representation of the semigroup is no longer applicable in our framework; instead, the proof relies on a modified  representation of the semigroup by commuting conditional expectations, together with the $K$-convexity of noncommutative $L_{p}$ spaces; see Lemma \ref{st-e-tl} and Lemma \ref{lem-st-2}. The holomorphic extension then gives us bounds for quantum Rademacher projections and for the functional calculus $w^{\mathcal L}$ on lens domains.

Chapters \ref{sec-qhse} and \ref{sec-bm} form the foundation of the semigroup theory developed in this monograph.  More precisely, Chapter \ref{sec-qhse} focuses on inequalities for observables lying in low-degree, or tail spaces. We prove quantum analogues of the Mendel--Naor heat smoothing estimate on $\mathcal{T}^{\geq d}(\mathbb{M}_{2^n})$, establish a quantum version of the Heilman--Mossel--Oleszkiewicz estimate in the first tail space, and derive reverse Bernstein--Markov inequalities for observables whose spectrum is localized between two degrees. Our arguments are inspired by \cite{MN2014,EI2020,HMO2017} and make essential use of the holomorphic functional calculus developed  in Chapter \ref{sec-hqs}, together with conformal mapping arguments on lens domains. Substantial additional difficulties arise from noncommutativity and the lack of pointwise estimates in the quantum setting. To overcome these obstacles, we also derive probabilistic representation of noncommutative semigroups. Chapter \ref{sec-bm} is devoted to observables lying in the low-degree regime.  We first prove lower bounds for the decay of $e^{-t\mathcal L}$ on low-degree subspace $\mathcal P^{\leq d}(\mathbb{M}_{2^{n}})$, which can be viewed as a ``dual'' version of Theorem \ref{mn-k-tail}.  The elementary Chebyshev bound gives a dimension-free estimate for all $1\leq p\leq \infty$, and the holomorphic functional calculus gives a sharper estimate when $1<p<\infty$.  We then prove Bernstein--Markov inequalities for the quantum Laplacian and related operators.  Our main estimates presented in this chapter use approximation-theoretic inequalities and a noncommutative Cowling-type theorem regarding the bound of imaginary Riesz potential.  The gradient estimates require a substantially different approach. More precisely, we first develop a transference principle, with no direct classical analogue, that reduces the problem to a vector-valued setting. A further key ingredient is a tensor-product representation of the relevant operator in terms of completely positive and trace-preserving maps. This representation allows us to combine noncommutative Khintchine inequalities, a quantum analogue of Naor's inequality, and Buser-type estimates to obtain the desired gradient bounds.

In Chapter \ref{sec-tala-tail}, we  focus on inequalities and covariance estimates in tail spaces.  We prove a Talagrand inequality for observables in $\mathcal{T}^{\geq d}(\mathbb{M}_{2^n})$. In particular, the geometric influence estimate obtained from Theorem \ref{lp-tala-d-tail} carries a full logarithmic denominator and improves the dependence on small influences in the corresponding result of \cite{RWZ2024}. We then develop covariance Poincar\'e and covariance Talagrand inequalities for pairs of quantum observables. Our approach replies on the covariance representation in terms of the semigroup and the partial derivative, combined with hypercontractivity and the additional decay property in tail spaces. Thus Chapter \ref{sec-tala-tail} converts spectral localization into refined influence and covariance bounds. 
It is also worthwhile to mention that the covariance versions of the Poincar\'e and Talagrand inequalities are even new for the classical Boolean analysis theory.

Chapter \ref{sec-tala-correlation} develops lower bounds for correlations on the quantum Boolean cube.  We introduce a notion of monotonically increasing element in $\mathbb{M}_{2^n}$ by requiring positivity of the three directional derivatives $\partial_{j,\ell}$ associated with each coordinate and each Pauli direction.  We then prove quantum analogues of the Harris--Kleitman, Talagrand \cite{Ta1996}, and Keller--Mossel--Sen \cite{KMS2014} correlation inequalities.  The quantitative Talagrand estimate is formulated in terms of the first-level quantity $\mathcal{W}_1(T_1,T_2)$, which records the interaction of matching Pauli directions, while the KMS-type estimate is expressed through the $L_1$-sizes of the coordinate derivatives. A key distinction from the classical setting is that each quantum coordinate carries three first-order Pauli coefficients rather than a single coefficient.  This gives rise to an essential obstruction: the natural product of the $L_{1}$-influences need not vanish even when the covariance does, which prevents a direct extension of the classical Talagrand correlation inequality. To address this obstruction, we introduce a similarity condition requiring the first-order Pauli coefficient vectors of the two observables to be proportional coordinate by coordinate.  This condition is automatic in the classical one-directional setting, but is genuinely additional in the quantum setting and is necessary for our Talagrand inequality, as demonstrated in Example \ref{example-2} and Theorem \ref{quantum-ta}.  Chapter \ref{sec-tala-correlation} thus extends monotone correlation inequalities to noncommutative observables while identifying the precise structural modification required by the three-directional Pauli structures.

Chapter \ref{sec-bks} investigates quantum noise sensitivity through the quantum BKS inequality and its converse. We first establish a dimension-free quantitative BKS estimate, showing that the quantitative BKS theorem relates the decay of $\mathrm{Var}(e^{-t\mathcal L}T)$ to the $L_{1}$-sizes of the quantum influences.  The main novelty is a noncommutative semigroup--martingale argument. The Pauli--Walsh expansion induces a dyadic filtration of matrix algebras, and the associated conditional expectations and martingale differences provide an orthogonal decomposition while contracting the coordinate derivatives. Combined with the commutation between the martingale differences, the quantum Ornstein--Uhlenbeck semigroup and  hypercontractivity, this yields a direct proof that avoids the random-restriction method. We then prove an inverse quantum BKS inequality for monotone increasing observables, providing a lower bound for the variance after the action of quantum noise in terms of the same influence quantities. We prove this converse estimate using the covariance representation from Chapter \ref{sec-tala-correlation}, together with two monotonicity facts: the semigroup preserves monotonicity, and the associated covariance terms are monotone. In the zero-noise limit, the estimate also gives a lower bound for the original variance. Hence, this chapter provides both upper and lower bounds relating quantum influence and noise sensitivity via noncommutative semigroup and martingale methods.

Chapter \ref{sec BH distance} develops coefficient estimates and approximation results for low-degree quantum observables.  We reprove  the quantum Bohnenblust--Hille inequality of Volberg and Zhang \cite{V-Z2024} with better constants. Our approach relies on the method of reducing a low-degree quantum observable to a classical polynomial on a Boolean cube with three copies of each coordinate and using a carefully chosen family of product density matrices.  The same reduction principle is then combined with a quantum Figiel-type estimate for Rademacher projections, duality arguments, and K-functional methods to study the $L_{1}$-distance of low-degree observables from quantum tail spaces. In particular, the chapter extends approximation results for Walsh tail spaces to the quantum setting, establishes a quantum counterpart of the two-sided Oleszkiewicz estimate in the degree-one case, and derives corresponding bounds for quantum symmetric polynomials, including lower estimates that reveal the necessity of the Pauli-direction factor. These results serve as Fourier-analytic tools used to explore the learning theory in the next chapter.

Chapter \ref{sec QLT} applies the preceding Fourier-analytic estimates to learning quantum observables.  We formulate quantum deterministic and random query models in terms of expectations against density matrices.  For bounded self-adjoint observables of degree at most $d$, we prove the sharp logarithmic asymptotic order of the random query complexity as $n\to\infty$.  The upper bound uses the quantum Bohnenblust--Hille inequality and the reduction from Chapter \ref{sec BH distance}. The lower bound is obtained from packing estimates for low-degree quantum Boolean observables, constructed through quantum analogues of decision trees. More precisely, the lower bound is based on two ingredients: a general counting principle that turns metric packings into query-complexity estimates, and explicit large packings built from quantum analogues of decision trees. The novelty on our construction of quantum decision trees use the three Pauli directions at each qubit to produce large families of pairwise separated, low-degree quantum Boolean observables. The packing argument also applies to all bounded self-adjoint contractions, not only to quantum Boolean observables, which gives a quantum counterpart of the classical sharpness result in \cite{EIS2023}. We also discuss exact learning, sparse recovery estimates, and metric entropy bounds. Thus, the final chapter shows that the analytic estimates developed earlier yield quantitative results for learning problems on the quantum Boolean cube.

\newpage
\section{List of notations}
Basic notations used throughout the monograph are summarized below.
\begin{itemize}
\item Let $\mathbb{N}$, $\mathbb{R}$, and $\mathbb{C}$ denote the sets of natural, real, and complex numbers, respectively.
\item $[n]\coloneqq \{1,2,3,\dots,n\}$.
\item Let $(\{-1,1\}^{n},\mu)$ be the $n$-dimensional Boolean cube equipped with the uniform probability measure $\mu$.
\item For each $A\subseteq [n]$ and $\varepsilon=\{\varepsilon_{j}\}_{j=1}^{n}\in \{-1,1\}^{n}$, define the Walsh function $w_A(\varepsilon)\coloneqq\prod_{j\in A}\varepsilon_j$.
\item $\sigma_{\s}=\sigma_{s_1}\otimes\cdots\otimes\sigma_{s_n}$ for multi-index $\s\in\{0,1,2,3\}^{n}$.
\item For $\s\in\{0,1,2,3\}^n$, define the support of $\s$ by $ \supp(\s)\coloneqq \{j\in[n]:s_j\neq0\}$, and $|s|=\#\supp(\s)$ is the cardinality of the support.
\item Let $\mathbb{M}_{2^{n}}\coloneqq \mathbb{M}^{\otimes n}_{2\times 2}(\mathbb{C})$ be the $n$-fold tensor product of $2\times 2$ complex matrices, and $\tr$ be the normalized trace on $\mathbb{M}_{2^{n}}$.
\item For $1\leq p\leq \infty$, let $L_{p}(\mathbb{M}_{2^{n}})$ be the noncommutative $L_{p}$ space generated by $(\mathbb{M}_{2^{n}},\tr)$.
\item For $d\in [n]$, the bounded degree space is defined by
\[
\mathcal{P}^{\leq d}(\mathbb{M}_{2^n})=\left\{T\in L_{2}(\mathbb{M}_{2^n}):T=\sum_{|\s|\leq d}\widehat{T}(\s)\sigma_{\s}\right\}.
\]
\item For $d\in [n]$, the tail space is defined by
\[
\mathcal{T}^{\geq d}(\mathbb{M}_{2^n})=\left\{T\in L_{2}(\mathbb{M}_{2^n}):T=\sum_{|\s|\geq d}\widehat{T}(\s)\sigma_{\s}\right\}.
\]
\item For $k\in [n]$, the Rademacher projection: $\mathrm{Rad}^{=k}(T) = \sum_{\substack{\mathbf{s} \in \{0,1,2,3\}^n \\ |\mathbf{s}| =k}} \widehat{T}(\mathbf{s}) \sigma_{\mathbf{s}}$, and $\mathrm{Rad}^{\leq k}(T) = \sum_{\substack{\mathbf{s} \in \{0,1,2,3\}^n \\ |\mathbf{s}|\leq k}} \widehat{T}(\mathbf{s}) \sigma_{\mathbf{s}}$.
%\item The quantum Ornstein-Uhlenbeck semigroup $e^{-t\mathcal{L}}(T) = \sum_{\s \in \{0,1,2,3\}^n} e^{-t|\s|} \widehat{T}(\s) \sigma_\s$.
%\item The infinitesimal generator/ Laplacian operator $\mathcal{L}(T) = \sum_{\s \in \{0,1,2,3\}^n} |\s| \widehat{T}(\s) \sigma_\s.$
%\item Given $T\in \M_{2^n}$, denote by $\widehat{T}$ the vector of Fourier coefficients $T$.
\end{itemize}

\chapter{Preliminaries}\label{sec-2}
\section{Quantum Boolean cube and Walsh system}\label{quantum-notations}

Throughout this monograph, $[n]=\{1,\ldots,n\}$.  Let $\M_{2}(\mathbb{C})$ be the $*$-algebra of all $2\times2$ complex matrices, equipped with the normalized trace $\tr_{2}$.  We use the following Pauli matrices
\[
\sigma_0 = \begin{pmatrix} 1 & 0 \\ 0 & 1 \end{pmatrix}, \quad
\sigma_1 = \begin{pmatrix} 0 & 1 \\ 1 & 0 \end{pmatrix}, \quad
\sigma_2 = \begin{pmatrix} 0 & -i \\ i & 0 \end{pmatrix}, \quad
\sigma_3 = \begin{pmatrix} 1 & 0 \\ 0 & -1 \end{pmatrix}.
\]
The Pauli matrices are Hermitian unitaries with $\tr_2(\sigma_j)=0$ for $j=1,2,3$, and
\begin{equation}\label{car-relation}
\sigma_j\sigma_k+\sigma_k\sigma_j=0,\quad  j\neq k \in \{1,2,3\}. 
\end{equation}
The quantum analogue of functions on the Boolean cubes is $\mathbb{M}_2(\mathbb{C})^{\otimes n}\cong\mathbb{M}_{2^n}(\mathbb{C})$ endowed with the normalized trace $\mathrm{tr}:=\mathrm{tr}_2^{\otimes n}$ and the unit $\mathbf{1}\coloneqq \sigma_{0}^{\otimes n}$. Following \cite[Definition 3.1]{MO2010}, an observable $A\in\mathbb{M}_{2^n}(\mathbb{C})$ is said to be a \emph{quantum Boolean function} if $A=A^*$ and $A^2=\mathbf{1}$. 

In the sequel, we may shortly write  $\mathbb{M}_2$ and $\mathbb{M}_{2^n}$ for $\mathbb{M}_2(\mathbb{C})$ and $\mathbb{M}_{2^n}(\mathbb{C})$, respectively. The identity maps on $\mathbb{M}_2$ and $\mathbb{M}_{2^n}$ will be denoted by $\mathbb{I}_2$ and $\mathbb{I}$, respectively. For $1\le p\le\infty$, let $L_p(\M_{2^n})$ be the noncommutative $L_p$-space associated with $(\M_{2^n},\tr)$.  Thus
\[
 \|x\|_{L_\infty(\M_{2^n})}=\|x\|_{\M_{2^n}}
\]
for $p=\infty$, while
\[
 \|x\|_{L_p(\M_{2^n})}=\bigl(\tr(|x|^p)\bigr)^{1/p},\qquad 1\le p<\infty,
\]
where $|x|=(x^*x)^{1/2}$.

For each multi-index $\s=(s_1,\ldots,s_n)\in\{0,1,2,3\}^n$, set
\[
 \sigma_{\s}=\sigma_{s_1}\otimes\cdots\otimes\sigma_{s_n}.
\]
The family $\{\sigma_{\s}:\s\in\{0,1,2,3\}^n\}$ forms an orthonormal basis of $L_2(\M_{2^n})$.  Consequently every $T\in\M_{2^n}$ has the Pauli--Walsh expansion
\[
 T=\sum_{\s\in\{0,1,2,3\}^n}\widehat T(\s)\sigma_{\s},
\]
where $ \widehat{T}(\mathbf{s})$ is the Fourier coefficient defined by $\widehat{T}(\mathbf{s}) = \operatorname{tr}(\sigma_{\mathbf{s}} T)$. For $\s\in\{0,1,2,3\}^n$, define the \emph{support} of $\s$ by
\[
 \supp(\s)\coloneqq \{j\in[n]:s_j\neq0\},
\]
and $|\s|=\#\supp(\s)$ is the cardinality of the support.

This is already explained in the notation section. Should we delete it here )For each $d\in [n]$, we define 
\[
\mathcal{T}^{\geq d}(\mathbb{M}_{2^n})=\left\{T\in L_{2}(\mathbb{M}_{2^n}):T=\sum_{|\s|\geq d}\widehat{T}(\s)\sigma_{\s}\right\}
\]
and 
\[
\mathcal{P}^{\leq d}(\mathbb{M}_{2^n})=\left\{T\in L_{2}(\mathbb{M}_{2^n}):T=\sum_{|\s|\leq d}\widehat{T}(\s)\sigma_{\s}\right\}.
\]
The space $\mathcal{T}^{\geq d}(\mathbb{M}_{2^n})$ is called the $d$-tail space in $\mathbb{M}_{2^n}$. When no confusion arises, for convenience, we may write $\mathcal{T}^{\geq d}$ and $\mathcal{P}^{\leq d}$ instead of $\mathcal{T}^{\geq d}(\mathbb{M}_{2^n})$ and $\mathcal{P}^{\leq d}(\mathbb{M}_{2^n})$, respectively.
\section{The quantum Ornstein-Uhlenbeck semigroup}
We next recall the product depolarizing semigroup, which will be denoted by the quantum Ornstein--Uhlenbeck semigroup. Let $L:\mathbb{M}_2 \rightarrow\mathbb{M}_2$ defined by
\begin{equation}\label{l-def}
L(A)\coloneqq A-\tr_2[A]\mathbf{1}_2,\quad A\in \M_{2}.
\end{equation}
Since the operator $L$ is indeed a projection on $\mathbb{M}_{2}$, it follows from the Taylor expansion that
\begin{equation}\label{etl-def}
e^{-tL}(A)=e^{-t}A+(1-e^{-t})\tr_2[A]\mathbf{1}_2,\quad A\in \M_{2}.
\end{equation}
In view of \eqref{l-def} and \eqref{etl-def}, for each $j \in \{0,1,2,3\}$, the following identities are obvious
\[
L(\sigma_j) = 
\begin{cases} 
\sigma_j, & j \neq 0; \\ 
0, & \text{otherwise,} 
\end{cases}
\]
and
\[
e^{-tL}(\sigma_j) = 
\begin{cases} 
e^{-t} \sigma_j, & j \neq 0; \\ 
\sigma_0, & \text{otherwise.} 
\end{cases}
\]
For each stabilizer operator $\sigma_{\s}$, let
\[
P_t(\sigma_{\s}) = e^{-tL}(\sigma_{{\s}_1}) \otimes \cdots \otimes e^{-tL}(\sigma_{{\s}_n}).
\]
Then $\{P_{t}\}_{t\geq 0}$ forms a semigroup on $L_{2}(\M_{2^{n}})$ with the \emph{infinitesimal generator} given by
\[
\mathcal{L}(\sigma_{\s}) := \lim_{t \to 0} \frac{\sigma_{\s} - (e^{-tL} \otimes \cdots \otimes e^{-tL})(\sigma_{\s})}{t} = |\s| \sigma_{\s},
\]
for each stabilizer operator $\sigma_\s$. In the sequel, we denote $\{ P_t \}_{t \geq 0}$ by $\left\{e^{-t\mathcal{L}}\right\}_{t \geq 0}$ to emphasize the generator $\mathcal{L}$. Equivalently, for $T \in \mathbb{M}_{2^n}$ we have
\begin{equation}\label{l-def-2}
\mathcal{L}(T) = \sum_{\s \in \{0,1,2,3\}^n} |\s| \widehat{T}(\s) \sigma_\s
\end{equation}
and
\begin{equation}\label{et-l-def}
e^{-t\mathcal{L}}(T) = \sum_{\s \in \{0,1,2,3\}^n} e^{-t|\s|} \widehat{T}(\s) \sigma_\s.
\end{equation}
For $j \in [n]$, the $j$-th partial differential operator (or, quantum bit-flip map) is given by
\[
\dc_j\coloneqq\mathbb{I}_2^{\otimes (j-1)} \otimes (\mathbb{I}_2 - \tr_2) \otimes \mathbb{I}_2^{\otimes (n-j)}.
\]
Due to the Pauli-Walsh expansion of $T\in \M_{2^{n}}$, it follows that each $\dc_j$ is a projection on $L_{2}(\M_{2^{n}})$ such that
\begin{equation}\label{dj-def}
\dc_j(T) = \sum_{\s \in \{0,1,2,3\}^n} \widehat{T}(\s) d_j(\sigma_\s) = \sum_{\substack{\s \in \{0,1,2,3\}^n \\ \s_j \neq 0}} \widehat{T}(\s) \sigma_\s,
\end{equation}
and $\mathcal{L} = \sum_{j=1}^n \dc_j$.

We recall the following \emph{optimal hypercontractivity} for the quantum Ornstein-Uhlenbeck semigroup $\{e^{-t\mathcal{L}}\}_{t\geq 0}$ as follows and  refer to \cite{K2014} (or, \cite[Theorem 8.4]{MO2010}) for more details.
\begin{theorem}\label{hc-1}
For each $1 < p < q < \infty$, we have
\[
\| e^{-t\mathcal{L}} \|_{L_p \to L_q} = 1 \quad \mbox{if and only if } t \geq \frac{1}{2} \log ( \frac{q-1}{p-1} ).
\]
\end{theorem}
The rest of this chapter is devoted to deriving and recalling technical lemmas that will be used throughout this monograph.

\begin{lemma}\label{intertwining}
For all $j\in[n]$ and $t\geq0$, we have
$$\dc_je^{-t\mathcal{L}}=e^{-t\mathcal{L}}\dc_j.$$
\end{lemma}
\begin{proof} For each $j\in[n]$, by the definitions of   \eqref{et-l-def} and \eqref{dj-def}, we have
$$\dc_je^{-t\mathcal{L}}(T)=e^{-t\mathcal{L}}\dc_j(T)=\sum_{\substack{\mathbf{s}\in\{0,1,2,3\}^n\\\s_j \neq 0}}e^{-t|\mathbf{s}|}\widehat{T}(\mathbf{s})\sigma_\mathbf{s},\quad \forall T\in\mathbb{M}_{2^n}.$$
\end{proof}

The following lemma is taken from \cite[Proposition 2.7 (i)]{JLLZ2025}.
\begin{lemma}\label{jlzz-lem}
For $T\in\mathbb{M}_{2^n}$ with $\|T\|_{L_\infty}\leq1$ and each $j\in[n]$, we have
\[
\|\dc_{j}(T)\|_{L_\infty}\leq1.
\]
\end{lemma}
\section{Auxiliary lemmas}
Set
\[
R_\rho:=\begin{pmatrix} 0 & i \\ -i & 0 \end{pmatrix}.
\]
Then,   $R_{\rho}=R^{*}_{\rho}=R^{-1}_{\rho}$ and for each $j \in \{0,1,2,3\}$, we have
\[
R_\rho\sigma_j^{\mathsf T}R_\rho = 
\begin{cases} 
-\sigma_j, & j \neq 0; \\ 
\sigma_j, & \text{otherwise,} 
\end{cases}
\]
where $\sigma_j^{\mathsf T}$ is the transpose of $\sigma_j$. The following map is inspired by Efraim and Lust-Piquard \cite{ELP2008}. Define the  map $\varPhi_\rho:\mathbb{M}_{2^n} \rightarrow \mathbb{M}_{2^n}$ by setting
\begin{equation}\label{var-def}
\varPhi_\rho(\sigma_{\mathbf{s}})\coloneqq R_\rho^{\otimes n}\sigma_{\mathbf{s}}^{\mathsf T}R_\rho^{\otimes n}=\otimes_{j=1}^nR_\rho\sigma_{s_j}^{\mathsf T}R_\rho.
\end{equation}
Equivalently,
\begin{equation}\label{var-def-2}
 \varPhi_\rho(\sigma_{\s})=R_\rho^{\otimes n}\sigma_{\s}^{\mathsf T}R_\rho^{\otimes n}
 =\bigotimes_{j=1}^n R_\rho\sigma_{s_j}^{\mathsf T}R_\rho
 =(-1)^{|\s|}\sigma_{\s}.
\end{equation}
Thus, $\varPhi_{\rho}$ is a linear bijection on $\mathbb{M}_{2^n}$.

\begin{lemma}\label{var-prop}
For $1 \leq p \leq \infty$ and $T \in \mathbb{M}_{2^n}$, the mapping $\varPhi_\rho$  given as in \eqref{var-def} satisfies the following properties:
\begin{enumerate}[{\rm (i)}]
\item $\varPhi_\rho$ is $*$-preserving, that is, $\varPhi_\rho(T^*)=\varPhi_\rho(T)^*$, for any $T\in\mathbb{M}_{2^n}$;
\item $\varPhi_\rho$ is an anti-isomorphism, that is, $\varPhi_\rho(T_1T_2)=\varPhi_\rho(T_2)\varPhi_\rho(T_1)$, for any $T_1,T_2\in\mathbb{M}_{2^n}$;
\item $\|\varPhi_\rho(T)\|_{L_{p}}=\|T\|_{L_{p}}$, for any $T\in\mathbb{M}_{2^n}$, and $1 \leq p \leq \infty$.
\end{enumerate}
\end{lemma}
\begin{proof} Since Pauli matrix $\sigma_{j}$  is a self adjoint matrix, we obtain
$$R_\rho(\sigma_{j}^*)R_\rho=[R_\rho\sigma_{j}R_\rho]^*= 
\begin{cases} 
-\sigma_j, & j \neq 0; \\ 
\sigma_j, & \text{otherwise.} 
\end{cases}$$
Therefore, for any $T\in\mathbb{M}_{2^n}$, by the definition and linearity of $\varPhi_\rho$, we have $\varPhi_\rho(T^*)=\varPhi_\rho(T)^*$.

We now prove item (ii). For any $T_1,T_2 \in \mathbb{M}_{2^n}$, we obtain
\begin{align}\label{anti-iso}
\varPhi_\rho(T_1T_2)&=\varPhi_\rho\left(\left(\sum_{\mathbf{s}\in\{0,1,2,3\}^n} \widehat{T_1}(\mathbf{s}) \sigma_{\mathbf{s}}\right)\left(\sum_{\mathbf{\ell}\in\{0,1,2,3\}^n} \widehat{T_2}(\mathbf{\ell}) \sigma_{\mathbf{\ell}}\right)\right)\nonumber\\
&=\sum_{\mathbf{s}\in\{0,1,2,3\}^n}\sum_{\mathbf{\ell}\in\{0,1,2,3\}^n}\widehat{T_1}(\mathbf{s})\widehat{T_2}(\mathbf{\ell})\varPhi_\rho(\sigma_{\mathbf{s}}\sigma_{\mathbf{\ell}})\\
&=\sum_{\mathbf{s}\in\{0,1,2,3\}^n}\sum_{\mathbf{\ell}\in\{0,1,2,3\}^n}\widehat{T_1}(\mathbf{s})\widehat{T_2}(\mathbf{\ell})\varPhi_\rho(\otimes_{j=1}^n\sigma_{\mathbf{s}_j}\sigma_{\mathbf{\ell}_j})\nonumber.\\\nonumber
\end{align}
By the definition of the operator $\varPhi_\rho$, it is clear that
\begin{align*}
\varPhi_\rho\left(\otimes_{j=1}^n\sigma_{\mathbf{s}_j}\sigma_{\mathbf{\ell}_j}\right)&=\otimes_{j=1}^nR_\rho(\sigma_{\mathbf{s}_j}\sigma_{\mathbf{\ell}_j})^{\mathsf T}R_\rho\\
&=\otimes_{j=1}^nR_\rho\sigma_{\mathbf{\ell}_j}^{\mathsf T}\sigma_{\mathbf{s}_j}^{\mathsf T}R_\rho\\
&=\otimes_{j=1}^n(R_\rho\sigma_{\mathbf{\ell}_j}^{\mathsf T}R_\rho) (R_\rho\sigma_{\mathbf{s}_j}^{\mathsf T}R_\rho)\\
&=\varPhi_\rho(\sigma_{\mathbf{\ell}})\varPhi_\rho(\sigma_{\mathbf{s}}).
\end{align*}
Then we get $\varPhi_\rho(T_1T_2)=\varPhi_\rho(T_2)\varPhi_\rho(T_1)$ by \eqref{anti-iso}.

Since $\|\cdot\|_{L_{p}}$ is unitary invariant, the following holds
\begin{align*}
\|\varPhi_\rho(T)\|_{L_{p}}&=\left\|\sum_{\mathbf{s} \in \{0,1,2,3\}^n} \widehat{T}(\mathbf{s}) R_\rho^{\otimes n}\sigma_{\mathbf{s}}^{\mathsf T}R_\rho^{\otimes n}\right\|_{L_{p}}\\
&=\left\|R_\rho^{\otimes n}\big(\sum_{\mathbf{s} \in \{0,1,2,3\}^n} \widehat{T}(\mathbf{s}) \sigma_{\mathbf{s}}\big)^{\mathsf T}R_\rho^{\otimes n}\right\|_{L_{p}}\\
&=\left\|\left(\sum_{\mathbf{s} \in \{0,1,2,3\}^n} \widehat{T}(\mathbf{s}) \sigma_{\mathbf{s}}\right)^{\mathsf T}\right\|_{L_{p}}\\
&=\|T\|_{L_{p}}.
\end{align*}
The proof is complete.
\end{proof}

In the sequel, we also need to apply the well-known \emph{noncommutative Khintchine inequality}, which was first proved by Lust-Piquard \cite{Lu1986} for the case $1<p<\infty$, and then extended to $0<p\leq 1$ by  Lust-Piquard, Pisier and Ricard \cite{LuP1991, PR2017}. For each finite sequence $\{T_{j}\}_{j=1}^{N}\subseteq L_p(\mathbb{M}_{2^n})$ ($0<p\leq \infty$), define 
$$\|(T_j)\|_{L_p(\mathbb{M}_{2^n}, \ell_2^c)}:=\Big\|\Big(\sum_{j=1}^{N}|T_j|^2\Big)^{1/2}\Big\|_{L_{p}},\quad \|(T_j)\|_{L_p(\mathbb{M}_{2^n}, \ell_2^r)}:=\Big\|\Big(\sum_{j=1}^{N}|T_j^*|^2\Big)^{1/2}\Big\|_{L_{p}}.$$
For $2\leq p\leq \infty$, define 
$$\|(T_j)\|_{L_p(\mathbb{M}_{2^n}, \ell_2^{cr})}\coloneqq \max\{\|(T_j)\|_{L_p(\mathbb{M}_{2^n}, \ell_2^c)},\|(T_j)\|_{L_p(\mathbb{M}_{2^n}, \ell_2^r)}\},$$
and for $0<p\leq 2$, define
$$\|(T_j)\|_{L_p(\mathbb{M}_{2^n}, \ell_2^{cr})}\coloneqq \inf_{T_j=A_j+B_j}\{\|(A_j)\|_{L_p(\mathbb{M}_{2^n}, \ell_2^c)}+\|(B_j)\|_{L_p(\mathbb{M}_{2^n}, \ell_2^r)}\}.$$
We only state the following special form of the noncommutative Khintchine inequality which will be used in the sequel.

\begin{lemma}[Lust-Piquard-Pisier]\label{nc-kh}
There exists a universal constant $C>0$ such that for each $2\leq p <\infty$, and any finite sequence $\{T_{j}\}_{j=1}^{N}\subseteq L_p(\mathbb{M}_{2^n})$ the following holds
\[
\left\|\{T_{j}\}_{j=1}^{N}\right\|_{L_p(\mathbb{M}_{2^n}, \ell_2^{cr})}\leq  \left(\mathbb{E}_{\delta}\left\|\sum_{j=1}^{N}\delta_{j} T_{j}\right\|_{L_{p}}^p\right)^{1/p}\leq C \sqrt{p}\left\|\{T_{j}\}_{j=1}^{N}\right\|_{L_p(\mathbb{M}_{2^n}, \ell_2^{cr})},
\]
where $\delta=\{\delta_{j}\}_{j=1}^{N}$ is a sequence of i.i.d. (unbiased) Rademacher random variables.
\end{lemma}
\chapter{Holomorphic quantum semigroups}\label{sec-hqs}

Recall that a bounded semigroup $\{S_{t}\}_{t\geq 0}$ on a real or complex Banach space $X$ is called holomorphic if there exists $\theta>0$ such that, for each $x\in X$ and $x^*\in X^*$
the function $t\rightarrow \left\langle x^*,S_t(x)\right\rangle$ has holomorphic extension in the sector $\{z\in\mathbb{C}:  \mathrm{Re}z\geq0,|\mathrm{arg}z|<\theta\}$. The purpose of this chapter is to prove that, for $1<p<\infty$, the quantum Ornstein-Uhlenbeck semigroup is holomorphic on $L_p(\M_{2^n})$ with constants independent of $n$; see Theorem \ref{e-zl-b}.
\section{Holomorphic quantum Ornstein-Uhlenbeck semigroup}

For $z\in\mathbb{C}$, define the Fourier multiplier
\begin{equation}\label{holo-extension}
e^{-z\mathcal L}(T)\coloneqq \sum_{\s\in\{0,1,2,3\}^{n}} \widehat T(\s)e^{-z|\s|}\sigma_{\s},
\end{equation}
where $T=\sum_{\s\in\{0,1,2,3\}^{n}}\widehat T(\s)\sigma_{\s}\in \M_{2^n}$. When $z=t\ge0$, this agrees with the quantum Ornstein--Uhlenbeck semigroup defined in \eqref{et-l-def}. The main result in this section is the following holomorphic extension theorem for quantum Ornstein-Uhlenbeck semigroups.

\begin{theorem}\label{e-zl-b}
For every $p \in(1, \infty)$, there exist constants $\theta=\theta_p \in(0, \frac{\pi}{2}]$ and $M_p \in[1, \infty)$ such that for every $n \in \mathbb{N}$ and for each $z\in \mathbb{C}$ with $|\arg z| \leq \theta$, one has
\begin{equation}\label{ezl-b}
\|e^{-z \mathcal{L}}\|_{L_{p} \rightarrow L_{p}} \leq M_p,\quad |\arg(z)|\leq \theta,
\end{equation}
where the operator $e^{-z\mathcal{L}}$ is defined in \eqref{holo-extension}.
\end{theorem}

Our proof follows Pisier's holomorphic-semigroup argument for $K$-convex Banach spaces (see \cite{Pi1982}), adapted to the product depolarizing semigroup. The key point is that the semigroup can be represented through commuting conditional expectations. More precisely, let $L:\mathbb{M}_2\to \mathbb{M}_2$ be as in  \eqref{l-def}. For each $j\in [n]$ and $t>0$, define the map $L_{j,t}:\mathbb{M}_{2^n}\rightarrow\mathbb{M}_{2^n}$ by setting
$$L_{j,t}=\mathbb{I}_2^{\otimes (j-1)} \otimes ((\mathbb{I}_2 - L)+e^{-t}L) \otimes \mathbb{I}_2^{\otimes (n-j)}$$
with $\mathbb{I}_2$ being the identity map over $\mathbb{M}_{2}$. It is not hard to see that $(L_{j,t})_{j=1}^n$ is a family of commuting projections on $L_2(\mathbb{M}_{2^n})$.
Set
\[
S_t\coloneqq \prod_{j=1}^nL_{j,t}.
\]
Moreover, the family of operators $\{S_{t}\}_{t>0}$ is exactly the quantum Ornstein-Uhlenbeck semigroup defined in \eqref{et-l-def}.
\begin{lemma}\label{st-e-tl}
For each $t>0$, we have
\begin{equation*}
S_t=e^{-t\mathcal{L}}.
\end{equation*}
\end{lemma}
\begin{proof}
Fix $t>0$, by the orthonormality of $\{\sigma_{\mathbf{s}}\}_{\s\in\{0,1,2,3\}^{n}}$, it suffices to verify that $S_{t}(\sigma_{\s})=e^{-t\mathcal{L}}(\sigma_{\s})$ for each $\s\in \{0,1,2,3\}^{n}$. Note that
\[
S_t(\sigma_{\mathbf{s}})=\prod_{j=1}^nL_{j,t}(\sigma_{\mathbf{s}})=\prod_{j=1}^nL_{j,t}(\otimes_{j=1}^n\sigma_{\mathbf{s}_j})=\otimes_{j=1}^n((\mathbb{I}_2-L)+e^{-t}L)(\sigma_{\mathbf{s}_j})=e^{-t|\s|}\sigma_\s,
\]
where we used the  definition of $L$ as in \eqref{l-def}. The assertion follows.
\end{proof}

We shall also use the following equivalent formula of $\{S_t\}_{t>0}$. For each $j\in [n]$, define $Q_j:\mathbb{M}_{2^n}\to \mathbb{M}_{2^n}$ by 
$$Q_j:=\mathbb{I}_2^{\otimes (j-1)} \otimes \mathcal{E}_2 \otimes \mathbb{I}_2^{\otimes (n-j)},$$
where $\mathcal{E}_2:\mathbb{M}_2\to\mathbb{M}_2$ is a conditional expectation given by $\mathcal{E}_2(A)=\tr_2(A)\mathbf{1}_2$, for all $A\in \mathbb{M}_2$.
By the construction, it follows that $\{Q_j\}_{j=1}^n$ forms a commuting family of conditional expectations. Thus $\{Q_{j}\}_{j=1}^n$ is a commuting family of norm-one projections on $L_{p}(\mathbb{M}_{2^n})$, $1\leq p\leq\infty$.

\begin{lemma}\label{lem-st-2}
For each $t>0$, we have
\[
S_t=\prod_{j=1}^{n}\left(Q_j+e^{-t}\left(\mathbb{I}_2^{\otimes_n}-Q_j\right)\right).
\]
\end{lemma}

\begin{proof} 
For any $t>0$ and $\s \in \{0,1,2,3\}^n$, it is not hard to check that for each $j\in [n]$, 
$$\big(Q_j+e^{-t}(\mathbb{I}_2^{\otimes_n}-Q_j)\big)(\sigma_{\s})=e^{-t}\sigma_{\s}$$
whenever $\s_j\neq 0$, and $\big(Q_j+e^{-t}(\mathbb{I}_2^{\otimes_n}-Q_j)\big)(\sigma_{\s})=\sigma_{\s}$ whenever $\s_j=0$.
It follows directly that 
$$\prod_{j=1}^n(Q_j+e^{-t}(\mathbb{I}_2^{\otimes_n}-Q_j))(\sigma_{\s})= e^{-t|\s|}\sigma_{\s}=e^{-t\mathcal{L}}(\sigma_{\s}).$$
Since $t>0$ and $\s\in \{0,1,2,3\}^n$ are arbitrary, the desired assertion follows from  Lemma  \ref{st-e-tl}.
\end{proof}

We next recall two standard facts from  semigroup theory in Banach space-valued setting. Given a Banach space $X$, recall that for each $k\in [n]$ and $f\in L_{p}(\{-1,1\}^n;X)$, the \emph{Rademacher projection} $\mathrm{Rad}^{=k}$ is defined by 
\begin{equation}\label{rad cls}
\mathrm{Rad}^{=k}(f)=\sum_{|A|=k}\widehat{f}(A)w_A.
\end{equation}
We say  a Banach space $X$ is $K$-convex if  
 $$\sup_{k\geq0}\|\mathrm{Rad}^{=k}\|_{L_{p}(\{-1,1\}^n;X)\rightarrow L_{p}(\{-1,1\}^n;X)}<\infty.$$ 
According to \cite[Theorem 2.1]{Pi1982}, a Banach space $X$ is $K$-convex is equivalent to that $X$ is non-trivial type. Hence, each noncommutative Lebesgue space $L_p(\mathcal{M})$ with $1<p<\infty$ is $K$-convex.

The following result is a restatement of the fundamental theorem of Pisier \cite{Pi1982}.

\begin{lemma}[{\cite[Lemma 1.5]{Pi1982}}]\label{rpo} 
Suppose that $X$ is a $K$-convex Banach space. For any  commuting norm-one projections $P_1,\dots,P_N \in B(X)\footnote{For a Banach space $E$, $B(E)$ stands for the space of all bounded linear operators from $E$ to itself.}$, there exists a real
number $0 <\rho < 2$ (depending on $N$ and $X$) such that
\[
\left\|\prod_{j=1}^N(\mathbb{I}-P_j)\right\|\leq \rho^N
\]
where $\mathbb{I}$ denotes the identity map on $X$.
\end{lemma}
Next lemma is due to  Kato and Beurling \cite{B1970,Ka1970} which is known as the \emph{Kato--Beurling criterion}.
\begin{lemma}[Kato--Beurling]\label{kato-burling}
Let $\{T_t\}_{t\geq0}$ be a strongly continuous semigroup on some Banach space $X$. Suppose that there exists a natural number $N$ such that
\[
\limsup_{t\to 0^{+}}\left\|(\mathbb{I}-T_t)^N\right\|_{X\to X}^{1/N} =\rho< 2,
\]
where $\mathbb{I}$ is the identity operator in $B(X)$. Then there exists $\theta>0$ such that $\{T_{t}\}_{t\geq 0}$ admits a holomorphic semigroup extension $z\mapsto T_{z}$ to the sector
\[
\Sigma_{\theta}\coloneqq\{z\in\mathbb{C}\setminus\{0\}: |\arg z|<\theta\}.
\]
Furthermore, after decreasing the angle if necessary, the extension is uniformly bounded on the closed smaller sector.
\end{lemma}  



We now are ready to prove Theorem \ref{e-zl-b}.
\begin{proof}[Proof of Theorem \ref{e-zl-b}]
For any $1<p<\infty$, by Lemma \ref{st-e-tl} and Lemma \ref{kato-burling}, it suffices to show that there exist $0<\rho<2$ and $N\in\mathbb{N}$ such that
\begin{equation}\label{st-rho}
\sup_{t>0}\|(\mathbb{I}-S_t)^{N}\|_{L_{p}\rightarrow L_{p}}\leq \rho^N.
\end{equation}

For a fixed $t>0$, let $\{\xi_j\}_{j=1}^n$ be a family of independent and identically distributed functions defined on $(\Omega,\mathcal{F},\mathbb{P}):=(\{0,1\}^n,\otimes_{i=1}^n\mu)$ satisfying
\[
\mathbb{P}(\{\xi_j=0\})=1-e^{-t}\quad\mbox{and}\quad \mathbb{P}(\{\xi_j=1\})=e^{-t}.
\]
Define $\pi:(\Omega,\mathcal{F},\mathbb{P})\rightarrow B(L_p(\mathbb{M}_{2^n}))$ by setting
\begin{equation}\label{pi-def}
\omega\mapsto\prod_{j=1}^n(Q_j+\xi_j(\omega)(\mathbb{I}_2^{\otimes_n}-Q_j)).
\end{equation}
By the independence of $\{\xi_j\}_{j=1}^n$, we set that 
\begin{align}\label{st-new}
\int_\Omega\pi(\omega)d\mathbb{P}(\omega)&=\int_\Omega\prod_{j=1}^n(Q_j+\xi_j(\omega)(\mathbb{I}_2^{\otimes_n}-Q_j))d\mathbb{P}(\omega)\nonumber\\
&=\prod_{j=1}^n\int_\Omega(Q_j+\xi_j(\omega)(\mathbb{I}_2^{\otimes_n}-Q_j))d\mathbb{P}(\omega)\nonumber\\
&=\prod_{j=1}^n[(1-e^{-t})Q_j+e^{-t}\mathbb{I}_2^{\otimes_n}]\\
&=\prod_{j=1}^n[Q_j+e^{-t}(\mathbb{I}_2^{\otimes_n}-Q_j)]\nonumber=S_t\nonumber,
\end{align}
where the last equality is due to Lemma \ref{lem-st-2}.


For $N\in\mathbb{N}$, let $(\Omega^N, \mathcal{F}^{\otimes_N},\mathbb{P}^{\otimes_N})$ be the $N$-Cartesian product of the probability space $(\Omega, \mathcal{F},\mathbb{P})$. For each $1\leq k\leq N$, let
$$P_k:(\Omega^N, \mathcal{F}^{\otimes_N},\mathbb{P}^{\otimes_N})\rightarrow(\Omega, \mathcal{F},\mathbb{P}),$$
that is, $P_k(\omega)=\omega_k$.
 We claim that, for each given $\omega\in \Omega^N$, $\{\pi\circ P_{\ell}(\mathbf{\omega})\}_{\ell=1}^N$ forms a commuting norm-one projections family of $B(L_p(\mathbb{M}_{2^n}))$. Once this claim is proved, applying Lemma \ref{rpo} and the $K$-convexity of $L_{p}(\M_{2^{n}})$ for all $1<p<\infty$,  there exists $0<\rho<2$ such that
\begin{equation}\label{eq-pi-1}
\left\|\prod_{\ell=1}^N(\mathbb{I}-\pi\circ P_{\ell}(\mathbf{\omega}))\right\|_{L_{p}\rightarrow L_{p}}\leq \rho^N.
\end{equation}
Furthermore, by \eqref{st-new}, we have
\begin{align*}
(\mathbb{I}-S_t)^N&=\left(\int_\Omega\mathbb{I}-\pi(\omega)d\mathbb{P}(\omega)\right)^N\\
&=\prod_{\ell=1}^N\int_{\Omega^N}(\mathbb{I}-\pi\circ P_{\ell}(\mathbf{\omega}))d\mathbb{P}^{\otimes_N}(\mathbf{\omega})\\
&=\int_{\Omega^N}\prod_{\ell=1}^N(\mathbb{I}-\pi\circ P_{\ell}(\mathbf{\omega}))d\mathbb{P}^{\otimes_N}(\mathbf{\omega}).
\end{align*}
Then the Jensen inequality entails that
$$\|(\mathbb{I}-S_t)^N\|_{L_{p}\rightarrow L_{p}}\leq \int_{\Omega^N}\|\prod_{\ell=1}^N(\mathbb{I}-\pi\circ P_{\ell}(\mathbf{\omega}))\|_{L_{p}\rightarrow L_{p}}d\mathbb{P}^{\otimes_N}.$$
which, together with \eqref{eq-pi-1}, gives
$$\|(\mathbb{I}-S_t)^N\|_{L_{p}\rightarrow L_{p}}\leq \int_{\Omega^N}\rho^Nd\mathbb{P}^{\otimes_N}=\rho^N.$$
This proves \eqref{st-rho}.

It now remains to verify the claim. Indeed, for any $\mathbf{\omega} \in\Omega^N$ and  $1\leq \ell\leq N$, by the definition of the map $\pi$, we have
$$\pi\circ P_{\ell}(\mathbf{\omega})=\pi(\omega_l)=\prod_{j=1}^n(Q_j+\xi_j(\omega_{\ell})(\mathbb{I}_2^{\otimes_n}-Q_j))=\prod_{\substack{\xi_j(\omega_{\ell})=0\\j\in[n]}}Q_j.$$
Since $(Q_j)_{j=1}^n$ is a family  norm-one projections on $L_p(\mathbb{M}_{2^n})$ with $1\leq p\leq\infty$, it follows that
$\pi\circ P_l(\mathbf{\omega})$ is also a norm-one projection.
Since the family $(Q_j)_{j=1}^n$ are  commuting,  for any $\ell_1,\ell_2\in[N]$ and $\ell_1\neq \ell_2$, it follows that
\begin{align*}
\pi\circ P_{\ell_1}(\mathbf{\omega})\cdot\pi\circ P_{\ell_2}(\mathbf{\omega})&=\left(\prod_{\substack{\xi_j(\omega_{\ell_1})=0\\j\in[n]}}Q_j\right)\left(\prod_{\substack{\xi_k(\omega_{\ell_2})=0\\k\in[n]}}Q_k\right)\\
&=\prod_{\substack{\xi_j(\omega_{\ell_1})=0,~\xi_k(\omega_{\ell_2})=0\\j,k\in[n]}}Q_j\cdot Q_k\\
&=\prod_{\substack{\xi_j(\omega_{\ell_1})=0,~\xi_k(\omega_{\ell_2})=0\\j,k\in[n]}}Q_k\cdot Q_j\\
&=\pi\circ P_{\ell_2}(\mathbf{\omega})\cdot\pi\circ P_{\ell_1}(\mathbf{\omega}).
\end{align*}
This verifies the claim, and
the proof is complete.
\end{proof}

\section{Quantum Rademacher projection}\label{sec-qr-projection}
The holomorphic estimate of Theorem \ref{e-zl-b} has two immediate consequences that will be used repeatedly in the sequel. First is a dimension-free bound for the spectral projections onto fixed Pauli--Fourier levels.  The second is a bounded functional calculus for the multiplier $w^{\mathcal L}$ on suitable lens domains inside the unit disc.  These two estimates are the analytic bridge between holomorphicity of the quantum Ornstein--Uhlenbeck semigroup and the degree sensitive estimates in the next two chapters.

For any $k\in [n]$, define the \emph{quantum Rademacher projection} $\mathrm{Rad}^{=k}_{\mathbb{M}_2}:\mathbb{M}_{2^n}\to\mathbb{M}_{2^n}$ of level $k$ is defined by
\begin{equation}
\mathrm{Rad}^{=k}_{\mathbb{M}_2}(T) = \sum_{\substack{\mathbf{s} \in \{0,1,2,3\}^n \\ |\mathbf{s}| =k}} \widehat{T}(\mathbf{s}) \sigma_{\mathbf{s}},
\end{equation}
where $T=\sum_{\mathbf{s} \in \{0,1,2,3\}^n} \widehat{T}(\mathbf{s}) \sigma_{\mathbf{s}}$. If the context is clear, we simply use $\mathrm{Rad}^{=k}$ instead of $\mathrm{Rad}^{=k}_{\mathbb{M}_2}$, and simply denote $\mathrm{Rad}^{=1}$ by $\mathrm{Rad}$. Notice that $\mathrm{Rad}^{=k}$ is the spectral projection of the number operator $\mathcal L$ corresponding to the eigenvalue $k$.

\begin{lemma}\label{radk-pp-ak}
Let $1<p<\infty$, and let $\theta=\theta_p\in(0,\pi/2]$ and $M_p$ be as in Theorem \ref{e-zl-b}. Then, for any $k \in[n]$, we have
$$\|\mathrm{Rad}^{=k}\|_{L_{p}\rightarrow L_{p}}\leq M_pe^{ak},\quad a=\frac{\pi}{\tan \theta}.$$
\end{lemma}

\begin{proof} Let $a=\frac{\pi}{\tan \theta}$.  For each $k \in[n]$ and every $-\pi \leq t\leq \pi$ (which ensures that $|\arg (a+it)|\leq \theta$),
applying the orthogonality of trigonometric series, we obtain
$$
\frac{1}{2 \pi} \int_{-\pi}^{\pi} e^{i k t} e^{-(a+i t) \mathcal{L}} \mathrm{d} t=\frac{1}{2 \pi} \int_{-\pi}^{\pi} \sum_{j=0}^{n} e^{i k t-(a+i t) j} \operatorname{Rad}^{=j} \mathrm{~d} t=e^{-k a} \mathrm{Rad}^{=k}.$$
Then, we have 
\begin{align*}
 \|\mathrm{Rad}^{=k}\|_{L_{p}\rightarrow L_{p}}
&=e^{k a}\|\frac{1}{2 \pi} \int_{-\pi}^{\pi} e^{i k t} e^{-(a+i t) \mathcal{L}}\mathrm{~d}t\|_{L_{p}\rightarrow L_{p}} \\
&\leq \frac{e^{k a}}{2 \pi} \int_{-\pi}^{\pi}\|e^{-(a+i t) \mathcal{L}}\|_{L_{p}\rightarrow L_{p}} \mathrm{~d}t \\
&\leq M_p e^{ka}, 
\end{align*}
where the last inequality is due to Theorem \ref{e-zl-b}. The proof is complete.
\end{proof}

We next explain why the following lens domains \eqref{omega-r} naturally enter the argument.  Estimates for heat smoothing and for low-degree lower bounds are not obtained from the positive semigroup parameter $w=e^{-t}$ alone.  They require applying scalar polynomial inequalities to the spectral variable $w$ and then transferring those polynomial estimates to the operator-valued multiplier $w^{\mathcal L}$.  In particular, the Hahn--Banach and conformal-mapping arguments used later require uniform bounds for $w^{\mathcal L}$ on domains that touch the unit circle at both $1$ and $-1$.  The point $1$ corresponds to the usual heat semigroup, while the point $-1$ is needed for Chebyshev-type and low-degree interpolation arguments.  The anti-isomorphism $\Phi_\rho$ from Lemma \ref{var-prop} supplies the required control of negative multipliers, since it implements the sign $(-1)^{|\s|}$ on each Pauli level.

Let $\mathbb{D}=\{z \in \mathbb{C}:|z|<1\}$ the open unit disc. For $r \in[1, \infty)$, define the lens domain
\begin{equation}\label{omega-r}
\Omega(r) \stackrel{\text { def }}{=}\left\{z \in \mathbb{C}: \max \left\{\left|z-i \sqrt{r^{2}-1}\right|,\left|z+i \sqrt{r^{2}-1}\right|\right\}<r\right\} \subseteq \mathbb{D}.
\end{equation}
Note that $\partial \Omega(r)$ is a Jordan curve with an interior angle 
\begin{equation}\label{theta-r}
\theta(r)=2 \arcsin (1 / r)
\end{equation} 
at $z= \pm 1$.

\begin{lemma}[{\cite[Lemma 27]{EI2020}}]\label{con-map}
For $\alpha \in(0, \infty)$, let $\phi_{\alpha}: \mathbb{C} \backslash(-\infty, 0] \rightarrow \mathbb{C}$ be a holomorphic branch of $z \mapsto z^{\alpha}$. Also, consider the M\"{o}bius transformations $\psi_{1}(z)=\frac{1-z}{1+z}$ and $\psi_{2}(z)=\frac{z+1}{z-1}$. Then, for every $r \in[1, \infty)$, the map $\psi_{1} \circ \phi_{\pi / \theta(r)} \circ \psi_{1}$ is a conformal equivalence between $\Omega(r)$ and the unit disc $\mathbb{D}$. Furthermore, the map $\psi_{2} \circ \phi_{\pi /(2 \pi-\theta(r))} \circ \psi_{2}$ is a conformal equivalence between the complement $\overline{\Omega(r)}^{c}$ of $\overline{\Omega(r)}$ and the complement $\overline{\mathbb{D}}^{c}$ of the closed unit disc.
\end{lemma}

\begin{lemma}\label{or-k} 
For any $p \in(1, \infty)$, there exist constants $r_p\in[1, \infty)$ and $K_p \in[1, \infty)$ such that
$$\sup _{w \in \Omega(r_p)} \|w^{\mathcal{L}}\|_{L_{p}\rightarrow L_{p}} \leq K_p.$$
\end{lemma}
\begin{proof} 
Let $\theta \in\left(0, \frac{\pi}{2}\right]$, and let  $a=\frac{\pi}{\tan \theta}$. Take $w\in \mathbb{D}$ with $|w|<e^{-2a}$. According to Lemma \ref{radk-pp-ak}, we have
\begin{align*}
 \|w^{\mathcal{L}}\|_{L_{p}\rightarrow L_{p}}
&\leq \sum_{k=0}^{n}|w|^{k}\left\|\mathrm{Rad}^{=k}\right\|_{L_{p}\rightarrow L_{p}}\\
&\leq M_p \sum_{k=0}^{n}(|w| e^{a})^{k},
\end{align*}
which implies that for $|w|<e^{-2a}$, we also have the estimate
$$
\|w^{\mathcal{L}}\|_{L_{p}\rightarrow L_{p}} \leq \frac{M_p}{1-e^{-a}}.
$$

For every $w \in\{e^{-z}:|\arg z|<\theta\}$, applying Lemma  \ref{var-prop} (iii) and  Theorem \ref{e-zl-b},  we have
\[
\left\|w^{\mathcal{L}}\right\|_{L_{p}\rightarrow L_{p}}=\left\|(-w)^{\mathcal{L}}\right\|_{L_{p}\rightarrow L_{p}} \leq M_p.
\]
By the above estimates, one can easily deduce that there exists $r_p \in[1, \infty)$ such that
\begin{equation}\label{or-def}
\Omega(r_p) \subseteq\left\{e^{-z}:|\arg z|<\theta\right\} \cup\left\{-e^{-z}:|\arg z|<\theta\right\} \cup\{w \in \mathbb{D}:|w|<e^{-2 a}\}.
\end{equation}
and the desired result follows with $K_p=\frac{M_p}{1-e^{-a}}$.
\end{proof}

\begin{lemma}\label{radk-ppb}
For each $k\in[n]$ and $T\in\mathbb{M}_{2^n}$, we have 
$$\sup_{k\geq0}\|\mathrm{Rad}^{=k}\|_{L_{p}\rightarrow L_{p}}^{\frac{1}{k}}\leq
\begin{cases} 
\sqrt{p-1}, & 2\leq p<\infty; \\ 
\sqrt{\frac{1}{p-1}}, & 1<p\leq2, 
\end{cases}.$$
\end{lemma}
\begin{proof} 
We first prove the case $2\leq p<\infty$ and use the duality argument to derive the case $1<p\leq 2$. It is clear that for $T\in \mathbb{M}_{2^{n}}$ we have
\[
\mathrm{Rad}^{=k}(T)=e^{tk}e^{-t\mathcal{L}}(\mathrm{Rad}^{=k}(T)).
\]
By the hypercontractivity (i.e., Theorem \ref{hc-1}), we obtain
\begin{align*}
\|\mathrm{Rad}^{=k}(T)\|_{L_{p}}&=\|e^{tk}e^{-t\mathcal{L}}(\mathrm{Rad}^{=k}(T))\|_{L_{p}}\\
&\leq e^{tk}\|\mathrm{Rad}^{=k}(T)\|_{L_2}\\
&\leq e^{tk}\|T\|_{L_2}\\
&\leq e^{tk}\|T\|_{L_{p}},
\end{align*}
where we used the fact $2\leq p<\infty$ and $(\mathbb{M}_{2^n},\mathrm{tr})$ is a probability space in the last inequality.
Taking $t=\frac{1}{2}\log(p-1)$,  it follows that
$$\sup_{k\geq0}\|\mathrm{Rad}^{=k}\|_{L_{p}\rightarrow L_{p}}\leq (p-1)^{k/2},\quad 2\leq p<\infty.$$
By the duality argument, we obtain the estimate for the case $1<p\leq2$, which completes our proof.
\end{proof}

\chapter{Quantum heat smoothing estimates}\label{sec-qhse}

This chapter develops the high-frequency side of the degree-sensitive semigroup theory on the quantum Boolean cube.  Throughout the chapter, $\mathcal{T}^{\geq d}(\M_{2^n})$ and $\mathcal{P}^{\leq d}(\M_{2^n})$ denote the corresponding Pauli--Walsh tail and low-degree subspaces, respectively.

The first section proves a quantum analogue of the Mendel--Naor heat smoothing estimate \cite{MN2014} for observables in $\mathcal{T}^{\geq d}(\M_{2^n})$.  The proof combines the holomorphic functional calculus from Chapter \ref{sec-hqs} with a conformal mapping argument on lens domains.  The second section gives a quantum analogy of the first-tail (i.e., $d=1$) Heilman--Mossel--Oleszkiewicz type strong contraction estimate \cite{HMO2017}. Our approach is inspired by \cite{HMO2017}, but the proof replaces the pointwise order arguments used in the commutative setting by a probabilistic representation for noncommutative Markov maps.  We conclude the chapter with reverse Bernstein--Markov inequalities for observables whose Pauli-Walsh spectrum is localized in the finite band $[d,d+m]$.
\section{Quantum Mendel-Naor heat smoothing estimates}

Heat smoothing estimates quantify how fast the heat semigroup damps the part of a function supported on high frequencies.  In the classical framework, this idea goes back to Meyer \cite{Me1984} in the Gaussian setting and was developed systematically by Mendel and Naor \cite{MN2014} in the Boolean cube setting.  The following theorem is the quantum counterpart needed later for the tail Talagrand and covariance estimates.

\begin{theorem}\label{mn-k-tail}
For $n\geq1, d \in [n]$ and $1<p<\infty$, there exist $C_p,c _p\in(0,\infty)$ and $A_p\in[1,\infty)$  depending only $p$ such that, for every  $T\in\mathcal{T}^{\geq d}(\mathbb{M}_{2^n})$,  we have
\begin{equation}\label{mn-k-tail-ineq}
\big\|e^{-t \mathcal{L}}(T)\big\|_{L_{p}} \leq C_p e^{-c_p d \min \{t, t^{A_p}\}}\|T\|_{L_{p}},\quad \forall t \geq 0.
\end{equation}
\end{theorem}

The proof is postponed until after Proposition \ref{cpro}. Recall that the $d$-th Chebyshev polynomial of the first kind, denoted by $T_{d}(x)$, is the unique polynomial of degree $d$ satisfying $T_{d}(\cos \theta)=\cos (d \theta)$ for all $\theta \in \mathbb{R}$. Chebyshev's inequality \cite[p.235]{BE1995} asserts that the $d$-th Chebyshev polynomial of the first kind is characterized by the extremal property
\begin{equation}\label{cp}
\left|T_{d}(x)\right|=\max \left\{|P(x)|: \operatorname{deg}(P) \leq d \text { and }\|P\|_{\mathcal{C}_{([-1,1])}}=1\right\},\; \forall x \in \mathbb{R} \backslash[-1,1].
\end{equation}
Here and what follows, for each polynomial $P\in \mathcal{P}(\mathbb{K})$, defined on a compact $\mathbb{K}\subseteq \mathbb{C}$, we denote $\|P\|_{\mathcal{C}(\mathbb{K})}=\sup_{z\in \mathbb{K}} |P(z)|$ its  supremum
norm in $\mathbb{K}$.
For each compact $\mathbb{K}\subseteq \mathbb{C}$ and $d\in \mathbb{N}$, denote 
$$\mathcal{P}^{\leq d}(\mathbb{K})=\{P:\mathbb{K}\rightarrow\mathbb{C}: \  \mbox{$P$ is a polynomial with $\mathrm{deg}(P)\leq d$}\}$$ 
and
\[
\mathcal{T}^{\geq d}(\mathbb{K})=\left\{P:\mathbb{K}\rightarrow\mathbb{C}: P(w)=\sum_{n=d}^{N}a_{n}w^{n},~N\in \mathbb{N}\right\}.
\] 

\begin{proposition}\label{cpro}
Let $n\geq1, d \in [n]$, $1<p<\infty$,  and $\Omega \subseteq \mathbb{D}$ be a simply connected domain bounded by a Jordan curve $\gamma$ such that $(-1,1) \subseteq \Omega$ and suppose that
\begin{equation}\label{cpro-k-i}
\sup _{w \in \Omega} \|w^{\mathcal{L}}\|_{L_{p}\rightarrow L_{p}} \leq K_p.
\end{equation}
Let $\phi_{\Omega}: \Omega \rightarrow \mathbb{D}$ be a conformal mapping of $\Omega$ onto $\mathbb{D}$ with $\phi_{\Omega}(0)=0$ which extends continuously as a homeomorphism between $\Omega \cup \gamma$ and $\overline{\mathbb{D}}$. Then, for  every $T\in \mathcal{T}^{\geq d}(\mathbb{M}_{2^n})$, we have
\[
\|e^{-t \mathcal{L}} (T)\|_{L_{p}} \leq K_p|\phi_{\Omega}(e^{-t})|^{d}\|T\|_{L_{p}},\quad \forall t>0.
\]
\end{proposition}
\begin{proof}
Consider the linear functional $\zeta_{t}:\left(\mathcal{T}^{\geq d}(\overline{\Omega}),\|\cdot\|_{\mathcal{C}(\overline{\Omega})}\right) \rightarrow \mathbb{C}$ defined by
\[
\zeta_{t}(P)=P(e^{-t}), \quad t\geq 0.
\]
For each $P\in \mathcal{T}^{\geq d}(\overline{\Omega})$, consider the function $h: \Omega \rightarrow \mathbb{C}$ given by
\[
h(w)\coloneqq\frac{P(w)}{\phi_{\Omega}(w)^{d}},\quad \forall w \in \Omega.
\]
Note that $\phi_{\Omega}$ does not vanish on $\Omega \backslash\{0\}$ and  has a single root at $0$. Therefore, the multiplicity of the root $0$ in the numerator is at least that in the denominator, ensuring that  $h$ is a holomorphic function on $\Omega$. Furthermore,  we have $$|\phi_{\Omega}(w)|=1,\quad w \in \gamma,$$ 
since $\phi_{\Omega}$ is a homeomorphism between $\gamma$ and $\partial \mathbb{D}$ (such a conformal map $\phi_{\Omega}$ always exists by Caratheodory's theorem, see \cite[Theorem 5.1.1]{K2006}). Therefore, by the maximum principle,
$$
\sup _{w \in \Omega}|h(w)|=\max _{w \in \gamma}|h(w)|=\max _{w \in \gamma} \frac{|P(w)|}{|\phi_{\Omega}(w)|^{d}}=\|P\|_ {\mathcal{C}(\gamma)}=\|P\|_{\mathcal{C}(\overline{\Omega})},$$
which implies 
$$|P(w)| \leq|\phi_{\Omega}(w)|^{d}\|P\|_{\mathcal{C}(\overline{\Omega})},\quad w \in \overline{\Omega}.$$
Applying the above inequality for $w=e^{-t} \in \Omega$, we deduce that the linear functional $\zeta_{t}$ satisfies
\begin{equation}\label{zeta-b}
\left|\zeta_{t}(P)\right| \leq\left|\phi_{\Omega}\left(e^{-t}\right)\right|^{d}\|P\|_{\mathbb{C}(\overline{\Omega})},\quad \forall P \in\mathcal{T}^{\geq d}(\overline{\Omega}) .
\end{equation}
Therefore, by the Hahn-Banach theorem and the Riesz representation theorem, there exists a complex measure $\tau_{t}$ on $\overline{\Omega}$ such that $$\left\|\tau_{t}\right\|_{V(\overline{\Omega})} \leq\left|\phi_{\Omega}\left(e^{-t}\right)\right|^{d}$$ 
and
\begin{equation}\label{p-zeta-i}
\int_{\overline{\Omega}} P(w) \mathrm{d} \tau_{t}(w)=\zeta_{t}(P)=P\left(e^{-t}\right), \quad P \in\mathcal{P}(\overline{\Omega}). 
\end{equation}
Since $\widehat{T}(\mathbf{s})=0$ when $|\mathbf{s}|<d$, by the Jensen inequality, we obtain
\begin{align*}
\|e^{-t \mathcal{L}}(T) \|_{L_{p}}
&=\Big\|\sum_{\mathbf{s}\in\{0,1,2,3\}^n,|\mathbf{s}|\geq d}e^{-t |\mathbf{s}|}\widehat{T}(\mathbf{s}) \sigma_{\mathbf{s}}\Big\|_{L_{p}}\\
&\stackrel{\eqref{p-zeta-i}}{=}\Big\|\int_{\overline{\Omega}}\sum_{\mathbf{s}\in\{0,1,2,3\}^n,|\mathbf{s}|\geq d}w^{|\mathbf{s}|}\widehat{T}(\mathbf{s}) \sigma_{\mathbf{s}}\mathrm{~d}\tau_t(w)\Big\|_{L_{p}}\\
&=\Big\|\int_{\overline{\Omega}} w^{\mathcal{L}}(T) \mathrm{~d}\tau_{t}(w)\Big\|_{L_{p}} \leq \int_{\overline{\Omega}}\|w^{\mathcal{L}} (T)\|_{L_{p}} \mathrm{d}|\tau_{t}|(w) \\
& \stackrel{\eqref{cpro-k-i}}{\leq} K_p \int_{\overline{\Omega}}\|T\|_{L_{p}} \mathrm{d}|\tau_{t}|(w)=K_p\left\|\tau_{t}\right\|_{V(\overline{\Omega})}\|T\|_{L_{p}} \\
& \leq K_p |\phi_{\Omega}(e^{-t})|^{d}\|T\|_{L_{p}},
\end{align*}
which concludes the proof.
\end{proof}

\begin{proof}[Proof of Theorem \ref{mn-k-tail}]
 Let $T\in \mathcal{T}^{\geq d}(\mathbb{M}_{2^{n}})$ and fix $1<p<\infty$. Choose a sectorial angle $0<\theta<\theta_p$, where $\theta_p$ is supplied by the holomorphic semigroup theorem of Chapter \ref{sec-hqs}, and put
\[
a\coloneqq\frac{\pi}{\tan\theta}>0 .
\]
We divide the proof into two cases.

\textbf{Case I}: For $t\geq 2a$, by Lemma \ref{radk-pp-ak},  we have the following estimate
\begin{align*}
\|e^{-t\mathcal{L}}(T)\|_{L_{p}} &\leq \sum_{k=d}^{n} e^{-t k}\|\mathrm{Rad}^{=k} (T)\|_{L_{p}}\\ 
&\leq M_p\sum_{k=d}^{n} e^{-t k}e^{ak}\|T\|_{L_{p}}
\leq K_p e^{-\frac{1}{2} t d}\|T\|_{L_{p}},
\end{align*}
where $K_p=\frac{M_p}{1-e^{-a}}$. 

\textbf{Case II}: For $t \leq2 a$, by Lemma \ref{or-k}, there exist constants $r_p \in[1, \infty)$ and  $K_p \in[1, \infty)$  such that
$$\sup _{\omega \in \Omega(r_p)} \sup _{n \in \mathbb{N}}\|\omega^{\mathcal{L}}\|_{L_{p} \rightarrow L_{p}} \leq K_p, $$
where $\Omega(r_p)$ is the domain \eqref{or-def}. By Proposition \ref{cpro},  we have
$$\|e^{-t \mathcal{L}} (T)\|_{L_{p}} \leq K_p|\phi_{\Omega(r_p)}(e^{-t})|^{d}\|T\|_{L_{p}},\quad \forall t \geq 0.
$$
By Lemma \ref{con-map}, $\theta=\theta(r_p)=2\arcsin\left(\frac{1}{r_{p}}\right)\in(0,\pi)$, and the function $\phi_{\Omega(r_p)}$ can be expressed as
$$
\phi_{\Omega(r_p)}(e^{-t})=\psi_1\circ \phi_{\pi / \theta} \circ \psi_{1}=\frac{(1+e^{-t})^{\pi / \theta}-(1-e^{-t})^{\pi / \theta}}{(1+e^{-t})^{\pi / \theta}+(1-e^{-t})^{\pi / \theta}} .
$$
To analyze the decay of $\phi_{\Omega(r_p)}$, we take the logarithm and use the inequality $\log(1-x)\leq-x$ on $(0,1)$:
\begin{align*}
\log |\phi_{\Omega(r_p)}(e^{-t})| & =\log \Big(1-\frac{2(1-e^{-t})^{\pi / \theta}}{(1+e^{-t})^{\pi / \theta}+(1-e^{-t})^{\pi / \theta}}\Big) \\
& \leq-\frac{2(1-e^{-t})^{\pi / \theta}}{(1+e^{-t})^{\pi / \theta}+(1-e^{-t})^{\pi / \theta}}.
\end{align*}
For all $t \leq 2 a$, we further simplify the bound
\begin{align*}
-\frac{2(1-e^{-t})^{\pi / \theta}}{(1+e^{-t})^{\pi / \theta}+(1-e^{-t})^{\pi / \theta}} & \leq-2^{1-\frac{\pi}{\theta}}(1-e^{-t})^{\pi / \theta} \\
& \leq \frac{-2^{1-\frac{\pi}{\theta}}(1-e^{-2 a})^{\pi / \theta}}{(2 a)^{\pi / \theta}} t^{\pi / \theta},
\end{align*}
where the last inequality  is because the function $f(x)=\frac{1-e^{-x}}{x}$ is monotonically decreasing on $x\in(0,\infty)$. Thus,  for all $0 \leq t \leq 2 a,$  we have
\begin{equation*}
|\phi_{\Omega(r)}\left(e^{-t}\right)| \leq e^{-c t^{\pi / \theta}}, 
\end{equation*}
where $c=\frac{2^{1-\frac{\pi}{\theta}}\left(1-e^{-2 a}\right)^{\pi / \theta}}{(2 a)^{\pi / \theta}}$. The proof is complete.
\end{proof}
\section{Quantum Heilman-Mossel-Oleszkiewicz heat smoothing estimate}

The preceding estimate gives heat smoothing on \(d\)-tail spaces for arbitrary \(d\).  We next prove a first-tail strong contraction estimate in the spirit of Heilman--Mossel--Oleszkiewicz \cite{HMO2017}. In the hilbertian case the spectral gap of the Laplacian operator is equivalent to the following exponential ergodicity of the semigroup, that is, for every $T\in \mathbb{M}_{2^{n}}$ with $\mathrm{tr}(T)=0$, we have
\begin{equation}\label{exponential-ergodic}
\|e^{-t\mathcal{L}}(T)\|_{L_{2}}\leq e^{-t}\|T\|_{L_{2}}.
\end{equation}
The point of this section is to derive  a general case of \eqref{exponential-ergodic}. After extensive investigations such as \cite{EI2020,EI2023,HMO2017,Vo2024}, the heat smoothing conjecture, that is, Conjecture \ref{conj-1}, remains open (even for $X=\mathbb{R}$ there). Particularly, Heilman, Mossel and Oleszkiewicz \cite{HMO2017} resolved Conjecture \ref{conj-1} with $X=\mathbb{R}$ and $d=1$. The next theorem is a quantum counterpart of the main estimate of Heilman, Mossel and Oleszkiewicz \cite{HMO2017}.

For $p\in(1,\infty)\setminus\{2\}$, set
\[
\kappa(p)\coloneqq\inf_{x>1}\frac{x^{p/(p-2)}+x^{-p/(p-2)}}{x-x^{-1}}.
\]


\begin{theorem}\label{heat-semigroup-f2828b}
For each $p\in (1,\infty)\setminus\{2\}$, the following heat smoothing inequality holds for each self-adjoint $T\in \mathbb{M}_{2^n}$ with $\mathrm{tr}(T)=0$,
\[
\|e^{-t\mathcal{L}}(T)\|_{L_p}\leq \exp\left(-\frac{(4p-4)t}{p^2(1+\kappa(p)^{-2})}\right)\|T\|_{L_{p}}.
\]
\end{theorem} 

We first isolate a spectral approximation bound for the heat semigroup.


\begin{lemma}\label{ete-22}
For any $\epsilon,t\in (0,\infty)$, we have
\[
\|e^{-t\mathcal{L}}-e^{-(\epsilon+t)\mathcal{L}}\|_{L_2\to L_2}\leq \frac{\epsilon}{t}.
\]
\end{lemma}
\begin{proof}
Without loss of generality, we assume that
\[
T=\sum_{\s\in\{0,1,2,3\}^{n}}\widehat{T}(\s)\sigma_{\s}.
\]

We have
\begin{equation}\label{basic-estimate-3-1}
e^{-ts}-e^{-(t+\epsilon)s}\leq \frac{\epsilon}{t},\quad\mbox{for all }s\geq 0.
\end{equation}
Indeed, let $F(s)\coloneqq e^{-ts}-e^{-(\epsilon+t)s}$ for all $s\geq 1$. It is easy to see that
\[
F^{\prime}(s)=e^{-ts}\left((t+\epsilon)e^{-\epsilon s}-t\right),\quad\mbox{for }s\geq 1.
\]
The function $(t+\epsilon)e^{-\epsilon s}-t$ is strictly decreasing and
vanishes at $s_{0}=\frac{1}{\epsilon}\log\left(1+\frac{\epsilon}{t}\right)$.  Hence, $F$ is increasing on $[0,s_0]$ and decreasing on
$[s_0,\infty)$, so it attains its maximum at $s_0$. Evaluating $F$ at
$s_0$, we obtain
\[
F(s_{0})=\left(1+\frac{\epsilon}{t}\right)^{-\frac{t}{\epsilon}}\left(1-\frac{t}{t+\epsilon}\right)\leq \frac{\epsilon}{t+\epsilon}\leq \frac{\epsilon}{t}.
\]
Hence, we apply  \eqref{basic-estimate-3-1} to get that
\begin{align*}
\left\|\big(e^{-t\mathcal{L}}-e^{-(\epsilon+t)\mathcal{L}}\big)(T)\right\|^{2}_{L_{2}}&=\sum_{\s\in \{0,1,2,3,\}^{n},\s\neq 0}(e^{-t|\s|}-e^{-(\epsilon+t)|\s|})^{2}\left|\widehat{T}(\s)\right|^{2}\\
&\leq \left(\frac{\epsilon}{t}\right)^{2}\|T-\mathrm{tr}(T)\|^{2}_{L_{2}}\leq \frac{\epsilon^{2}}{t^{2}}\|T\|^{2}_{L_{2}},
\end{align*}
which is the desired inequality.
\end{proof}


The next proposition is a probabilistic representation of positive operators, established in \cite[Lemma 4.5]{Xi2017}, which enables us to bypass pointwise estimate heavily relied upon in the approach of \cite{HMO2017}.

\begin{proposition}\label{tr-var-psi}
Suppose that $\mathcal{N}$ is a $\sigma$-weakly dense $C^*$-subalgebra of $(\mathcal{M},\tau)$. Let $x \in \mathcal{N}$ be a normal element and $P : \mathcal{N} \to \mathcal{N}$ be a positive map. Then, there exists a positive measure $\mu$ on $\sigma(x) \times \sigma(x)$ with $\|\mu\| \leq \|P\|$ such that for all $\varphi, \psi \in C(\sigma(x))$,
\[
\tau\bigl(\varphi(x)P(\psi(x))\bigr)=\int_{\sigma(x) \times \sigma(x)} \varphi(t)\psi(s)\mathrm{~d}\mu(t,s).
\]
\end{proposition}


For  $a\in \mathbb{R}$ and $q\in [1,\infty)$, we use the shorthand notation 
\begin{equation}\label{phiq}
\phi_{q}(a)\coloneqq |a|^{q-1}\mathrm{sign}(a)
\end{equation}
 in the sequel. The well-known Stroock-Varopulos inequality is recalled as follows; see, for instance, \cite[Lemma 2.4]{HMO2017}.

\begin{lemma}[Stroock--Varopoulos]\label{svi}
Let \(a, b \in \mathbb{R}\) and \(p > 1\). Then
\[ 
\bigl(\phi_{p-1}(a) - \phi_{p-1}(b)\bigr)(a - b) \geq \frac{4(p-1)}{p^2} \bigl(\phi_{\frac{p}{2}}(a) - \phi_{\frac{p}{2}}(b)\bigr)^2.
\]
\end{lemma}

Combining Proposition \ref{tr-var-psi} with the Stroock-Varopoulos inequality, we get the following inequality.
\begin{lemma}\label{p-var}
For every $1<p<\infty$ and $T\in L_{p}(\mathbb{M}_{2^{n}})$ with $T=T^{*}$, we have
\[
(p-2)^{2}\|T\|^{p}_{L_{p}}+4(p-1)\mathrm{tr}(\phi_{\frac{p}{2}}(T)\cdot e^{-t\mathcal L}(\phi_{\frac{p}{2}}(T)))\geq p^2\mathrm{tr}(\phi_{p-1}(T)\cdot e^{-t\mathcal L}(T))
\]
\end{lemma}
\begin{proof} 
By Lemma \ref{svi}, the following inequality holds for all $a$, $b\in \mathbb{R}$
\begin{equation*}
p^2\bigl(\phi_{p-1}(a) - \phi_{p-1}(b)\bigr)(a - b) \geq 4(p-1) \bigl(\phi_{\frac{p}{2}}(a) - \phi_{\frac{p}{2}}(b)\bigr)^{2},
\end{equation*}
which can be rearranged to the follows
\begin{equation}\label{ab1}
(p-2)^2(|a|^p+|b|^p)+8(p-1)\phi_{\frac p2}(a)\phi_{\frac p2}(b)\geq p^2(a\phi_{p-1}(b)+\phi_{p-1}(a)b).
\end{equation}
Applying Proposition \ref{tr-var-psi} to \eqref{ab1} and taking the integration with respect to the positive measure $\mu$ on $\sigma(T)\times \sigma(T)$ yields the desired inequality.
\end{proof}




\begin{lemma}\label{basic-estimate-1-1}
Let  $a$, $b\in\mathbb{R}$. For each $2<p<\infty$, we have
\begin{equation}\label{large-range}
|a|^{p}+(p-1)|b|^{p}-pa|b|^{p-1}\mathrm{sign}(b)\leq p(p-1)|a-b|^{2}(|a|^{p-2}+|b|^{p-2}),
\end{equation}
and for each $1<p\leq 2$, we have
\begin{equation}\label{small-range}
|a|^{p}+(p-1)|b|^{p}-pa|b|^{p-1}\mathrm{sign}(b)\leq 2p|a-b|^{p}.
\end{equation}
\end{lemma}
\begin{proof}
The proofs of both inequalities are similar to each other which relies on specific convexity of the function
\[
F(t)\coloneqq |b+t(a-b)|^{p},\quad\mbox{for }t\in [0,1].
\]
It is easy to see that $F$ is a differentiable convex function on $[0,1]$ such that $F(0)=|b|^{p}$, $F(1)=|a|^{p}$ and
\[
F^{\prime}(t)=p(a-b)|b+t(a-b)|^{p-1}\mathrm{sign}(b+t(a-b)),
\]
and
\[
F^{\prime\prime}(t)=p(p-1)(a-b)^2|b+t(a-b)|^{p-2},
\]
for all $t\in [0,1]$. Thanks to the Newton-Leibniz formula, we have
\begin{equation}\label{n-l-expansion-1}
\begin{split}
|a|^{p}-|b|^{p}&=F(1)-F(0)=\int_{0}^1F^{\prime}(t)dt\\
&=p(a-b)\int_{0}^{1} |b+t(a-b)|^{p-1}\mathrm{sign}(b+t(a-b))dt=\int_{0}^{1}\int_{0}^{t}F^{\prime\prime}(u)dudt+F^{\prime}(0)\\
&=p(p-1)|a-b|^{2}\int_{0}^{1}\int_{0}^{t}|b+u(a-b)|^{p-2}dudt+p(a-b)|b|^{p-1}\mathrm{sign}(b).
\end{split}
\end{equation}
Rearranging \eqref{n-l-expansion-1} yields the following
\begin{equation}\label{n-l-expansion-2}
|a|^{p}+(p-1)|b|^{p}-pa|b|^{p-1}\mathrm{sign}(b)=p(p-1)|a-b|^{2}\int_{0}^{1}\int_{0}^{t}|b+u(a-b)|^{p-2}dudt,
\end{equation}
where $t\in [0,1]$.

For the inequality \eqref{large-range}, 
we note that, for every $u\in[0,1]$,
\[	|b+u(a-b)|
	=|(1-u)b+ua|
	\leq (1-u)|b|+u|a|
\leq \max\{|a|,|b|\}.
\]
Since $p>2$, it follows that
\[
|b+u(a-b)|^{p-2}
\leq \max\{|a|^{p-2},|b|^{p-2}\}
\leq |a|^{p-2}+|b|^{p-2}.
\]
Therefore, by \eqref{n-l-expansion-2} we have
\[
|a|^{p}+(p-1)|b|^{p}-pa|b|^{p-1}\mathrm{sign}(b)\leq p(p-1)|a-b|^{2}(|a|^{p-2}+|b|^{p-2}),
\]
which completes the proof of \eqref{large-range}.

If $1<p\leq 2$, we note that the convexity of $F$ implies that $F^{\prime}$ is an increasing function. Hence,
\begin{equation}\label{n-l-expansion-3}
\begin{split}
&|a|^{p}+(p-1)|b|^{p}-pa|b|^{p-1}\mathrm{sign}(b)=|a|^p-|b|^p+p(b-a)|b|^{p-1}\mathrm{sign}(b)\\
=&F(1)-F(0)-F^{\prime}(0)=\int_0^1F^{\prime}(t)dt-F^{\prime}(0)\leq F^{\prime}(1)-F^{\prime}(0)\\
\leq &p|a-b|\left||a|^{p-1}\mathrm{sign}(a)-|b|^{p-1}\mathrm{sign}(b)\right|.
\end{split}
\end{equation}
Since $1<p\leq 2$, it follows that $\left||a|^{p-1}-|b|^{p-1}\right|\leq |a-b|^{p-1}$ for all $a$, $b\in \mathbb{R}$ and $\mathrm{sign}(a)\mathrm{sign}(b)\geq0$. Thus, if $\mathrm{sign}(a)\mathrm{sign}(b)\geq0$, \eqref{n-l-expansion-3} can be estimated by
\begin{equation}\label{small-range-1}
|a|^{p}+(p-1)|b|^{p}-pa|b|^{p-1}\mathrm{sign}(b)\leq p|a-b|^{p}.
\end{equation}
If $\mathrm{sign}(a)\mathrm{sign}(b)<0$, we have
\[
\left||a|^{p-1}+|b|^{p-1}\right|\leq 2||a|+|b||^{p-1}=2|a-b|^{p-1}.
\]
Thus, \eqref{n-l-expansion-3} can be estimated by
\begin{equation}\label{small-range-2}
|a|^{p}+(p-1)|b|^{p}-pa|b|^{p-1}\mathrm{sign}(b)\leq 2p|a-b|^{p}.
\end{equation}
Combining \eqref{small-range-1} and \eqref{small-range-2} yields the desired inequality \eqref{small-range}.
\end{proof}

\begin{definition}\label{alpha-p-me}
For $M$, $\epsilon>0$, define
\[
\mathcal{C}(M,\epsilon)\coloneqq \{T: T=e^{-\epsilon\mathcal{L}}(S), S=S^*,\|
S\|_{L_{\infty}}<M\}.
\]
For every $1<p<\infty$, we define the parameter
\[
\alpha_{p, M, \epsilon}(t)\coloneqq\sup_{T \in \mathcal{C}(M, \epsilon)} (p - 1)\|T\|^{p}_{L_{p}} +\|e^{-t\mathcal{L}}(T)\|^{p}_{L_{p}}-p\mathrm{tr}(\phi_{p-1}(T)\cdot e^{-t\mathcal L}(T)).
\]
\end{definition}

\begin{lemma}\label{alpha-0}
For any \(1<p <\infty\) and \(M, \epsilon > 0\), we have
\[
\alpha_{p, M, \epsilon}(t)/t \to 0,\quad \mbox{as}\quad t \to 0^{+}.
\]
\end{lemma}
\begin{proof}
Combining Lemma \ref{basic-estimate-1-1} and Proposition \ref{tr-var-psi}, we get that
\[
\alpha_{p,M,\epsilon}(t)\leq
p(p-1)\sup_{T\in \mathcal{C}(M,\epsilon)}\left(\|T\|^{p-2}_{L_{\infty}}+\|e^{-t\mathcal{L}}(T)\|^{p-2}_{L_{\infty}}\right)\int_{\sigma(T)\times \sigma(T)}|a-b|^{2}d\mu(a,b),
\]
if $2<p<\infty$, and
\[
\alpha_{p,M,\epsilon}(t)\leq 2p\sup_{T\in \mathcal{C}(M,\epsilon)}\int_{\sigma(T)\times \sigma(T)}|a-b|^{p}d\mu(a,b),
\]
if $1<p\leq 2$. It is clear that
\[
\int_{\sigma(T)\times \sigma(T)}|a-b|^{2}d\mu(a,b)=\|e^{-t\mathcal{L}}(T)\|^{2}_{L_{2}}-2\mathrm{tr}\left(e^{-t\mathcal{L}}(T)T\right)+\|T\|^{2}_{L_{2}}=\|e^{-t\mathcal{L}}(T)-T\|^{2}_{L_{2}}.
\]
Hence, for the case $2<p<\infty$, we have
\begin{align*}
\alpha_{p,M,\epsilon}(t)&\leq p(p-1)\sup_{T\in \mathcal{C}(M,\epsilon)}(\|T\|^{p-2}_{L_{\infty}}+\|e^{-t\mathcal{L}}(T)\|^{p-2}_{L_{\infty}})\|e^{-t\mathcal{L}}(T)-T\|^{2}_{L_{2}}\\
&\leq 2p(p-1)M^{p-2}\sup_{T\in \mathcal{C}(M,\epsilon)}\|e^{-t\mathcal{L}}(T)-T\|^{2}_{L_{2}},
\end{align*}
where we used $T\in \mathcal{C(M,\epsilon)}$ and the fact $e^{-t\mathcal{L}}$ is a contraction in the last inequality.  Applying Lemma \ref{ete-22} to the above inequality yields that
\begin{align*}
	\alpha_{p,M,\varepsilon}(t)&\leq 2p(p-1)M^{p-2}\sup_{\|\varphi\|_{L_{\infty}}\leq M}\|e^{-(t+\epsilon)\mathcal{L}}(\varphi)-e^{-\epsilon\mathcal{L}}(\varphi)\|^{2}_{L_{2}}\\
	&\leq 2p(p-1)M^{p}\left(\frac{t^{2}}{\epsilon^{2}}\right),
\end{align*}
which further entails that $\alpha_{p,M,\epsilon}(t)/t\to 0^{+}$ as $t\to 0^{+}$.

For the case $1<p\leq 2$,  Jensen's inequality gives
\begin{align*}
\int_{\sigma(T)\times \sigma(T)}|a-b|^{p}d\mu(a,b)&=\int_{\sigma(T)\times \sigma(T)}\big(|a-b|^{2}\big)^{p/2}d\mu(a,b)\\
&\leq\left(\int |a-b|^2\,\mathrm d\mu(a,b)
\right)^{p/2}\\
&= \|e^{-t\mathcal{L}}(T)-T\|_{L_{2}}^p.
\end{align*}
%Therefore, we have
%\begin{equation}\label{estimate-for-alpha}
%\alpha_{p,M,\varepsilon}(t)\leq2p(p+1)M^{p-2}\sup_{T\in \mathcal{C}(M,\varepsilon)}\|e^{-t\mathcal{L}}(T)-T\|^{2}_{L_{2}}.
%\end{equation}
It follows that
\[
\alpha_{p,M,\epsilon}(t)
\leq
2pM^p\left(\frac{t}{\epsilon}\right)^p,
\]
and hence
\[
0\leq
\frac{\alpha_{p,M,\epsilon}(t)}{t}
\leq
\frac{2pM^p}{\epsilon^p}t^{p-1}
\longrightarrow0
\qquad\text{as }t\to0^+.
\]
The proof is complete.
\end{proof}

\begin{lemma}[{\cite[Lemma 2.2 and Remark 3.2]{HMO2017}}]\label{ab-k-p} 
Let $p\in(1,\infty)\setminus\{2\}$. For any $a~,b>0$, we have
\[
\bigl(1 + \kappa(p)^2\bigr)(a - b)^2 \leq a^2 + b^2 + a^{\frac{2}{2-p}}b^{\frac{2p-2}{p-2}} + b^{\frac{2}{2-p}}a^{\frac{2p-2}{p-2}}.
\]
\end{lemma}

\begin{lemma}\label{t-p}
For each $p\in(1,\infty)\setminus\{2\}$ and a non-zero self-adjoint $ T\in L_p(\mathbb{M}_{2^n})$ with $\mathrm{tr}(T)=0$, we have
\begin{align*}
(\mathrm{tr}(T_+^{p/2}))^2 + (\mathrm{tr}(T_-^{p/2}))^2&+ \big(\mathrm{tr}(T_+^{p/2})\big)^{-\frac{2}{p-2}} \bigl(\mathrm{tr}(T_-^{p/2})\bigr)^{\frac{2p-2}{p-2}}\\
&+ \bigl(\mathrm{tr}(T_-^{p/2})\bigr)^{-\frac{2}{p-2}} \bigl(\mathrm{tr}(T_+^{p/2})\bigr)^{\frac{2p-2}{p-2}} \leq \mathrm{tr}(|T|^p).
\end{align*}
\end{lemma}
\begin{proof}
%Put
%\[
%A=\operatorname{tr}(T_+^{p/2}),\qquad
%B=\operatorname{tr}(T_-^{p/2}),\qquad
%R=\operatorname{tr}(|T|^p),\qquad
%m=\frac{\operatorname{tr}(|T|)}2,
%\]
%and set 
%\[
%\alpha=\frac{2p-2}{p-2},\qquad
%\beta=-\frac2{p-2}.
%\]
%Then $\alpha+\beta=2$, so the desired left-hand side factors as
%\[
%(A^\alpha+B^\alpha)(A^\beta+B^\beta).
%\]
%It is therefore enough to prove the two auxiliary bounds
%\begin{equation}\label{t-1}
%A^\alpha+B^\alpha\le m^{p/(p-2)}R 
%\end{equation}
%and
%\begin{equation}\label{t-2}
%A^\beta+B^\beta\le m^{-p/(p-2)} .
%\end{equation}
%
%By the spectral calculus,
%\[
%T=T_+-T_-,\qquad |T|=T_++T_-,\qquad T_+T_-=0.
%\]
%Since $\operatorname{tr}(T)=0$,
%\[
%\operatorname{tr}(T_+)=\operatorname{tr}(T_-)=m,
%\qquad
%\operatorname{tr}(|T|)=2m.
%\]
%Let $e_+=1_{(0,\infty)}(T)$ and
%$e_-=1_{(-\infty,0)}(T)$ be the support projections of $T_+$ and
%$T_-$.  With normalized trace,
%\[
%r_+:=\operatorname{tr}(e_+),\qquad r_-:=\operatorname{tr}(e_-),
%\qquad r_++r_-\le1.
%\]
%
%\emph{Case $p>2$.}
%Use H\"older's inequality with conjugate exponents
%\[
%\alpha=\frac{2p-2}{p-2},\qquad \alpha'=\frac{2p-2}{p}.
%\]
%For the positive part, we have
%\begin{align*}
%A
%&=\operatorname{tr}\left(
%T_+^{\frac{p(p-2)}{2p-2}}
%T_+^{\frac{p}{2p-2}}\right)\leq
%\operatorname{tr}(T_+^p)^{\frac{p-2}{2p-2}}
%\operatorname{tr}(T_+)^{\frac{p}{2p-2}},
%\end{align*}
%which implies
%\[
%A^\alpha\leq
%\operatorname{tr}(T_+^p)m^{p/(p-2)}.
%\]
%The same argument for $T_-$ gives
%\[
%B^\alpha\leq
%\operatorname{tr}(T_-^p)m^{p/(p-2)}.
%\]
%Adding and using
%$\operatorname{tr}(T_+^p)+\operatorname{tr}(T_-^p)=R$ proves \eqref{t-1}.
%For \eqref{t-2}, H\"older's inequality with exponents $p/2$ and $p/(p-2)$ gives
%\[
%m=\operatorname{tr}(T_+)=\operatorname{tr}(T_+e_+)
%\leq A^{2/p}r_+^{(p-2)/p}.
%\]
%After raising to $p/(p-2)$ and rearranging,
%\[
%A^\beta\leq m^{-p/(p-2)}r_+.
%\]
%Similarly, we have
%$B^\beta\leq m^{-p/(p-2)}r_-$.  Adding these estimates and using
%$r_++r_-\leq1$ proves \eqref{t-2}.
%
%\emph{Case $1<p<2$.}
%Now $\alpha<0$ and $\beta>0$.  First use H\"older's inequality with
%exponents $2/p$ and $2/(2-p)$:
%\[
%A=\operatorname{tr}(T_+^{p/2}e_+)
%\leq m^{p/2}r_+^{(2-p)/2}.
%\]
%Raising to $\beta=2/(2-p)$ yields
%\[
%A^\beta\leq m^{p/(2-p)}r_+.
%\]
%The analogous estimate for the negative part and $r_++r_-\leq1$ give
%\[
%A^\beta+B^\beta\leq m^{p/(2-p)}=m^{-p/(p-2)},
%\]
%which is \eqref{t-2}.
%For \eqref{t-1}, use the conjugate exponents
%\[
%q=\frac{p}{2p-2},\qquad q'=\frac{p}{2-p}.
%\]
%Since $T_+=T_+^{p-1}T_+^{2-p}$ on the support of $T_+$,
%\begin{align*}
%m
%&=\operatorname{tr}(T_+^{p-1}T_+^{2-p})\\
%&\leq A^{(2p-2)/p}\operatorname{tr}(T_+^p)^{(2-p)/p}.
%\end{align*}
%Raising to $p/(2-p)$ and rearranging gives
%\[
%A^{-2(p-1)/(2-p)}
%\leq m^{-p/(2-p)}\operatorname{tr}(T_+^p).
%\]
%Because $\alpha=-2(p-1)/(2-p)$, the corresponding negative-part estimate
%and addition yield
%\[
%A^\alpha+B^\alpha
%\leq m^{-p/(2-p)}R=m^{p/(p-2)}R,
%\]
%which is (I).


For each $T\in L_p(\mathbb M_{2^n})$,  we write $T=T_+-T_-$. Then $|T|=T_++T_{-}$ and $|T|^p=T_+^p+T_{-}^p$. By the fact $\mathrm{tr}(T)=0$, we have
\begin{equation}\label{t-tr}
\mathrm{tr}(T_+)=\mathrm{tr}(T_{-})=\frac{\mathrm{tr}(|T|)}{2}.
\end{equation}
If we prove the following two inequalities,
\begin{equation}\label{t-1}
\mathrm{tr}\big(T_+^{\frac{p}{2}}\big)^{\frac{2p-2}{p-2}} + \mathrm{tr}\big(T_-^{\frac{p}{2}}\big)^{\frac{2p-2}{p-2}} \leq 2^{-\frac{p}{p-2}}\mathrm{tr}(|T|)^{\frac{p}{p-2}}\mathrm{tr}(|T|^p)
\end{equation}
and
\begin{equation}\label{t-2}
\mathrm{tr}\big(T_+^{\frac{p}{2}}\big)^{-\frac{2}{p-2}} + \mathrm{tr}\big(T_-^{\frac{p}{2}}\big)^{-\frac{2}{p-2}}\leq 2^{\frac{p}{p-2}} \mathrm{tr}(|T|)^{-\frac{p}{p-2}},
\end{equation}
multiplying them together yields the desired estimate.
	
For $p>2$, by the H\"older inequality with indexes $\left(\frac{2p-2}{p-2},\frac{2p-2}{p}\right)$ and \eqref{t-tr}, we have
\begin{equation}\label{positive part}
\mathrm{tr}\left(T_+^{\frac{p}{2}}\right)=\mathrm{tr}\left(T_+^{\frac{p^2-2p}{2p-2}} T_+^{\frac{p}{2p-2}}\right) \leq \mathrm{tr}\left(T_+^p\right)^{\frac{p-2}{2p-2}} \mathrm{tr}\left(\frac{|T|}{2}\right)^{\frac{p}{2p-2}},
\end{equation}
and
\begin{equation}\label{negative part}
\mathrm{tr}\left(T_-^{\frac{p}{2}}\right)
\leq\mathrm{tr}(T_-^p)^{\frac{p-2}{2p-2}}\mathrm{tr}\big(|T|/2\big)^{\frac{p}{2p-2}}.
\end{equation}
Taking the power of $(2p-2)/(p-2)$ to both \eqref{positive part} and \eqref{negative part}, and summing the resulting inequalities together, we obtain \eqref{t-1}.  In addition, by the H\"older inequality with $(\frac{p}{2},\frac{p}{p-2})$ and \eqref{t-tr},
\begin{equation}\label{basic inequality 1}
\mathrm{tr}(|T|/2) = \mathrm{tr}(T_+)
\leq \mathrm{tr}\big(T_+^{\frac{p}{2}}\big)^{\frac{2}{p}} \tr(\chi_{(0,\infty)}(T))^{\frac{p-2}{p}}.
\end{equation}
Then, taking the power $p/(p-2)$ to the both sides of \eqref{basic inequality 1} yields the inequality
\[
\mathrm{tr}\left(T_+^{\frac{p}{2}}\right)^{-\frac{2}{p-2}}
\leq 2^{\frac{p}{p-2}} \mathrm{tr}(|T|)^{-\frac{p}{p-2}}  \cdot \tr(\chi_{(0,\infty)}(T)).
\]
Similarly, we have
\[
\mathrm{tr}\left(T_{-}^{\frac{p}{2}}\right)^{-\frac{2}{p-2}}
\leq 2^{\frac{p}{p-2}}\mathrm{tr}(|T|)^{-\frac{p}{p-2}}  \cdot \tr(\chi_{(-\infty,0)}(T)).
\]
Summing the two inequalities above gives \eqref{t-2}.	 
	
For $1<p<2$, by the H\"older inequality with $(\frac{2}{p},\frac{2}{2-p})$ and \eqref{t-tr},  
\begin{equation*}
\mathrm{tr}\big(T_+^{\frac{p}{2}}\big)\leq \mathrm{tr}(T_+)^{\frac{p}{2}} \mathbb \tau(\chi_{(0,\infty)}(T))^{\frac{2-p}{2}}=2^{-\frac {p}2}\mathrm{tr}(|T|)^{\frac{p}{2}}\tr\left(\chi_{(0,\infty)}(T)\right)^{\frac{2-p}{2}},
\end{equation*}
and 
\[
\mathrm{tr}\big(T_{-}^{\frac{p}{2}}\big)\leq 2^{-\frac p2}\mathrm{tr}(|T|)^{\frac{p}{2}}\tr(\chi_{(-\infty,0)}(T))^{\frac{2-p}{2}}.
\]
Exponentiating both sides of the two inequalities above to the power $2/(2-p)$ and summing the resulting relations gives inequality \eqref{t-2}. The H\"older inequality with indexes $\left(\frac{p}{2p-2},\frac{p}{2-p}\right)$ and \eqref{t-tr} give
\[
\mathrm{tr}(|T|/2) = \mathrm{tr}(T_+)=\mathrm{tr}\left(T_+^{p-1}T_+^{2-p}\right)\leq \mathrm{tr}\big(T_+^{\frac{p}{2}}\big)^{\frac{2p-2}{p}} \mathrm{tr}\big(T_+^{p}\big)^{\frac{2-p}{p}}.
\]
Then, taking both sides to the \(p/(2-p)\)-th power, we obtain
\[
\mathrm{tr}\big(T_+^{\frac{p}{2}}\big)^{-\frac{2p-2}{2-p}}\leq 2^{\frac{p}{2-p}} \mathrm{tr}(|T|)^{-\frac{p}{2-p}} \mathrm{tr}\big(T_+^{p}\big).
\]
Similarly,
\[
\mathrm{tr}\big(T_{-}^{\frac{p}{2}}\big)^{-\frac{2p-2}{2-p}}\leq 2^{\frac{p}{2-p}} \mathrm{tr}(|T|)^{-\frac{p}{2-p}} \mathrm{tr}\big(T_{-}^{p}\big).
\]
Adding the two inequalities together yields \eqref{t-1}, which completes the proof.
\end{proof}


\begin{lemma}\label{t-kp-trt}
For each $p\in(1,\infty)\setminus\{2\}$ and  self-adjoint $T\in L_p(\mathbb{M}_{2^n})$ with $\mathrm{tr}(T)=0$, 
\[
\bigl(1 + \kappa(p)^2\bigr)\big(\mathrm{tr}(T_+^{\frac{p}{2}})-\mathrm{tr}(T_-^{\frac{p}{2}})\big)^2\leq \mathrm{tr}(|T|^p).
\]
\end{lemma}
\begin{proof}
Lemma \ref{ab-k-p} applies only to two positive real numbers $a = \mathrm{tr}\big(T_+^{\frac{p}{2}}\big)$ and $b = \mathrm{tr}\big(T_-^{\frac{p}{2}}\big)$. Combining this with Lemma \ref{t-p}, we complete the proof.
\end{proof}

Before proving Theorem \ref{heat-semigroup-f2828b}, we note that for self-adjoint $T$,  $\phi_{q}(T)= u_{T}|T|^{q-1}$, where $\phi_q$ is as in \eqref{phiq}, $q\geq 1$,  and $u_{T}$ is a partial isometry coming from the polar decomposition $T=u_{T}|T|$.
\begin{proof}[Proof of Theorem \ref{heat-semigroup-f2828b}]
Denote  for simplicity that
\[
\mathcal{C}_0(M,\epsilon)=\{T\in \mathcal{C}(M,\epsilon):\mathrm{tr}(T)=0\},
\]
 and
\[
L_{p,0}(\mathbb{M}_{2^{n}})=\{T\in L_{p}(\mathbb{M}_{2^{n}}):\mathrm{tr}(T)=0\}.
\]
It is clear that
\[
\bigcup_{M,\varepsilon>0}\mathcal{C}_{0}(M,\varepsilon)=L_{p,0}(\mathbb{M}_{2^{n}}).
\]
Indeed, for each $T\in L_{p,0}(\mathbb{M}_{2^{n}})$, we have $T=e^{-\epsilon\mathcal{L}}\left(e^{\epsilon\mathcal{L}}(T)\right)$ with
\[
\begin{split}
\left\|e^{\epsilon\mathcal{L}}(T)\right\|_{L_{\infty}}&\leq 2^{n/2}\|e^{\epsilon\mathcal{L}}(T)\|_{L_{2}}
\leq 2^{n/2}e^{\epsilon n}\|T\|_{L_{2}}\leq 2^{n}e^{\varepsilon n}\|T\|_{L_{\infty}}.
\end{split}
\]

For each $T\in \mathcal{C}_0(M,\epsilon)$ and $M,\epsilon>0$, it follows from Definition \ref{alpha-p-me} and Lemma \ref{p-var} that
\begin{equation*}\label{tr-etp}
\begin{aligned}
\|e^{-t\mathcal{L}}(T)\|^{p}_{L_{p}}&\leq \alpha_{p,M,\epsilon}(t)+p\mathrm{tr}(\phi_{p-1}(T)\cdot e^{-t\mathcal L}(T))-(p-1)\|T\|^{p}_{L_{p}}\\
&\leq\alpha_{p,M,\epsilon}(t)+\frac{1}{p}\left[(p-2)^2\|T\|^{p}_{L_{p}}+4(p-1)\mathrm{tr}(\phi_{\frac{p}{2}}(T)\cdot e^{-t\mathcal L}(\phi_{\frac{p}{2}}(T)))\right]\\
&\quad-(p-1)\|T\|^{p}_{L_{p}}.
\end{aligned}
\end{equation*}
By the exponential ergodicity of $\{e^{-t\mathcal{L}}\}_{t\geq 0}$, we have
\begin{align*}
\mathrm{tr}(\phi_{\frac{p}{2}}(T)e^{-t\mathcal L}(\phi_{\frac{p}{2}}(T)))-(\mathrm{tr}(\phi_{\frac p2}(T)))^2&=\mathrm{Var}\left(e^{-t\mathcal{L}/2}(\phi_{\frac{p}{2}}(T))\right)\\
&\leq e^{-t}\mathrm{Var}(\phi_{\frac{p}{2}}(T))=e^{-t}\left(\|T\|^{p}_{L_{p}}-\mathrm{tr}(\phi_{\frac{p}{2}}(T))^{2}\right).
\end{align*}
By Lemma \ref{t-kp-trt}, it follows that 
\begin{align*}
\mathrm{tr}(\phi_{\frac{p}{2}}(T))^{2}=\big(\mathrm{tr}(T_+^{\frac{p}{2}})-\mathrm{tr}(T_-^{\frac{p}{2}})\big)^2\leq \frac{1}{\bigl(1 + \kappa(p)^2\bigr)}\mathrm{tr}(|T|^p)
\end{align*}
Combining the above estimates, we have
\begin{equation*}\label{k-1}
\begin{aligned}
\|e^{-t\mathcal{L}}(T)\|^{p}_{L_{p}}&\leq \alpha_{p,M,\epsilon}(t)+\frac1p\big[\left((p-2)^2+4(p-1)e^{-t}\right)\|T\|^{p}_{L_{p}}\\
	&\quad+4(p-1)(1-e^{-t})(\mathrm{tr}(\phi_{\frac p2}(T)))^2\big]-(p-1)\|T\|^{p}_{L_{p}}\\
&\leq \alpha_{p,M,\epsilon}(t)+\Big[\frac{(p-2)^2+4(p-1)e^{-t}}{p}+\frac{4(p-1)(1-e^{-t})}{(1+\kappa(p)^2)p}-(p-1)\Big]\|T\|^{p}_{L_{p}}\\
	&=\alpha_{p,M,\epsilon}(t)+\Big[1-\frac{4(p-1)(1-e^{-t})}{p}+\frac{4(p-1)(1-e^{-t})}{(1+\kappa(p)^2)p}\Big]\|T\|^{p}_{L_{p}}\\
	&=\alpha_{p,M,\epsilon}(t)+\Big[1-\frac{4(p-1)(1-e^{-t})}{(1+\kappa(p)^{-2})p}\Big]\|T\|^{p}_{L_{p}}.
\end{aligned}
\end{equation*}
For arbitrary $m\in \mathbb{N}$ and $k\in\mathbb{N}$, standard induction argument gives
\begin{equation}\label{kn}
\left\|e^{-kt\mathcal{L}/m}(T)\right\|^{p}_{L_{p}}\leq k \alpha_{p,M,\epsilon}\left(\frac{t}{m}\right)+\left(1-\frac{4(p-1)(1-e^{-\frac t m})}{(1+\kappa(p)^{-2})p}\right)^k\|T\|^{p}_{L_{p}}.
\end{equation}
%Indeed, for $k=1$, the desired estimate follows from \eqref{k-1}. We now assume from now on that inequality \eqref{kn} holds for $m$ with $m\leq k-1$. By \eqref{kn} again, we have
%\begin{align*}
%&\left\|e^{-kt\mathcal{L}/n}(T)\right\|^{p}_{L_{p}}\\
%=&\left\|e^{-t\mathcal{L}/n}\circ e^{-(k-1)t\mathcal{L}/n}(T)\right\|^{p}_{L_{p}}\\
%\le&\alpha_{p,M,\epsilon}\left(\frac{t}{n}\right)+\left(1-\frac{4(p-1)(1-e^{-\frac t n})}{(1+\kappa(p)^{-2})p}\right)\left\|e^{-\frac{(k-1)t}{n}\mathcal{L}}(T)\right\|^{p}_{L_{p}}\\
%\le&\left(1-\frac{4(p-1)(1-e^{-\frac t n})}{(1+\kappa(p)^{-2})p}\right)\left[\left(1-\frac{4(p-1)(1-e^{-\frac t n})}{(1+\kappa(p)^{-2})p}\right)^{k-1}+(k-1)\alpha_{p,M,\epsilon}\left(\frac{t}{n}\right)\right]\|T\|^{p}_{L_{p}}\\
%&\quad+\alpha_{p,M,\epsilon}\left(\frac{t}{n}\right)\\
%=&\left(1-\frac{4(p-1)(1-e^{-\frac t n})}{(1+\kappa(p)^{-2})p}\right)^{k}\|T\|^{p}_{L_{p}}+\Big[\Big(1-\frac{4(p-1)(1-e^{-\frac t n})}{(1+\kappa(p)^{-2})p}\Big)(k-1)+1\Big]\alpha_{p,M,\epsilon}\left(\frac{t}{n}\right)\\
%\leq&\left(1-\frac{4(p-1)(1-e^{-\frac t n})}{(1+\kappa(p)^{-2})p}\right)^{k}\|T\|^{p}_{L_{p}}+k\alpha_{p,M,\epsilon}\left(\frac{t}{n}\right).
%\end{align*}
It follows from Lemma \ref{alpha-0} that
\begin{equation*}\label{limit-1}
m\alpha_{p,M,\varepsilon}(t/m)\to 0,\quad
\mbox{as }m\to\infty.
\end{equation*}
Also note that
\[
\lim_{m\to\infty}\left(1-\frac{4(p-1)(1-e^{-\frac t m})}{(1+\kappa(p)^{-2})p}\right)^{m}=\exp\left(-\frac{(4p-4)t}{p(1+\kappa(p)^{-2})}\right)
\]
for some constant $c>0$. The desired inequality follows by taking $k=m$ and $m\to\infty$ in \eqref{kn}.
\end{proof} 

\begin{remark}\label{rem-hmo}

In the classical treatments of establishing the Heilman--Mossel--Oleszkiewicz inequality \cite{HMO2017}, the proof of Lemmas \ref{p-var} and \ref{alpha-0} relies on pointwise estimates for the heat semigroup. However, such pointwise arguments are not available in the quantum Boolean cube setting. To overcome this difficulty, we applied the probabilistic representation of  the noncommutative semigroups (see Proposition \ref{tr-var-psi})  to prove the desired results.
\end{remark}

We now provide a simpler, although weaker, contraction estimate.  It will be useful when only an explicit dimension-free contraction factor is needed. The proof presented below is motivated by \cite[Proposition 4.1]{HMO2017}. 
\begin{theorem}\label{heat-inter}
For each $1<p<\infty$, let $T\in L_{p}(\mathbb{M}_{2^n})$ with $\mathrm{tr}(T)=0$. Then
$$\|e^{-t\mathcal{L}}(T)\|_{L_p}\leq \big(1-2^{2-p^*}(1-e^{-2t})\big)^{\frac{1}{p*}}\|T\|_{L_p},$$
where $p^*=\max\{p,p'\}$.
\end{theorem} 

\begin{proof}
Let $c>0$ be the constant that will be specified later. We define a positive measure on the two-point set $\{0,1\}$ as follows
\[
\mu_{c}(\{0\})=1\quad\mbox{and}\quad \mu_{c}(\{1\})=c^{2}(1-e^{-2t}).
\]
Let $\mathcal{N}\coloneqq L_{\infty}(\{0,1\},\mu_{c})\bar\otimes \mathbb{M}_{2^{n}}$ be the tensor product of two von Neumann algebras, equipped with the trace $\tau\coloneqq \mu_{c}\otimes \mathrm{tr}$. Let $F:L_{1}(\mathbb{M}_{2^{n}})\to L_{1}(\mathcal{N})$ be the linear operator defined by
\[
F(T)(0)\coloneqq c e^{-t\mathcal{L}}(T)\quad\mbox{and}\quad F(T)(1)\coloneqq T-\mathrm{tr}(T),
\]
for each $T\in L_{1}(\mathbb{M}_{2^{n}})$. We now verify the boundedness of the operator $F$ as follows. By the contraction of the semigroup, it follows that
\begin{align*}
\|F(T)\|_{L_1(\mathcal{N})}&=\|F(T)(0)\|_{L_1(\mathbb{M}_{2^n})}+c^2(1-e^{-2t})\|F(T)(1)\|_{L_1(\mathbb{M}_{2^n})}\\
&=c\|e^{-t\mathcal{L}}(T)\|_{L_1(\mathbb{M}_{2^n})}+c^2(1-e^{-2t})\|T-\mathrm{tr}(T)\|_{L_1(\mathbb{M}_{2^n})}\\
&\leq c\|T\|_{L_1(\mathbb{M}_{2^n})}+2c^2(1-e^{-2t})\cdot\|T\|_{L_1(\mathbb{M}_{2^n})}.
\end{align*}
Hence, we have
\begin{equation}\label{f-11}
\|F\|_{L_1(\mathbb{M}_{2^n})\to L_1(\mathcal{N})}\leq c+2c^2(1-e^{-2t}).
\end{equation}


Given $T\in L_2(\M_{2^n})$, set $T_0\coloneqq T-\mathrm{tr}(T)$. Then
\begin{equation*}\label{f-22-1}
\begin{split}
\|F(T)\|_{L_2(\mathcal{N})}^{2}&=c^{2}\|e^{-t\mathcal{L}}(T)\|_{L_2(\mathbb{M}_{2^n})}^2+c^2(1-e^{-2t})\cdot\|T_0\|_{L_2(\mathbb{M}_{2^n})}^{2}\\
&=c^{2}\|e^{-t\mathcal{L}}(T_0)\|^{2}_{L_{2}}+c^{2}|\mathrm{tr}(T)|^{2}+c^{2}(1-e^{-2t})\|T_0\|^{2}_{L_{2}}\\
&\leq c^{2}e^{-2t}\|T_0\|^{2}_{L_{2}}+c^{2}|\mathrm{tr}(T)|^{2}+c^{2}(1-e^{-2t})\|T_0\|^{2}_{L_{2}}\\
&=c^{2}\left(|\mathrm{tr}(T)|^{2}+\|T_0\|^{2}_{L_{2}}\right)=c^{2}\|T\|^{2}_{L_{2}},
\end{split}
\end{equation*}
where we used the exponential ergodic of the semigroup in the inequality. Thus, 
\begin{equation}\label{f-22}
	\|F\|_{L_2(\mathbb{M}_{2^n})\to L_2(\mathcal{N})}\leq c.
\end{equation}
Moreover, it is clear that
\begin{align*}
\|F(T)\|_{L_\infty(\mathcal{N})}&=\max\{c\|e^{-t\mathcal{L}}(T)\|_{L_\infty(\mathbb{M}_{2^n})},\|T-\mathrm{tr}(T)\|_{L_\infty(\mathbb{M}_{2^n})}\}\\
	&\leq \max\{c,2\}\|T\|_{L_\infty(\mathbb{M}_{2^n})}.
\end{align*}
Hence, 
\begin{equation}\label{f-infty}
\|F\|_{L_\infty(\mathbb{M}_{2^n})\to L_\infty(\mathcal{N})}\leq \max\{c,2\}.
\end{equation}

We now apply the Riesz-Thorin interpolation theorem to conclude the proof. For the case $1<p\leq2$, applying the Riesz-Thorin interpolation theorem to \eqref{f-11} and \eqref{f-22} yields
\[
\|F\|_{L_p(\mathbb{M}_{2^n})\to L_p(\mathcal{N})}\leq \big(c+2c^2(1-e^{-2t})\big)^{\frac{2-p}{p}}\cdot c^{\frac{2p-2}{p}}=c\big(1+2c(1-e^{-2t})\big)^{\frac{2-p}{p}}.
\]
Putting $c=\left(2^{\frac{1}{p-1}}-2(1-e^{-2t})\right)^{-1}$, it follows that
\begin{equation*}\label{f-p12}
\begin{aligned}
\|e^{-t\mathcal{L}}(T)\|_{L_{p}}&\leq \Big(\big(1+2c(1-e^{-2t})\big)^{2-p}-c^{2-p}(1-e^{-2t})\Big)^{\frac{1}{p}}\|T\|_{L_{p}}\\
&=\Big(\frac{c^{1-p}}{2}\Big)^{\frac{1}{p}}\|T\|_{L_{p}}\\
&=\Big(\frac{(2^{\frac{1}{p-1}}-2(1-e^{-2t}))^{\frac{1}{p^{\prime}}}}{2^{\frac{1}{p}}}\Big)\|T\|_{L_{p}}\\
&=(1-2^{2-p^{\prime}}(1-e^{-2t}))^{\frac{1}{p^{\prime}}}\|T\|_{L_{p}}.
\end{aligned}
\end{equation*}
For the case $2\leq p<\infty$, we put $c=2$. Then, combining  \eqref{f-22}, \eqref{f-infty} with interpolation theorem, we have $\|F\|_{L_{p}\to L_p(\mathcal{N})}\leq 2$. Therefore, for each $T\in L_{p}(\mathbb{M}_{2^{n}})$ with $\mathrm{tr}(T)=0$, it follows from the definition of $F$ that
\[
\|e^{-t\mathcal{L}}(T)\|_{L_{p}}\leq (1-2^{2-p}(1-e^{-2t}))^{\frac{1}{p}}\|T\|_{L_{p}},
\]
which completes our proof.
\end{proof}
\section{A reverse Bernstein-Markov type inequality}
The previous results control the decay property of observables under the heat semigroup.  We now convert this heat control into reverse Bernstein--Markov inequalities for observables whose spectrum lies in the band $[d,d+m]$.  Such estimates are the high-frequency counterpart of the low-degree Bernstein--Markov inequalities proved in Chapter \ref{sec-bm}.

\begin{theorem}\label{d1-2-dtail-1} Let $1<p<\infty$, and let $n, d, m \in \mathbb{N}$  satisfy $d+m \leq n$.  For every $T\in\mathcal{P}^{\leq d+m}(\mathbb{M}_{2^n})\cap \mathcal{T}^{\geq d}(\mathbb{M}_{2^n})$, there exists $c_p \in(0, \infty)$ such that
	$$
	\|\mathcal{L}^{1 / 2} (T)\|_{L_{p}} \geq c_p \sqrt{\frac{d}{m}}\|T\|_{L_{p}}.$$
\end{theorem}


\begin{theorem}\label{rbm-d-tail}  Let $1<p<\infty$, and let $n, d, m \in \mathbb{N}$  satisfy $d+m \leq n$.  For every $T\in\mathcal{P}^{\leq d+m}(\mathbb{M}_{2^n})\cap \mathcal{T}^{\geq d}(\mathbb{M}_{2^n})$, there exists $c_p \in(0, \infty)$ such that
$$\|\mathcal{L}(T)\|_{L_{p}}\geq c_p \frac{d}{m}\|T\|_{L_{p}}.
$$
\end{theorem}

We now derive Theorem \ref{rbm-d-tail} from Theorem \ref{d1-2-dtail-1} as follows. 

\begin{proof}[Proof of Theorem \ref{rbm-d-tail}]  By Theorem \ref{d1-2-dtail-1} twice, we have
\begin{align*}
\|\mathcal{L}(T)\|_{L_{p}} & =\|\mathcal{L}^{1/2} \mathcal{L}^{1 / 2} (T)\|_{L_{p}} \\
& \geq c_p \sqrt{\frac{d}{m}}\|\mathcal{L}^{1 / 2} (T)\|_{L_{p}}  \geq c_p^{2} \frac{d}{m}\|T\|_{L_{p}}.
\end{align*}
\end{proof}


To prove Theorem \ref{d1-2-dtail-1}, we shall use the following standard fractional-power estimate for generators of contraction semigroups which can be found from \cite{EL2008}.

\begin{lemma}[{\cite[Lemma 5.6]{EL2008}}]\label{nsi}
Let $E$ be a Banach space, $\{T_{t}\}_{t\geq 0}$ be a strongly continuous semigroup of contractions of $E$, and let $A$ be its infinitesimal generator. Then for $0 <\beta< 1$ and every $f$ in the domain of $A$
\[
\|A^\beta (f)\|_E\leq 4\|f\|_E^{1-\beta}\|A(f)\|_E^\beta.
\]
\end{lemma}


\begin{proposition}\label{d1-2-dtail-2}
Let $1<p<\infty$. For  every  $T\in \mathbb{M}_{2^n}$, there exists $K_p \in(0, \infty)$ such that
\[
\|\mathcal{L}^{1 / 2} e^{-t \mathcal{L}} (T)\|_{L_{p}} \leq \frac{K_p}{\sqrt{e^{2 t}-1}}\|T\|_{L_{p}},\quad \forall t>0.
\]
\end{proposition}
\begin{proof}
For $1<p<\infty$, we take $\theta=\theta_p\in (0,\frac{\pi}{2}]$ as in Theorem \ref{e-zl-b}.  Set $a=\frac{\pi}{\tan(\theta)}>0$. We divide the proof into following two cases.

\textbf{Case I}: For $t \geq 2 a$, we apply Lemma \ref{radk-pp-ak} to obtain the following estimate
\begin{align*}
\|\mathcal{L}^{1/2} e^{-t \mathcal{L}} (T)\|_{L_{p}} & =\left\|\sum_{k=1}^{n} \sqrt{k} e^{-t k}\mathrm{Rad}^{=k} (T)\right\|_{L_{p}} \\
&\leq \sum_{k=1}^{n} \sqrt{k} e^{-t k}\|\mathrm{Rad}^{=k} (T)\|_{L_{p}} \\
& \leq M_p \sum_{k=1}^{n} \sqrt{k} e^{-(t-a) k}\|T\|_{L_{p}}  \\
& \leq\left(M_p e^{2a} \sum_{k=1}^{\infty} \sqrt{k} e^{-a k}\right) e^{-t}\|T\|_{L_{p}},
\end{align*}
where the last inequality follows from the assumption $t\geq 2a$. Indeed, for every $k\geq1$ we have
\[
e^{-(t-a)k}
=e^{-ak}e^{-(t-2a)k}
\leq e^{-ak}e^{-(t-2a)}
=e^{2a}e^{-t}e^{-ak}.
\]
Moreover, $
e^{-t}\leq\frac{1}{\sqrt{e^{2t}-1}}$ for $ t>0.$ Consequently,
$$
\|\mathcal{L}^{1/2} e^{-t \mathcal{L}} (T)\|_{L_{p}}\leq \frac{K_p}{\sqrt{e^{2 t}-1}}\|T\|_{L_{p}},
$$
where $K_p=M_p e^{2a} \sum_{k=1}^{\infty} \sqrt{k} e^{-a k}$. 

\textbf{Case II}: For $0<t <2 a$, let $r=t \sin \theta$. Notice that the closed disc $\overline{D}(t, r)$ of radius $r$ centered at $t$ is contained in $\{z \in \mathbb{C}:|\arg z| \leq \theta\}$. By the Cauchy integral formula for the derivative, we have
$$
-\mathcal{L} e^{-t \mathcal{L}}=\frac{1}{2 \pi i} \int_{\partial D(t, r)} \frac{e^{-\zeta \mathcal{L}}}{(\zeta-t)^{2}} \mathrm{d} \zeta,
$$
Hence,
\begin{align*}
\|\mathcal{L} e^{-t \mathcal{L}} (T)\|_{L_{p}}&=\Big\|\frac{1}{2 \pi i} \int_{\partial D(t, r)} \frac{e^{-\zeta \mathcal{L}}T}{(\zeta-t)^{2}} \mathrm{~d} \zeta\Big\|_{L_{p}}\\
&\leq\left|\frac{1}{2 \pi i} \int_{\partial D(t, r)} \frac{1}{(\zeta-t)^{2}} \mathrm{~d} \zeta\right|\sup _{\zeta \in \partial D(t, r)}\|e^{-\zeta \mathcal{L}}(T)\|_{L_{p}}\\
& \leq \frac{1}{r} M_p\|T\|_{L_{p}}\\ &= \frac{M_p}{t \sin \theta}\|T\|_{L_{p}},
\end{align*}
where the second inequality is due to Theorem \ref{e-zl-b}. Applying Lemma \ref{nsi} (with $f=e^{-t \mathcal{L}} (T)$ and $\beta=\frac{1}{2}$ there) and Lemma \ref{hc-1} (with $p=q$), we have
\begin{align*}
\|\mathcal{L}^{1/2} e^{-t \mathcal{L}} (T)\|_{L_{p}} &\leq 4\|\mathcal{L}e^{-t \mathcal{L}} (T)\|_{L_{p}}^{\frac{1}{2}}\|e^{-t \mathcal{L}} (T)\|_{L_{p}}^{\frac{1}{2}}\leq\frac{4 \sqrt{M_p}}{\sqrt{t\sin \theta}}\|T\|_{L_{p}}.
\end{align*}
Since for all $t\in (0,2a)$, we have $\sqrt{e^{2 t}-1} \leq c_{a} \sqrt{t}$, where   the constant $c_{a}>0$ depends only on $p$, because $a=\pi/\tan(\theta_p)$, it gives the desired estimate.
\end{proof}

We shall also use the following reverse Markov inequality in the proof of Theorem \ref{d1-2-dtail-1}.
\begin{lemma}[{\cite[Theorem 2.2]{E2020}}]\label{er}
Let $d, m \in \mathbb{N}$, and let $P(x)$ be a polynomial of the form
\[
P(x)=a_{d} x^{d}+a_{d+1} x^{d+1}+\cdots+a_{d+m} x^{d+m}.
\]
Then, the following estimate holds
\[
|P(1)| \leq 6 \sqrt{\frac{m}{d}}\left\|\sqrt{1-x^{2}} P^{\prime}(x)\right\|_{\mathcal{C}([0,1])}.
\]
The order of this estimate is sharp up to the universal constant.
\end{lemma}

\begin{proof}[Proof of Theorem \ref{d1-2-dtail-1}] 
Let $x=e^{-t}\in[0,1]$. Rewrite Proposition \ref{d1-2-dtail-2}  as follows 
\begin{equation}\label{1x}
\left\|\sqrt{1-x^{2}} \mathcal{L} x^{\mathcal{L}-1} (T)\right\|_{L_{p}} =	\left\|\sqrt{x^{-2}-1} \mathcal{L}^{1/2} x^{\mathcal{L}} (\mathcal{L}^{1/2} T)\right\|_{L_{p}} 	\leq K_p\|\mathcal{L}^{1 / 2} (T)\|_{L_{p}}. 
\end{equation}
Consider the linear functional
$$
\xi:\left(\operatorname{span}\left\{\sqrt{1-x^{2}} x^{d-1}, \ldots, \sqrt{1-x^{2}} x^{d+m-1}\right\},\|\cdot\|_{\mathcal{C}([0,1])}\right) \rightarrow \mathbb{C}
$$
given by
\[
\xi\left(\sqrt{1-x^{2}} \sum_{k=d}^{d+m} k a_{k} x^{k-1}\right)\coloneqq\sum_{k=d}^{d+m} a_{k} =P(1),
\]
where $P(x)=\sum_{k=d}^{d+m}a_kx^k$ is the polynomial defined as in Lemma \ref{er}. Thus, by lemma \ref{er}, it follows that for each
\[
Q(x)=\sqrt{1-x^{2}}P'(x)
=\sqrt{1-x^{2}}
\sum_{k=d}^{d+m}ka_kx^{k-1}\in \operatorname{span}\left\{\sqrt{1-x^{2}} x^{d-1}, \ldots, \sqrt{1-x^{2}} x^{d+m-1}\right\},
\]
we have 
$$
|\xi(Q)|=|P(1)|\leq 6\sqrt{\frac{m}{d}}\|\sqrt{1-x^{2}}P^{\prime}(x)\|_{\mathcal{C}([0,1])}=6 \sqrt{\frac{m}{d}}\|Q\|_{\mathcal{C}([0,1])}.
$$
Therefore, by the Hahn-Banach theorem and the Riesz representation theorem, there exists a complex measure $\sigma$ on $[0,1]$ such that $$\|\sigma\|_{\mathcal{M}([0,1])} \leq 6 \sqrt{\frac{m}{d}}$$ and
$$\int_{0}^{1} Q(x) \operatorname{d} \sigma(x)=\xi(Q),\quad\forall Q \in \operatorname{span}\left\{\sqrt{1-x^{2}} x^{d-1}, \ldots, \sqrt{1-x^{2}} x^{d+m-1}\right\}. $$
Since $\widehat{T}(\mathbf{s})=0$ when $|\mathbf{s}| \notin\{d, \ldots, d+m\}$, by the above equation, Jensen inequality and \eqref{1x}, we obtain that
\begin{align*}
 \|T\|_{L_{p}}&=\left\|\sum_{\substack{\mathbf{s}\in\{0,1,2,3\}^n\\d\leq|\mathbf{s}|\leq d+m}}\widehat{T}(\mathbf{s})\sigma_\mathbf{s}\right\|_{L_{p}}\\
%&=\|P(1)\|_{L_{p}(\mathbb{M}_{2^n})}\\
%&= \|\xi(p)\|_{L_{p}(\mathbb{M}_{2^n})}\\
&=\left\|\int_{0}^{1} \sqrt{1-x^{2}} \mathcal{L} x^{\mathcal{L}-1} (T) \mathrm{~d} \sigma(x)\right\|_{L_{p}}\\
& \leq \int_{0}^{1}\left\|\sqrt{1-x^{2}} \mathcal{L} x^{\mathcal{L}-1} (T)\right\|_{L_{p}}\mathrm{d}|\sigma|(x) \\
& \leq K_{p}\int_{0}^{1}\left\|\mathcal{L}^{1/2}(T)\right\|_{L_{p}} \mathrm{d}|\sigma|(x)  \\
&=K_{p}\|\sigma\|_{\mathcal{M}([0,1])}\left\|\mathcal{L}^{1/2}(T)\right\|_{L_{p}} \leq 6 K_p \sqrt{\frac{m}{d}}\left\|\mathcal{L}^{1 / 2}(T)\right\|_{L_{p}}.
\end{align*}
\end{proof}


\chapter[BM Inequalities for Low Degree Observables]{Bernstein-Markov type inequalities for low degree observables}\label{sec-bm}

In this chapter, the observables of interest have a degree of at most $d \in [n]$, which shifts our focus from the preceding chapter. We first investigate a lower bound on the decay of the quantum heat semigroup, which may be viewed as the ``dual'' formulation of Theorem \ref{mn-k-tail}. Furthermore, we establish Bernstein-Markov-type inequalities for low-degree quantum functions with respect to both the quantum Laplacian and its associated gradient. 

\section{Lower bound on the  decay of the quantum heat semigroup}
In the classical setting, the study of lower bounds on the decay of the heat semigroup
acting on functions of low degree goes back at least to \cite{FHKL2016}. Later, Eskenazis  and Ivanisvili \cite{EI2020} extended it to the vector-valued setting and found interesting applications to moment comparison for functions of low degree and Pisier's remarkable inequality \eqref{eq-pisier-ineq}. The main results here, namely Theorems \ref{LB-C} and \ref{MN-K-most-d}, serve as the quantum counterparts of theirs. The Chebyshev polynomial was defined above (see the explanation after Theorem \ref{mn-k-tail}).


\begin{theorem}\label{LB-C}
For every $ d \in [n]$, $1\leq p \leq\infty$,  and   $T\in \mathcal{P}^{\leq d}(\mathbb{M}_{2^n})$, we have
\begin{equation}\label{LBI-C}
\|e^{-t \mathcal{L}} (T)\|_{L_{p}} \geq \frac{1}{T_{d}(e^{t})}\|T\|_{L_{p}}, \quad \forall t \geq0,
\end{equation}
where $T_{d}$ is the $d$-th Chebyshev polynomial of the first kind.
\end{theorem}
\begin{proof} 
Let $1\leq p\leq \infty$. The contractivity of the semigroup can be expressed as: 
$$\big\|x^{\mathcal{L}} (T)\big\|_{L_{p}} \leq\|T\|_{L_{p}},\quad \forall x \in[0,1],T\in\mathbb{M}_{2^n}.$$  
As for $x \in [-1,0]$, we have
\begin{align*}
\big\|(-x)^{\mathcal{L}} (T)\big\|_{L_{p}}&=\left\|\sum_{\mathbf{s} \in \{0,1,2,3\}^n} (-x)^{|\s|}\widehat{T}(\mathbf{s}) \sigma_{\mathbf{s}}\right\|_{L_{p}}\\
&=\left\|\sum_{\mathbf{s} \in \{0,1,2,3\}^n} x^{|\s|}\widehat{T}(\mathbf{s}) \varPhi_\rho(\sigma_{\s})\right\|_{L_{p}}\\
&=\left\|x^{\mathcal{L}} \varPhi_\rho(T)\right\|_{L_{p}}\\
&\leq \left\| \varPhi_\rho(T)\right\|_{L_{p}}=\| T\|_{L_{p}},
\end{align*}
where we used \eqref{var-def-2} and Lemma  \ref{var-prop} (iii). Consequently, we conclude from the above argument that
\begin{equation}\label{Cont-[-1,1]}
\big\|x^{\mathcal{L}} (T)\big\|_{L_{p}} \leq\|T\|_{L_{p}},\quad\forall x \in[-1,1].
\end{equation}
Let $\mu$ be any complex measure on $[-1,1]$. Then, averaging over \eqref{Cont-[-1,1]}, and applying the Jensen inequality, we get
\begin{align}\label{AC}
\left\|\int_{-1}^{1} x^{\mathcal{L}} (T) d\mu(x)\right\|_{L_{p}} & \leq \int_{-1}^{1}\|x^{\mathcal{L}} (T)\|_{L_{p}} \mathrm{d}|\mu|(x) \nonumber\\
& \stackrel{\eqref{Cont-[-1,1]}}{\leq} \int_{-1}^{1}\|T\|_{L_{p}} \mathrm{d}|\mu|(x)\\
& =\|\mu\|_{V([-1,1])}\|T\|_{L_{p}},\nonumber
\end{align}
where  $\|\mu\|_{V([-1,1])}$ denotes the total variation of $\mu$. Consider the linear functional
$$
\varphi_{t}:(\mathcal{P}^{\leq d}[-1,1],\|\cdot\|_{\mathcal{C}([-1,1])}) \rightarrow \mathbb{C},\quad \mbox{for every $t\geq0$}
$$
given by
\begin{equation*}
\varphi_t(P)=P(e^t),\quad \forall P\in \mathcal{P}^{\leq d}[-1,1].
\end{equation*}
Then, Chebyshev's inequality \eqref{cp} can be rewritten as
\begin{equation*}\label{Var-I1}
|\varphi_{t}(P)| =|P(e^t)|\leq T_{d}(e^t)\|P\|_{\mathcal{C}([-1,1])}, \quad\forall P \in \mathcal{P}^{\leq d}[-1,1]. 
\end{equation*}
Therefore, by the Hahn-Banach theorem and the Riesz representation theorem, there exists a complex measure $\mu_{t}$ on $[-1,1]$ such that 
\begin{equation}\label{mu-T}
\|\mu_{t}\|_{V([-1,1])}=\|\varphi_t\| \leq T_{d}(e^{t}),
\end{equation}
and 
\begin{equation*}\label{RP-C}
\int_{-1}^{1} P(x) \mathrm{d} \mu_{t}(x)=\varphi_t(P)=P(e^{t}),\quad \forall P\in\mathcal{P}^{\leq d}[-1,1].
\end{equation*}
Since $\widehat{T}(\mathbf{s})=0$ when $|\mathbf{s}|>d$, we deduce that for all $t \geq 0$,
\begin{align*}
\|e^{t \mathcal{L}} (T)\|_{L_{p}}
&=\left\|\sum_{\substack{\mathbf{s}\in\{0,1,2,3\}^n\\|\mathbf{s}|\leq d}}e^{t |\mathbf{s}|}\widehat{T}(\mathbf{s}) \sigma_{\mathbf{s}}\right\|_{L_{p}}\\
&\stackrel{\eqref{RP-C}}{=}\left\|\int_{-1}^1\sum_{\substack{\mathbf{s}\in\{0,1,2,3\}^n\\ |\mathbf{s}|\leq d}}x^{|\mathbf{s}|}\widehat{T}(\mathbf{s}) \sigma_{\mathbf{s}}d\mu_t(x)\right\|_{L_{p}}\\
&=\left\|\int_{-1}^{1} x^{\mathcal{L}} (T) \mathrm{~d} \mu_{t}(x)\right\|_{L_{p}} \\
& \stackrel{\eqref{AC},\eqref{mu-T}}{\leq} T_{d}(e^{t})\|T\|_{L_{p}}.
\end{align*}
We can get the desired estimate \eqref{LBI-C} by taking $T=e^{-t\mathcal{L}}(T)$.
\end{proof}


We should note that even though Theorem \ref{MN-K-most-d} can be understood as the ``dual'' of Theorem \ref{mn-k-tail}, it can not be deduced from Theorem \ref{mn-k-tail} via duality argument. This is merely a formal duality relation.

We now establish an improvement  of Theorem \ref{LB-C} whenever $1<p<\infty$. 

\begin{theorem}\label{MN-K-most-d}
Let $ d \in [n]$ and $1<p<\infty$. for every  $T\in \mathcal{P}^{\leq d}(\mathbb{M}_{2^n})$, there exist $C_p,c_p\in(0,\infty)$ and $\eta_p\in(\frac{1}{2},1]$ such that
\begin{equation}\label{MN-K-most-d-Ineq}
\big\|e^{-t \mathcal{L}} (T)\big\|_{L_{p}} \geq C_p e^{-c_p d \min \{t, t^{\eta_p}\}}\|T\|_{L_{p}},\quad\forall t \geq 0.
\end{equation}
\end{theorem}

To prove this theorem, we begin by stating a fundamental result from complex analysis, which will be instrumental in proving the subsequent proposition.
\begin{lemma}[{\cite[Lemma 29]{EI2020}}]\label{con-map-2}
Let $\Omega \subseteq \mathbb{C}$ be a bounded simply connected domain and let $\phi_{\Omega^{c}}: \overline{\Omega}^{c} \rightarrow \overline{\mathbb{D}}^{c}$  be  a conformal mapping of the complement of $\overline{\Omega}$ onto the complement of the closed unit disc,  satisfying $\phi_{\Omega^{c}}(\infty)=\infty$. Then
\begin{equation*}
\lim _{z \rightarrow \infty} \frac{\phi_{\Omega^{c}}(z)}{z}=\beta \in \mathbb{C} \backslash\{0\}. 
\end{equation*}
\end{lemma}


\begin{proposition}\label{Cpro-Reverse}
Let $ d \in [n]$, $1<p<\infty$, $K_p \in(0, \infty)$, and let $\Omega \subseteq \mathbb{D}$ be a simply connected domain bounded by a Jordan curve $\gamma$. Suppose that
\begin{equation}\label{omega-K}
\sup _{w \in \Omega} \|w^{\mathcal{L}}\|_{L_{p}\rightarrow L_{p}} \leq K_p.
\end{equation}
Also, let $\phi_{\Omega^{c}}: \overline{\Omega}^{c} \rightarrow \overline{\mathbb{D}}^{c}$ be a conformal mapping of the complement of $\overline{\Omega}$ onto the complement of the closed unit disc, which extends continuously as a homeomorphism between $\Omega^{c}$ and $\mathbb{D}^{c}$. Then, for  $T\in \mathcal{P}^{\leq d}(\mathbb{M}_{2^n})$,  we have
$$\|e^{-t \mathcal{L}} (T)\|_{L_{p}} \geq \frac{1}{K_p|\phi_{\Omega^{c}}(e^{t})|^{d}}\|T\|_{L_{p}},\quad \forall t \geq 0.
$$
\end{proposition}
\begin{proof} Consider the linear functional
$$\varphi_{t}:(\mathcal{P}(\overline{\Omega}),\|\cdot\|_{\mathcal{C}(\overline{\Omega})}) \rightarrow \mathbb{C}
$$
defined by
$$
\varphi_{t}(P)=\varphi_{t}(\sum_{k=0}^{d} a_{k} w^{k}) \stackrel{\text { def }}{=} \sum_{k=0}^{d} a_{k} e^{t k}=P(e^t),\quad P\in \mathcal{P}^{\leq d}(\overline{\Omega}).
$$
 For such a polynomial $P\in \mathcal{P}^{\leq d}(\overline{\Omega})$, define the function $h:\overline{\Omega}^{c}\rightarrow \mathbb{C}$ by
$$h(z)=\frac{P(z)}{(\phi_{\Omega^c})^d},\quad z\in\overline{\Omega}^{c},$$
where $\phi_{\Omega^{c}}$ is a conformal map that extends to a homeomorphism between $\Omega^{c}$ and $\mathbb{D}^{c}$, satisfying $\phi_{\Omega^{c}}(\infty)=\infty$. By Lemma \ref{con-map-2}, it follows  that $\lim _{z \rightarrow \infty} h(z) \in \mathbb{C}$, and thus $h$ is a holomorphic function on $\overline{\Omega}^{c} \cup\{\infty\}$. For $z \in \gamma$, we have $|\phi_{\Omega^{c}}(z)|=1$ since $\phi_{\Omega^{c}}$ is a homeomorphism between $\gamma$ and $\partial \mathbb{D}$. By the maximum principle,
$$
\sup _{z \in \Omega^{c}}|h(z)|=\max _{z \in \gamma}|h(z)|=\max _{z \in \gamma} \frac{|P(z)|}{|\phi_{\Omega^{c}}(z)|^{d}}=\|P\| _{\mathcal{C}(\gamma)}=\|P\|_{\mathcal{C}(\overline{\Omega})},$$
which implies that 
$$
|P(z)| \leq|\phi_{\Omega^{c}}(z)|^{d}\|P\|_{\mathcal{C}(\overline{\Omega})},\quad z \in \Omega^{c}.
$$
For  $z=e^t\in\Omega^{c}$, we obtain the following estimate
$$
|\varphi_{t}(P)| \leq|\phi_{\Omega^{c}}(e^{t})|^{d}\|P\|_{\mathcal{C}(\overline{\Omega})},\quad \forall P \in \mathcal{P}^{\leq d}(\overline{\Omega}). $$
Therefore, by the Hahn-Banach theorem and the Riesz representation theorem, there exists a complex measure $\rho_{t}$ on $\overline{\Omega}$ such that $\|\rho_{t}\|_{V(\overline{\Omega})} \leq|\phi_{\Omega^{c}}(e^{t})|^{d}$ and
\begin{equation}\label{pet-I}
\int_{\overline{\Omega}} P(w) \mathrm{d}\rho_{t}(w)=P(e^{t}),\quad P \in \mathcal{P}^{\leq d}(\overline{\Omega}). 
\end{equation}
Since $\widehat{T}(\mathbf{s})=0$ when $|\mathbf{s}|>d$, combining the above estimates, the Jensen inequality, and the given assumptions, it follows that
\begin{align*}
\|e^{t \mathcal{L}} (T)\|_{L_{p}} 
&=\Big\|\sum_{\mathbf{s}\in\{0,1,2,3\}^n,|\mathbf{s}|\leq d}e^{t |\mathbf{s}|}\widehat{T}(\mathbf{s}) \sigma_{\mathbf{s}}\Big\|_{L_{p}}\\
&\stackrel{\eqref{pet-I}}{=}\Big\|\int_{\overline{\Omega}}\sum_{\mathbf{s}\in\{0,1,2,3\}^n,|\mathbf{s}|\leq d}w^{|\mathbf{s}|}\widehat{T}(\mathbf{s}) \sigma_{\mathbf{s}}\mathrm{~d} \rho_{t}(w)\Big\|_{L_{p}}\\
&=\Big\|\int_{\overline{\Omega}}w^{\mathcal{L}} (T) \mathrm{~d} \rho_{t}(\omega)\Big\|_{L_{p}} \\
&\leq \int_{\overline{\Omega}}\|w^{\mathcal{L}} (T)\|_{L_{p}} \mathrm{d}|\rho_{t}|(w) \\
& \stackrel{\eqref{omega-K}}{\leq} K_p \int_{\overline{\Omega}}\|T\|_{L_{p}} \mathrm{d}\left|\rho_{t}\right|(w)=K_p\|\rho_{t}\|_{V(\overline{\Omega})}\|T\|_{L_{p}} \\
& \leq K_p |\phi_{\Omega^{c }}(e^{t})|^{d}\|T\|_{L_{p}},
\end{align*}
which is equivalent to the desired estimate by taking $T=e^{-t\mathcal{L}}(T)$.
\end{proof}


Now we are ready to prove Theorem \ref{MN-K-most-d}.
\begin{proof}[Proof of Theorem \ref{MN-K-most-d}] Let $d \in[n]$, $1<p<\infty$,   and $T\in \mathcal{P}^{\leq d}(\mathbb{M}_{2^n})$. By Theorem \ref{LB-C}, for every $t \geq 0$,
$$
\|e^{-t \mathcal{L}} (T)\|_{L_{p}} \geq \frac{1}{T_{d}(e^{t})}\|T\|_{L_{p}},
$$
where
$$T_{d}(e^t)=\frac{(e^t+\sqrt{e^{2t}-1})^{d}+(e^t-\sqrt{e^{2t}-1})^{d}}{2}\leq (2e^t)^d\leq e^{2td}.$$
The desired assertion for $1<t<\infty$ follows. 

It remains to consider the case $0<t\leq1$.  By Lemma \ref{or-k}, there exist $r_p \in[1, \infty)$ and  $K_p \in[1, \infty)$ such that
$$\sup _{w\in \Omega(r_p)} \|w^{\mathcal{L}}\|_{L_{p} \rightarrow L_{p}} \leq K_p, $$
where $\Omega(r_p)$ is the domain \eqref{or-def}. Therefore, by Proposition \ref{Cpro-Reverse},  we have
\begin{equation}\label{E-tl-I}
\|e^{-t\mathcal{L}}T\|_{L_{p}} \geq \frac{1}{K_p|\phi_{\Omega(r_p)^{c}}(e^{t})|^{d}}\|T\|_{L_{p}},\quad \forall t \geq 0.
\end{equation}
Let $\theta=\theta(r_p)=2 \arcsin (1 / r_p) \in(0, \pi]$. By Lemma \ref{con-map}, the  conformal map $\phi_{\Omega(r_p)^{c}}$ is given by
$$
\phi_{\Omega(r_p)^{c}}(e^{t})=\psi_2\circ\phi_{\frac{\pi}{2\pi-\theta}}\circ\psi_2(e^{t})=\frac{\left(e^{t}+1\right)^{\frac{\pi}{2 \pi-\theta}}+\left(e^{t}-1\right)^{\frac{\pi}{2 \pi-\theta}}}{\left(e^{t}+1\right)^{\frac{\pi}{2 \pi-\theta}}-\left(e^{t}-1\right)^{\frac{\pi}{2 \pi-\theta}}} . 
$$
Taking the logarithm of its modulus, we obtain
\begin{align*}
\log \left|\phi_{\Omega(r_p)^{c}}\left(e^{t}\right)\right| & =\log \left(1+\frac{2\left(e^{t}-1\right)^{\frac{\pi}{2 \pi-\theta}}}{\left(e^{t}+1\right)^{\frac{\pi}{2 \pi-\theta}}-\left(e^{t}-1\right)^{\frac{\pi}{2 \pi-\theta}}}\right) \\
& \leq \frac{2\left(e^{t}-1\right)^{\frac{\pi}{2 \pi-\theta}}}{\left(e^{t}+1\right)^{\frac{\pi}{2 \pi-\theta}}-\left(e^{t}-1\right)^{\frac{\pi}{2 \pi-\theta}}},
\end{align*}
where the inequality is due to  the fact $\log(1+x)\leq x$ for any $x\in(0,\infty)$. 
We further have
\begin{align*}
\frac{2(e^{t}-1)^{\frac{\pi}{2 \pi-\theta}}}{(e^{t}+1)^{\frac{\pi}{2 \pi-\theta}}-(e^{t}-1)^{\frac{\pi}{2 \pi-\theta}}} & \leq \frac{2(e^{t}-1)^{\frac{\pi}{2 \pi-\theta}}}{(e+1)^{\frac{\pi}{2 \pi-\theta}}-(e-1)^{\frac{\pi}{2 \pi-\theta}}} \\
& \leq \frac{2(e-1)^{\frac{\pi}{2 \pi-\theta}}}{(e+1)^{\frac{\pi}{2 \pi-\theta}}-(e-1)^{\frac{\pi}{2 \pi-\theta}}} t^{\frac{\pi}{2 \pi-\theta}}, 
\end{align*}
where the first inequality is because that the denominator is monotonically decreasing in $t$, and the second inequality is due to the fact $e^t-1\leq (e-1)t$ for $t\leq 1$. By \eqref{E-tl-I}, we obtain the desired estimate with $C=\frac{2(e-1)^{\frac{\pi}{2 \pi-\theta}}}{(e+1)^{\frac{\pi}{2 \pi-\theta}}-(e-1)^{\frac{\pi}{2 \pi-\theta}}}$ and $\eta_p=\frac{\pi}{2\pi-\theta}\in(\frac{1}{2},1]$. The proof is complete.
\end{proof}


\section{Bernstein-Markov type inequalities for Laplacian}
Recall that partial answers to Conjecture \ref{conj-2-2} have been provided in \cite{EI2020} and \cite{DDFKOSW2025}. In this section, we extend the corresponding Bernstein-Markov type inequalities in \cite{EI2020} and \cite{DDFKOSW2025} to the quantum setting. Although we adopt similar ideas to \cite{EI2020} and \cite{DDFKOSW2025}, we still need to apply the newly developed tools in noncommutative analysis, such as the noncommutative Cowling theorem. By the definition of \eqref{l-def-2}, when $p=2$, it follows directly from Parseval's identity
\begin{equation}\label{BM-2}
\|\mathcal{L}(T)\|_{L_{2}}\leq d\|T\|_{L_{2}},
\end{equation}
whenever $T\in \mathcal{P}^{\leq d}(\mathbb{M}_{2^n})$. 
The first purpose of this section is to establish a variant of \eqref{BM-2} on $L_p$ spaces when $p\neq 2$. 

\begin{theorem}\label{BM-d}
Let $d\in[n]$ and $1<p<\infty$. For  $T\in \mathcal{P}^{\leq d}(\mathbb{M}_{2^n})$,  there exist $\alpha_p \in[1,2)$ and $C_p \in(0, \infty)$ such that
$$
\|\mathcal{L} (T)\|_{L_{p}} \leq C_p d^{\alpha_p}\|T\|_{L_{p}}.
$$
\end{theorem}

We shall use the  classical approximation theoretic result of Szeg\"{o} \cite{S1925}, indeed, we just apply the following Szeg\"{o} theorem specifically for lens domains of the form \eqref{omega-r}. 

\begin{lemma}[{\cite[Proposition 15]{EI2020-JAT}}]\label{S-1} Let $r\in [1,\infty)$, and let lens domain $\Omega(r) \subseteq \mathbb{C}$ (defined in \eqref{omega-r}) be a domain bounded by a Jordan curve $\gamma$. The exterior angle of $\gamma$ at $\pm1$ is $2\pi-\theta(r)$ ($\theta(r)$ is given in \eqref{theta-r}). Then, 
for every $d \in \mathbb{N}$ and every polynomial $P\in \mathcal{P}^{\leq d}(\overline{\Omega})$, we have
$$
|P^{\prime}(1)| \leq 10 d^{2\pi-\theta(r) / \pi}\|P\|_{\mathcal{C}(\overline{\Omega})}.
$$
\end{lemma}




\begin{proof}[Proof of Theorem \ref{BM-d}]
Let $T\in \mathcal{P}^{\leq d}(\mathbb{M}_{2^n})$.
Consider the linear functional
$$\psi:(\mathcal{P}^{\leq d}(\overline{\Omega}),\|\cdot\|_{\mathcal{C}(\overline{\Omega})}) \rightarrow \mathbb{C}
$$
given by
$$
\psi(P) \stackrel{\text { def }}{=} P'(1),
$$
where 
$$\Omega=\{e^{-z}:|\arg z|<\theta\} \subseteq \mathbb{D},$$ 
where $\theta$ is given by Theorem \ref{e-zl-b}.
 Notice that $\Omega$ is a domain bounded by a Jordan curve $\gamma$ with $1 \in \gamma$ which forms an exterior angle $2 \pi-2 \theta$ at $1$. Then, Lemma \ref{S-1} implies that 
$$|\psi(P)| \leq 10 d^{2-\frac{2 \theta}{\pi}}\|p\|_{\mathcal{C}(\overline{\Omega})},\quad \forall P \in \mathcal{P}^{\leq d}(\overline{\Omega}).
$$
Therefore, by the Hahn-Banach theorem and the Riesz representation theorem, there exists a complex measure $v$ on $\overline{\Omega}$ such that 
$$\|v\|_{\mathcal{M}(\overline{\Omega})} \leq 10 d^{2-\frac{2 \theta}{\pi}}$$ 
and 
\begin{equation}\label{P'1-I}
\int_{\overline{\Omega}} P(w) \mathrm{d} v(w)=P^{\prime}(1),\quad P\in\mathcal{P}^{\leq d}(\overline{\Omega}).
\end{equation}
Since $\widehat{T}(\mathbf{s})=0$ when $|\mathbf{s}|>d$, we use \eqref{P'1-I}, Jensen inequality and Theorem \ref{e-zl-b} to derive
\begin{align*}
\|\mathcal{L}(T)\|_{L_{p}} 
&=\Big\|\sum_{\mathbf{s}\in\{0,1,2,3\}^n,|\mathbf{s}|\leq d} |\mathbf{s}|\widehat{T}(\mathbf{s}) \sigma_{\mathbf{s}}\Big\|_{L_{p}}\\
&\stackrel{\eqref{P'1-I}}{=}\Big\|\int_{\overline{\Omega}}\sum_{\mathbf{s}\in\{0,1,2,3\}^n,|\mathbf{s}|\leq d}w^{|\mathbf{s}|}\widehat{T}(\mathbf{s}) \sigma_{\mathbf{s}}\mathrm{~d} v(w)\Big\|_{L_{p}}\\
&=\Big\|\int_{\overline{\Omega}}w^{\mathcal{L}} (T) \mathrm{~d} v(w)\Big\|_{L_{p}} \\
&\leq \int_{\overline{\Omega}}\|w^{\mathcal{L}} (T)\|_{L_{p}} \mathrm{d}|v|(w) \\
& \leq M_p \int_{\overline{\Omega}}\|T\|_{L_{p}} \mathrm{d}|v|(w) \\
& \leq 10 M_pd^{2-\frac{2\theta}{\pi}}\|T\|_{L_{p}}.
\end{align*}
 The proof is complete.
\end{proof}


Next, we aim to apply the three-lines lemma to derive a bound for $\mathcal{L}^{k}$ when acting on low-degree polynomial. For each $p\in (1,\infty)$ and $\epsilon>0$, we put
\[
p_{\epsilon}\coloneqq
\begin{cases}
p+\epsilon,&\mbox{if }p\in[2,\infty);\\
p-\epsilon,&\mbox{if }1<p<2,~\epsilon\in(0,p-1).
\end{cases}
\]

\begin{theorem}\label{lower degree estimate}
Let $d\in [n]$, and let $p\in(1,\infty)$,  $\theta\in(0,1)$  be such that
\[
\frac{1}{p}=\frac{\theta}{p_{\epsilon}}+\frac{1-\theta}{2},\quad \epsilon\in (0,\infty).
\]
There exists a constant $K_{p,\epsilon}>0$ (depending only on $p$ and $\epsilon>0$) such that for each $k\in \mathbb{N}$ we have
\begin{equation}\label{bound for lower degree}
\left\|\mathcal{L}^{k}(T)\right\|_{L_p}\leq K_{p,\epsilon}^{k}d^{k}\|T\|^{1-\theta}_{L_2}\|T\|^{\theta}_{L_{p_{\epsilon}}},\quad T\in \mathcal{P}^{\leq d}(\mathbb{M}_{2^n}).
\end{equation}
\end{theorem}

The proof of Theorem \ref{lower degree estimate} is a consequence of the three-lines lemma and the following noncommutative Cowling inequality  regarding the bound of imaginary Riesz potential.
\begin{theorem}[{\cite{JW2017}}]\label{imaginary power}
Suppose that $\Delta$ is the infinitesimal generator of the Markov semigroup $\{e^{-t\Delta}\}_{t\geq 0}$ acting on a noncommutative probability space. For each $1<p<\infty$, there exists a universal constant $K>0$ such that
\[
\left\|\Delta^{iz}\right\|_{L_{p}\to L_{p}}\leq \frac{Kp^{2}}{p-1}(1+|z|^{3}\log^{2}(|z|))^{|1/p-1/2|}\exp\{\pi|1/p-1/2||z|\},\quad z\in \mathbb{R}.
\]
\end{theorem}

We now proceed to prove Theorem \ref{lower degree estimate}.

\begin{proof}[Proof of Theorem \ref{lower degree estimate}]
For each $z\in\mathbb{C}$ and $T\in \M_{2^{n}}$, we write
\[
\mathcal{L}^{z}(T)=\sum_{\s\in \{0,1,2,3\}^{n}}\widehat{T}(\s)|\s|^{z}\sigma_{\s}.
\]
Since the summation above is finite, it follows that $\mathcal{L}^{z}(T):\mathbb{C}\to \M_{2^{n}}$ is analytic.

In order to derive \eqref{bound for lower degree}, it suffices to show that there exists a constant $K_{p,\epsilon}>0$ such that
\begin{equation}\label{dual version}
\left|\langle \mathcal{L}^{k}(T),G \rangle\right|\leq K_{p,\epsilon}^{k}d^{k}\|T\|^{1-\theta}_{L_2}\|T\|^{\theta}_{L_{p_{\epsilon}}},
\end{equation}
for all $G\in L_{p^{\prime}}(\M_{2^{n}})$ with $\|G\|_{L_{p^{\prime}}}=1$ where $p^{\prime}$ is the conjugate index of $p$, i.e., $\frac{1}{p}+\frac{1}{p^{\prime}}=1$. We now apply the three-lines lemma to prove \eqref{dual version}. For any $G\in L_{p^{\prime}}(\M_{2^{n}})$ with $\|G\|_{L_{p^{\prime}}}=1$, we define $\psi:S\coloneqq\{z\in \mathbb{C}:0\leq \Re(z)\leq 1\}\to \mathbb{C}$ as follows
\[
\psi(z)\coloneqq \exp\left\{z^{2}\pi k|1/p_{\epsilon}-1/2|/(1-\theta)\right\}\left\langle \mathcal{L}^{(1-z)k/(1-\theta)}(T),G^{*}_{z}\right\rangle,
\]
where $G_{z}\coloneqq u(g)|G|^{\frac{p^{\prime}(1-z)}{2}+\frac{p^{\prime}z}{p_{\epsilon}^{\prime}}}$ and $u(G)$ is the partial isometry coming from the polar decomposition $G=u(G)|G|$. Following from the spectral decomposition of $|G|$, it is clear that $G_{z}$ is analytic, and hence $\psi$ is an analytic function on the open strip $\{z\in \mathbb{C}:0<\Re(z)<1\}$ and continuous on closed unit strip $S$. It is clear that $G_{\theta}=G$. Moreover, notice that $u(G)^*u(G)$ is an orthogonal projection onto $\mathrm{range}(|G|)$, it follows that
\begin{equation}
|G_{z}|=
\begin{cases}
|G|^{\frac{p^{\prime}}{2}},& \mbox{if }\Re(z)=0,\\
|G|^{\frac{p^{\prime}}{p_{\epsilon}^{\prime}}},& \Re(z)=1.
\end{cases}
\end{equation}

We first consider the case $z\in S$ with $\Re(z)=0$. In this case, we note that $\|G_z\|_{2}=1$ and $\|\mathcal{L}^{-\frac{zk}{1-\theta}}\|=1$ for each  $z\in S$. By the Cauchy-Schwarz inequality, we have
\begin{equation}\label{imaginary part}
\begin{split}
|\psi(z)|&=\left|\exp\left\{z^{2}\pi k|1/p_{\epsilon}-1/2|/(1-\theta)\right\}\left\langle \mathcal{L}^{(1-z)k/(1-\theta)}(T),G^{*}_{z}\right\rangle\right|\\
&=\left|\left\langle \mathcal{L}^{-\frac{zk}{1-\theta}}\left(\mathcal{L}^{\frac{k}{1-\theta}}(T)\right),G^{*}_{z}\right\rangle\right|\\
&\leq\left\| \mathcal{L}^{\frac{k}{1-\theta}}(T)\right\|_{L_2}\leq d^{\frac{k}{1-\theta}}\|T\|_{L_2},
\end{split}
\end{equation}
where we used \eqref{BM-2} in the last line.

We now turn to consider the case $z\in S$ with $\Re(z)=1$. Let $z=1+iv$ with $v\in \mathbb{R}$.  Applying Theorem \ref{imaginary power}, we get
\begin{equation}\label{real part 1}
\begin{split}
|\psi(z)|&=\left|\exp\left\{z^{2}\pi k|1/p_{\epsilon}-1/2|/(1-\theta)\right\}\left\langle \mathcal{L}^{(1-z)k/(1-\theta)}(T),G^{*}_{z}\right\rangle\right|\\
&\leq \exp\left\{(1-v^{2})\pi k|1/p_{\epsilon}-1/2|/(1-\theta)\right\} \left\|\mathcal{L}^{ikv/(1-\theta)}(T)\right\|_{L_{p_{\epsilon}}}\\
&\leq K_{p_{\epsilon}} \left(1+\left(\frac{k|v|}{1-\theta}\right)^{3}\log^{2}\left(\frac{k|v|}{1-\theta}\right)\right)^{{|1/p_{\epsilon}-1/2|}}\\
&\quad\times  \exp\left\{(1+|v|-v^{2})\pi k|1/p_{\epsilon}-1/2|/(1-\theta)\right\}\|T\|_{L_{p_{\epsilon}}}.
\end{split}
\end{equation}
Here we used the H\"{o}lder inequality and the fact $\|G_z\|_{p_{\epsilon}'}=1$. Note that $1+t^{3}\log^{2}(t)\leq 6e^{t}$ for all $t\geq 0$. Therefore, 
\begin{equation}\label{real part 2}
\begin{split}
|\psi(z)|&\leq K_{p_{\epsilon}}\exp\left\{\frac{\pi k|1/p_{\epsilon}-1/2|\left(1+(1+\frac{1}{\pi})|v|-v^{2}\right)}{1-\theta}\right\}\|T\|_{L_{p_{\varepsilon}}}\\
&\leq K_{p_{\epsilon}}\exp\left\{\frac{\pi k|1/p_{\epsilon}-1/2|\left(\frac{4\pi^{2}+(1+\pi)^{2}}{4\pi^{2}}\right)}{1-\theta}\right\}\|T\|_{L_{p_{\epsilon}}}.
\end{split}
\end{equation}
Combining \eqref{imaginary part} and \eqref{real part 2} with the three-lines lemma \cite[Lemma 1.3.5]{Gr2014}, we get that
\[
|\psi(\theta)|\leq K_{p,\epsilon}d^{k}\exp\left\{\frac{\pi\theta k|1/p_{\epsilon}-1/2|\left(\frac{4\pi^{2}+(1+\pi)^{2}}{4\pi^{2}}\right)}{1-\theta}\right\}\|T\|^{1-\theta}_{L_2}\|T\|^{\theta}_{L_{p_{\epsilon}}},
\]
where $\theta=\frac{(2-p)p_{\epsilon}}{(2-p_{\epsilon})p}$. The proof is complete.
\end{proof}


\begin{corollary}\label{BM-d-p+e}
Let $d\in [n]$, $p\in(1,\infty)$ and $\epsilon\in (0,\infty)$. 
There exists a constant $K_{p,\epsilon}>0$ (depending only on $p$ and $\epsilon>0$) such that
$$
\left\|\mathcal{L}(T)\right\|_{L_p}\leq K_{p,\epsilon}d\|T\|_{L_{p+\epsilon}},\quad T\in \mathcal{P}^{\leq d}(\mathbb{M}_{2^n}).
$$
\end{corollary}
\begin{proof} {\bf Case 1:}  $p\in [2,\infty)$. The assertion follows from  
 Theorem \ref{lower degree estimate} with $k=1$ and the fact $\|T\|_{L_2}\leq \|T\|_{L_{p+\epsilon}}$. 

{\bf Case 2:}  $p\in(1,2)$. By duality,  we have
\begin{equation}\label{L-dual-12}
\begin{split}
\left\|\mathcal{L}(T)\right\|_{L_p}&=\sup_{\|G\|_{L_{p'}}\leq1} | \langle\mathcal{L}(T),G\rangle|\\
&=\sup_{\|G\|_{L_{p'}}\leq1} |\langle T,\mathcal{L}(G)\rangle |\\
&\leq \|T\|_{L_{p+\epsilon}}\left\|\mathcal{L}(G)\right\|_{L_{(p+\epsilon)'}}.
\end{split}
\end{equation}
One can choose $\epsilon>0$ small enough so that $p<p+\epsilon<2$ and $(p+\epsilon)'+\epsilon <p'$. Applying the result of {\bf Case 1}, to \eqref{L-dual-12}, it follows that
$$\left\|\mathcal{L}(T)\right\|_{L_p}\leq K_{p,\epsilon} d\|T\|_{L_{p+\epsilon}}\left\|G\right\|_{L_{(p+\epsilon)'+\epsilon}}\leq K_{p,\epsilon} d \|T\|_{L_{p+\epsilon}}\left\|G\right\|_{L_{p'}}\leq K_{p,\epsilon} d \|T\|_{L_{p+\epsilon}}.$$
This finishes the proof.
\end{proof}

We have a further consequence of Theorem \ref{lower degree estimate}. 

\begin{corollary}
Let $d\in [n]$, $p\in(1,\infty)$, and let $T\in \mathcal{P}^{\leq d}(\mathbb{M}_{2^n})$. 
If $T$ is an unitary or a projection operator, then there exists a constant $K_{p}>0$  such that
$$
\left\|\mathcal{L}(T)\right\|_{L_p}\leq K_{p}d\|T\|_{L_p}.
$$
\end{corollary}
\begin{proof}
Note that, for any $r\in (0,\infty)$, $\|T\|_{L_r}=1$ whenever $T\in \M_{2^n}$ is unitary. The assertion follows immediately from Theorem \ref{lower degree estimate} (with $k=1$ and any given $\epsilon>0$). 

Now let us assume that $T$ is a projection, that is, $T$ is self-adjoint and $T^2=T$. In this case, we have
$$\|T\|_{L_2}^2=\|T\|_{L_p}^p=\|T\|_{L_{p_\epsilon}}^{p_\epsilon},\quad \forall \epsilon>0.$$
Consequently,
\begin{align*}
\|T\|^{1-\theta}_{L_2}\|T\|^{\theta}_{L_{p_{\epsilon}}}&=\|T\|^{(1-\theta)\frac{p}{2}}_{L_p}\|T\|^{\theta\frac{p}{p_{\epsilon}}}_{L_p}=\|T\|_{L_p}.
\end{align*}
Then the desired assertion for projection $T$ follows immediately from Theorem \ref{lower degree estimate} (with $k=1$ and any given $\epsilon>0$). 
\end{proof}

\begin{remark}
The last corollary means that quantum version of Conjecture \ref{conj-2-2} holds for  unitary or  projection elements which are of at most of degree $d$.  Its commutative version  can be found in Domelevo et al.  \cite{DDFKOSW2025}. 
\end{remark}


\section[Quantum inequalities for gradients]{Quantum Naor's inequality and  Bernstein-Markov type inequalities for gradients}
This section is devoted to deriving gradient estimate inequality and the corresponding Riesz transform-type inequality in quantum Boolean cubes.  To begin with, for $T\in \M_{2^{n}}$, let
\[
\|\nabla(T)\|_{L_{p}}\coloneqq \left\|(\dc_{j}(T))_{j=1}^{n}\right\|_{L_{p}(\M_{2^{n}};\ell^{cr}_{2})},
\]
where the norm given in the right-hand side comes from the noncommutative Khintchine inequality (see more details in Lemma \ref{nc-kh}). We first establish the gradient estimates presented in Theorem \ref{gradient estimate1} and Theorem \ref{Riesz-type 1}, which are quantum analogues of the results due to Naor, Eskenazis and Ivanisvili (see for example \cite[Proposition 43]{EI2020} and  \cite[Lemma 44]{EI2020}), and then we prove Bernstein-Markov type inequalities for gradient in Theorem \ref{BM-improve-gradient}. 

\begin{theorem}\label{gradient estimate1}
For each $p\in(1,\infty)$ and each $T\in L_{p}(\mathbb{M}_{2^n})$, there exists a constant $K_{p}>0$ depending only on $p$ such that  the following inequality holds
\begin{equation}\label{gradient estimate1-con}
\mathbb{E}_{\delta}\left\|\sum_{j=1}^{n}\delta_{j}\dc_{j}(e^{-t\mathcal{L}}(T))\right\|_{L_{p}}\leq
\frac{K_{p}}{(e^{\gamma_pt}-1)^{1/{\gamma_p}}}\|T\|_{L_{p}},\quad \forall t>0,
\end{equation}
where $\delta=\{\delta_{j}\}_{j=1}^{n}\in \{-1,1\}^n$ and  $\gamma_p=\min\{p,2\}$.
\end{theorem}
\begin{proof} Fix $t>0$,
and let $p\in [2,\infty)$ and $T\in\mathbb{M}_{2^n}$.	Recall that the following quantum (dual) Buser-type inequality was recently proved in \cite[Lemma 5.6]{JLLZ2025} (see also its proof there with changing $e^{2t}-1$ to $2t$):
\begin{equation}\label{Q-Buser}
\left\|\left(\sum_{j=1}^{n}|\dc_{j}(e^{-t\mathcal{L}}(T))|^2\right)^{1/2}\right\|_{L_p}\leq \frac{1}{\sqrt{e^{2t}-1}}\|T\|_{L_p}, \quad 2\leq p\leq\infty.
\end{equation}
Then, by the noncommutative Khintchine inequality (Lemma \ref{nc-kh}), we have
\begin{equation}\label{Byser}
\mathbb{E}_{\delta}\left\|\sum_{j=1}^{n}\delta_{j}\dc_{j}(e^{-t\mathcal{L}}(T))\right\|_{L_{p}}\leq C\sqrt{p}\|\nabla(e^{-t\mathcal{L}}(T))\|_{L_{p}}\leq \frac{K_p}{\sqrt{e^{2t}-1}}\|T\|_{L_{p}}.
\end{equation}	

It remains to deal with the case $p\in(1,2]$. We divide its proof into three steps.

{\bf Step I.} For each $T\in \M_{2^{n}}$ and $t>0$, we write
\[
T=\sum_{\s\in\{0,1,2,3\}^{n}}\widehat{T}(\s)\sigma_{\s},
\]
and define $H_t: L_p(\mathbb{M}_{2^n})\to L_p(\{-1,1\}^n, L_p(\mathbb{M}_{2^n}))$ by setting
\begin{equation}
H_{t}(T)(\delta)\coloneqq \sum_{\s\in\{0,1,2,3\}^{n}}\widehat{T}(\s)\sigma_{\s}\prod_{j\in \supp(\s)}(e^{-t}+(\frac{1-e^{-t}}3)\delta_{j}),
\end{equation}
where $\delta=\{\delta_{j}\}_{j=1}^{n}\in \{-1,1\}^n$.

 Recall that for each $A\subset [n]$, the Walsh function $w_A:\{-1,1\}^n\to \{-1,1\}$ is given by $w_A(\delta)=\prod_{j\in A}\delta_j$, $\delta=\{\delta_{j}\}_{j=1}^{n}\in \{-1,1\}^n$. By the definition of $H_{t}(T)$,
\begin{equation}\label{basic formula 1}
H_{t}(T)=\sum_{\s\in\{0,1,2,3\}^{n}}\sum_{B\subseteq \supp(\s)}e^{-t|B|}(\frac{1-e^{-t}}3)^{|\s|-|B|}\widehat{T}(\s)\sigma_{\s} w_{\supp(\s)\setminus B}.
\end{equation}


We now make the connection between the term $\sum_{j=1}^{n}\delta_{j}\dc_{j}\left(e^{-t\mathcal{L}}(T)\right)$ with $H_{t}(T)$ as follows.
Due to \eqref{basic formula 1}, it is clear that
\begin{equation}\label{basic formula 3}
\begin{aligned}
&\left(\mathrm{Id}_{\M_{2^{n}}}\otimes \mathrm{Rad}^{=1}_{\{-1,1\}^{n}}\right)(H_{t}(T))=\sum_{j=1}^{n}\delta_{j}\left(\mathrm{Id}_{\M_{2^{n}}}\otimes \mathbb{E}_{\delta}\right)\left(\delta_{j}H_{t}(T)\right)\\
&\qquad=\sum_{j=1}^{n}\delta_{j}\sum_{\s\in\{0,1,2,3\}^{n}}\sum_{B\subseteq\supp(\s)}\widehat{T}(\s)\sigma_{\s}e^{-t|B|}(\frac{1-e^{-t}}3)^{|\s|-|B|}\mathbb{E}(\delta_{j}w_{\supp(\s)\setminus B})\\
&\qquad=e^{t}(\frac{1-e^{-t}}3)\sum_{j=1}^{n}\delta_{j}\sum_{\substack{\s\in\{0,1,2,3\}^{n}\\s_{j}\neq 0}}e^{-t|\s|}\widehat{T}(\s)\sigma_{\s}\\
&\qquad=\frac{e^{t}-1}3\sum_{j=1}^{n}\delta_{j}\dc_{j}\left(e^{-t\mathcal{L}}(T)\right).
\end{aligned}
\end{equation}
Here, $\mathrm{Rad}^{=1}_{\{-1,1\}^{n}}$ is the classical Rademacher projection (see \eqref{rad cls} or \cite{Pi1982}). 
Since $\mathbb E_{\delta} w_A=0$ whenever $A\neq\varnothing$,
	only the term $B=\supp(\s)$ survives after taking expectation.
	Therefore,
	\begin{equation}\label{mean-of-modified-H}
		\left(
		\mathrm{Id}_{\mathbb M_{2^n}}\otimes\mathbb E_\delta
		\right)(H_t(T))
		=
		\sum_{\s}
		e^{-t|\s|}
		\widehat T(\s)\sigma_{\s}
		=
		e^{-t\mathcal L}(T).
	\end{equation}

{\bf Step II.} For each $t>0$ and $T\in \M_{2^{n}}$, let
\[
S_{t}(T)\coloneqq H_{t}(T)-e^{-t\mathcal{L}}(T).
\]
Note that
\[
\|S_{t}(T)\|_{L_{1}}\leq \|H_{t}(T)\|_{L_{1}}+\|e^{-t\mathcal{L}}(T)\|_{L_{1}}\leq \|H_{t}(T)\|_{L_{1}}+\|T\|_{L_{1}}.
\]
To get an upper bound for $\|S_{t}(T)\|_{L_{1}}$, it remains to estimate the term $\|H_{t}(T)\|_{L_{1}}$. 
For
	$-1/3\leq\rho\leq1$, define the single-qubit Pauli channel
	$D_\rho:\mathbb M_2\to\mathbb M_2$ by
	\[
	D_\rho(T)
	\coloneqq
	\frac{1+3\rho}{4}T
	+
	\frac{1-\rho}{4}
	\sum_{\ell=1}^3\sigma_\ell T\sigma_\ell.
	\]
	The coefficients in the preceding expression are nonnegative and
	their sum is equal to one. Thus, $D_\rho$ is completely positive
	and trace-preserving. Moreover, by \eqref{car-relation}, 
	\[
	D_\rho(\sigma_0)=\sigma_0,
	\qquad
	D_\rho(\sigma_\ell)=\rho\sigma_\ell,
	\quad \ell=1,2,3.
	\]
For $\delta\in{-1,1}^{n}$, set
	\[
	\rho_j(\delta)\coloneqq e^{-t}+\frac{1-e^{-t}}{3}\delta_j.
	\]
we have
	\[-\frac13\leq\rho_j(\delta)\leq1.\]
By the definition of $H_t$,
	\[
	H_t(T)(\delta)
	=
	\left(
	D_{\rho_1(\delta)}
	\otimes\cdots\otimes
	D_{\rho_n(\delta)}
	\right)(T).
	\]
The tensor product in the preceding identity is completely positive
	and trace-preserving and is therefore contractive on $L_1$.
	Consequently,
\[\|H_t(T)\|_{L_1(\Omega_n;L_1(\mathbb M_{2^n}))}=
		\mathbb E_{\delta}
		\|H_t(T)(\delta)\|_{L_1}\leq
		\|T\|_{L_1}.\]
	Since $e^{-t\mathcal L}$ is also contractive on $L_1$, it follows that $\|S_{t}(T)\|_{L_{1}}\leq 2\|T\|_{L_{1}}$.

Since $\{\sigma_{\s}\otimes w_{A}\}_{\s\in\{0,1,2,3\}^{n},A\subseteq[n]}$ forms an orthonormal basis in $L_{2}(\M_{2^{n}})\otimes L_{2}(\{-1,1\}^{n})$, it follows from \eqref{basic formula 1} that
\begin{align*}
\|S_{t}(T)\|^{2}_{L_{2}}&=\left\|H_{t}(T)-e^{-t\mathcal{L}}(T)\right\|^{2}_{L_{2}}\\
&=\left\|\sum_{\s\in\{0,1,2,3\}^{n}}\sum_{B\subsetneq \supp(\s)}e^{-t|B|}(\frac{1-e^{-t}}{3})^{|\s|-|B|}\widehat{T}(\s)\sigma_{\s} w_{\supp(\s)\setminus B}\right\|^{2}_{L_{2}}\\
&=\sum_{\s\in\{0,1,2,3\}^{n}}\sum_{B\subsetneq \supp(\s)}|\widehat{T}(\s)|^{2}e^{-2t|B|}(\frac{1-e^{-t}}{3})^{2(|\s|-|B|)}\\
&=\sum_{\s\in\{0,1,2,3\}^{n}}|\widehat{T}(\s)|^{2}\left(\sum_{l=0}^{|\s|-1}\binom{|\s|}{l}e^{-2tl}(\frac{1-e^{-t}}{3})^{2(|\s|-l)}\right).
\end{align*}
Note that for each $\s\in \{0,1,2,3\}^n$, 
\begin{align*}
	\sum_{l=0}^{|\s|-1}\binom{|\s|}{l}e^{-2tl}(\frac{1-e^{-t}}{3})^{2(|\s|-l)}&=(e^{-2t}+(\frac{1-e^{-t}}{3})^{2})^{|\s|}-e^{-2t|\s|}\\
	&\leq(e^{-2t}+(1-e^{-t})^{2})^{|\s|}-e^{-2t|\s|}\leq 1-e^{-t},
\end{align*}
where we used the elementary fact $(u^{2}+(1-u)^{2})^{l}\leq u^{2l}+(1-u)$ whenever $u\in(0,1)$ and $l\in \mathbb{N}$. Indeed,  it suffices to note that
\begin{align*}
\sum_{j=1}^{l}\binom{l}{j}u^{2(l-j)}(1-u)^{2j}&\leq(1-u)\left(\sum_{j=1}^{l}\binom{l}{j}u^{2(l-j)}(1-u)^{j}\right)\\
&\leq (1-u)\left(\sum_{j=0}^{l}\binom{l}{j}u^{l-j}(1-u)^{j}\right)=1-u.
\end{align*}
Now we conclude from the above argument that 
$$\|S_{t}(T)\|_{L_{2}}\leq (1-e^{-t})^{1/2} \|T\|_{L_2}.$$
For general $p\in(1,2)$, we choose $\theta\in(0,1)$ such that $\frac{1}{p}=\frac{\theta}{1}+\frac{1-\theta}{2}$ and apply the complex interpolation theorem (see e.g. \cite[Theorem 1.1.1]{BL1976}) to get that
\[
\|S_{t}\|_{L_{p}\to L_{p}}\leq 2^{\frac{2}{p}-1}(1-e^{-t})^{1-\frac{1}{p}}.
\]

{\bf Step III.} For general $p\in(1,2]$ and $\mathfrak{T}\in  L_{p}(\{-1,1\}^{n};L_{p}(\M_{2^{n}}))$, we have
%\begin{align*}
%\left\|\left(\mathrm{Id}_{\M_{2^{n}}}\otimes \mathrm{Rad}^{=1}_{\{-1,1\}^{n}}\right)(\mathfrak{T})\right\|_{L_{p}}&\leq \left\|\left(\mathrm{Id}_{\M_{2^{n}}}\otimes \mathrm{Rad}^{=1}_{\{-1,1\}^{n}}\right)(\mathfrak{T})\right\|_{L_{2}}\\
%&=e^{t}\left\|\left(\mathrm{Id}_{\M_{2^{n}}}\otimes \mathrm{Rad}^{=1}_{\{-1,1\}^{n}}\right)\circ\left(\mathrm{Id}_{\M_{2^{n}}}\otimes e^{-t\Delta}\right)(\mathfrak{T})\right\|_{L_{2}}\\
%&\leq e^{t}\left\|\left(\mathrm{Id}_{\M_{2^{n}}}\otimes e^{-t\Delta}\right)(\mathfrak{T})\right\|_{L_{2}},
%\end{align*}
%and
%\[
%\left\|\left(\mathrm{Id}_{\M_{2^{n}}}\otimes \mathrm{Rad}^{=1}_{\{-1,1\}^{n}}\right)(\mathfrak{T})\right\|_{L_{p}}\leq e^{t}\left\|\left(\mathrm{Id}_{\M_{2^{n}}}\otimes e^{-t\Delta}\right)(\mathfrak{T})\right\|_{L_{2}}\leq{\frac{1}{\sqrt{p-1}}}\|\mathfrak{T}\|_{L_{p}}.
%\]
\[\left\|\left(\mathrm{Id}_{\M_{2^{n}}}\otimes \mathrm{Rad}^{=1}_{\{-1,1\}^{n}}\right)(\mathfrak{T})\right\|_{L_{p}}\leq C_p\|\mathfrak{T}\|_{L_{p}}.\] 
Therefore, combining the previous estimates with \eqref{basic formula 3} and taking the expectation with respect to $\{\delta_{j}\}_{j=1}^{n}$ yields that
\begin{equation}\label{basic formula 4}
\begin{split}
\mathbb{E}_{\delta}\left\|\sum_{j=1}^{n}\delta_{j}\dc_{j}(e^{-t\mathcal{L}}(T))\right\|_{L_{p}}&=\frac{3}{e^{t}-1}\mathbb{E}_{\delta}\left\|\left(\mathrm{Id}_{\M_{2^{n}}}\otimes \mathrm{Rad}^{=1}_{\{-1,1\}^{n}}\right)(H_{t}(T))\right\|_{L_{p}}\\
&=\frac{3}{e^{t}-1}\mathbb{E}_{\delta}\left\|\left(\mathrm{Id}_{\M_{2^{n}}}\otimes \mathrm{Rad}^{=1}_{\{-1,1\}^{n}}\right)(S_{t}(T))\right\|_{L_{p}}\\
&\leq \frac{3C_p}{(e^{t}-1)}\left\|S_{t}(T)\right\|_{L_{p}}\\
&\leq \frac{3C_p2^{\frac{2}{p}-1}(1-e^{-t})^{1-\frac{1}{p}}}{(e^{t}-1)}\|T\|_{L_{p}}\\
&\leq K_p\frac{1}{(e^{tp}-1)^{1/p}}\|T\|_{L_{p}},
\end{split}
\end{equation}
where the last inequality follows from
\[
\frac{(1-e^{-t})^{1-\frac1p}}{e^t-1}
\leq
\frac{p^{1/p}}{(e^{pt}-1)^{1/p}}.
\]
The proof is complete.
\end{proof}
%\begin{proof} Fix $t>0$,
%and let $p\in [2,\infty)$ and $T\in\mathbb{M}_{2^n}$.	Recall that the following quantum (dual) Buser-type inequality was recently proved in \cite[Lemma 5.6]{JLLZ2025} (see also its proof there with changing $e^{2t}-1$ to $2t$):
%\begin{equation}\label{Q-Buser}
%\left\|\left(\sum_{j=1}^{n}|\dc_{j}(e^{-t\mathcal{L}}(T))|^2\right)^{1/2}\right\|_{L_p}\leq \frac{1}{\sqrt{e^{2t}-1}}\|T\|_{L_p}, \quad 2\leq p\leq\infty.
%\end{equation}
%Then, by the noncommutative Khintchine inequality (Lemma \ref{nc-kh}), we have
%\begin{equation}\label{Byser}
%\mathbb{E}_{\delta}\left\|\sum_{j=1}^{n}\delta_{j}\dc_{j}(e^{-t\mathcal{L}}(T))\right\|_{L_{p}}\leq C\sqrt{p}\|\nabla(e^{-t\mathcal{L}}(T))\|_{L_{p}}\leq \frac{K_p}{\sqrt{e^{2t}-1}}\|T\|_{L_{p}}.
%\end{equation}	
%
%It remains to deal with the case $p\in(1,2]$. We divide its proof into three steps.
%
%{\bf Step I.} For each $T\in \M_{2^{n}}$ and $t>0$, we write
%\[
%T=\sum_{\s\in\{0,1,2,3\}^{n}}\widehat{T}(\s)\sigma_{\s},
%\]
%and define $H_t: L_p(\mathbb{M}_{2^n})\to L_p(\{-1,1\}^n, L_p(\mathbb{M}_{2^n}))$ by setting
%\begin{equation}
%H_{t}(T)(\delta)\coloneqq \sum_{\s\in\{0,1,2,3\}^{n}}\widehat{T}(\s)\sigma_{\s}\prod_{j\in \supp(\s)}(e^{-t}+(1-e^{-t})\delta_{j}),
%\end{equation}
%where $\delta=\{\delta_{j}\}_{j=1}^{n}\in \{-1,1\}^n$.
%
%There exist two explicit formulas for $H_{t}(T)$. Recall that for each $A\subset [n]$, the Walsh function $w_A:\{-1,1\}^n\to \{-1,1\}$ is given by $w_A(\delta)=\prod_{j\in A}\delta_j$, $\delta=\{\delta_{j}\}_{j=1}^{n}\in \{-1,1\}^n$. By the definition of $H_{t}(T)$,
%\begin{equation}\label{basic formula 1}
%H_{t}(T)=\sum_{\s\in\{0,1,2,3\}^{n}}\sum_{B\subseteq \supp(\s)}e^{-t|B|}(1-e^{-t})^{|\s|-|B|}\widehat{T}(\s)\sigma_{\s} w_{\supp(\s)\setminus B}.
%\end{equation}
%On the other hand, we have
%\begin{equation}\label{basic formula 2}
%\begin{split}
%H_{t}(T)&=\sum_{\s\in\{0,1,2,3\}^{n}}\widehat{T}(\s)\sigma_{\s}\prod_{j=1}^{n}\left(e^{-t}+(1-e^{-t})\delta_j^{\mathds{1}_{\supp(\s)}(j)}\right)\\
%&=\sum_{\s\in\{0,1,2,3\}^{n}}\widehat{T}(\s)\sigma_{\s}\sum_{B\subseteq [n]}e^{-t|B|}(1-e^{-t})^{n-|B|}w_{\supp(\s)\setminus B}\\
%&=\sum_{B\subseteq [n]}\sum_{\s\in\{0,1,2,3\}^{n}}e^{-t|B|}(1-e^{-t})^{n-|B|}\widehat{T}(\s)\sigma_{\s}w_{\supp(\s)\setminus B}.
%\end{split}
%\end{equation}
%
%We now make the connection between the term $\sum_{j=1}^{n}\delta_{j}\dc_{j}\left(e^{-t\mathcal{L}}(T)\right)$ with $H_{t}(T)$ as follows.
%Due to \eqref{basic formula 1}, it is clear that
%\begin{equation}\label{basic formula 3}
%\begin{aligned}
%&\left(\mathrm{Id}_{\M_{2^{n}}}\otimes \mathrm{Rad}^{=1}_{\{-1,1\}^{n}}\right)(H_{t}(T))=\sum_{j=1}^{n}\delta_{j}\left(\mathrm{Id}_{\M_{2^{n}}}\otimes \mathbb{E}_{\{-1,1\}^{n}}\right)\left(\delta_{j}H_{t}(T)\right)\\
%&\qquad=\sum_{j=1}^{n}\delta_{j}\sum_{\s\in\{0,1,2,3\}^{n}}\sum_{B\subseteq\supp(\s)}\widehat{T}(\s)\sigma_{\s}e^{-t|B|}(1-e^{-t})^{|\s|-|B|}\mathbb{E}(\delta_{j}w_{\supp(\s)\setminus B})\\
%&\qquad=e^{t}(1-e^{-t})\sum_{j=1}^{n}\delta_{j}\sum_{\substack{\s\in\{0,1,2,3\}^{n}\\s_{j}\neq 0}}e^{-t|\s|}\widehat{T}(\s)\sigma_{\s}\\
%&\qquad=(e^{t}-1)\sum_{j=1}^{n}\delta_{j}\dc_{j}\left(e^{-t\mathcal{L}}(T)\right).
%\end{aligned}
%\end{equation}
%Here, $\mathrm{Rad}^{=1}_{\{-1,1\}^{n}}$ is the classical Rademacher projection (see \eqref{rad cls} or \cite{Pi1982}). 
%
%
%{\bf Step II.} For each $t>0$ and $T\in \M_{2^{n}}$, let
%\[
%S_{t}(T)\coloneqq H_{t}(T)-e^{-t\mathcal{L}}(T).
%\]
%Note that
%\[
%\|S_{t}(T)\|_{L_{1}}\leq \|H_{t}(T)\|_{L_{1}}+\|e^{-t\mathcal{L}}(T)\|_{L_{1}}\leq \|H_{t}(T)\|_{L_{1}}+\|T\|_{L_{1}}.
%\]
%To get an upper bound for $\|S_{t}(T)\|_{L_{1}}$, it remains to estimate the term $\|H_{t}(T)\|_{L_{1}}$. Thanks to \eqref{basic formula 2}, it follows that
%\begin{align*}
%\|H_{t}(T)\|_{L_{1}}&=\left\|\sum_{B\subseteq [n]}\sum_{\s\in\{0,1,2,3\}^{n}}e^{-t|B|}(1-e^{-t})^{n-|B|}\widehat{T}(\s)\sigma_{\s}w_{\supp(\s)\setminus B}\right\|_{L_{1}}\\
%&\leq \sum_{l=0}^{n}\binom{n}{l}e^{-tl}(1-e^{-t})^{n-l}\left\|\sum_{\s\in\{0,1,2,3\}^{n}}\widehat{T}(\s)\sigma_{\s}\right\|_{L_{1}}=\|T\|_{L_{1}}.
%\end{align*}
%Thus, we have $\|S_{t}(T)\|_{L_{1}}\leq 2\|T\|_{L_{1}}$.
%
%Since $\{\sigma_{\s}\otimes w_{A}\}_{\s\in\{0,1,2,3\}^{n},A\subseteq[n]}$ forms an orthonormal basis in $L_{2}(\M_{2^{n}})\otimes L_{2}(\{-1,1\}^{n})$, it follows from \eqref{basic formula 1} that
%\begin{align*}
%\|S_{t}(T)\|^{2}_{L_{2}}&=\left\|H_{t}(T)-e^{-t\mathcal{L}}(T)\right\|^{2}_{L_{2}}\\
%&=\left\|\sum_{\s\in\{0,1,2,3\}^{n}}\sum_{B\subsetneq \supp(\s)}e^{-t|B|}(1-e^{-t})^{|\s|-|B|}\widehat{T}(\s)\sigma_{\s} w_{\supp(\s)\setminus B}\right\|^{2}_{L_{2}}\\
%&=\sum_{\s\in\{0,1,2,3\}^{n}}\sum_{B\subsetneq \supp(\s)}|\widehat{T}(\s)|^{2}e^{-2t|B|}(1-e^{-t})^{2(|\s|-|B|)}\\
%&=\sum_{\s\in\{0,1,2,3\}^{n}}|\widehat{T}(\s)|^{2}\left(\sum_{l=0}^{|\s|-1}\binom{|\s|}{l}e^{-2tl}(1-e^{-t})^{2(|\s|-l)}\right).
%\end{align*}
%Note that for each $\s\in \{0,1,2,3\}^n$, 
%$$\sum_{l=0}^{|\s|-1}\binom{|\s|}{l}e^{-2tl}(1-e^{-t})^{2(|\s|-l)}=(e^{-2t}+(1-e^{-t})^{2})^{|\s|}-e^{-2t|\s|}\leq 1-e^{-t},$$
% where we used the elementary fact $(u^{2}+(1-u)^{2})^{l}\leq u^{2l}+(1-u)$ whenever $u\in(0,1)$ and $l\in \mathbb{N}$. Indeed,  it suffices to note that
%\begin{align*}
%\sum_{j=1}^{l}\binom{l}{j}u^{2(l-j)}(1-u)^{2j}&\leq(1-u)\left(\sum_{j=1}^{l}\binom{l}{j}u^{2(l-j)}(1-u)^{j}\right)\\
%&\leq (1-u)\left(\sum_{j=0}^{l}\binom{l}{j}u^{l-j}(1-u)^{j}\right)=1-u.
%\end{align*}
%Now we conclude from the above argument that 
%$$\|S_{t}(T)\|_{L_{2}}\leq (1-e^{-t})^{1/2} \|T\|_{L_2}.$$
%
%
%For general $p\in(1,2)$, we choose $\theta\in(0,1)$ such that $\frac{1}{p}=\frac{\theta}{1}+\frac{1-\theta}{2}$ and apply the complex interpolation theorem (see e.g. \cite[Theorem 1.1.1]{BL1976}) to get that
%\[
%\|S_{t}\|_{L_{p}\to L_{p}}\leq 2^{\frac{2}{p}-1}(1-e^{-t})^{1-\frac{1}{p}}.
%\]
%
%{\bf Step III.} For general $p\in(1,2]$ and $\mathfrak{T}\in L_{p}(\M_{2^{n}})\otimes L_{p}(\{-1,1\}^{n})$, we have
%\begin{align*}
%\left\|\left(\mathrm{Id}_{\M_{2^{n}}}\otimes \mathrm{Rad}^{=1}_{\{-1,1\}^{n}}\right)(\mathfrak{T})\right\|_{L_{p}}&\leq \left\|\left(\mathrm{Id}_{\M_{2^{n}}}\otimes \mathrm{Rad}^{=1}_{\{-1,1\}^{n}}\right)(\mathfrak{T})\right\|_{L_{2}}\\
%&=e^{t}\left\|\left(\mathrm{Id}_{\M_{2^{n}}}\otimes \mathrm{Rad}^{=1}_{\{-1,1\}^{n}}\right)\circ\left(\mathrm{Id}_{\M_{2^{n}}}\otimes e^{-t\Delta}\right)(\mathfrak{T})\right\|_{L_{2}}\\
%&\leq e^{t}\left\|\left(\mathrm{Id}_{\M_{2^{n}}}\otimes e^{-t\Delta}\right)(\mathfrak{T})\right\|_{L_{2}},
%\end{align*}
%and
%\[
%\left\|\left(\mathrm{Id}_{\M_{2^{n}}}\otimes \mathrm{Rad}^{=1}_{\{-1,1\}^{n}}\right)(\mathfrak{T})\right\|_{L_{p}}\leq e^{t}\left\|\left(\mathrm{Id}_{\M_{2^{n}}}\otimes e^{-t\Delta}\right)(\mathfrak{T})\right\|_{L_{2}}\leq{\frac{1}{\sqrt{p-1}}}\|\mathfrak{T}\|_{L_{p}}.
%\]
%Therefore, combining the previous estimates with \eqref{basic formula 3} and taking the expectation with respect to $\{\delta_{j}\}_{j=1}^{n}$ yields that
%\begin{equation}\label{basic formula 4}
%\begin{split}
%\mathbb{E}_{\delta}\left\|\sum_{j=1}^{n}\delta_{j}\dc_{j}(e^{-t\mathcal{L}}(T))\right\|_{L_{p}}&=\frac{1}{e^{t}-1}\mathbb{E}_{\delta}\left\|\left(\mathrm{Id}_{\M_{2^{n}}}\otimes \mathrm{Rad}^{=1}_{\{-1,1\}^{n}}\right)(H_{t}(T))\right\|_{L_{p}}\\
%&=\frac{1}{e^{t}-1}\mathbb{E}_{\delta}\left\|\left(\mathrm{Id}_{\M_{2^{n}}}\otimes \mathrm{Rad}^{=1}_{\{-1,1\}^{n}}\right)(S_{t}(T))\right\|_{L_{p}}\\
%&\leq \frac{1}{(p-1)^{1/2}(e^{t}-1)}\left\|S_{t}(T)\right\|_{L_{p}}\\
%&\leq \frac{2^{\frac{2}{p}-1}(1-e^{-t})^{1-\frac{1}{p}}}{(p-1)^{1/2}(e^{t}-1)}\|T\|_{L_{p}}\\
%&=C\frac{2^{\frac{2}{p}-1}}{(p-1)^{1/2}(e^{tp}-1)^{1/p}}\|T\|_{L_{p}}.
%\end{split}
%\end{equation}
%The proof is complete.
%\end{proof}

\begin{remark}\label{gradient estimate 2}
For each $p\in(1,\infty]$, there exists a constant $K_{p}>0$ depending only on $p$ such that
\begin{equation}
\|\nabla(e^{-t\mathcal{L}}(T))\|_{L_{p}}\leq \begin{cases} 
\frac{K_{p}}{(e^{pt}-1)^{1/p}}\|T\|_{L_{p}}, & 1<p<2, \\ 
\frac{1}{\sqrt{e^{2t}-1}}\|T\|_{L_{p}}, & {2\leq p\leq\infty,}
\end{cases}\quad\forall~t>0.
\end{equation}
\end{remark}
\begin{proof}
In the case of $2\leq p\leq\infty$, we can obtain it directly from  \eqref{Q-Buser}. Applying the noncommutative Khintchine inequality Lemma \ref{nc-kh} to the left-hand side of \eqref{gradient estimate1-con} of Theorem \ref{gradient estimate1} yields a constant $K_{p}>0$ depending only on $p\in(1,2)$ such that
\[
\left\|\nabla(e^{-t\mathcal{L}}(T))\right\|_{L_{p}}\leq\frac{K_{p}}{(e^{tp}-1)^{1/p}}\|T\|_{L_{p}},\quad\mbox{for all }t>0,
\]
which is the desired inequality.
\end{proof}

The following result,  a noncommutative analogy of Naor's Riesz transform inequality \cite[Proposition 43]{EI2020},  can be deduced from Theorem \ref{gradient estimate1} and the functional calculus of the number operator.

\begin{theorem}\label{Riesz-type 1}
For each $p\in(1,\infty)$ and $\epsilon\in(\max\{0,\frac{1}{2}-\frac{1}{p}\},1-\frac{1}{p})$, there exists a constant $K_{p}>0$ depending only on $p$ such that for each $n\in \mathbb{N}$, we have
\begin{equation}\label{Riesz-type 2}
\left\|\nabla(T)\right\|_{L_{p}}\leq \left(\frac{K_{p}}{\epsilon}\right)\left\|\mathcal{L}^{\frac{1}{p}+\epsilon}(T)\right\|_{L_{p}},\quad\forall~T\in L_{p}(\M_{2^{n}}).
\end{equation}
\end{theorem}
\begin{proof}
By the functional calculus of $\mathcal{L}$, we have
\begin{equation}\label{Riesz-type 3}
\mathcal{L}^{-(\frac{1}{p}+\epsilon)}=\frac{1}{\Gamma(\frac{1}{p}+\epsilon)}\int_{0}^{\infty}t^{\frac{1}{p}+\epsilon-1}e^{-t\mathcal{L}}dt.
\end{equation}
By Lemma \ref{nc-kh}, the Kahane-Khintchine inequality (see for instance \cite{Pi2016})  and \eqref{Riesz-type 3}, we have
\begin{equation}\label{eq Delta 1}
\begin{aligned}
\left\|\nabla \mathcal{L}^{-(\frac{1}{p}+\epsilon)}(T)\right\|_{L_{p}}&\leq A^{(1)}_{p}\mathbb{E}_{\delta}\left\|\sum_{j=1}^{n}\delta_{j}\dc_{j}\left(\mathcal{L}^{-(\frac{1}{p}+\epsilon)}(T)\right)\right\|_{L_p}\\
&\leq A^{(2)}_{p}\int_{0}^{\infty}t^{\frac{1}{p}+\epsilon-1}\mathbb{E}_{\delta}\left\|\sum_{j=1}^{n}\delta_{j}\dc_{j}(e^{-t\mathcal{L}}(T))\right\|_{L_{p}}dt,
\end{aligned}
\end{equation}
which, together with Theorem \ref{gradient estimate1} for $1<p\leq 2$, implies
$$\left\|\nabla \mathcal{L}^{-(\frac{1}{p}+\epsilon)}(T)\right\|_{L_{p}}\leq A^{(3)}_{p}\|T\|_{L_{p}}\int_{0}^{\infty}\frac{t^{\frac{1}{p}+\epsilon-1}}{(e^{tp}-1)^{1/p}}dt.$$
Note that (see \cite[eq. (214)]{EI2020})
$$\int_{0}^{\infty}\frac{t^{\frac{1}{p}+\epsilon-1}}{(e^{tp}-1)^{1/p}}dt=O(\frac{1}{\epsilon}),\quad p\in (1,2].$$
The desired inequality for $p\in(1,2]$ follows. 

As for the case $p\in (2,\infty)$, combining \eqref{eq Delta 1} and Theorem \ref{gradient estimate1} for $2<p<\infty$, we have
$$\left\|\nabla \mathcal{L}^{-(\frac{1}{p}+\epsilon)}(T)\right\|_{L_{p}}\leq A^{(4)}_{p}\|T\|_{L_{p}}\int_{0}^{\infty}\frac{t^{\frac{1}{p}+\epsilon-1}}{\sqrt{e^{2t}-1}}dt.$$
Hence, it suffices to check 
$$I:=\int_{0}^{\infty}\frac{t^{\frac{1}{p}+\epsilon-1}}{\sqrt{e^{2t}-1}}dt<\infty.$$ 
We write
$$I=\int_{0}^{1} \frac{t^{\frac{1}{p}+\epsilon-1}}{\sqrt{e^{2t}-1}} dt+\int_{1}^{\infty} \frac{t^{\frac{1}{p}+\epsilon-1}}{\sqrt{e^{2t}-1}}dt=:I_1+I_2.$$
For  $I_1$, as $t \to 0^+$, we have  $e^{2t} - 1 \approx 2t$, the integrand becomes:
$$
\frac{t^{\frac{1}{p}+\epsilon-1}}{\sqrt{e^{2t}-1}} \approx \frac{t^{\frac{1}{p}+\epsilon-1}}{\sqrt{2t}} = \frac{t^{\frac{1}{p}-\frac{1}{2}+\epsilon-1}}{\sqrt{2}}.
$$
Thus, the integrand behaves like $t^{\frac{1}{p}-\frac{1}{2}+\epsilon-1}$, and the integral converges since $ \epsilon > \frac{1}{2}-\frac{1}{p}$. The exact result is:
$$
\int_{0}^{1} \frac{t^{\frac{1}{p}-\frac{1}{2}+\epsilon-1}}{\sqrt{2}} dt = \frac{1}{\sqrt{2}} \int_{0}^{1} t^{\frac{1}{p}-\frac{1}{2}+\epsilon-1}dt = \frac{1}{\sqrt{2}} \left[ \frac{t^{\frac{1}{p}-\frac{1}{2}+\epsilon}}{\frac{1}{p}-\frac{1}{2}+\epsilon} \right]_0^1 =: \frac{K_p}{\epsilon}.
$$
In the second part $I_2$, $e^{pt}$ grows very rapidly as $t\to \infty$, so the denominator $\sqrt{e^{2t} - 1} \approx e^{t}$. The integrand becomes:
$$
\frac{t^{\frac{1}{p}+\epsilon-1}}{\sqrt{e^{2t} - 1}} \approx t^{\frac{1}{p}+\epsilon-1} e^{-t}.
$$
Due to the exponential decay $e^{-t}$, the integral always converges as $t \to \infty$. The exact result is:
$$
\int_{1}^{\infty} t^{\frac{1}{p}+\epsilon-1} e^{-t} dt = \Gamma\left(\frac{1}{p}+\epsilon, 1\right),
$$
where $\Gamma(a, x)$ is the incomplete Gamma function.
Combining the results of the two parts, we obtain the desired result.
\end{proof}





The following result is motivated by \cite[Theorem 9.22]{ODo2014}. 

\begin{lemma}\label{2-1}
	For every  $T\in\mathcal{P}^{\leq d}(\mathbb{M}_{2^n})$, we have
	$$\|T\|_{L_2}\leq e^{d}\|T\|_{L_1}.$$
\end{lemma}
\begin{proof} For $\epsilon>0$, and let $0<\theta<1$ be the solution of $\frac{1}{2}=\frac{\theta}{1}+\frac{1-\theta}{2+\epsilon}$, that is $\theta=\frac{\epsilon}{2(1+\epsilon)}$. Applying  H\"{o}lder’s inequality, we get
\begin{equation}\label{2-1-theta}
\|T\|_{L_2}\leq \|T\|_{L_{2+\epsilon}}^{1-\theta}\|T\|_{L_1}^\theta.
\end{equation}
In addition, for each $2\leq q<\infty$ and $T\in\mathcal{P}^{\leq d}(\mathbb{M}_{2^n})$, we have
\begin{equation}\label{q-2}
\|T\|_{L_q}\leq \sqrt{q-1}^d\|T\|_{L_2}.
\end{equation}
Indeed, using Lemma \ref{hc-1}, we have
$$\|T\|_{L_q}^2=\|e^{-\log\sqrt{q-1}}(e^{\log\sqrt{q-1}}(T))\|_{L_q}^2\leq \|e^{\log\sqrt{q-1}}(T)\|_{L_2}^2$$
and
$$\|e^{\log\sqrt{q-1}}(T)\|_{L_2}^2=\sum_{k=0}^d\sum_{|\s|=k}(q-1)^k|\widehat{T}(\s)|^2\leq (q-1)^d\|T\|_{L_2}^2.$$
The inequality \eqref{q-2} now follows from the last two inequalities. Then, by \eqref{2-1-theta} and \eqref{q-2} with $q=2+\epsilon$,  we get
$$\|T\|_{L_2}\leq (1+\epsilon)^\frac{d(1-\theta)}{2}\|T\|_{L_2}^{1-\theta}\|T\|_{L_1}^\theta,$$
which implies
\begin{align*}
\|T\|_{L_2}&\leq (1+\epsilon)^\frac{d(1-\theta)}{2\theta}\|T\|_{L_1}=\Big((1+\epsilon)^{\frac{1}{\epsilon}+\frac{1}{2}}\Big)^d\|T\|_{L_1}.
\end{align*}
Since $\lim_{\epsilon\to0}(1+\epsilon)^{\frac{1}{\epsilon}+\frac{1}{2}}=e$, the desired result follows.
\end{proof}


Now we can prove the following quantum Bernstein–Markov type inequality for the gradient  as a
consequence of Theorem \ref{Riesz-type 1} and Theorem \ref{BM-d}.

\begin{theorem}\label{BM-improve-gradient} Let $n\geq1$, $d\in[n]$, $1\leq p\leq\infty$ and $\theta_p\in(0,\frac{\pi}{2}]$ as Theorem \ref{e-zl-b}. Then there exist $\alpha_p \in[1,2)$ and $C_p \in(0, \infty)$ such that, for every  $T\in\mathcal{P}^{\leq d}(\mathbb{M}_{2^n})$, we have
	$$
	\|\nabla (T)\|_{L_{p}} \leq 
	\begin{cases} 
		C_p d^{\frac{2}{p}-\frac{\theta_p}{p\pi}}\log(d+1)\|T\|_{L_{p}}, & 1< p<\infty; \\ 
		2d\|T\|_{L_{p}}, & p=\infty;\\
		e^d\sqrt{d}\|T\|_{L_{p}}& p=1.
	\end{cases}
	$$
\end{theorem}
\begin{proof}  For each $1<p<\infty$ and $\epsilon\in(\max\{0,\frac{1}{2}-\frac{1}{p}\},1-\frac{1}{p})$, 
	by Theorem \ref{Riesz-type 1} and Lemma \ref{nsi} with $\beta=\frac{1}{p}+\epsilon$, we get
	\begin{align*}
		\|\nabla (T)\|_{L_{p}} &\leq \frac{K_p}{\epsilon}\|\mathcal{L}^{\frac{1}{p}+\epsilon}(T)\|_{L_p}\leq \frac{4K_p}{\epsilon}\|\mathcal{L}(T)\|_{L_p}^{\frac{1}{p}+\epsilon}\|T\|_{L_p}^{1-(\frac{1}{p}+\epsilon)}.
	\end{align*}
	Using Theorem \ref{BM-d}, there exist $\alpha_p=2-\frac{\theta_p}{\pi}\in[1,2)$ (where $\theta_p\in(0,\frac{\pi}{2}]$ is as in Theorem \ref{e-zl-b}) and $c_p>0$ such that 
	$$\|\nabla (T)\|_{L_{p}}\leq\left(\frac{4c_pK_p}{\epsilon}\right)d^{(2-\frac{\theta_p}{\pi})(\frac{1}{p}+\epsilon)}\|T\|_{L_{p}}.$$
	Choosing $\epsilon=\frac{1}{\log(d+1)}$, we have
\begin{align*}
\|\nabla (T)\|_{L_{p}}&\leq 4c_pK_pd^{\frac{2}{p}-\frac{\theta_p}{p\pi}}d^{(2-\frac{\theta_p}{\pi})\frac{1}{\log(d+1)}}\log(d+1)\|T\|_{L_{p}}\\
&\leq 4c_pK_pd^{\frac{2}{p}-\frac{\theta_p}{p\pi}}d^{\frac{2}{\log(d+1)}}\log(d+1)\|T\|_{L_{p}}\\
&\leq C_pd^{\frac{2}{p}-\frac{\theta_p}{p\pi}}\log(d+1)\|T\|_{L_{p}},
\end{align*}
where the last inequality follows from $d^{\frac{2}{\log(d+1)}}\leq e^2$.

	For the case $p=\infty$, by the quantum (dual) Buser inequality \eqref{Q-Buser} and Theorem \ref{LB-C} with $T_d(e^t)\leq e^{td^2}$,  we have
	$$\|\nabla (T)\|_{L_{\infty}}\leq \frac{1}{\sqrt{e^{2t}-1}}\|e^{t\mathcal{L}}(T)\|_{L_{\infty}}\leq \frac{e^{td^2}}{\sqrt{e^{2t}-1}}\|T\|_{L_{\infty}}.$$
	Taking $t=\frac{1}{d^{2}}$ and by the inequality $\sqrt{e^{2t}-1}\geq \sqrt{2t}$, we obtain
	$$\|\nabla (T)\|_{L_{\infty}}\leq \frac{e}{\sqrt{2t}}\|T\|_{L_{\infty}}\leq 2d\|T\|_{L_{\infty}}.$$
	
	For the final case where $p=1$, the result follows from orthogonality, $\widehat{T}(\s)=0$ whenever $|\s|>d$ and Lemma \ref{2-1}. Indeed, it is clear that
\begin{equation*}\label{gradient estimate 1}
\begin{split}
\|\nabla (T)\|_{L_{1}}&\leq \|\nabla (T)\|_{L_{2}}\leq\left\|\sum_{j=1}^n|d_j(T)|^2\right\|_{L_{1}}^{\frac{1}{2}}\\
&\leq \left(\sum_{j=1}^n\||d_j(T)|^2\|_{L_{1}}\right)^{\frac{1}{2}}\leq\left(\sum_{j=1}^n\||d_j(T)|\|_{L_{2}}^2\right)^{\frac{1}{2}},
\end{split}
\end{equation*}
which, together with Lemma \ref{2-1}, implies
\[
\begin{split}
\|\nabla (T)\|_{L_{1}}\leq \left(\sum_{\substack{\s\in\{0,1,2,3\}^n\\|\s|\leq d}}|\s||\widehat{T}(\s)|^2\right)^{\frac{1}{2}}\leq\sqrt{d}\|T\|_{L_{2}}\leq e^d\sqrt{d}\|T\|_{L_{1}}.
\end{split}
\]
\end{proof}
\begin{remark}\label{p12infty} Theorem \ref{BM-improve-gradient}  can be traced back to \cite[Theorem 15]{EI2020} and \cite[Proposition 46]{EI2020}; however, for $2\leq p<\infty$,  \cite[Theorem 15]{EI2020} was proved by using  Lust-Piquard's  Riesz transform inequality (\cite{BL2008,Lu1998}), and here we apply the quantum (dual) Buser inequality established recently in \cite{JLLZ2025}.
\end{remark}

\chapter[Covariance tail quantum inequalities]{Covariance quantum Poincar\'e and quantum Talagrand-type inequalities in tail spaces}\label{sec-tala-tail}

This chapter develops Talagrand and covariance inequalities for
quantum observables whose Pauli--Fourier spectrum is supported on the
tail space $\mathcal{T}^{\geq d}(\mathbb{M}_{2^n})$. The main point is that the lower spectral cutoff degree $d$ yields specific decay property under the quantum Ornstein--Uhlenbeck semigroup and therefore strengthens the corresponding influence-based estimates.

We first establish tail Talagrand inequalities for $1\leq p<2$.  In particular, the geometric
influence estimate obtained from Theorem \ref{lp-tala-d-tail} has a full logarithmic denominator and improves the dependence on small influences in the corresponding result of \cite{RWZ2024}.  We then derive tail-sensitive covariance Poincar\'e and covariance Talagrand
inequalities.  These results are obtained within a common semigroup framework and make the improvement arising from spectral localization explicit.

\section{Quantum tail Talagrand inequalities }\label{improved-talagrand}

For $j\in[n]$, $1\leq p<\infty$, and $T\in\mathbb{M}_{2^n}$, define the
$j$-th $L_p$-influence of $T$ by
\[
\mathrm{Inf}_j^p(T)\coloneqq \|\dc_j(T)\|_{L_p}^p.
\]
The case $p=1$ is commonly referred to as the \emph{geometric influence}; see,
for example, \cite{RWZ2024}.  When $p=2$, we write
\[
\mathrm{Inf}_j(T)\coloneqq \mathrm{Inf}_j^2(T),
\qquad
\mathrm{Inf}(T)\coloneqq \sum_{j=1}^n \mathrm{Inf}_j(T).
\]
By \eqref{dj-def} and Parseval's identity,
\[
\mathrm{Inf}_j(T)=\sum_{\substack{\s\in\{0,1,2,3\}^n\\ s_j\neq0}}|\widehat T(\s)|^2.
\]
We use the variance convention
\[
\mathrm{Var}(T)
 \coloneqq \mathrm{tr}(|T|^2)-|\mathrm{tr}(T)|^2
 =\|T-\mathrm{tr}(T)\mathbf{1}\|_{L_2}^2.
\]
The next result is tail-version of the \emph{Talagrand inequality}, which quantifies the variance decay property for $T$ whose Fourier support is contained in levels at least $d$.

\begin{theorem}\label{lp-tala-d-tail}
Let $p\in[1,2)$ and $d\in [n]$. For every $T\in \mathcal{T}^{\geq d}(\M_{2^n})$ with $\|T\|_{L_\infty}\leq1$,  we have
\[\mathrm{Var}(T)
\leq\left(\frac{4}{2-p}\right)\sum_{j=1}^{n}\frac{\|d_{j}(T)\|^{p}_{L_p}}{d+\log\left(1/\|d_{j}(T)\|^{p}_{L_p}\right)}.\]
\end{theorem}

To prove Theorem \ref{lp-tala-d-tail}, we begin by recalling some elementary estimates.  The first is the spectral decay on tail spaces.

\begin{lemma}\label{etl-t-l2-pd}
For every $T \in \mathcal{T}^{\geq d}(\M_{2^n})$, we have
$$\|e^{-t \mathcal{L}}(T)\|_{L_{2}} \leq e^{-td}\|T\|_{L_{2}}.$$
\end{lemma}

%\begin{lemma}[{\cite[Proposition 2.7]{JLLZ2025}}]\label{t-infty} For $T\in\mathbb{M}_{2^n}$ with $\|T\|_{L_\infty}\leq1$ and each $j\in[n]$, we have
%\[
%\|\dc_{j}(T)\|_{L_\infty}\leq1.
%\]
%\end{lemma}

The second estimate is a scalar integral bound that will retain both the tail parameter and the logarithmic influence term.

\begin{lemma}\label{tail-kernel-estimate}
For every $d\geq1$ and $\lambda\geq0$,
\[
\int_{0}^\infty e^{-2dt-\lambda\tanh(t)}\,\mathrm{d}t\leq \frac{1}{d+\lambda}.
\]
\end{lemma}

\begin{proof}
Set $y=\tanh(t)=\frac{e^{t}-e^{-t}}{e^{t}+e^{-t}}$.  Since
\[
e^{-2t}=\frac{1-y}{1+y}\qquad\text{and}\qquad \mathrm{d}t=\frac{\mathrm{d}y}{1-y^2},
\]
changing the variable the integral equals
\[
\int_{0}^{1}\frac{(1-y)^{d-1}}{(1+y)^{d+1}}e^{-\lambda y}\,\mathrm{d}y.
\]
For $0\leq y<1$, let
\[
F_d(y)\coloneqq dy+(d-1)\log(1-y)-(d+1)\log(1+y).
\]
It is easy to see that $F_d(0)=0$ and
\[
F_{d}^{\prime}(y)=\frac{2y-d(1+y^2)}{1-y^2}\leq0,
\]
where we used $d\geq1$ and $2y\leq1+y^2$.  Hence
\[
\frac{(1-y)^{d-1}}{(1+y)^{d+1}}\leq e^{-dy}.
\]
It follows that
\[
\int_0^\infty e^{-2dt-\lambda\tanh t}\,\mathrm{d}t \leq\int_0^1e^{-(d+\lambda)y}\,\mathrm{d}y \leq\frac{1}{d+\lambda}.
\]
\end{proof}


\begin{proof}[Proof of Theorem \ref{lp-tala-d-tail}]
Let $d\in [n]$ and  $T\in \mathcal{T}^{\geq d}$.
By Lemma \ref{cov-semigroup-0-infty} and Lemma \ref{etl-t-l2-pd}, we have 
\begin{equation}\label{semigroup-approach-lp}
\begin{aligned}
\mathrm{Var}(T)
&=\sum_{j=1}^{n}\int_{0}^{\infty}\mathrm{tr}((e^{-t\mathcal{L}}\dc_{j}(T))^*\cdot \dc_{j}(T)) \mathrm{d}t\\
&=4\sum_{j=1}^{n}\int_{0}^{\infty}\mathrm{tr}((e^{-2t\mathcal{L}}\dc_{j}(T))^*\cdot e^{-2t\mathcal{L}}\dc_{j}(T)) \mathrm{d}t\\
&=4\sum_{j=1}^{n}\int_{0}^{\infty}\left\|e^{-2t\mathcal{L}}(\dc_{j}(T))\right\|^{2}_{L_{2}} \mathrm{d}t\\
&\leq 4\sum_{j=1}^{n}\int_{0}^{\infty}e^{-2td}\|e^{-t\mathcal{L}}(\dc_{j}(T))\|^{2}_{L_{2}}\mathrm{d}t\\
&\leq 4\sum_{j=1}^{n}\int_{0}^{\infty}e^{-2td}\|\dc_{j}(T)\|^{2}_{L_{1+e^{-2t}}} \mathrm{d}t,
\end{aligned}
\end{equation}
where the last inequality is due to the (optimal) hypercontractivity given in Theorem \ref{hc-1}.
Hence, to get the desired inequality it suffices to estimate
\begin{equation}\label{key-term}
\int_{0}^{\infty}e^{-2td}\|\dc_{j}(T)\|^{2}_{L_{1+e^{-2t}}}\mathrm{d}t,\quad j\in[n].
\end{equation}

 Assume first that $1<p<2$. Denote that $X_{j}\coloneqq \dc_{j}(T)$ for each $j\in [n]$. Put
\[
a_k=\frac{k}{2}\log\Bigl(\frac{1}{p-1}\Bigr),
\qquad k\geq0,
\]
and, for a fixed $j$, set
\[
I_{j}\coloneqq\int_0^\infty e^{-2td}\|X_j\|_{L_{1+e^{-2t}}}^2~\mathrm{d}t.
\]
Writing the integral as a sum over the intervals $[a_k,a_{k+1}]$ and
translating each interval to $[0,a_1]$, we obtain
\begin{align*}
I_{j}&=\sum_{k=0}^\infty (p-1)^{kd}\int_0^{a_1}e^{-2ud}\|X_j\|_{L_{1+(p-1)^ke^{-2u}}}^2~\mathrm{d}u \\
&\leq\sum_{k=0}^\infty (p-1)^k\int_0^{a_1}e^{-2ud}\|X_j\|_{L_{1+e^{-2u}}}^2~\mathrm{d}u \\
&=\frac{1}{2-p}\int_0^{a_1}e^{-2ud}\|X_j\|_{L_{1+e^{-2u}}}^2~\mathrm{d}u.
\end{align*}
Here we used $d\geq1$ and the monotonicity of the $L_{r}$-norms with respect to
$r$ for the normalized trace.
By Lemma \ref{jlzz-lem}, $\|X_j\|_{L_\infty}\leq1$ for each $j\in [n]$.  Set
\[
b_j=\|X_j\|_{L_p}^p
\qquad\text{and}\qquad
\lambda_j=\log(1/b_j),\quad j\in [n]
\]
whenever $X_j\neq0$.  For $0\leq u\leq a_1$, the exponent
$r(u)=1+e^{-2u}$ belongs to $[p,2]$.  Interpolation between $L_p$ and
$L_\infty$ therefore gives
\[
\|X_j\|_{L_{r(u)}}^2
 \leq b_j^{2/r(u)}
 =b_j^{1+\tanh(u)}
 =b_j e^{-\lambda_j\tanh(u)}.
\]
Consequently, Lemma \ref{tail-kernel-estimate} yields
\begin{equation}\label{tail-Ij-power}
I_{j}\leq \frac{b_j}{(2-p)(d+\lambda_j)}=\left(\frac{1}{2-p}\right)\frac{\|X_{j}\|_{L_p}^p}{d+\log\left(1/\|X_j\|_{L_p}^p\right)}.
\end{equation}
Combining this estimate with \eqref{semigroup-approach-lp} proves the claim
for $1<p<2$.

It remains to consider $p=1$.  In this case interpolation between $L_1$ and
$L_\infty$ applies for every $t\geq0$.  If
$a_j=\|X_j\|_{L_1}\neq0$ and $\mu_j=\log(1/a_j)$, then
\[
\|X_j\|_{L_{1+e^{-2t}}}^2
 \leq a_j^{2/(1+e^{-2t})}
 =a_j e^{-\mu_j\tanh(t)}.
\]
Combining Lemma \ref{tail-kernel-estimate} and \eqref{semigroup-approach-lp}, we obtain 
\[
\mathrm{Var}(T)
 \leq4\sum_{j=1}^n
 \frac{\|\dc_{j}(T)\|_{L_1}}
 {d+\log\bigl(1/\|\dc_{j}(T)\|_{L_1}\bigr)},
\]
which is the asserted estimate at $p=1$.
\end{proof}

A direct consequence of Theorem \ref{lp-tala-d-tail} is the following geometric influence Talagrand inequality, which was established in \cite[Theorem 3.2]{RWZ2024} and serves as the key ingredient for quantum KKL-type inequality invoking geometric influence; see \cite[Theorem 3.9]{RWZ2024}.

\begin{corollary}\label{cov-l1-tala-d-tail-d-1}
For each $T\in \mathbb{M}_{2^n}$ with $\|T\|_{L_\infty}\leq1$, we have
\[
\mathrm{Var}(T)\leq4\sum_{j=1}^{n}\frac{\|\dc_{j}(T)\|_{L_{1}}}{1+\log\left(1/\|\dc_{j}(T)\|_{L_{1}}\right)}.
\]
\end{corollary}
\begin{proof}
Choosing $p=1$ and $d=1$ in Theorem \ref{lp-tala-d-tail} yields the desired result.
\end{proof}

We conclude this section by comparing Theorem \ref{lp-tala-d-tail} with the geometric-influence inequality obtained in \cite[Theorem 3.2]{RWZ2024}.  The latter asserts that, for every
$A\in\mathbb{M}_{2^n}$ with $\|A\|_{L_\infty}\leq1$,
\[
\mathrm{Var}(A)\leq C\sum_{j=1}^{n}
 \frac{\|\dc_{j}(A)\|_{L_1}\bigl(1+\|\dc_{j}(A)\|_{L_1}\bigr)}
 {\bigl[1+\log^+\bigl(1/\|\dc_{j}(A)\|_{L_1}\bigr)\bigr]^{1/2}},
\]
where $C>0$ is universal constant.  By contrast, the case $p=d=1$
of Theorem \ref{lp-tala-d-tail} gives
\[
\mathrm{Var}(A)
 \leq 8\sum_{j=1}^n
 \frac{\|\dc_{j}(A)\|_{L_1}}
 {1+\log\bigl(1/\|\dc_{j}(A)\|_{L_1}\bigr)}.
\]
Since $\|\dc_{j}(A)\|_{L_1}\leq1$, the latter estimate improves the
dependence on each individual geometric influence: it replaces the
square-root logarithmic denominator in \cite{RWZ2024} by a full
logarithmic denominator and removes the additional factor
$1+\|\dc_{j}(A)\|_{L_1}$.  In particular, writing
$a_j=\|\dc_{j}(A)\|_{L_1}$, one has, for $0<a_j\leq1$,
\[
\frac{a_j}{1+\log(1/a_j)}
 \leq
\frac{a_j(1+a_j)}
 {\bigl[1+\log^+(1/a_j)\bigr]^{1/2}}.
\]
More generally, Theorem \ref{lp-tala-d-tail} provides a family of estimates for every
$1\leq p<2$ and retains the lower Fourier-level parameter $d$ through
the denominator $d+\log\left(1/\|\dc_{j}(T)\|_{L_p}\right)$. Thus the estimate becomes stronger when the Fourier support of $T$ is pushed farther into the tail space. 

\section{Covariance Poincar\'e and Talagrand-type inequalities}\label{covariance-poincare-for-tail}

This section aims to derive covariance estimates for observables whose Pauli--Fourier spectra are supported on the tail space $\mathcal{T}^{\geq d}(\M_{2^{n}})$. More precisely, covariance Poincar\'e inequalities and covariance Talagrand-type inequalities will be established via semigroup tools. The semigroup representation below reduces covariance to coordinatewise pairings. For non-Hilbertian norms, we combine this representation with the heat smoothing estimate of Theorem \ref{mn-k-tail}; in Hilbertian case, the exact spectral decay yields the sharper factor $d^{-1}$.

For $T_1,T_2\in\mathbb{M}_{2^n}$, the covariance of For $T_1,T_2$ is defined by
\begin{equation}\label{cov-def}
\mathrm{Cov}(T_1,T_2)=\mathrm{tr}((T_1-\mathrm{tr}(T_1))^*)(T_2-\mathrm{tr}(T_2))=\mathrm{tr}(T_1^*T_2)-\overline{\mathrm{tr}(T_1)}\mathrm{tr}(T_2).
\end{equation}
Thus, $\mathrm{Cov}(T,T)=\mathrm{Var}(T)$ for each $T\in \M_{2^{n}}$.
\begin{lemma}\label{cov-semigroup-0-infty}
For $T_1,T_2\in\mathbb{M}_{2^n}$, we have
$$\mathrm{Cov}(T_1,T_2)=\sum_{j=1}^n\int_{0}^\infty\mathrm{tr}[((e^{-t\mathcal{L}}\dc_{j}(T_1)))^*\dc_{j}(T_2)]dt.$$
\end{lemma}
\begin{proof} For $T_1,T_2\in\mathbb{M}_{2^n}$, according to the definition \eqref{cov-def}  and  the interpolation of semigroups, we have
\begin{equation}\label{cov-1}
\begin{aligned}
\mathrm{Cov}(T_1,T_2)&=\mathrm{tr}(T_1^*T_2)-\overline{\mathrm{tr}(T_1)}\mathrm{tr}(T_2)\\
&=-\int_{0}^\infty\mathrm{tr}((\partial_te^{-t\mathcal{L}}(T_1))^*T_2)~\mathrm{d}t,
\end{aligned}
\end{equation}
where we used $\lim_{t\to 0^{+}}e^{-t\mathcal{L}}(T)=T$ and $\lim_{t\to \infty}e^{-t\mathcal{L}}(T)=\tr(T)$.
Since for each $j\in[n]$ and  each $T\in\mathbb{M}_{2^n}$, we have
\begin{align*}
\partial_te^{-t\mathcal{L}}(T)&=\partial_t\left(\sum_{\mathbf{s}\in\{0,1,2,3\}^n}e^{-t|\mathbf{s}|}\widehat{T}(\mathbf{s})\sigma_\mathbf{s}\right)\\
&=-\sum_{\mathbf{s}\in\{0,1,2,3\}^n}|\mathbf{s}|e^{-t|\mathbf{s}|}\widehat{T}(\mathbf{s}|)\sigma_\mathbf{s}\\
&=-\sum_{j=1}^n\sum_{\substack{\mathbf{s}\in\{0,1,2,3\}^n\\j\in\mathrm{supp}(\mathbf{s})}}e^{-t|\mathbf{s}|}\widehat{T}(\mathbf{s})\sigma_\mathbf{s}\\
&=-\sum_{j=1}^n\dc_{j}e^{-t\mathcal{L}}(T).
\end{align*}
Combining the above equation with Lemma \ref{intertwining} and  the orthogonality of $\{\sigma_{\s}\}_{\s\in\{0,1,2,3\}^{n}}$ in $L_{2}(\M_{2^{n}})$, we obtain that
\begin{align*}
\mathrm{Cov}(T_1,T_2)&=\sum_{j=1}^n\int_{0}^\infty\mathrm{tr}((\dc_{j}e^{-t\mathcal{L}}(T_1))^*T_2)dt\\
&=\sum_{j=1}^n\int_{0}^\infty\mathrm{tr}[(e^{-t\mathcal{L}}\dc_{j}(T_1))^*\dc_{j}(T_2)]dt.
\end{align*}
The proof is complete.
\end{proof}

The following lemma provides consequences of tail heat smoothing that will be used here and in the next section. The first estimate is a dimension-free exponential decay estimate on the first tail space (i.e., tail space with $d=1$); the second is an estimate on the $d$-tail space.

\begin{lemma}\label{etl-pp-theta}
Fix $1<p<\infty$, and let $C_{p},c_{p}>0$ and $A_{p}\geq1$ be the constants in Theorem \ref{mn-k-tail}.  Put $\alpha_p\coloneqq A_{p}^{-1}$.
\begin{enumerate}[{\rm (i)}]
\item There exist constants $B_{r}\geq1$ and $b_{p}>0$, depending only on $p$, such that
\begin{equation}\label{mean-zero-lp-decay}
\|e^{-t\mathcal{L}}T\|_{L_p}
 \leq B_p e^{-b_pt}\|T\|_{L_p},
\qquad t\geq0,
\end{equation}
whenever $\tr(T)=0$.

\item There exists a constant $K_p>0$, depending only on $p$, such that
\begin{equation}\label{integrated-tail-decay}
\int_0^\infty\|e^{-t\mathcal{L}}T\|_{L_p}\,\mathrm{d}t \leq K_p d^{-\alpha_p}\|T\|_{L_p}
\end{equation}
for every $d\in[n]$ and every
$T\in\mathcal{T}^{\geq d}(\mathbb{M}_{2^n})$.
\end{enumerate}
\end{lemma}

\begin{proof}
If $\tr(T)=0$, that is, $T\in\mathcal{T}^{\geq 1}(\M_{2^{n}})$.  For $0\leq t\leq1$,
contractivity gives $\|e^{-t\mathcal{L}}T\|_{L_p}\leq\|T\|_{L_p}$,
whereas for $t\geq1$, Theorem \ref{mn-k-tail} with $d=1$ gives
\[
\|e^{-t\mathcal{L}}T\|_{L_p}
 \leq C_p e^{-c_pt}\|T\|_{L_p}.
\]
Thus \eqref{mean-zero-lp-decay} holds, for example, with
$b_p=c_p$ and $B_p=\max\{e^{c_p},C_p\}$.

For $T\in\mathcal{T}^{\geq d}(\M_{2^{n}})$, combining $A_p\geq1$ with Theorem
\ref{mn-k-tail} yields
\begin{align*}
\int_0^\infty\|e^{-t\mathcal{L}}T\|_{L_p}\,\mathrm{d}t&\leq C_p\|T\|_{L_p}\left(\int_0^1 e^{-c_pdt^{A_p}}\,\mathrm{d}t+\int_1^\infty e^{-c_pdt}\,\mathrm{d}t\right)\\
&\leq C_p\|T\|_{L_p}\left(\frac{\Gamma(\alpha_p)}{(c_pd)^{\alpha_p}}+\frac{e^{-c_pd}}{c_pd}\right).
\end{align*}
Since $d^{-1}\leq d^{-\alpha_{p}}$ for $d\geq 1$ and $0<\alpha_{p}\leq 1$, it follows that there exists a constant $K_{p}>0$ depending only on $p$ such that
\[
\frac{\Gamma(\alpha_p)}{(c_pd)^{\alpha_p}}+\frac{e^{-c_pd}}{c_pd}\leq K_{p}d^{-\alpha_{p}},
\]
which yields the desired estimate.
\end{proof}

\begin{theorem}\label{ci-pq}
Let $d\in[n]$,  $2<p<\infty$, and let $q=p/(p-1)$.  With
$\alpha_p$ and $K_p$ as in Lemma \ref{etl-pp-theta}, one has
\begin{equation}\label{cov-poincare-pq}
\big|\mathrm{Cov}(T_1,T_2)\big|
 \leq K_p d^{-\alpha_p}
 \sum_{j=1}^n
 \big(\mathrm{Inf}_j^p(T_1)\big)^{1/p}
 \big(\mathrm{Inf}_j^q(T_2)\big)^{1/q}
\end{equation}
for all $T_1,T_2\in\mathcal{T}^{\geq d}(\M_{2^{n}})$.
\end{theorem}

\begin{proof}
For each $j\in[n]$, the operator $\dc_j(T_1)$ still belongs to $\mathcal{T}^{\geq d}(\M_{2^{n}})$.  Lemma \ref{cov-semigroup-0-infty}, H\"older's inequality, and \eqref{integrated-tail-decay} therefore give
\begin{align*}
\big|\mathrm{Cov}(T_1,T_2)\big|&\leq \sum_{j=1}^n\int_0^\infty \|e^{-t\mathcal{L}}\dc_j(T_1)\|_{L_p}\|\dc_j(T_2)\|_{L_q}\,\mathrm{d}t\\
&\leq K_{p}d^{-\alpha_p}\sum_{j=1}^n\|\dc_{j}(T_1)\|_{L_p}\|\dc_j(T_2)\|_{L_q}.
\end{align*}
The conclusion follows from the definition of the $L_{p}$-influences.
\end{proof}

In the Hilbertian case, we apply Parseval's identity to give the optimal dependence on the lower Fourier level.

\begin{corollary}\label{cov-poincare-l2-tail}
For every $d\in[n]$ and all $T_{1}$, $T_{2}\in\mathcal{T}^{\geq d}(\M_{2^{n}})$,
\begin{equation}\label{cov-poincare-l2}
\big|\mathrm{Cov}(T_{1},T_{2})\big|\leq\left(\frac{1}{d}\right)\sum_{j=1}^n\sqrt{\mathrm{Inf}_{j}(T_1)\mathrm{Inf}_{j}(T_2)}.
\end{equation}
In particular,
\[
\mathrm{Var}(T)\leq \left(\frac{1}{d}\right)\mathrm{Inf}(T),
\quad T\in\mathcal{T}^{\geq d}(\M_{2^{n}}).
\]
\end{corollary}

\begin{proof}
By Lemmas \ref{cov-semigroup-0-infty} and \ref{etl-t-l2-pd},
\begin{align*}
\big|\mathrm{Cov}(T_1,T_2)\big|
&\leq\sum_{j=1}^n\int_0^\infty
 \|e^{-t\mathcal{L}}\dc_j(T_1)\|_{L_2}
 \|\dc_j(T_2)\|_{L_2}\,\mathrm{d}t\\
&\leq\sum_{j=1}^n\int_0^\infty e^{-dt}\,\mathrm{d}t\,
 \|\dc_j(T_1)\|_{L_2}\|\dc_j(T_2)\|_{L_2},
\end{align*}
which is \eqref{cov-poincare-l2}.  Taking $T_1=T_2=T$ gives the
variance estimate.
\end{proof}

The covariance Poincar\'e inequalities stated previously control the
bilinear form by the sizes of the coordinate derivatives.  The purpose of the rest of the section is to refine those estimates into Talagrand-type inequalities invoking $L_{1}$ and $L_{p}$ influences. The next theorem is a covariance form of Talagrand's influence inequality.  It
may be viewed as a logarithmic refinement of the covariance Poincar\'e
estimate at level $d=1$.

\begin{theorem}\label{cov-taragrand-pq}
For every $1<p,q<\infty$ with $\frac{1}{p}+\frac{1}{q}=1$, there exist a constant $C_p$ such that for any $T_{1}$,  and $T_{1}\in \M_{2^{n}}$, we have
\begin{equation}\label{cov-talagrand-pq-estimate}
\begin{split}
\big|\mathrm{Cov}(T_{1},T_{2})\big|\leq C_{p}\sum_{j=1}^{n}\frac{\|\dc_{j}(T_{1})\|_{L_{p}}\|\dc_{j}(T_{2})\|_{L_{q}}}{\sqrt{\left(1+\log\left(\frac{\|\dc_{j}(T_{1})\|_{L_{p}}}
                                  {\|\dc_{j}(T_{1})\|_{L_{1}}}\right)\right)
       \left(1+\log\left(\frac{\|\dc_{j}(T_{2})\|_{L_{q}}}
                                  {\|\dc_{j}(T_{2})\|_{L_{1}}}\right)\right)}}
\end{split}
\end{equation}
\end{theorem}
\begin{proof}
By conjugate symmetry of covariance, after interchanging $T_{1}$ and $T_{2}$
if necessary, it suffices to consider $2\leq p<\infty$ and $1<q\leq2$.
For $j\in[n]$, set
\[
X_{j}=\dc_{j}(T_{1}), \qquad Y_{j}=\dc_{j}(T_{2}),
\]
and it is clear that both $X_{j}$ and $Y_{j}$ are trace zero.  Let $B_{p},b_{p}>0$ be the
constants in \eqref{mean-zero-lp-decay}, and put $a_{p}=2/b_{p}$.  Applying
Lemma \ref{cov-semigroup-0-infty}, changing variables according to
$t=(a_{p}+2)u$, and using self-adjointness of the semigroup, we obtain
\begin{align*}
\big|\mathrm{Cov}(T_{1},T_{2})\big|
&\leq (a_{p}+2)\sum_{j=1}^{n}\int_{0}^{\infty}
 \left|\tr\!\left(
  (e^{-a_{p}u\mathcal{L}}e^{-u\mathcal{L}}X_{j})^{*}
   e^{-u\mathcal{L}}Y_{j}
 \right)\right|\,\mathrm{d}u\\
&\leq (a_{p}+2)B_{p}\sum_{j=1}^{n}\int_{0}^{\infty}e^{-2u}
 \|e^{-u\mathcal{L}}X_{j}\|_{L_{p}}
 \|e^{-u\mathcal{L}}Y_{j}\|_{L_{q}}\,\mathrm{d}u.
\end{align*}
Indeed, the second inequality follows from H\"older's inequality and
\eqref{mean-zero-lp-decay}, applied to the trace-zero operator
$e^{-u\mathcal{L}}X_{j}$ for the time $a_{p}u$.

Define
\[
p(u)=1+(p-1)e^{-2u},
\qquad
q(u)=1+(q-1)e^{-2u}.
\]
Hypercontractivity gives
\[
\|e^{-u\mathcal{L}}X_{j}\|_{L_{p}}
 \leq \|X_{j}\|_{L_{p(u)}},
\qquad
\|e^{-u\mathcal{L}}Y_{j}\|_{L_{q}}
 \leq \|Y_{j}\|_{L_{q(u)}}.
\]
Interpolation between $L_{1}$ and $L_{p}$, and between $L_{1}$ and $L_{q}$,
respectively, yields
\begin{align*}
\|X_{j}\|_{L_{p(u)}}
&\leq \|X_{j}\|_{L_{1}}^{\theta_{p}(u)}
       \|X_{j}\|_{L_{p}}^{1-\theta_{p}(u)},\\
\|Y_{j}\|_{L_{q(u)}}
&\leq \|Y_{j}\|_{L_{1}}^{\theta_{q}(u)}
       \|Y_{j}\|_{L_{q}}^{1-\theta_{q}(u)},
\end{align*}
where
\[
\frac{1}{p(u)}=\theta_{p}(u)+\frac{1-\theta_{p}(u)}{p},
\qquad
\frac{1}{q(u)}=\theta_{q}(u)+\frac{1-\theta_{q}(u)}{q}.
\]
Fix a summand for which $X_{j}$ and $Y_{j}$ are nonzero, and write
\[
x_{j}=\frac{\|X_{j}\|_{L_{1}}}{\|X_{j}\|_{L_{p}}},
\qquad
y_{j}=\frac{\|Y_{j}\|_{L_{1}}}{\|Y_{j}\|_{L_{q}}}.
\]
Since the trace is normalized, $0<x_{j},y_{j}\leq1$.  The preceding estimates
and the Cauchy--Schwarz inequality reduce the remaining integral to
$A_{j}K_{j}$, where
\[
A_{j}^{2}=\int_{0}^{\infty}e^{-2u}x_{j}^{2\theta_{p}(u)}\,\mathrm{d}u,
\qquad
K_{j}^{2}=\int_{0}^{\infty}e^{-2u}y_{j}^{2\theta_{q}(u)}\,\mathrm{d}u.
\]
Changing variables from $u$ to $\nu=p(u)$ gives
\[
A_{j}^{2}
 =\frac{1}{2(p-1)}\int_{1}^{p}x_{j}^{2\theta_{p}(\nu)}\,\mathrm{d}\nu,
\qquad
\theta_{p}(\nu)=\frac{p-\nu}{\nu(p-1)}.
\]
A second change of variables, $s=2\theta_{p}(\nu)$, yields
\[
A_{j}^{2}\leq\frac{p}{4}\int_{0}^{2}x_{j}^{s}\,\mathrm{d}s \leq\left(\frac{p}{2}\right)\frac{1}{1+\log(1/x_{j})}.
\]
Here the last inequality follows from the elementary estimate
\[
\int_{0}^{2}x^{s}~\mathrm{d}s=\frac{1-x^{2}}{\log(1/x)}\leq\frac{2}{1+\log(1/x)}, \quad 0<x\leq1,
\]
with the value at $x=1$ understood by continuity.  The same argument gives
\[
K_{j}^{2}\leq\left(\frac{q}{2}\right)\frac{1}{1+\log(1/y_{j})}.
\]
Substitution proves \eqref{cov-talagrand-pq-estimate}, for example with
\[
C_{p}=\frac{(a_{p}+2)B_{p}\sqrt{pq}}{2}.
\]
%For any $T_1\in L_p(\mathbb{M}_{2^n})$  and $T_2\in L_q(\mathbb{M}_{2^n})$, by H\"{o}lder inequality, Lemma \ref{etl-pp-theta} and changing the variable $t=(p+2)u$,  we have
%\begin{align*}
%\mathrm{Cov}(T_1,T_2)&=\sum_{j=1}^n\int_{0}^\infty\mathrm{tr}(e^{-t\mathcal{L}}d_j(T_1)\cdot d_j(T_2))dt\\
%&= (p+2)\sum_{j=1}^n\int_{0}^\infty \mathrm{tr}(e^{-pu\mathcal{L}}e^{-u\mathcal{L}}d_j(T_1)e^{-u\mathcal{L}}d_j(T_2))du\\
%&\leq (p+2)\sum_{j=1}^n\int_{0}^\infty\|e^{-pu\mathcal{L}}e^{-u\mathcal{L}}d_j(T_1)\|_{L_{p}}\|e^{-u\mathcal{L}}d_j(T_2)\|_{L_{q}}du\\
%&\leq(p+2)\sum_{j=1}^n\int_{0}^\infty e^{-pu\frac{2}{p}}\|e^{-u\mathcal{L}}d_j(T_1)\|_{L_{p}}\|e^{-u\mathcal{L}}d_j(T_2)\|_{L_{q}}du.
%\end{align*}
%Applying Theorem \ref{hc-1}  with $2<p<\infty$ and $p(u)=1+e^{-2 u}(p-1)$, $1<q\leq2$ and $q(u)=1+e^{-2 u}(q-1)$, and using the H\"{o}lder inequality $\frac{1}{p(u)}=\frac{\theta(p(u))}{1}+\frac{1-\theta(p(u))}{p}$ and $\frac{1}{q(u)}=\frac{\theta(q(u))}{1}+\frac{1-\theta(q(u))}{p}$,  we obtain
%\begin{align*}
%\mathrm{Cov}(T_1,T_2)&\leq(p+2)\sum_{j=1}^n\int_{0}^\infty e^{-2u}\|d_j(T_1)\|_{L_{p(u)}}\|d_j(T_2)\|_{L_{q(u)}}du\\
%&\leq(p+2)\sum_{j=1}^n\int_{0}^\infty e^{-2u}\|d_j(T_1)\|_{L_{1}}^{\theta(p(u))}\|d_j(T_1)\|_{L_{p}}^{1-\theta(p(u))}\\
%&\quad \cdot\|d_j(T_2)\|_{L_{1}}^{\theta(q(u))}\|d_j(T_2)\|_{L_{q}}^{1-\theta(q(u))}du\\
%&=(p+2)\sum_{j=1}^n\|d_j(T_1)\|_{L_{p}}\|d_j(T_2)\|_{L_{q}}\int_{0}^\infty e^{-2u}\big(\frac{\|d_j(T_1)\|_{L_{1}}}{\|d_j(T_1)\|_{L_{p}}}\big)^{\theta(p(u))}\\ &\quad\cdot\big(\frac{\|d_j(T_2)\|_{L_{1}}}{\|d_j(T_2)\|_{L_{q}}}\big)^{\theta(q(u))}du.
%\end{align*}
%By Cauchy-Schwartz inequality,  we have
%\begin{equation}\label{cov-ab}
%\mathrm{Cov}(T_1,T_2)\leq(p+2)\sum_{j=1}^n\|d_j(T_1)\|_{L_{p}}\|d_j(T_2)\|_{L_{q}}\cdot AB.
%\end{equation}
%where 
%$$A=\big(\int_{0}^\infty e^{-2u}\big(\frac{\|d_j(T_1)\|_{L_{1}}}{\|d_j(T_1)\|_{L_{p}}}\big)^{2\theta(p(u))}\big)^{\frac{1}{2}}$$
%and 
%$$B=\big(\int_{0}^\infty e^{-2u}\big(\frac{\|d_j(T_2)\|_{L_{1}}}{\|d_j(T_2)\|_{L_{q}}}\big)^{2\theta(q(u))}\big)^{\frac{1}{2}}.$$
%Similar to the treatment in the preceding Talagrand estimates, changing the variable with $p(u)=\nu$ and $s=2\theta(\nu)$, 
%\begin{equation}\label{a}
%\begin{aligned}
%A&=\big(\frac{1}{2(p-1)}\int_1^pa_j^{2\theta(\nu)}d\nu\big)^{\frac{1}{2}}
%\leq \big(\frac{p}{4}\int_0^2a_j^sds\big)^{\frac{1}{2}}\\
%&=\big(\frac{p}{4}\frac{a_j^2-1}{\log a_j}\big)^{\frac{1}{2}}
%\leq \big(\frac{p}{4}\frac{2}{1+ \log\frac{1}{a_j}}\big)^{\frac{1}{2}}=\big(\frac{p}{2}\frac{1}{1+ \log\frac{1}{a_j}}\big)^{\frac{1}{2}},
%\end{aligned}
%\end{equation}
%where $a_j=\frac{\|d_j(T_1)\|_{L_{1}}}{\|d_j(T_1)\|_{L_{p}}}$, and the last inequality is due to the fact $\frac{a^2-1}{\log a}\leq \frac{2}{1+\log \frac{1}{a}}$ for $a\in(0,2)$. Similarly, we obtain
%$$B\leq\big(\frac{q}{2}\frac{1}{1+ \log\frac{1}{b_j}}\big)^{\frac{1}{2}},$$
%where $b_j=\frac{\|d_j(T_2)\|_{L_{1}}}{\|d_j(T_2)\|_{L_{q}}}$. Combining \eqref{cov-ab}, \eqref{a} and the above inequality,  it follows that
%$$\mathrm{Cov}(T_1,T_2)\leq\frac{(p+2)\sqrt{pq}}{2}\sum_{j=1}^n\frac{\|d_j(T_1)\|_{L_{p}}\|d_j(T_2)\|_{L_{q}}}{\sqrt{\big(1+\log (\frac{\|d_j(T_1)\|_{L_{p}}}{\|d_j(T_1)\|_{L_{1}}})\big)\big(1+\log (\frac{\|d_j(T_2)\|_{L_{q}}}{\|d_j(T_2)\|_{L_{1}}})\big)}}.$$
%The proof is complete.
\end{proof}

The preceding theorem does not yet record the location of the Fourier support.
In $L_{2}$, the exact spectral decay on $\mathcal{T}^{\geq d}(\M_{2^{n}})$ can be inserted
before the hypercontractive step. The next theorem provides a \emph{covariance Talagrand $L_{1}$-$L_{2}$ inequality for tail spaces}, which captures the decay property depending the lower Fourier level.

\begin{theorem}\label{cov-tala-d-tail-l2}
%There exists a universal constant $K>0$ such that 
For each $n\in \mathbb{N}$, $d\in [n]$, and $T_{1}$, $T_{2}\in \mathbb{M}_{2^{n}}$, we have
\begin{equation}\label{cov-tail-l2-estimate}
\big|\mathrm{Cov}(T_{1},T_{2})\big|\leq4\sum_{j=1}^{n}
\frac{\|\dc_{j}(T_{1})\|_{L_{2}}\|\dc_{j}(T_{2})\|_{L_{2}}} {\sqrt{\left(d+\log\left(\frac{\|\dc_{j}(T_{1})\|_{L_{2}}}{\|\dc_{j}(T_{1})\|_{L_{1}}}\right)\right)\left(d+\log\left(\frac{\|\dc_{j}(T_{2})\|_{L_{2}}}{\|\dc_{j}(T_{2})\|_{L_{1}}}\right)\right)}}
\end{equation}
\end{theorem}

\begin{proof}
Set $X_{j}=\dc_{j}(T_{1})$ and $Y_{j}=\dc_{j}(T_{2})$. By Lemma
\ref{cov-semigroup-0-infty}, the change of variables $t=4u$, and
self-adjointness of the semigroup give
\begin{align*}
\big|\mathrm{Cov}(T_{1},T_{2})\big|
&\leq4\sum_{j=1}^{n}\int_{0}^{\infty}
 \left|\tr\!\left((e^{-3u\mathcal{L}}X_{j})^{*}
                         e^{-u\mathcal{L}}Y_{j}\right)\right|\,\mathrm{d}u\\
&\leq4\sum_{j=1}^{n}\int_{0}^{\infty}e^{-2du}
 \|e^{-u\mathcal{L}}X_{j}\|_{L_{2}}
 \|e^{-u\mathcal{L}}Y_{j}\|_{L_{2}}\,\mathrm{d}u.
\end{align*}
In the second line we used Lemma \ref{etl-t-l2-pd} in the form
\[
\|e^{-3u\mathcal{L}}X_{j}\|_{L_{2}}
 =\|e^{-2u\mathcal{L}}(e^{-u\mathcal{L}}X_{j})\|_{L_{2}}
 \leq e^{-2du}\|e^{-u\mathcal{L}}X_{j}\|_{L_{2}}.
\]
Hypercontractivity and interpolation between $L_{1}(\M_{2^{n}})$ and $L_{2}(\M_{2^{n}})$ imply
\[
\|e^{-u\mathcal{L}}X_{j}\|_{L_{2}}
 \leq\|X_{j}\|_{L_{1+e^{-2u}}}
 \leq\|X_{j}\|_{L_{2}}
       \left(\frac{\|X_{j}\|_{L_{1}}}{\|X_{j}\|_{L_{2}}}\right)^{\tanh(u)},
\]
and the same estimate holds for $Y_{j}$.  For a nonzero derivative, put
\[
\lambda_{1,j}=\log\left(\frac{\|X_{j}\|_{L_{2}}}{\|X_{j}\|_{L_{1}}}\right),
\qquad
\lambda_{2,j}=\log\left(\frac{\|Y_{j}\|_{L_{2}}}{\|Y_{j}\|_{L_{1}}}\right).
\]
The Cauchy--Schwarz inequality therefore gives
\begin{align*}
\big|\mathrm{Cov}(T_{1},T_{2})\big|
&\leq4\sum_{j=1}^{n}\|X_{j}\|_{L_{2}}\|Y_{j}\|_{L_{2}}
 I_{1,j}^{1/2}I_{2,j}^{1/2},
\end{align*}
where
\[
I_{i,j}=\int_{0}^{\infty}
 e^{-2du-2\lambda_{i,j}\tanh(u)}\,\mathrm{d}u,
\qquad i\in\{1,2\}.
\]
Recall that Lemma \ref{tail-kernel-estimate} yields
\[
I_{i,j}\leq\frac{1}{d+2\lambda_{i,j}}
 \leq\frac{1}{d+\lambda_{i,j}}.
\]
Substituting these bounds yields \eqref{cov-tail-l2-estimate}.
%Note  that
%\begin{align*}
%\mathrm{Cov}(T_{1},T_{2})&=\sum_{j=1}^n\int_{0}^\infty\mathrm{tr}((e^{-t\mathcal{L}}d_j(T_1))d_j(T_2))dt\\
%%\langle T_{1},f\rangle&=\langle f, P_{0}(f)-P_{\infty}(f)\rangle=\int_{0}^{\infty}\langle f,\Delta e^{-t\Delta}(f)\rangle dt\\
%&=4\sum_{j=1}^n\int_{0}^\infty\mathrm{tr}(e^{-3t\mathcal{L}}d_j(T_1)\cdot e^{-t\mathcal{L}}d_j(T_2))dt\\
%&\leq 4\sum_{j=1}^{n}\int_{0}^{\infty}\|e^{-3t\mathcal{L}}d_j(T_1)\|_{L_{2}}\|e^{-t\mathcal{L}}d_j(T_2)\|_{L_{2}}dt\\
%&\leq4\sum_{j=1}^{n}\int_{0}^{\infty}e^{-2td}\|e^{-t\mathcal{L}}d_j(T_1)\|_{L_{2}}\|e^{-t\mathcal{L}}d_j(T_2)\|_{L_{2}}dt\\
%&\leq 4\sum_{j=1}^{n}\int_{0}^{\infty}e^{-2td}\|d_j(T_1)\|_{L_{1+e^{-2t}}}\|d_j(T_2)\|_{L_{1+e^{-2t}}}dt\\
%&\leq 4\sum_{j=1}^{n}\|d_j(T_1)\|_{L_{2}}\|d_j(T_2)\|_{L_{2}}\int_{0}^{\infty}e^{-2td}\left(\frac{\|d_{j}(T_1)\|_{L_{1}}}{\|d_{j}(T_1)\|_{L_{2}}}\right)^{\frac{1-e^{-2t}}{1+e^{-2t}}}\left(\frac{\|d_{j}(T_2)\|_{L_{1}}}{\|d_{j}(T_2)\|_{L_{2}}}\right)^{\frac{1-e^{-2t}}{1+e^{-2t}}}dt,
%\end{align*}
%where we used the H\"older inequality in the last line with parameters
%\[
%\frac{1}{1+e^{-2t}}=\frac{1-e^{-2t}}{1+e^{-2t}}+\frac{e^{-2t}}{1+e^{-2t}},\quad t\geq 0.
%\]
%Let $b_1\coloneqq \frac{\|d_{j}(T_1)\|_{L_{1}}}{\|d_{j}(T_1)\|_{L_{2}}}, b_2\coloneqq \frac{\|d_{j}(T_2)\|_{L_{1}}}{\|d_{j}(T_2)\|_{L_{2}}}\in (0,1)$. By the Cauchy-Schwarz inequality, we have
%\begin{align}\label{semigroup-approach}
%\mathrm{Cov}(T_{1},T_{2})&\leq4\sum_{j=1}^{n}\big(\int_{0}^{\infty}e^{-2td}b_1^{2(\frac{1-e^{-2t}}{1+e^{-2t}})}dt\big)^{\frac{1}{2}}\big(\int_{0}^{\infty}e^{-2td}b_2^{2(\frac{1-e^{-2t}}{1+e^{-2t}})}dt\big)^{\frac{1}{2}}\\
%&=:4\sum_{j=1}^{n}A^{\frac{1}{2}}B^{\frac{1}{2}}.\nonumber
%\end{align}
%Since A and B are handled similarly, it now remains to estimate the following term
%\begin{equation}\label{key-term-1}
%A=\int_{0}^{\infty}e^{-2td}b_1^{2(\frac{1-e^{-2t}}{1+e^{-2t}})}dt.
%\end{equation}
%Changing the variable $1+e^{-2t}\mapsto u$ yields that
%\begin{equation}\label{key-term-2}
%A=\frac{1}{2}\int_{1}^{2}(u-1)^{d-1}b_1^{2(\frac{2-u}{u})}du.
%\end{equation}
%Changing variable $2(\frac{2-u}{u})\mapsto t$ yields that
%\begin{equation}\label{key-term-3}
%A=\frac{1}{2\times4}\int_{0}^{2}(\frac{4}{2+t})^2\frac{(2-t)^{d-1}}{(2+t)^{d-1}}\cdot b_1^{t}dt=2\int_{0}^{2}\frac{(2-t)^{d-1}}{(2+t)^{d+1}}\cdot b_1^{t}dt.
%\end{equation}
%Changing the variable $t\mapsto \frac{2\nu}{d}$ yields that
%\begin{equation}\label{key-term-4}
%\begin{split}
%A&=\frac{2\times2}{d}\int_{0}^{d}\frac{2^{-2}d^{2}(d-\nu)^{d-1}}{(d+\nu)^{d+1}}\cdot b_{1}^{2\nu/d} d\nu\\
%&\leq \frac{1}{d}\int_{0}^{d}b_{1}^{2\nu/d}d\nu,
%\end{split}
%\end{equation}
%the last inequality is due to  the decreasing of the function $\nu\mapsto \frac{d^{2}(d-\nu)^{d-1}}{(d+\nu)^{d+1}}$ for $\nu\in[1,d]$.
%And $b\in (0,1)$, we have that
%\begin{align*}
%A&\leq\frac{1}{d}\int_{0}^{d}b_{1}^{2\nu/d}d\nu\\
%&=\left(\frac{1}{d}\right)\frac{b_1^{2}-1}{\log\left(b_1^{2/d}\right)}\\
%&\leq\left(\frac{1}{d}\right)\frac{b_1^{2/d}-1}{\log\left(b_1^{2/d}\right)}=\left(\frac{1}{2d}\right)\frac{b_1^{2/d}-1}{\log\left(b_1^{1/d}\right)}.
%\end{align*}
%Thanks to the numerical inequality $\frac{t^2-1}{\log(t)}\leq \frac{2}{1-\log(t)}$, we have
%\begin{equation}\label{key-term-9}
%A\leq \left(\frac{1}{d}\right)\frac{1}{1+\log\left(b_1^{-1/d}\right)}.
%\end{equation}
%Similarly, we obtain
%\begin{equation}\label{key-term-9-b}
%B\leq \left(\frac{1}{d}\right)\frac{1}{1+\log\left(b_2^{-1/d}\right)}.
%\end{equation}
%Combing \eqref{semigroup-approach}, \eqref{key-term-9}  and \eqref{key-term-9-b} yields that for each $T_1,T_2\in \mathcal{T}^{\geq d}$
%\begin{equation}\label{key-term-10}
%\begin{split}
%\mathrm{Cov}(T_{1},T_{2})&\leq \left(\frac{4}{d}\right)\sum_{j=1}^{n}\frac{\|d_{j}(T_1)\|_{L_{2}}\|d_{j}(T_2)\|_{L_{2}}}{\sqrt{\Big(1+\frac{1}{d}\log\big(\frac{\|d_{j}(T_1)\|_{L_{2}}}{\|d_{j}(T_1)\|_{L_{1}}}\big)\Big)\Big(1+\frac{1}{d}\log\big(\frac{\|d_{j}(T_2)\|_{L_{2}}}{\|d_{j}(T_2)\|_{L_{1}}}\big)\Big)}}\\
%&=4\sum_{j=1}^{n}\frac{\|d_{j}(T_1)\|_{L_{2}}\|d_{j}(T_2)\|_{L_{2}}}{\sqrt{\Big(d+\log\big(\frac{\|d_{j}(T_1)\|_{L_{2}}}{\|d_{j}(T_1)\|_{L_{1}}}\big)\Big)\Big(d+\log\big(\frac{\|d_{j}(T_2)\|_{L_{2}}}{\|d_{j}(T_2)\|_{L_{1}}}\big)\Big)}},
%\end{split}
%\end{equation}
%which is the desired inequality.
\end{proof}

Taking the two observables to be equal gives the following tail Talagrand $L_{1}$-$L_{2}$ inequality.
\begin{corollary}\label{tala-d-tail}
For each $d\in [n]$ and $T\in \M_{2^{n}}$, we have
\[
\|T\|^{2}_{L_{2}}\leq 4\sum_{j=1}^{n}\frac{\|\dc_{j}(T)\|^{2}_{L_{2}}}{d+\log\left(\|\dc_{j}(T)\|_{L_{2}}/\|\dc_{j}(T)\|_{L_{1}}\right)}.
\]
\end{corollary}

We conclude this section with the following \emph{covariance Talagrand $L_{1}$-$L_{p}$ inequality for tail spaces}, which is the most general covariance inequality established in this chapter.
\begin{theorem}\label{cov-tala-d-tail-pq}
Let $1<p<\infty$, let $q=p/(p-1)$, and set
$r=\max\{p,q\}$.  Let $A_{r}\geq1$ be the exponent in Theorem
\ref{mn-k-tail}, and put $\alpha_{r}=A_{r}^{-1}$.  There exists a constant
$C_{p}>0$, depending only on $p$, such that
\begin{equation*}\label{cov-tail-pq-estimate}
\big|\mathrm{Cov}(T_{1},T_{2})\big|\leq C_{p}\sum_{j=1}^{n}
\frac{\|\dc_{j}(T_{1})\|_{L_{p}}\|\dc_{j}(T_{2})\|_{L_{q}}}
{\sqrt{\left(d^{\alpha_{r}}+\log\left(\frac{\|\dc_{j}(T_{1})\|_{L_{p}}}{\|\dc_{j}(T_{1})\|_{L_{1}}}\right)\right)\left(d^{\alpha_{r}}+\log\left(\frac{\|\dc_{j}(T_{2})\|_{L_{q}}}{\|\dc_{j}(T_{2})\|_{L_{1}}}\right)\right)}}
\end{equation*}
for all $d\in[n]$ and all $T_{1}$, $T_{2}\in\mathcal{T}^{\geq d}(\mathbb{M}_{2^{n}})$.
\end{theorem}

\begin{proof}
By the symmetry of conjugation, it is enough to treat the case $p=r\geq2$ and $1<q\leq2$.  Let $C_{p}^{(0)},c_{p}>0$ and $A_{p}\geq1$ be constants for
which Theorem \ref{mn-k-tail} holds.  Set
$X_{j}=\dc_{j}(T_{1})$ and $Y_{j}=\dc_{j}(T_{2})$.  After the change of
variables $t=3u$, Lemma \ref{cov-semigroup-0-infty} and self-adjointness of the
semigroup yield
\begin{align*}
\big|\mathrm{Cov}(T_{1},T_{2})\big|
&\leq3\sum_{j=1}^{n}\int_{0}^{\infty}
 \|e^{-2u\mathcal{L}}X_{j}\|_{L_{p}}
 \|e^{-u\mathcal{L}}Y_{j}\|_{L_{q}}\,\mathrm{d}u.
\end{align*}
Since $e^{-u\mathcal{L}}(X_{j})\in\mathcal{T}^{\geq d}$, Theorem
\ref{mn-k-tail}, applied for the time $u$, gives
\[
\|e^{-2u\mathcal{L}}X_{j}\|_{L_{p}}
\leq C_{p}^{(0)}
 e^{-c_{p}d\min\{u,u^{A_{p}}\}}
 \|e^{-u\mathcal{L}}X_{j}\|_{L_{p}}.
\]
Consequently,
\begin{align}\label{cov-tail-pq-kernel-reduction}
\big|\mathrm{Cov}(T_{1},T_{2})\big|
&\leq C_{p}^{(1)}\sum_{j=1}^{n}
 \|X_{j}\|_{L_{p}}\|Y_{j}\|_{L_{q}}
 \int_{0}^{\infty}K_{d}(u)
 x_{j}^{\theta_{p}(u)}y_{j}^{\theta_{q}(u)}\,\mathrm{d}u,
\end{align}
where
\[
K_{d}(u)=e^{-c_{p}d\min\{u,u^{A_{p}}\}},
\qquad
x_{j}=\frac{\|X_{j}\|_{L_{1}}}{\|X_{j}\|_{L_{p}}},
\qquad
y_{j}=\frac{\|Y_{j}\|_{L_{1}}}{\|Y_{j}\|_{L_{q}}},
\]
and
\[
\theta_{s}(u)=\frac{1-e^{-2u}}{1+(s-1)e^{-2u}},
\qquad s\in\{p,q\}.
\]
Here we used the same hypercontractive and interpolation estimates as in the
proof of Theorem \ref{cov-taragrand-pq}.

We now claim the following estimate needed to bound the integral in
\eqref{cov-tail-pq-kernel-reduction}.  Fix $1<s<\infty$, $A\geq1$, and
$c>0$.  There exists $D_{s,A,c}>0$ such that, for every $d\geq1$ and
$\lambda\geq0$,
\begin{equation}\label{tail-logarithmic-kernel-pq}
\int_{0}^{\infty}
 e^{-cd\min\{u,u^{A}\}-2\lambda\theta_{s}(u)}\,\mathrm{d}u
\leq\frac{D_{s,A,c}}{d^{1/A}+\lambda}.
\end{equation}
Indeed, with $\kappa_{s}=(1-e^{-2})/s$, one has
\[
\theta_{s}(u)\geq\kappa_{s}\min\{u,1\},
\qquad u\geq0.
\]
Writing $R=d^{1/A}$ and splitting the integral at $u=1$, we obtain
\begin{align*}
\int_{0}^{1}e^{-cdu^{A}-2\kappa_{s}\lambda u}\,\mathrm{d}u
&\leq\min\left\{C_{A,c}R^{-1},(2\kappa_{s}\lambda)^{-1}\right\}
 \leq\frac{C_{s,A,c}}{R+\lambda},
\end{align*}
where the second term in the minimum is omitted when $\lambda=0$.
Moreover, by
\[
\int_{1}^{\infty}e^{-cdu-2\kappa_{s}\lambda}\,\mathrm{d}u \leq\frac{C_{s,c}}{d(1+\lambda)} \leq\frac{C_{s,c}}{R+\lambda}.
\]
the desired inequality \eqref{tail-logarithmic-kernel-pq} follows.

Put
\[
\lambda_{1,j}\coloneqq \log(1/x_{j}),\qquad \lambda_{2,j}\coloneqq\log(1/y_{j}).
\]
The Cauchy--Schwarz inequality and \eqref{tail-logarithmic-kernel-pq}, first
with $s=p$ and then with $s=q$, give
\begin{align*}
&\int_{0}^{\infty}K_{d}(u)
 x_{j}^{\theta_{p}(u)}y_{j}^{\theta_{q}(u)}\,\mathrm{d}u\\
\leq &\left(\int_{0}^{\infty}K_{d}(u)x_{j}^{2\theta_{p}(u)}\,\mathrm{d}u\right)^{1/2}
\left(\int_{0}^{\infty}K_{d}(u)y_{j}^{2\theta_{q}(u)}\,\mathrm{d}u\right)^{1/2}\\
\leq&\frac{C_{p}}
{\sqrt{\left(d^{\alpha_{p}}+\lambda_{1,j}\right)
       \left(d^{\alpha_{p}}+\lambda_{2,j}\right)}}.
\end{align*}
Substituting this estimate into \eqref{cov-tail-pq-kernel-reduction} proves
the desired inequality.  If $p<2$, the same argument is applied after
interchanging the two observables; then $r=q$ and the exponent is
$\alpha_{r}=A_{r}^{-1}$.
\end{proof}

\begin{remark}
Theorems \ref{cov-tala-d-tail-l2} and \ref{cov-tala-d-tail-pq} distinguish two
features of tail covariance estimates.  In $L_{2}$, spectral orthogonality
produces the exact lower-level parameter $d$; for general conjugate exponents,
the verified dimension-free gain is governed by the heat-smoothing scale
$d^{\alpha_{r}}$.
\end{remark}







%\section{The Talagrand influence inequality in quantum  functions}
%
%{\color{red}It seems that we will delete all the section!}
%\begin{theorem}\label{talagrand-d-tail} Let $T \in \mathbb{M}_{2^n}$ be in the $d$-th tail space with $\mathrm{tr}(T)$. Then we have
%$$
%\|T\|_{L_{2}}^{2} \leq \frac{(1+d)^{\frac{1}{d}+1}}{d} \sum_{j=1}^{n} \frac{\|d_j(T)\|_{L_{2}}^{2}}{1+\log (\|d_j(T)\|_{L_{2}}/\|d_j(T)\|_{L_{1}})}.
%$$
%\end{theorem}
%\begin{proof}  For every $t \geq 0$ and $d\in[n]$, and $T \in \mathbb{M}_{n}$ is in the $d$-th tail space. 
%\begin{align*}
%\|T\|_{L_2}^{2}&=\|e^{-t \mathcal{L}}(T)\|_{L_2}^{2}+\|T\|_{L_{2}(\mathbb{M}_{2^n})}^{2}-\|e^{-t \mathcal{L}}(T)\|_{L_2}^{2}\\
%&=\sum_{\substack{\s\in\{0,1,2,3\}^n \\ |\s|\geq d} }e^{-2t|\s|}|\widehat{T}(\s)|^2+\|T\|_{L_2}^{2}-\|e^{-t \mathcal{L}}(T)\|_{L_2}^{2}\\
%&\leq e^{-2td} \|T\|_{L_2}^{2}+\|T\|_{L_2}^{2}-\|e^{-t \mathcal{L}}(T)\|_{L_2}^{2}.
%\end{align*}
%Then
%\begin{equation}\label{t2-etd}
%\|T\|_{L_2}^2 \leq \frac{1}{1-e^{-2 td}}(\|T\|_{L_2}^{2}-\|e^{-t \mathcal{L}}(T)\|_{L_2}^{2}). 
%\end{equation}
%Similar to the proof process of \cite[Theorem 1.10 (i)]{JLL2025}, we can obtain the desired estimate. For the convenience of the readers, we provide the detailed proof process here. By Fourier expansion of such $T \in \mathbb{M}_{2^n}$, it follows that
%
%\begin{align*}
%\|T\|_{L_2}^{2}-\| e^{-t \mathcal{L}}(T) \|_{L_2}^{2} & =\sum_{\substack{\s \in\{0,1,2,3\}^n\\ |\s|\geq d}}(1-e^{-2 t|\s|})|\widehat{T}(\s)|^{2} \\
%&=-\sum_{\substack{\s \in\{0,1,2,3\}^n\\ |\s|\geq d}}\int_{0}^{t} \frac{\partial}{\partial y} e^{-2 y|\s|} |\widehat{T}(\s)|^{2}dy \\
%&=2 \int_{0}^{t} \sum_{\substack{\s \in\{0,1,2,3\}^n\\ |\s|\geq d}}|\s| e^{-2 y|\s|}|\widehat{T}(\s)|^{2} dy\\
%&=2 \int_{0}^{t} \sum_{j=1}^n\sum_{\substack{\s \in\{0,1,2,3\}^n\\ \s_j\neq 0\\ |\s|\geq d}} e^{-2 y|\s|}|\widehat{T}(\s)|^{2} dy\\
%&=2 \int_{0}^{t} \sum_{j=1}^n\|d_je^{-t \mathcal{L}}(T)\|_{L_2}^2dt.
%\end{align*}
%Combining the above equation \eqref{t2-etd},  $d_je^{-t \mathcal{L}}=e^{-t \mathcal{L}}d_j$, the hypercontractivity \eqref{hc-1} with $q=2$ and $p(y)=e^{-2y}+1$ together, we get
%\begin{align}\label{t2-etd-integral}
%\|T\|_{L_2}^{2} &\leq \frac{2}{1-e^{-2 td}} \int_{0}^{t} \sum_{j=1}^{n}\|e^{-y\mathcal{L}} d_j(T)\|_{L_2}^{2} d y\\
%&\leq \frac{2}{1-e^{-2 td}} \int_{0}^{t} \sum_{j=1}^{n}\|d_j(T)\|_{L_{p(y)}}^{2} dy\nonumber, 
%\end{align}
%where $1<p(y)\leq2$. 
%Changing the variable with $p(y)=\nu$ entails that
%$$
%\|T\|_{L_2}^{2} \leq \min _{t\geq0}\{\frac{e^{2 t}}{1-e^{-2 td}}\} \int_{1}^{2} \sum_{j=1}^{n}\left\|d_j(T)\right\|_{L_\nu}^{2} d \nu.
%$$
%It is clear that
%\begin{equation}\label{min-t}
%\min _{t\geq0}\{\frac{e^{2 t}}{1-e^{-2 td}}\}=\frac{(1+d)^{\frac{1}{d}+1}}{d}\approx \mathcal{O}(1),\quad \mbox{as $d\rightarrow\infty$}.
%\end{equation}
%Let $\theta(\nu)=\frac{2-\nu}{\nu}$ for $\nu \in[1,2]$. Then $\frac{1}{\nu}=\frac{\theta(\nu)}{1}+\frac{1-\theta(\nu)}{2}$. 
%Applying the H\"{o}lder inequality, we have
%\begin{align*}
%\|T\|_{L_2}^{2} & \leq\frac{(1+d)^{\frac{1}{d}+1}}{d} \int_{1}^{2} \sum_{j=1}^{n}\|d_j(T)\|_{L_1}^{2 \theta(\nu)} \cdot\|d_j(T)\|_{L_2}^{2(1-\theta(\nu))} d \mu \\
%& =\frac{(1+d)^{\frac{1}{d}+1}}{d} \sum_{j=1}^{n}\|d_j(T)\|_{L_2}^{2} \int_{1}^{2}(\frac{\|d_j(T)\|_{L_1}}{\|d_j(T)\|_{L_{2}}})^{2 \theta(\nu)} d \nu.
%\end{align*}
% Changing variables $y=2\theta(\nu)$ again, let $b_{j}=\|d_j(T)\|_{L_1}/\|d_j(T)\|_{L_{2}}$. we can get the further estimate 
%\begin{align*}
%\|T\|_{L_2} & \leq \frac{(1+d)^{\frac{1}{d}+1}}{d} \sum_{j=1}^{n}\|d_j(T)\|_{L_{2}}^{2} \int_{1}^{2} b_{j}^{2 \theta(\nu)} d \nu \\
%& \leq\frac{(1+d)^{\frac{1}{d}+1}}{d} \sum_{j=1}^{n}\|d_j(T)\|_{L_{2}}^{2} \int_{0}^{2} b_{j}^{s} d s \\
%& =\frac{(1+d)^{\frac{1}{d}+1}}{d} \sum_{j=1}^{n}\|d_j(T)\|_{L_{2}}^{2} \cdot \frac{b_{j}^{2}-1}{\log \left(b_{j}\right)} \\
%& \leq 2\frac{(1+d)^{\frac{1}{d}+1}}{d} \sum_{j=1}^{n} \frac{\|d_j(T)\|_{L_{2}}^{2}}{1+\log \left(\|d_j(T)\|_{L_{2}(\mathbb{M})} /\|d_j(T)\|_{L_{1}}\right)},
%\end{align*}
%where we used the fact $\frac{b^{2}-1}{\log (b)} \leq \frac{2}{1+\log (1 / b)}$ for $b \in(0,2)$ in the last inequality. This completes our proof.
%\end{proof}
%
%\begin{remark}\label{covd-talad} In Theorem \ref{ci-d-tail}, we observe that the estimate for the covariance inequality of the two quantum Boolean function versions related to $d$-tail spaces $\frac{1}{d}\rightarrow0$, when $d\rightarrow\infty$.
%However, when applying a similar approach to Talagrand's method for quantum Boolean functions in $d$ tail spaces, the constant obtained \eqref{min-t} is $\frac{(1+d)^{\frac{1}{d}+1}}{d}\approx \mathcal{O}(1)$. We know that without considering $d$-tail spaces, Talagrand's inequality implicitly implies the Poincar\'{e} inequality. However, the results related to $d$-tail probability spaces are significantly different.
%\end{remark}
%In order to obtain constant estimates similar to those of the Poincar\'{e} inequality, we  establish the following  result analogous to \cite[Theorem 1.8]{JLZ2025}.
%\begin{theorem}\label{talagrand-d-tail-improve} Let $T \in \mathbb{M}_{2^n}$ be in the $d$-th tail space with $\mathrm{tr}(T)$. Then we have
%$$
%\|T\|_{L_{2}}^{2} \leq \frac{(1+2d)^{\frac{1}{2d}+1}}{d\sqrt{2d+1}} \sum_{j=1}^{n} \frac{\|d_j(T)\|_{L_{2}}^{2}}{\sqrt{1+2\log (\|d_j(T)\|_{L_{2}}/\|d_j(T)\|_{L_{1}})}}.
%$$
%\end{theorem}
%\begin{proof} Let $T \in \mathbb{M}_{2^n}$ be  in the $d$-th tail space. By the inequality \eqref{t2-etd-integral} and the hapercontractivity \eqref{hc-1} with $q=2$,$p(y)=e^{-y}+1$,  we have 
%\begin{align*}
%\|T\|_{L_2}^{2} &\leq \frac{2}{1-e^{-2 td}} \int_{0}^{t} \sum_{j=1}^{n}\|e^{-y\mathcal{L}} d_j(T)\|_{L_2}^{2} d y\\
%&=\frac{2}{1-e^{-2 td}} \int_{0}^{t} \sum_{j=1}^{n}\|e^{-y/2\mathcal{L}}e^{-y/2\mathcal{L}} d_j(T)\|_{L_2}^{2} d y\\
%&\leq \frac{2}{1-e^{-2 td}} \sum_{j=1}^{n} \int_{0}^{t} e^{-yd}\|e^{-y/2\mathcal{L}} d_j(T)\|_{L_2}^{2} d y\\
%&\leq\frac{2}{1-e^{-2 td}} \sum_{j=1}^{n} \int_{0}^{t} e^{-yd}\|d_j(T)\|_{L_{p(y)}}^{2} d y.
%\end{align*}
%Note that $1<p(y)\leq2$. By changing the variable with $\mu=p(y)$ and applying H\"{o}lder inequality, we obtain the following result
%\begin{align}\label{t2-v-1}
%\|T\|_{L_2}^{2} &\leq \min _{t\geq0}\{\frac{2e^ t}{1-e^{-2 td}}\}\sum_{j=1}^{n} \int_{1}^{2} (\nu-1)^d\left\|d_j(T)\right\|_{L_\nu}^{2} d \nu\\
%&\leq\min _{t\geq0}\{\frac{2e^ t}{1-e^{-2 td}}\}\sum_{j=1}^{n} \big(\int_{1}^{2} (\nu-1)^{2d} d \nu\big)^{\frac{1}{2}}\big(\int_{1}^{2} \left\|d_j(T)\right\|_{L_\nu}^{4} d \nu\big)^{\frac{1}{2}}\nonumber,
%\end{align}
%where 
%\begin{equation}\label{min-td}
%\min_{t\geq0}\{\frac{2e^ t}{1-e^{-2 td}}\}=\frac{(1+2d)^{\frac{1}{2d}+1}}{d}\approx\mathcal{O}(1),\quad \mbox{as $d\rightarrow\infty$},
%\end{equation}
%and 
%\begin{equation}\label{v-2d}
%\big(\int_{1}^{2} (\nu-1)^{2d} d \nu\big)^{\frac{1}{2}}=\frac{1}{\sqrt{2d+1}}.
%\end{equation}
%The subsequent treatment is similar to the proof process of Theorem \ref{talagrand-d-tail}.
%Let $\theta(\nu)=\frac{2-\nu}{\nu}$ for $\nu \in[1,2]$. Then $\frac{1}{\nu}=\frac{\theta(\nu)}{1}+\frac{1-\theta(\nu)}{2}$.  By the H\"{o}lder inequality, we have
%\begin{align*}
%\big(\int_{1}^{2} \left\|d_j(T)\right\|_{L_\nu}^{4} d \nu\big)^{\frac{1}{2}}&\leq\big(\int_{1}^{2}\|d_j(T)\|_{L_1}^{4 \theta(\nu)} \cdot\|d_j(T)\|_{L_2}^{4(1-\theta(\nu))} d \mu\big) ^{\frac{1}{2}}\\
%&=\big(\int_{1}^{2}\|d_j(T)\|_{L_2}^{4} \cdot\big(\frac{\|d_j(T)\|_{L_1}}{\|d_j(T)\|_{L_{2}}}\big)^{4 \theta(\nu)} d \nu\big)^{\frac{1}{2}}.
%\end{align*}
%Changing variables $y=2\theta(\mu)$ again, set $b_{j}=\|d_j(T)\|_{L_1}/\|d_j(T)\|_{L_{2}}$. we can get the further estimate 
%\begin{align*}
%\big(\int_{1}^{2} \left\|d_j(T)\right\|_{L_\nu}^{4} d \nu\big)^{\frac{1}{2}}&\leq\|d_j(T)\|_{L_2}^{2}\big(\int_{0}^{2}  b_j^{2y} dy\big)^{\frac{1}{2}}\\
%&=\|d_j(T)\|_{L_2}^{2}\big(\frac{b_j^2-1}{\log b_j^2}\big)^{\frac{1}{2}}\\
%&\leq\|d_j(T)\|_{L_2}^{2}\big(\frac{1}{1+\log\frac{1}{b_j} }\big)^{\frac{1}{2}},
%\end{align*}
%where we used the fact $\frac{b^{2}-1}{\log (b)} \leq \frac{2}{1+\log (1 / b)}$ for $b \in(0,2)$ in the last inequality. Combining \eqref{t2-v-1}, \eqref{min-td}, \eqref{v-2d} and the above estimate, we get
%$$\|T\|_{L_{2}}^{2} \leq \frac{(1+2d)^{\frac{1}{2d}+1}}{d\sqrt{2d+1}} \sum_{j=1}^{n} \frac{\|d_j(T)\|_{L_{2}}^{2}}{\sqrt{1+2\log (\|d_j(T)\|_{L_{2}}/\|d_j(T)\|_{L_{1}})}}.$$
%This completes our proof, and $\frac{(1+2d)^{\frac{1}{2d}+1}}{d\sqrt{2d+1}}\rightarrow0$, as $d\rightarrow\infty$.
%\end{proof}
%Here we present an alternative proof of the quantum Talagrand inequality.
%\begin{theorem}\label{taragrand-ano-proof} For all $T\in\mathbb{M}_{2^n}$ one has
%$$
%\|T\|_{L_{2}}^{2} \leq 4 \sum_{j=1}^{n} \frac{\|d_j(T)\|_{L_{2}}^{2}}{1+\log (\|d_j(T)\|_{L_{2}}/\|d_j(T)\|_{L_{1}})}.
%$$
%\end{theorem}
%\begin{proof} For each $T\in\mathbb{M}_{2^n}$, by Lemma \ref{cov-semigroup-0-infty}, H\"{o}lder inequality and changing the variable $t=3u$,  we have
%\begin{align*}
%\|T\|_{L_{2}}&=\sum_{j=1}^n\int_{0}^\infty\mathrm{tr}(e^{-t\mathcal{L}}d_j(T)\cdot d_j(T))dt\\
%&=4\sum_{j=1}^n\int_{0}^\infty\mathrm{tr}(e^{-2t\mathcal{L}}d_j(T)\cdot e^{-2t\mathcal{L}}d_j(T))dt\\
%&\leq4\sum_{j=1}^n\int_{0}^\infty \|e^{-2t\mathcal{L}}d_j(T)\|^2_{L_{2}}dt\\
%&\leq 4\sum_{j=1}^n\int_{0}^\infty e^{-2t} \|e^{-t\mathcal{L}}d_j(T)\|^2_{L_{2}}dt.
%\end{align*}
%Applying Theorem \ref{hc-1}  with $q=2$ and $p(t)=1+e^{-2 t}$ and changing the variable with $p(t)=\nu$, we obtain
%\begin{align*}\|T\|_{L_{2}}&\leq  4\sum_{j=1}^n\int_{0}^\infty e^{-2t} \|d_j(T)\|^2_{L_{p(t)}}dt\\
%&=2\sum_{j=1}^n\|d_j(T))\|_{L_{2}}\int_{1}^2  \|d_j(T)\|_{L_{\nu}}d\nu.
%\end{align*}
%Let $\theta(\nu)=\frac{2-\nu}{\nu}$ for $\nu \in[1,2]$. Then $\frac{1}{\nu}=\frac{\theta(\nu)}{1}+\frac{1-\theta(\nu)}{2}$. Applying the H\"{o}lder inequality, we obtain
%\begin{align*}
%\|T\|_{L_{2}\left(\mathbb{M}_{n}\right)}^{2} & \leq 2 \sum_{j=1}^{n}\int_{1}^{2} \|d_{j}(T)\|_{L_{1}}^{2\theta(\nu)} \cdot\|d_{j}(T)\|_{L_{2}}^{2(1-\theta(\nu))} d\nu \\
%& =2 \sum_{j=1}^{n}\|d_j(T))\|_{L_{2}}^2\int_{1}^{2} \big(\frac{\|d_{j}(T)\|_{L_{1}}}{\|d_{j}(T)\|_{L_{2}}}\big)^{2\theta(\nu)}d\nu\\
%&=2 \sum_{j=1}^{n}\|d_j(T))\|_{L_{2}}^2\int_{1}^{2} b_j^{2\theta(\nu)}d\nu,
%\end{align*}
%where  $b_j=\frac{\|d_{j}(T)\|_{L_{1}}}{\|d_{j}(T)\|_{L_{2}}}\in(0,1)$. Changing the variable $s=2\theta(\nu)$ and using the fact $\frac{b^2-1}{\log b}\leq \frac{2}{1+\log \frac{1}{b}}$ for $b\in(0,2)$, we have
%\begin{align*}
%\|T\|_{L_{2}\left(\mathbb{M}_{n}\right)}^{2} 
%&\leq2\sum_{j=1}^{n}\|d_j(T))\|_{L_{2}}^2\int_{0}^{2} b_j^{s}ds\\
%&=2\sum_{j=1}^{n}\|d_j(T))\|_{L_{2}}^2\frac{b_j^2-1}{\log b_j}\\
%&\leq 4\sum_{j=1}^{n}\|d_j(T))\|_{L_{2}}^2\frac{1}{1+\log \frac{1}{b_j}},
%\end{align*}
%the proof is complete.
%\end{proof}
%
%
%
%
%
%\begin{theorem}\label{l1-tala-d-tail}
%For each $n\in \mathbb{N}$ and $d\in [n]$, we have
%\[
%\|T\|^2_{L_2}\leq 8\sum_{j=1}^{n}\frac{\|d_{j}(T)\|_{L_{1}}}{d+\log\left(1/\|d_{j}(T)\|_{L_{1}}\right)}\]
%for all $T\in \mathcal{T}^{\geq d}$ with $\|T\|_{L_\infty}\leq1$.
%\end{theorem}
%\begin{proof}
%By Lemma \ref{cov-def} and \ref{hc-1}, we obtain
%\begin{equation}\label{semigroup-approach-l1}
%\begin{split}
%\|T\|_{L_2}&=\sum_{j=1}^n\int_{0}^\infty\mathrm{tr}((e^{-t\mathcal{L}}d_j(T))\cdot d_j(T))dt\\
%&=4\sum_{j=1}^n\int_{0}^\infty\mathrm{tr}((e^{-2t\mathcal{L}}d_j(T))\cdot e^{-2t\mathcal{L}}d_j(T))dt\\
%&\leq 4\sum_{j=1}^{n}\int_{0}^{\infty}\|e^{-2t\mathcal{L}}d_{j}(T)\|^2_{L_{2}}dt\\
%&=4\sum_{j=1}^{n}\int_{0}^{\infty}e^{-2td}\|e^{-t\mathcal{L}}d_{j}(T)\|^2_{L_{2}}dt\\
%&\leq 4\sum_{j=1}^{n}\int_{0}^{\infty}e^{-2td}\|d_{j}(T)\|^2_{L_{p(t)}}dt\\
%&\leq 4\sum_{j=1}^{n}\int_{0}^{\infty}e^{-2td}\|d_{j}(T)\|^{\frac{2}{p(t)}}_{L_{1}}\|d_{j}(f)\|_{L_{\infty}}^{2(1-\frac{1}{p(t)})}dt\\
%&\leq4\sum_{j=1}^{n}\|d_{j}(T)\|_{L_{1}}\int_{0}^{\infty}e^{-2td}\|d_{j}(T)\|^{\frac{2}{p(t)}-1}_{L_{1}}dt\\
%&=4\sum_{j=1}^{n}\|d_{j}(T)\|_{L_{1}}\int_{0}^{\infty}e^{-2td}\|d_{j}(T)\|^{\frac{1-e^{-2t}}{1+e^{-2t}}}_{L_{1}}dt
%\end{split}
%\end{equation}
%where $p(t)=1+e^{-2t}$, we used the H\"older inequality and Lemma \ref{t-infty} in the last two inequalities. %line with parameters
%%\[
%%\frac{1}{1+e^{-2t}}=\frac{1-e^{-2t}}{1+e^{-2t}}+\frac{e^{-2t}}{1+e^{-2t}},\quad t\geq 0.
%%\]
%By Lemma \ref{t-infty}, $\|d_{j}(T)\|_{L_{1}}\leq\|d_{j}(T)\|_{L_{\infty}}\leq1$. Let $b=\|d_{j}(T)\|_{L_{1}}\in(0,1)$. It now remains to estimate the following term
%$$\int_{0}^{\infty}e^{-2td}b^{\frac{1-e^{-2t}}{1+e^{-2t}}}dt.$$
%Changing the variable $1+e^{-2t}\mapsto u$ yields that
%$$\int_{0}^{\infty}e^{-2td}b^{\frac{1-e^{-2t}}{1+e^{-2t}}}dt=\frac{1}{2}\int_{1}^{2}(u-1)^{d-1}b^{\frac{2-u}{u}}du.$$
%Changing variable $\frac{2-u}{u}\mapsto t$ yields that
%$$
%\int_{0}^{\infty}e^{-2td}b^{\frac{1-e^{-2t}}{1+e^{-2t}}}dt=\int_{0}^{1}\frac{(1-t)^{d-1}}{(t+1)^{d+1}}\cdot b^{t}dt.$$
%Changing the variable $t\mapsto \frac{\nu}{d}$ yields that
%\begin{equation}\label{key-term-4-l1}
%\begin{split}
%\int_{0}^{\infty}e^{-2td}b^{\frac{1-e^{-2t}}{1+e^{-2t}}}dt&=\frac{1}{d}\int_{0}^{d}\frac{d^{2}(d-\nu)^{d-1}}{(d+\nu)^{d+1}}\cdot b^{\nu/d} d\nu\leq \frac{1}{d}\int_{0}^{d}b^{\nu/d}d\nu,
%\end{split}
%\end{equation}
%where 
%the last inequality is due to that the function $\nu\mapsto \frac{d^{2}(d-\nu)^{d-1}}{(d+\nu)^{d+1}}$ for $\nu\in[1,d]$ is decreasing. 
%And $b\in (0,1)$, we have that {\color{red} something is missing here}
%%\begin{equation}\label{key-term-5-l1}
%%\begin{split}
%%\int_{1}^{d}\frac{d^{2}(d-\nu)^{d-1}}{(d+\nu)^{d+1}}\cdot b^{\nu/d}d\nu&=\sum_{j=1}^{d-1}\int_{j}^{j+1}\frac{d^{2}(d-\nu)^{d-1}}{(d+\nu)^{d+1}}\cdot b^{\nu/d}d\nu\\
%%&\leq\sum_{j=1}^{d-1}\frac{d^{2}(d-j)^{d-1}}{(d+j)^{d+1}}\int_{0}^{1}b^{\nu/d}d\nu,
%%\end{split}
%%\end{equation}
%%where we used the monotone decreasing of $\nu\mapsto b^{\nu/d}$ for $\nu\geq 0$.
%%It is clear that
%%\[
%%\sum_{j=1}^{d-1}\frac{d^{2}(d-j)^{d-1}}{(d+j)^{d+2}}\leq \sum_{j=1}^{d-1}\frac{d^{d+1}}{(d+j)^{d+1}}.
%%\]
%%Since $x\mapsto \frac{1}{(d+x)^{d+1}}$ is a monotone decreasing function on $x\geq 0$, draw a picture and note that
%%\begin{equation}\label{key-term-6-l1}
%%\begin{split}
%%\sum_{j=1}^{d-1}\frac{d^{d+1}}{(d+j)^{d+1}}&\leq d^{d+1}\int_{0}^{d}(d+x)^{-d-1} dx\\
%%&=-d^{d}(d+x)^{-d}|_{0}^{d}\leq 1.
%%\end{split}
%%\end{equation}
%Substituting estimate given as above, we get that
%%\begin{equation}\label{key-term-8-l1}
%%\frac{1}{d}\int_{1}^{d}\frac{d^{2}(d-\nu)}{(d+\nu)^{d+1}}\cdot b^{\nu/d}d\nu\leq \frac{1}{d}\int_{0}^{1}b^{\nu/d}d\nu.
%%\end{equation}
%\begin{align*}
%\int_{0}^{\infty}e^{-2td}b^{\frac{1-e^{-2t}}{1+e^{-2t}}}dt&\leq\big(\frac{1}{d}\big)\frac{b-1}{\log b^{\frac{1}{d}}}\\
%&\leq \big(\frac{1}{d}\big)\frac{b^{\frac{1}{d}}-1}{\log b^{\frac{1}{d}}}.
%\end{align*}
%Thanks to the numerical inequality $\frac{t-1}{\log(t)}\leq \frac{2}{1-\log(t)}$ by Lemma \ref{log-b}, we have
%\begin{equation}\label{key-term-9-l1}
%\int_{0}^{\infty}e^{-2td}b^{\frac{1-e^{-2t}}{1+e^{-2t}}}dt\leq \left(\frac{2}{d}\right)\frac{1}{1+\log\left(b^{-1/d}\right)}.
%\end{equation}
%Combing \eqref{semigroup-approach-l1} and \eqref{key-term-9-l1} yields that for each $T\in \mathcal{T}^{\geq d}$
%\begin{equation}\label{key-term-10-1}
%\begin{split}
%\|T\|^{2}_{L_{2}}&\leq \left(\frac{8}{d}\right)\sum_{j=1}^{n}\frac{\|d_{j}(T)\|_{L_{1}}}{1+\frac{1}{d}\log\left(1/\|d_{j}(T)\|_{L_{1}}\right)}\\
%&=8\sum_{j=1}^{n}\frac{\|d_{j}(T)\|_{L_{1}}}{d+\log\left(1/\|d_{j}(T)\|_{L_{1}}\right)},
%\end{split}
%\end{equation}
%which is the desired inequality.
%\end{proof}
%\begin{lemma}\label{log-b} For any $b\in(0,1)$,  we have
%$$\frac{b-1}{\log(b)}\leq \frac{2}{1+\log(\frac{1}{b})}$$
%\end{lemma}
%\begin{proof}
%The fact that $\frac{b-1}{\log(b)}\leq \frac{2}{1+\log(\frac{1}{b})}$ is equivalent to show $$\frac{e^{-x}-1}{-x}\leq\frac{2}{1+x},\quad \mbox{for all $x\in (0,\infty)$}.$$ 
%Rearranging the above  inequality,  it suffices to show that 
%$$x-1+e^{-x}+xe^{-x},\quad \mbox{for all $x\in (0,\infty)$}.$$ 
%Rearranging the inequality yields
%\[
%(x-1)e^{x}+1+x\geq 0,\quad\mbox{for all }x\in(0,\infty),
%\]
%which also holds trivially.
%\end{proof}




%From Theorem \ref{l1-tala-d-tail}, we can obtain an optimized version of the quantum L1-Talagrand inequality derived from \cite[Theorem 3.2]{RWZ2024}.


\chapter{Quantum correlation inequalities with lower bound}\label{sec-tala-correlation}
A function $f: \{-1,1\}^n\to \mathbb{R}$ is said monotone  increasing if $f(x)\leq f(y)$ whenever $x=(x_j)_{j=1}^n\leq y=(y_j)_{j=1}^n$, $x,y\in \{-1,1\}^n$. Here $x\leq y$ means that $x_j\leq y_j$ for all $j\in [n]$. Recall that $f: \{-1,1\}^n\to \mathbb{R}$ is motonone  increasing if and only if 
\begin{equation}\label{basic-mi}
w_{\{j\}}D_j(f)\geq 0,\quad \forall j\in [n].
\end{equation}
The remarkable Harris-Kleitman correlation inequality (\cite{Ha1960,Kl1966}) states that for monotone increasing functions $f,g: \{-1,1\}^n\to \{-1,1\}$, their correlation is always positive, that is, 
\begin{equation}\label{hk-ine}
\mathrm{Cov}(f,g):=\mathbb{E}(fg)-\mathbb{E}(f)\mathbb{E}(g)\geq 0.
\end{equation}
Talagrand \cite{Ta1996} improved the Harris--Kleitman inequality via establishing the following quantitative inequality: for monotone increasing functions 
$f,g: \{-1,1\}^n\to \{-1,1\}$, there exists an absolute constant $C>0$ such that 
\begin{equation}\label{ta-coi}
\mathrm{Cov}(f,g):=\mathbb{E}(fg)-\mathbb{E}(f)\mathbb{E}(g)\geq C \Psi\Big(\sum_{j\in [n]} \|D_j(f)\|_{L_1}\|D_j(g)\|_{L_1}\Big),
\end{equation}
where $\Psi(t)=t/\log(\frac{e}{t})$. 
Talagrand's method and results inspired a number of important
works in combinatorics and probability theory.
An alternative correlation lower bound was recently proved  by
Keller, Mossel, and Sen in \cite{KMS2014}. For monotone increasing functions 
$f,g: \{-1,1\}^n\to \{-1,1\}$, there exists an absolute constant $K>0$ such that
	\begin{equation}\label{kms-correlation}
	\mathrm{Cov}(f,g)\geq K\sum_{j=1}^{n}\frac{\|D_{j}(f)\|_{L_{1}}\|D_{j}(g)\|_{L_{1}}}{\sqrt{\log(e/\|D_{j}(f)\|_{L_{1}})\log(e/\|D_{j}(g)\|_{L_{1}})}}.
\end{equation}
Note that the lower bounds in \eqref{ta-coi} and \eqref{kms-correlation} are uncomparable. To combine the advantages of \eqref{ta-coi} and \eqref{kms-correlation}, Kalai, Keller and Mossel \cite{KKM2016} studied  stronger correlation lower bounds for monotone increasing functions 
under additional natural regularity or symmetry conditions.
Recently, the famous Chv\'{a}tal's conjecture  was found by Friedgut et al. 
 \cite{FKKK2018} to be equivalent to a certain strengthening of Talagrand's correlation inequality. Motivated by this conjecture, Eldan \cite{El2022} improved \eqref{ta-coi} and \eqref{kms-correlation} when one of the monotone increasing functions is so-called ``antipodal". We also refer the reader to \cite{CC2025} for  recent developments on this problem.
 
 
In this chapter, we aim to extend the above results \eqref{hk-ine}, \eqref{ta-coi},  and \eqref{kms-correlation} to the quantum Boolean cube setting. The main results of this chapter are Theorem  \ref{cov-0}, Theorem \ref{quantum-ta} and Theorem \ref{quantum-kms}.
\section[Quantum Monotone Functions]{Quantum monotonically increasing function and Harris-Kleitman inequality}

We first introduce the definition of monotone elements in $\mathbb{M}_{2^n}$. To this end, for each $j\in [n]$ and  $\ell\in\{1,2,3\}$,  we define
\[
\dc_{j,\ell}(T)\coloneqq\sum_{\substack{\s\in\{0,1,2,3\}^{n}\\ s_{j}=\ell}}\widehat{T}(\s)\sigma_{\s},
\]
\[
\partial_{j,\ell}(T)\coloneqq\sigma_{\ell e_{j}}\dc_{j,\ell}(T),
\]
and
\[
\partial_{j}(T)\coloneqq \sum_{\ell=1}^3\partial_{j,\ell}(T).
\]
Here, 
	$e_j=(0,0,\cdots, 1,0,0)$, and $1$ only appears in the $j$-th position. Similarly,  $\ell e_j=(0,0,\cdots, \ell,0,0)$ and $\ell$ only appears in the $j$-th position.

\begin{definition} \label{mo-in-o}  
	An element $T\in \mathbb{M}_{2^n}$ is said to be monotonically increasing if
	$$\partial_{j,\ell}(T)\geq0,\quad\mbox{for each $j\in[n], \ell\in\{1,2,3\}$}.$$
\end{definition}


This definition is of course motivated by \eqref{basic-mi}. 
We here provide a basic example of monotonically increasing
elements in $\M_{2^n}$. 
\begin{example}\label{ex-mo-in-o}  Let  $T\in \mathbb{M}_{2^n}$ and assume that the Fourier coefficients of $T$ satisfy $\widehat{T}(\s)\in\mathbb{R}$ and 
	$$\widehat{T}(\ell e_j)\geq \sum_{\ell\in\{1,2,3\}}\sum_{\substack{\s_j=\ell \\ |\s|\geq2}} |\widehat{T}(\s)|.$$
	Then $T$ is monotonically increasing.
\end{example}
\begin{proof} According to the choice of $T$, it follows that $T$ is self-adjoint. For each $j\in[n]$ and $\ell\in\{1,2,3\}$, to show $\partial_{j,\ell}(T)\geq0$, it suffices to verify that $\mathrm{tr}[p\cdot\partial_{j,\ell}(T)]\geq0$ for each projection $p\in\mathcal{P}(\mathbb{M}_{2^n})$. Indeed, we have
	\begin{align*}
		\mathrm{tr}[p\cdot\partial_{j,\ell}(T)]&=\sum_{\s_j=\ell}\widehat{T}(\s)\mathrm{tr}(p\cdot\sigma_{\ell e_j}\sigma_\s)\\
		&=\widehat{T}(\ell e_j)\mathrm{tr}(p)+\sum_{\substack{\s_j=\ell\\ |\s|\geq2}}\widehat{T}(\s)\mathrm{tr}(p\cdot\sigma_{\ell e_j}\sigma_\s)\\
		&\geq \widehat{T}(\ell e_j)\mathrm{tr}(p)-\sum_{\substack{\s_j=\ell \\ |\s|\geq2}}|\widehat{T}(\s)| \cdot |\mathrm{tr}(p\cdot\sigma_{\ell e_j}\sigma_\s)|\\
		&\geq\big(\widehat{T}(\ell e_j)-\sum_{\substack{\s_j=\ell  \\ |\s|\geq2}}|\widehat{T}(\s)|\big)\mathrm{tr}(p)\geq 0,
	\end{align*}
	where we used the fact $|\mathrm{tr}(p\cdot\sigma_{\ell e_j}\sigma_\s)|\leq \mathrm{tr}(p)$ and the assumption. The desired assertion follows.
\end{proof}


We  collect useful properties of the partial differential operators $\partial_{j,\ell}$ and $\partial_j$.

\begin{lemma}\label{comparison}
	Let $T_1, T_2\in \M_{2^{n}}$.
	\begin{enumerate}[{\rm (i)}]
		\item For each $j\in [n]$ and $1\leq p\leq\infty$, $\|\partial_{j}(T)\|_{L_{p}}\leq 3\|\dc_{j}(T)\|_{L_{p}}$.
		\item For all $j\in [n]$ and $\ell\in\{1,2,3\}$, we have $$\sigma_{\ell e_j}e^{-t\mathcal{L}}\dc_{j,\ell}(T)=\partial_{j,\ell}\left(e^{-t\La}(T)\right)=e^{-t}e^{-t\La}\left(\partial_{j,\ell}(T)\right).$$

		\item We always have the following covariance identity
	\begin{equation}\label{key-identity}
		\mathrm{Cov}(T_1,T_2)=\sum_{j=1}^{n}\sum_{\ell=1}^{3}\int_{0}^{\infty}e^{-t}\mathrm{tr}\left(\partial_{j,\ell}(T_1)^*\cdot e^{-t\La}\left(\partial_{j,\ell}(T_2)\right)\right)dt.
	\end{equation}
	\end{enumerate}
\end{lemma}

\begin{proof}
	(i) Fix $j\in [n]$. It is clear that  $$\|\dc_{j,\ell}(T)\|_{L_{p}}=\left\|\mathcal{E}_{\M_{2^{n}}\setminus\{\sigma_{me_{j}}:m\neq \ell\}}\dc_{j}(T)\right\|_{L_{p}}\leq \|\dc_{j}(T)\|_{L_{p}}\qquad \ell=\{1,2,3\}.$$ 
	The definition of conditional expectation is referred to Chapter \ref{sec BKS 1}. By the contractivity of conditional expectations on \(L_p\), we have
	$$\|\partial_{j}(T)\|_{L_{p}}\leq \sum_{\ell=1}^{3}\|\sigma_{\ell e_{j}}\dc_{j,\ell}(T)\|_{L_{p}}\leq 3\|\dc_{j}(T)\|_{L_{p}}.$$
	
	(ii) The first equality is obvious. We here only show the second equality. Fix $j\in [n]$ and $\ell\in \{1,2,3\}$. 
According to the definitions, we write
\begin{equation}\label{fourier-expansion-1}
	\partial_{j,\ell}\left(e^{-t\La}(T)\right)=\partial_{j,\ell}\left(\sum_{\s\in\{0,1,2,3\}^{n}}\widehat{T}(\s)e^{-t|\s|}\sigma_{\s}\right)=\sum_{\s_{j}=\ell}\widehat{T}(\s)e^{-t|\s|}\sigma_{\s\ominus \ell e_{j}}
\end{equation}
and
\begin{equation}\label{fourier-expansion-2}
	e^{-t\La}\left(\partial_{j,\ell}(T)\right)=\sum_{\s_{j}=\ell}\widehat{T}(\s)e^{-t(|\s|-1)}\sigma_{\s\ominus \ell e_{j}}.
\end{equation}
Comparing \eqref{fourier-expansion-1} and \eqref{fourier-expansion-2} yields the desired result. 

	(iii) Combining Lemma \ref{cov-semigroup-0-infty} and item (ii) above, we have	\begin{equation*}\label{covariance-1}
	\begin{split}
		\mathrm{Cov}(T_1,T_2)
		&=\sum_{j=1}^{n}\int_{0}^{\infty}\mathrm{tr}\left(\dc_{j}(T_1)^*\cdot \dc_{j}(e^{-t\La}(T_2))\right)dt\\
		&=\sum_{j=1}^{n}\sum_{\ell=1}^{3}\int_{0}^{\infty}\mathrm{tr}\left(\dc_{j,\ell}(T_1)^*\cdot \dc_{j,\ell}\left(e^{-t\La}(T_2)\right)\right)dt\\
		&=\sum_{j=1}^{n}\sum_{\ell=1}^{3}\int_{0}^{\infty}\mathrm{tr}\left(\partial_{j,\ell}(T_1)^*\cdot \partial_{j,\ell}\left(e^{-t\La}(T_2)\right)\right)dt\\
		&=\sum_{j=1}^{n}\sum_{\ell=1}^{3}\int_{0}^{\infty}e^{-t}\mathrm{tr}\left(\partial_{j,\ell}(T_1)^*\cdot e^{-t\La}\left(\partial_{j,\ell}(T_2)\right)\right)dt.
	\end{split}
\end{equation*}
The proof is complete.
\end{proof}




Useful properties of monotonically increasing elements in $\mathbb{M}_{2^n}$ are stated below.

\begin{lemma}\label{basic}
Let $T\in \M_{2^{n}}$ be monotonically increasing. 
	\begin{enumerate}[{\rm (i)}]
			\item For each $j\in [n]$, $ \|\dc_{j}(T)\|_{L_{1}}\leq \|\partial_{j}(T)\|_{L_{1}}$.
\item  For each $j\in [n]$ and $\ell\in \{1,2,3\}$, 
$$\|\dc_{j,\ell}(T)\|_{L_{1}}=\widehat{T}(\ell e_j).$$
		\item For each $j\in [n]$ and $\ell\in \{1,2,3\}$, 
		\[
		\partial_{j,\ell}\left(e^{-t\La}(T)\right)\geq 0,
		\]
that is, $e^{-t\La}(T)$ is also monotonically increasing.
	\end{enumerate}
\end{lemma}
\begin{proof}
		(i) Since $T$ is monotonically increasing, it follows that $\partial_{j,\ell}(T)\geq 0$ for all $j\in [n]$ and $\ell\in \{1,2,3\}$. Then, for each $j\in [n]$, we have
	\begin{align*}
		\|\partial_{j}(T)\|_{L_{1}}&=\sum_{\ell=1}^{3}\mathrm{tr}\left(\partial_{j,\ell}(T)\right)=\sum_{\ell=1}^{3}\|\dc_{j,\ell}(T)\|_{L_{1}}\\
		&\geq \left\|\sum_{\ell=1}^{3}\dc_{j,\ell}(T)\right\|_{L_{1}}=\|\dc_{j}(T)\|_{L_{1}}.
	\end{align*}
	
	
(ii) Since $\|\cdot\|_{L_1}$ is unitary variant, for each $j\in [n]$ and $\ell\in \{1,2,3\}$, we have
\begin{align*}
	\|\dc_{j,\ell}(T)\|_{L_{1}}&=\|\sigma_{\ell e_j}\dc_{j,\ell}(T)\|_{L_{1}}=\|\partial_{j,\ell}(T)\|_{L_1}\\
	&=\mathrm{tr}(\partial_{j,\ell}(T))=\widehat{T}(\ell e_j),
\end{align*}
where we used the fact that $T\in \M_{2^{n}}$ is monotonically increasing in the third equality, and the last equality is due to orthogonality of the quantum Walsh series. 
	
	(iii)  By Lemma \ref{comparison} (ii) and the fact that $e^{-t\La}$ is positive preserving, it follows that
	\[
	\partial_{j,\ell}\left(e^{-t\La}(T)\right)=e^{-t}e^{-t\La}\left(\partial_{j,\ell}(T)\right)\geq 0, \quad \forall j\in [n], \ell\in \{1,2,3\}.
	\]
\end{proof}







Now we are ready to prove quantum version of the Harris--Kleitman inequality \eqref{hk-ine}.

\begin{theorem}\label{cov-0}  
Let $T_1,T_2\in\mathbb{M}_{2^n}$ be monotonically increasing operators. Then
	$$\mathrm{Cov}(T_1,T_2)\geq0.$$
\end{theorem}
\begin{proof} 
For $t>0$, set $h(t)=\mathrm{tr}(e^{-t\mathcal{L}}(T_1)^*T_2)$. Then we find that
	$$h(0)=\mathrm{tr}(T_1^*T_2)\quad \mbox{and \quad $h(\infty)=\mathrm{tr}(T_1)^*\mathrm{tr}(T_2)$}.$$
By the definition of $h(t)$ and Lemma \ref{basic} (iii), we get
\begin{align*}
h'(t)&=-\mathrm{tr}((\mathcal{L}e^{-t\mathcal{L}}(T_1))^*T_2)\\
&=-\sum_{j=1}^n\mathrm{tr}(e^{-t\mathcal{L}}\dc_j(T_1)^*\dc_j(T_2))\\
&=-\sum_{j=1}^n\sum_{\ell=1}^3\mathrm{tr}(e^{-t\mathcal{L}}\dc_{j,\ell}(T_1)^*\dc_{j,\ell}(T_2))\\
&=-\sum_{j=1}^n\sum_{\ell=1}^3\mathrm{tr}((\sigma_{\ell e_j}e^{-t\mathcal{L}}\dc_{j,\ell}(T_1))^*\dc_{j,\ell}(T_2)\sigma_{\ell e_j})\\
&=-\sum_{j=1}^n\sum_{\ell=1}^3e^{-t}\mathrm{tr}\Big(e^{-t\mathcal{L}}(\partial_{j,\ell}(T_1))^* \partial_{j,\ell}(T_2)\Big)\\
&\leq 0,
\end{align*}
where the last equality is due to Lemma \ref{comparison} (ii).
Hence, the function $h$ is decreasing, and consequently, we have
$$\mathrm{Cov}(T_1,T_2)=h(0)-h(\infty)\geq0.$$
The proof is complete.
\end{proof}


\section{The Talagrand correlation inequality}
We next aim to establish quantum version of the Talagrand quantitative correlation inequality \eqref{ta-coi}. 
Given $T_1,T_2\in\mathbb{M}_{2^n}$, $d\in [n]$ and $J\subset[n]$, we set $$\mathcal{W}_d(T_1,T_2)=\sum_{|\s|=d}\widehat{T_1}(\s)\widehat{T_2}(\s),$$
and
$$\mathcal{W}_{1,J}(T_1,T_2)=\sum_{j\in J}\sum_{\ell=1}^3\widehat{T_1}(\ell e_j)\widehat{T_2}(\ell e_j).$$
If  $T_1=T_2=T$, the preceding symbols will be abbreviated as $\mathcal{W}_d(T)$, $\mathcal{W}_{1,J}(T)$,  respectively.

For each $T_1$, $T_2\in \M_{2^{n}}$ with $\max\{\|T_1\|_{L_{\infty}},\|T_2\|_{L_{\infty}}\}\leq 1$, 
it is clear that $\left|\mathcal{W}_{d}(T_1,T_2)\right|\leq 1$ for each $d\in [n]$. Indeed, by the Cauchy-Schwarz inequality we have
\begin{equation}\label{upper-bound-1}
	\begin{split}
		|\mathcal{W}_{d}(T_1,T_2)|&=\left|\sum_{|\s|=d}\widehat{T_1}(\s)\widehat{T_2}(\s)\right|\\
		&\leq\Big(\sum_{|\s|=d}|\widehat{T_1}(\s)|^{2}\Big)^{1/2}\Big(\sum_{|\s|=d}|\widehat{T_2}(\s)|^{2}\Big)^{1/2}\\
		&\leq\|T_1\|_{L_{2}}\|T_2\|_{L_{2}}\leq \|T_1\|_{L_{\infty}}\|T_2\|_{L_{\infty}}\leq 1.
	\end{split}
\end{equation}
Next, we will establish  Talagrand's lower bound correlation inequality for the covariance between monotonically increasing elements $T_1,T_2\in\mathbb{M}_{2^n}$. Due to the following counterexample, the lower bound we obtain is of the $\mathcal{W}_{1}(T_1,T_2)$-form rather than 
$M(T_1,T_2)$-form lower bound,
where 
$$M(T_1,T_2):=\sum_{j=1}^n\|\dc_{j}(T_1)\|_{L_{1}}\|\dc_{j}(T_2)\|_{L_{1}}.$$





\begin{example}\label{example-2} 
Let $T_1=\sum_{j=1}^n\sigma_{e_{j}}$ and $T_2=\sum_{j=1}^n\sigma_{2e_{j}}$. Then we have $\mathrm{Cov}(T_1,T_2)=0$, but $M(T_1,T_2)>0$.
\end{example}
\begin{proof}
	It is clear that both $T_1$ and $T_2$ are monotonically increasing. By the definition of the covariance and for any $i\in\{1,2,3\}$, $\mathrm{tr}(\sigma_i)=0$, we have
	$$\mathrm{Cov}(T_1,T_2)=\mathrm{tr}(T_1T_2)-\mathrm{tr}(T_1)\mathrm{tr}(T_2)=0.$$
	In addition,  
	$$M(T_1,T_2)=\sum_{j=1}^n\|d_{j}(T_1)\|_{L_{1}}\|d_{j}(T_2)\|_{L_{1}}=\sum_{j=1}^n\|\sigma_{e_j}\|_{L_{1}}\|\sigma_{2e_j}\|_{L_{1}}=n>0.$$
	Therefore, we cannot obtain the following correlation inequality
	$$\mathrm{Cov}(T_1,T_2)\geq c M(T_1,T_2)\cdot\left(\log\left(\frac{e}{M(T_1,T_2)}\right)\right)^{-1}.$$
\end{proof}



For each $j\in[n]$ and  any $f:\{-1,1\}^n\rightarrow\mathbb{R}$,   the classical framework involves only a single Fourier coefficient $\widehat{f}(\{j\})$. For each $j\in[n]$ and  $T\in\mathbb{M}_{2^n}$, whereas the noncommutative setting yields three first-order Fourier coefficients $\widehat{T}(\ell e_j), \ell\in\{1,2,3\}$. This fundamental difference directly leads to the fact that there is no inherent relation between $M(T_1,T_2)$ and $\mathcal{W}_1(T_1,T_2)$, when 
$T_1,T_2$ are monotone increasing matrices. Motivated by the role of structural conditions in classical inequalities \cite[Definition 1.6]{KKM2016}, we introduce the following  similarity condition in the noncommutative setting. 
\begin{definition}\label{weak-similar-def}
Let  $T_1,T_2\in \mathbb{M}_{2^n}$ be monotonically increasing. We say $T_1,T_2$ satisfy the  similarity condition if 
for each $j\in[n]$, there exists a constant $c_j\in (0,\infty)$ such that
\begin{equation}\label{weak-similar}
	\frac{\widehat{T_1}(e_j)}{\widehat{T_2}(e_j)}= \frac{\widehat{T_1}(2e_j)}{\widehat{T_2}(2e_j)}= \frac{\widehat{T_1}(3e_j)}{\widehat{T_2}(3e_j)}=c_j.
\end{equation}
\end{definition} 

\begin{remark}\label{rem-similarity}
The above similarity condition holds automatically in the classical setting. Indeed, for each $j\in[n]$ and any functions $f,g:\{-1,1\}^n\rightarrow\mathbb{R}$, the pointwise Fourier coefficient reduces to a single term, and it suffices to set $c_j=\frac{\widehat{f}(\{j\})}{\widehat{g}(\{j\})}$, which in effect means that no additional  condition is imposed. 
\end{remark}



\begin{remark}\label{weak-similar-def2}  If we replace the above condition \eqref{weak-similar} with the following weak similarity condition 
\begin{equation}\label{weak-similar-def2-con}
\max_{\ell\in\{1,2,3\}}\left\{\frac{\widehat{T_1}(\ell e_j)}{\widehat{T_2}(\ell e_j)}\right\}\leq c_j \min_{\ell\in\{1,2,3\}}\left\{\frac{\widehat{T_1}(\ell e_j)}{\widehat{T_2}(\ell e_j)}\right\},\quad c_j>0,
\end{equation}
all the corresponding results (like Theorem \ref{quantum-ta}, Proposition \ref{tala-correlation-pro1} and Proposition \ref{tala-correlation-pro2}) remain valid.
\end{remark}
This section is devoted to establishing the following \emph{quantum Talagrand correlation inequality}. Note that it follows from Lemma \ref{basic} (ii) that 
\begin{equation}
\mathcal{W}_{1}(T_1,T_2)=\	\sum_{j=1}^{n}\sum_{\ell=1}^{3}{\|\dc_{j,\ell}(T_1)\|_{L_{1}}\|\dc_{j,\ell}(T_2)\|_{L_{1}}}
\end{equation}
provided both  $T_1,T_2\in \M_{2^{n}}$ are  monotonically increasing.

\begin{theorem}[quantum Talagrand correlation inequality]\label{quantum-ta}
 Suppose that monotonically increasing elements $T_1,T_2\in \M_{2^{n}}$ satisfy condition \eqref{weak-similar} and $$\max\{\|T_1\|_{L_{\infty}},\|T_2\|_{L_{\infty}}\}\leq 1.$$ Then exists a universal constant $K>0$ such that 
	\begin{equation}
		\mathrm{Cov}(T_1,T_2)\geq K\frac{\mathcal{W}_{1}(T_1,T_2)}{\log\left(e/\mathcal{W}_{1}(T_1,T_2)\right)}.
	\end{equation}
\end{theorem}



Before proving Theorem \ref{quantum-ta}, we need following auxiliary lemmas regarding the basic properties related to the partial differential operators and Rademacher projections (see Chapter \ref{sec-hqs}).
%\begin{lemma}\label{auxiliary-1}
%	For each $T\in \M_{2^{n}}$, $j$, $k\in [n]$ and $\dag\in\{=,\neq,\leq,\geq\}$, we have
%	\begin{equation}\label{commuting-relations}
%		\dc_{j}\circ\mathrm{Rad}^{\dag k}(T)=\mathrm{Rad}^{\dag k}\circ \dc_{j}(T).
%	\end{equation}
%{\color{blue}we may not need this part, no use}
%\end{lemma}
%\begin{proof}
%	The proof of the lemma is elementary, which follows from the Fourier expansion. We leave the proof for interested readers.
%\end{proof}

\begin{lemma}\label{monotone-lemma-1}
	Let $T_1$, $T_2\in \M_{2^{n}}$. For each $j\in [n]$ and $\ell\in\{1,2,3\}$, we have
	\begin{equation*}
		\left\langle \dc_{j,\ell}(T_1),\dc_{j,\ell}(T_2)\right\rangle=\left\langle\partial_{j,\ell}(T_1),\partial_{j,\ell}(T_2)\right\rangle,
	\end{equation*}
and consequently,
	\begin{equation}\label{transfer}
		\left\langle \dc_{j}(T_1),\dc_{j}(T_2)\right\rangle=\sum_{\ell=1}^3\left\langle\partial_{j,\ell}(T_1),\partial_{j,\ell}(T_2)\right\rangle.
	\end{equation}
\end{lemma}

\begin{proof}
We write
	\begin{align*}
		\left\langle \dc_{j,\ell}(T_1),\dc_{j,\ell}(T_2)\right\rangle&=\mathrm{tr}\left(\dc_{j,\ell}(T_1)^{*}\dc_{j,\ell}(T_2)\right)\\
		&=\mathrm{tr}\left(\left(\sigma_{\ell e_{j}}\dc_{j, \ell}(T_1)\right)^{*}\left(\sigma_{\ell e_{j}}\dc_{j,\ell}(T_2)\right)\right)\\
		&=\left\langle\partial_{j,\ell}(T_1),\partial_{j,\ell}(T_2)\right\rangle,
	\end{align*}
	which completes the proof.
\end{proof}

\begin{lemma}\label{l1-norm}
If $T\in \mathbb{M}_{2^n}$ is monotonically increasing, then, for each $j\in [n]$ and $\ell\in\{1,2,3\}$, we have
\begin{equation*}
	\partial_{j,\ell}\left(\mathrm{Rad}^{=1}(T)\right)=\left\|\partial_{j,\ell}(T)\right\|_{L_{1}}=\widehat{T}(\ell e_j).
\end{equation*}	
\end{lemma}


\begin{proof}Fix $j\in [n]$ and  $\ell\in\{1,2,3\}$.  Since  $T\in \mathbb{M}_{2^n}$ is monotonically increasing, it follows that
	\begin{align*}
	\|\partial_{j,\ell}(T)\|_{L_{1}}
		=\mathrm{tr}(\sigma_{\ell e_j}d_{j,\ell}(T))
		=\sum_{\s:\s_j=\ell} \widehat{T}(\s) \mathrm{tr}(\sigma_{\ell e_j}\sigma_{\s})
		=\widehat{T}(\ell e_j).
	\end{align*}
On the other hand, for a given $j\in [n]$ and $\ell\in\{1,2,3\}$, we note that
	\begin{equation*}\label{positive-1}
		\begin{split}
			\partial_{j,\ell}\left(\mathrm{Rad}^{=1}(T)\right)&=\sum_{|\s|=1}\widehat{T}(\s)\partial_{j,\ell}(\sigma_{\s})
=\widehat{T}(\ell e_j).
		\end{split}
	\end{equation*}
The desired result follows.
\end{proof}


\begin{lemma}\label{h-xi}
	 Suppose that $T_1, T_2\in \M_{2^{n}}$ are monotonically increasing. Then the function
	 $$H(\xi)\coloneqq\mathrm{Cov}\left(e^{-\xi\mathcal L}(T_1),T_2\right),\quad \xi\geq 0$$
	 is decreasing. 
\end{lemma}

\begin{proof}
 By Lemma \ref{cov-semigroup-0-infty}, we have
	\[
	\mathrm{Cov}\left(e^{-\xi\La}(T_1),T_2\right)=\sum_{j=1}^{n}\int_{0}^{\infty}\left\langle\dc_{j}e^{-\xi \La}(T_1),\dc_{j}e^{-t\La}(T_2)\right\rangle dt.
	\]
	Hence, by the fact that $\{\dc_{j}\}_{j=1}^{n}$ are orthogonal projections \eqref{transfer} and Lemma \ref{comparison} (ii), we have
	\begin{equation}\label{derivative-1}
		\begin{split}
			\frac{d}{d\xi}H(\xi)&=-\sum_{j=1}^{n}\int_{0}^{\infty}\left\langle\La\dc_{j}e^{-\xi\La}(T_1),\dc_{j}e^{-t\La}(T_2)\right\rangle dt\\
			&=-\sum_{j=1}^{n}\int_{0}^{\infty}\left\langle\dc_{j}e^{-\xi\La}(T_1),\dc_{j}e^{-t\La}(T_2)\right\rangle dt\\
			&=-\sum_{j=1}^{n}\sum_{\ell=1}^{3}\int_{0}^{\infty}\left\langle\dc_{j,\ell}e^{-\xi\La}(T_1),\dc_{j,\ell}e^{-t\La}(T_2)\right\rangle dt\\
			&=-\sum_{j=1}^{n}\sum_{\ell=1}^{3}\int_{0}^{\infty}\left\langle\partial_{j,\ell}e^{-\xi\La}(T_1),\partial_{j,\ell}e^{-t\La}(T_2)\right\rangle dt\\
			&=-\sum_{j=1}^{n}\sum_{\ell=1}^{3}\int_{0}^{\infty} e^{-(\xi+t)}\left\langle e^{-\xi\La}\partial_{j,\ell}(T_1),e^{-t\La}\partial_{j,\ell}(T_2)\right\rangle dt\leq 0,
		\end{split}
	\end{equation}
	where the last inequality is due to the fact that $T_1, T_2\in \M_{2^{n}}$ are monotonically increasing and the semigroup $\{e^{-t\La}\}_{t\geq 0}$ is positive preserving.
\end{proof}

To prove Theorem \ref{quantum-ta}, we need to prove the following proposition. We will provide its proof after we prove Theorem \ref{quantum-ta}.

\begin{proposition}\label{tala-correlation-pro1}  For each $2\leq d\in [n]$, suppose that monotonically increasing elements $T_1$ and $T_2\in \M_{2^{n}}$ satisfy condition \eqref{weak-similar}, $\max\{\|T_1\|_{L_{\infty}},\|T_2\|_{L_{\infty}}\}\leq 1$ and $\mathcal{W}_1(T_1,T_2)\leq e^{-2(d-1)}$.  Then we have
	$$|\mathcal{W}_d(T_1,T_2)|\leq \frac{26e}{d}\Big(\frac{4e}{d-1}\Big)^{d-1}\mathcal{W}_1(T_1,T_2)\Big(\log(\frac{d}{\mathcal{W}_1(T_1,T_2)})\Big)^{d-1}.$$
\end{proposition}


We here provide the proof of  Theorem \ref{quantum-ta}.

\begin{proof}[Proof of Theorem \ref{quantum-ta}]
Take monotonically increasing elements $T_1$, $T_2\in \M_{2^{n}}$. Choose specific $\xi>0$ such that
\[
\rho\coloneqq e^{-\xi}=\frac{\epsilon}{\log(e/\mathcal{W}_1(T_1,T_2))}
\]
with $0<\epsilon<\frac{1}{208e^{2}}$. According to Lemma 
\ref{h-xi}, 
$$\mathrm{Cov}(T_1,T_2)\geq \mathrm{Cov}\left(e^{-\xi\La}(T_1),T_2\right).$$
Hence, it suffices to estimate $\mathrm{Cov}\left(e^{-\xi\La}(T_1),T_2\right)$.
By Lemma \ref{cov-semigroup-0-infty}, we have
	\begin{equation*}\label{covariance}
		\begin{aligned}
			\mathrm{Cov}\left(e^{-\xi\La}(T_1),T_2\right)&=\sum_{j=1}^n\int_{0}^{\infty}\left\langle \dc_{j}e^{-\xi\La}(T_1),e^{-t\La}\dc_{j}(T_2)\right\rangle dt\\
			&=\sum_{j=1}^{n}\int_{0}^{\infty}\left\langle \dc_{j}\left(e^{-\xi\La}\left(\mathrm{Rad}^{=1}(T_1)\right)\right),e^{-t\La}\left(\dc_{j}(\mathrm{Rad}^{=1}(T_2))\right)\right\rangle dt\\
			&\quad +\sum_{j=1}^{n}\int_{0}^{\infty}\left\langle \dc_{j}\left(e^{-\xi\La}\left(\mathrm{Rad}^{\geq 2}(T_1)\right)\right),e^{-t\La}\left(\dc_{j}(\mathrm{Rad}^{\geq 2}(T_2))\right)\right\rangle dt\\
			&\coloneqq \mathbf{A}+\mathbf{B}.
		\end{aligned}
	\end{equation*}
	We claim that 
	\begin{equation}\label{claim-a}
		\mathbf{A}=\rho \mathcal{W}_1(T_1,T_2)
	\end{equation}
	and 
\begin{equation}\label{claim-b}
	\mathbf{B}\leq \frac{\rho}{2} \mathcal{W}_1(T_1,T_2).
\end{equation}	
Once these claims have been proved, we deduce that
\[
\mathrm{Cov}(T_1,T_2)\geq \mathrm{Cov}\left(e^{-\xi\La}(T_1),T_2\right)\geq \big(\frac{3\epsilon}{2}\big)\frac{\mathcal{W}_1(T_1,T_2)}{\log\Big(e/\mathcal{W}_1(T_1,T_2)\Big)},
\]
which is the desired inequality.


	
	We next prove the claims \eqref{claim-a} and \eqref{claim-b} respectively.



	
	{\bf Estimate of $\mathbf{A}$}. For the term $\mathbf{A}$, by Lemma \ref{monotone-lemma-1} and Lemma \ref{l1-norm}, we have
	\begin{align*}
			\mathbf{A}
			&=\sum_{j=1}^{n}\int_{0}^{\infty}\left\langle \dc_{j}\left(e^{-\xi\La}\left(\mathrm{Rad}^{=1}(T_1)\right)\right),e^{-t\La}\left(\dc_{j}(\mathrm{Rad}^{=1}(T_2))\right)\right\rangle dt\\
			&=\sum_{j=1}^{n}e^{-\xi}\left(\int_{0}^{\infty}e^{-t}dt\right)\left\langle \dc_{j}\mathrm{Rad}^{=1}(T_1),\dc_{j}\mathrm{Rad}^{=1}(T_2)\right\rangle\\
			&=\sum_{j=1}^{n}\sum_{\ell=1}^{3}e^{-\xi}\left\langle \dc_{j,\ell}\mathrm{Rad}^{=1}(T_1),\dc_{j,\ell}\mathrm{Rad}^{=1}(T_2)\right\rangle\\
            &=\sum_{j=1}^{n}\sum_{\ell=1}^{3}e^{-\xi}\left\langle \partial_{j,\ell}\mathrm{Rad}^{=1}(T_1),\partial_{j,\ell}\mathrm{Rad}^{=1}(T_2)\right\rangle\\
			&=e^{-\xi}\mathcal{W}_1(T_1,T_2)=\rho \mathcal{W}_1(T_1,T_2).
		\end{align*}
This verifies the claim \eqref{claim-a}.

	
	{\bf Estimate of $\mathbf{B}$}. For the term $\mathbf{B}$, by the orthogonality of \(\{\sigma_{\s}\}_{\s}\), we have
\begin{align*}
\mathbf{B}&=\sum_{j=1}^{n}\int_{0}^{\infty}\left\langle\dc_{j}\mathrm{Rad}^{\geq 2}\left(e^{-\xi\La}(T_1)\right),e^{-t\La}\left(\dc_{j}\mathrm{Rad}^{\geq 2}(T_2)\right)\right\rangle dt\\
			&=\sum_{j=1}^{n}\sum_{\substack{|\s|\geq 2\\\s_{j\neq 0}}}\sum_{\substack{|\mathbf{w}|\geq 2\\\mathbf w_{j}\neq 0}}\int_{0}^{\infty}e^{-\xi|\s|}e^{-t|\mathbf{w}|}\widehat{T_1}(\s)\widehat{T_2}(\mathbf{w})\langle\sigma_{\s},\sigma_{\mathbf{w}}\rangle dt\\
			&=\sum_{j=1}^{n}\sum_{\substack{|\s|\geq 2\\\s_{j\neq 0}}}e^{-\xi|\s|}\left(\int_{0}^{\infty}e^{-t|\s|}dt\right)\widehat{T_1}(\s)\widehat{T_2}(\s)\\
&=\sum_{j=1}^{n}\sum_{\substack{|\s|\geq 2\\\s_{j\neq 0}}}\frac{e^{-\xi|\s|}}{|\s|}\widehat{T_1}(\s)\widehat{T_2}(\s).
\end{align*}
 After interchanging the order of summation, we write
	\begin{equation}\label{split-term-1}
		\begin{split}
\mathbf{B}
        &=\sum_{|\s|\geq 2}|\s|\frac{e^{-\xi|\s|}}{|\s|}\widehat{T_1}(\s)\widehat{T_2}(\s)\\
			&=\sum_{\ell:=|\s|=2}^{n}\rho^{\ell}\mathcal{W}_{\ell}(T_1,T_2)\\
			&=\sum_{\ell=2}^{\frac 12\log\left(\frac{e}{\mathcal{W}_{1}(T_1,T_2)}\right)}\rho^{\ell}\mathcal{W}_{\ell}(T_1,T_2)+\sum_{\ell>\frac 12\log\left(\frac{e}{\mathcal{W}_{1}(T_1,T_2)}\right)}^{n}\rho^{\ell}\mathcal{W}_{\ell}(T_1,T_2)\\
			&\coloneqq \mathbf{B_{I}}+\mathbf{B_{II}},
		\end{split}
	\end{equation}
	where $\mathcal{W}_{\ell}(T_1,T_2)\coloneqq \sum_{|\s|=\ell}\widehat{T_1}(\s)\widehat{T_2}(\s)$.
	
	We  assume from now on that $\frac{1}{2}\log\left(\frac{e}{\mathcal{W}_{1}(T_1,T_2)}\right)\geq 2$, otherwise the desired inequality is trivial. By this assumption, we have
	\begin{equation}\label{upper-bound-to-main-1}
		0<\mathcal{W}_{1}(T_1,T_2)\leq e^{-3}.
	\end{equation}
	
We proceed to estimate $\mathbf{B_{I}}$. In this case, $2\leq \ell\leq \frac 12\log\left(\frac{e}{\mathcal{W}_{1}(T_1,T_2)}\right)$, and hence, $\mathcal{W}_{1}(T_1,T_2)\leq e^{-2\ell+1}\leq e^{-2(\ell-1)}$. 
Note that for each $\ell \geq 2$, we have
\begin{equation}\label{basic-estimate-1}
	\left(\frac{4e}{\ell-1}\right)\log\left(\frac{\ell}{\mathcal{W}_{1}(T_1,T_2)}\right)\leq 4e\cdot\log\left(\frac{e}{\mathcal{W}_{1}(T_1,T_2)}\right).
\end{equation}
Indeed, rearranging \eqref{basic-estimate-1}, this is equivalent to
\[
\ell \mathcal{W}_{1}(T_1,T_2)^{\ell-2}\leq e^{\ell-1},
\]
which is trivial since  by \eqref{upper-bound-to-main-1},  we have
\[
\ell \mathcal{W}_{1}(T_1,T_2)^{\ell-2}\leq \ell e^{-3(\ell-2)}\leq \ell\leq e^{\ell-1},
\quad \ell\geq 2.\]
Combining Proposition \ref{tala-correlation-pro1} and \eqref{basic-estimate-1}, we get
\begin{equation*}\label{estimate-on-b1-2}
	\begin{split}
		|\mathbf{B_{I}}|&\leq \sum_{\ell=2}^{\frac 12\log\left(\frac{e}{\mathcal{W}_{1}(T_1,T_2)}\right)}\rho^{\ell}\cdot\frac{26e}{\ell}\left(\frac{4e}{\ell-1}\right)^{\ell-1}\mathcal{W}_{1}(T_1,T_2)\left(\log\left(\frac{\ell}{\mathcal{W}_{1}(T_1,T_2)}\right)\right)^{\ell-1}\\
&\leq \sum_{\ell=2}^{\frac 12\log\left(\frac{e}{\mathcal{W}_{1}(T_1,T_2)}\right)}\rho^{\ell}\cdot\frac{26e}{\ell}({4e})^{\ell-1}\left(\log\left(\frac{e}{\mathcal{W}_{1}(T_1,T_2)}\right)\right)^{\ell-1}\mathcal{W}_{1}(T_1,T_2)\\
		&= \rho \mathcal{W}_{1}(T_1,T_2)\sum_{\ell=2}^{\frac 12\log\left(\frac{e}{\mathcal{W}_{1}(T_1,T_2)}\right)}\epsilon^{\ell-1}\left(\frac{26e}{\ell}\right)(4e)^{\ell-1}.
	\end{split}
\end{equation*}
Note that for  $0<\epsilon<\frac{1}{208e^{2}}$, we have
\begin{equation*}
	\epsilon^{\ell-1}\left(\frac{26e}{\ell}\right)(4e)^{\ell-1}\leq 2^{-\ell},\quad \ell\geq 2.
\end{equation*}
Thus, we further have
\begin{equation}\label{estimate-on-b1-3}
	|\mathbf{B_{I}}|\leq \rho \mathcal{W}_{1}(T_1,T_2)\sum_{\ell=2}^{\frac 12\log\left(\frac{e}{\mathcal{W}_{1}(T_1,T_2)}\right)}2^{-\ell}.
\end{equation}

It remains to estimate $\mathbf{B_{II}}$. Note here that (see also \eqref{upper-bound-1}) 
	\begin{equation}\label{estimate-on-b2-1}
		\begin{split}
			|\mathbf{B_{II}}|&\leq \sum_{\ell>\frac 12\log\left(\frac{e}{\mathcal{W}_1(T_1,T_2)}\right)}\rho^{\ell}|\mathcal{W}_{\ell}(T_1,T_2)|\\
			&\leq \sum_{\ell>\frac 12\log\left(\frac{e}{\mathcal{W}_1(T_1,T_2)}\right)}\rho^{\ell}\\
			&=\rho\sum_{\ell>\frac 12\log\left(\frac{e}{\mathcal{W}_1(T_1,T_2)}\right)} \left(\frac{\epsilon^{\ell-1}}{\left(\log\frac{e}{\mathcal{W}_1(T_1,T_2)}\right)^{\ell-1}}\right).
		\end{split}
	\end{equation}
We now show that
	\begin{equation}\label{basic-estimate-3}
		\left(\log\frac{e}{\mathcal{W}_1(T_1,T_2)}\right)^{-(\ell-1)}\leq e^2\mathcal{W}_1(T_1,T_2),
	\end{equation}
	whenever $\ell>\frac 12\log\left(\frac{e}{\mathcal{W}_1(T_1,T_2)}\right)$ and $0<\mathcal{W}_1(T_1,T_2)\leq e^{-3}$. Indeed, changing the variable $\mathcal{W}_1(T_1,T_2)\mapsto e^{-u}$ with $u\geq 3$, \eqref{basic-estimate-3} is equivalent to
$$		1\leq e^{2-u}(1+u)^{\ell-1},$$
whenever $u\geq 3$ and $\ell> \frac{u+1}{2}\geq 2$. Since $e^{2-u}(1+u)^{\frac{u-1}{2}}\leq e^{2-u}(1+u)^{\ell-1}$ for $\ell-1\geq \frac{u-1}{2}$. It now remains to show that
	\begin{equation}\label{basic-estimate-5}
		1\leq e^{2-u}(1+u)^{\frac{u-1}{2}},
	\end{equation}
	for $u\geq 3$. Rearranging \eqref{basic-estimate-5} yields that the desired inequality is equivalent to
	\begin{equation}\label{basic-estimate-6}
		u-\left(\frac{u-1}{2}\right)\log(1+u)\leq 2,
	\end{equation}
	for $u\geq 3$. We define $F(u)\coloneqq u-\left(\frac{u-1}{2}\right)\log(1+u)$. We now show that $F(u)\leq 2$ for $u\geq 3$.
For $3\leq u\leq 4$, we have
\[
F(u)\leq
u-\frac{u-1}{2}\log 4\leq
4-\frac{3}{2}\log 4
<2.
\]
For $u\geq 4$, we compute
\[
F'(u)
=
\frac{1}{2}
-\frac{1}{2}\log(1+u)
+\frac{1}{1+u}
\]
and
\[
F''(u)
=
-\frac{1}{2(1+u)}
-\frac{1}{(1+u)^2}
<0.
\]
Since $F'(4)<0$, it follows that $F$ is decreasing on
$[4,\infty)$. Hence,
\[
F(u)\leq F(4)<2,
\qquad u\geq 4.
\]
This proves \eqref{basic-estimate-6}, and therefore
\eqref{basic-estimate-3}.
 Substituting \eqref{basic-estimate-3} to \eqref{estimate-on-b2-1} yields that
	\begin{equation}\label{estimate-on-b2-2}
	\begin{split}
|\mathbf{B_{II}}|&\leq e^2\mathcal{W}_1(T_1,T_2)\rho\sum_{\ell>\frac 12\log\left(\frac{e}{\mathcal{W}_1(T_1,T_2)}\right)}\epsilon^{\ell-1} \\
&\leq \rho \mathcal{W}_1(T_1,T_2)\sum_{\ell>\frac 12\log\left(\frac{e}{\mathcal{W}_1(T_1,T_2)}\right)}2^{-\ell},
\end{split}	
	\end{equation}
where we used the observation
\[
e^2\epsilon^{\ell-1}\leq2^{-\ell},
\quad \mbox{for}\quad  \ell>\frac 12\log\left(\frac{e}{\mathcal{W}_1(T_1,T_2)}\right)\geq2.
\]
Indeed, this observation holds for $\ell=2$, since
\[
e^2\epsilon<\frac1{208}<\frac14,
\]
and the general case follows by induction because
$\epsilon<\frac{1}{208e^{2}}<1/2$.	  Combining \eqref{estimate-on-b2-2} and \eqref{estimate-on-b1-3}, we arrive at \eqref{claim-b}. 
 The proof of the theorem is complete.
\end{proof}


It remains to prove Proposition \ref{tala-correlation-pro1}. To this end, we first establish another useful result in Proposition \ref{tala-correlation-pro2}. 
For each $T_1, T_2\in \mathbb{M}_{2^n}$, $J\subset [n]$, $d\geq2$ and $t>0$, set  
\[L_{J,d,t}(T_1,T_2)=\left\{j\in J:\sum_{\substack{\supp(\mathbf{s})\subset J^c \\ |\mathbf{s}|=d-1,\alpha\in\{1,2,3\}}}\left|\widehat{T_1}(\s\oplus \alpha e_j)\widehat{T_2}(\s\oplus \alpha e_j)\right|>t\sum_{\ell=1}^3\left|\widehat{T_1}(\ell e_j)\widehat{T_2}(\ell e_j)\right|\right\}.\]
If $T_1,T_2$ are  monotonically increasing operators  and satisfy condition \eqref{weak-similar},  then, by Lemma \ref{basic} (ii) and by  \eqref{weak-similar}, for each $\ell\in\{1,2,3\}$ and $j\in [n]$, $\widehat{T_1}(\ell e_j),\widehat{T_2}(\ell e_j)\geq0$, and we have 
$$\sum_{\ell=1}^3\widehat{T_1}(\ell e_j)\widehat{T_2}(\ell e_j)=\big(\sum_{\ell=1}^3\widehat{T_1}(\ell e_j)^2\big)^{\frac{1}{2}}\cdot\big(\sum_{\ell=1}^3\widehat{T_2}(\ell e_j)^2\big)^{\frac{1}{2}},$$ 
and consequently, by the Cauchy--Schwarz inequality, we have
\begin{equation}\label{lj}
L_{J,d,t}(T_1,T_2)\subset L_{J,d,t}(T_1,T_1)\cup L_{J,d,t}(T_2,T_2).
\end{equation}

We need the following lemma. 

\begin{lemma}\label{tala-correlation-lemma1-lemma}
Let $T\in \mathbb{M}_{2^n}$ be monotonically increasing with $\|T\|_{L_{\infty}}\leq 1$.	For each   $J\subset [n]$, $d\geq2$ and $t\geq 12(4e)^{d-1}$, we have
	\begin{equation}\label{tala-correlation-lemma1}
		\sum_{j\in L_{J,d,t}(T,T)}\sum_{\ell=1}^3\widehat{T}(\ell e_j)^2\leq \frac{5e}{3}(\frac{t}{12})^{-\frac{1}{d-1}}\exp\left\{-\frac{(d-1)(\frac{t}{12})^\frac{1}{d-1}}{2e}\right\}.
	\end{equation} 
\end{lemma}

\begin{proof}
Take monotonically increasing $T\in \mathbb{M}_{2^n}$,  $J\subset [n]$, $d\geq2$ and $t\geq 12(4e)^{d-1}$. During the proof, let us write $L_t^T$ instead of $ L_{J,d,t}(T,T)$.
Applying Lemma \ref{basic} (ii), for each $j\in [n]$, we have
$$\|\dc_j(T)\|_{L_1}\leq \sum_{\ell=1}^3\|\dc_{j,\ell}(T)\|_{L_1}=\sum_{\ell=1}^3\widehat{T}(\ell e_{j}),$$
which further implies 
$$\|\dc_j(T)\|_{L_1}^2\leq 3 \sum_{\ell=1}^3\widehat{T}(\ell e_{j})^2.$$
Thus, we have
	\begin{equation*}\label{j-t2-1}
		\begin{aligned}
	\sum_{j\in J}\sum_{\substack{\supp(\mathbf{s})\subset J^c \\ |\mathbf{s}|=d-1,\alpha\in\{1,2,3\}}}&|\widehat{T}(\s\oplus \alpha e_j)|^2\\ &\leq2t_0^2\sum_{j\in J}\|\dc_j(T)\|_{L_1}^2+10et_0^{2-\frac{2}{d-1}}\exp\left\{-\frac{(d-1)t_0^{\frac{2}{d-1}}}{2e}\right\}\\
	&\leq 6t_0^2\sum_{j\in J}\sum_{\ell=1}^3\widehat{T}(\ell e_{j})^2+10et_0^{2-\frac{2}{d-1}}\exp\left\{-\frac{(d-1)t_0^{\frac{2}{d-1}}}{2e}\right\},
		\end{aligned}
	\end{equation*}
where the first inequality is taken from the proof of \cite[Proposition 5.2]{JLLZ2025} (provided $t_0\geq (4e)^{\frac{d-1}{2}}$).
	According to the definition of $L_t^T$,  taking $t_0=\frac{\sqrt{t}}{2\sqrt{3}}$,  it follows that
	\begin{align*}
		t\sum_{j\in L_t^{T}}\sum_{\ell=1}^3\widehat{T}(le_{j})^2&\leq \sum_{j\in L_t^{T}}\sum_{\substack{\supp(\mathbf{s})\subset J^c \\ |\mathbf{s}|=d-1,\alpha\in\{1,2,3\}}}|\widehat{T}(\s\oplus \alpha e_j)|^2\\
		&\leq\frac{t}{2}\sum_{j\in J}\sum_{\ell=1}^3\widehat{T}(\ell e_{j})^2+10e(\frac{t}{12})^{1-\frac{1}{d-1}}\exp\left\{-\frac{(d-1)(\frac{t}{12})^{\frac{1}{d-1}}}{2e}\right\}.
	\end{align*}
From the above inequality,  by subtracting the first term on the right-hand side from the left-hand side and subsequently dividing both sides by 
$\frac{t}{2}$, we have
\begin{align*}
		\sum_{j\in L_t^{T}}\sum_{\ell=1}^3\widehat{T}(\ell e_{j})^2&\leq \sum_{j\in L_t^{T}}\sum_{\substack{\supp(\mathbf{s})\subset J^c \\ |\mathbf{s}|=d-1,\alpha\in\{1,2,3\}}}|\widehat{T}(\s\oplus \alpha e_j)|^2\\
		&\leq\frac{2}{t}\cdot10e(\frac{t}{12})^{1-\frac{1}{d-1}}\exp\left\{-\frac{(d-1)(\frac{t}{12})^{\frac{1}{d-1}}}{2e}\right\}\\
        &=\frac{5e}{3}(\frac{t}{12})^{-\frac{1}{d-1}}\exp\left\{-\frac{(d-1)(\frac{t}{12})^\frac{1}{d-1}}{2e}\right\}.
	\end{align*}
The proof is complete.
\end{proof}



\begin{proposition}\label{tala-correlation-pro2} Let $T_1,T_2\in\mathbb{M}_{2^n}$ be  monotonically increasing operators satisfying the condition \eqref{weak-similar},  $\max\{\|T_1\|_{L_\infty},\|T_2\|_{L_\infty}\}\leq1$ and $\mathcal{W}_1(T_1,T_2)\leq e^{-2(d-1)}$. Then, for each $J\subset [n]$ and  $d\geq2$, we have
	$$\sum_{j\in J}\sum_{\substack{\supp(\mathbf{s})\subset J^c \\ |\mathbf{s}|=d-1,\alpha\in\{1,2,3\}}}\left|\widehat{T_1}(\s\oplus \alpha e_j)\widehat{T_2}(\s\oplus \alpha e_j)\right|\leq 26\left(\frac{4e}{d-1}\right)^{d-1}\mathcal{W}_{1,J}(T_1,T_2)\left(\log\left(\frac{1}{\mathcal{W}_{1,J}(T_1,T_2)}\right)\right)^{d-1}.$$
\end{proposition}
\begin{proof} 
	We divide the proof into three steps. We write for convenience $L_t$ (resp. $L_t^{T_1}$, $L_t^{T_2}$) instead of $L_{J,d,t}(T_1,T_2)$ (resp.  $L_{J,d,t}(T_1,T_1)$, $L_{J,d,t}(T_2,T_2)$).
	
	
	\textbf{Step I.} For each $j\in J$, we define the following function on $J$ by setting
	$$f_1(j)=\sum_{\ell=1}^3\widehat{T_1}(\ell e_j)\widehat{T_2}(\ell e_j).$$
We further introduce the function $f_2(\cdot)$ on $J$ by setting
	$$\sum_{\substack{\supp(\mathbf{s})\subset J^c \\ |\mathbf{s}|=d-1,\alpha\in\{1,2,3\}}}\left|\widehat{T_1}(\s\oplus \alpha e_j)\widehat{T_2}(\s\oplus \alpha e_j)\right|=f_1(j)f_2(j).$$
	By  \cite[Lemma 11]{KK2013}, we get
	\begin{align*}
		\sum_{j\in J}\sum_{\substack{\supp(\mathbf{s})\subset J^c \\ |\mathbf{s}|=d-1,\alpha\in\{1,2,3\}}}&\left|\widehat{T_1}(\s\oplus \alpha e_j)\widehat{T_2}(\s\oplus \alpha e_j)\right|=\sum_{j\in J} f_1(j)f_2(j)\\
		&=\int_0^\infty\big(\mathbb{E}_J\big(f_1(j)\cdot\chi_{\{j:f_2(j)>t\}}(j)\big)\big)dt=\int_0^\infty \sum_{j\in L_t}\sum_{\ell=1}^3\widehat{T_1}(\ell e_j)\widehat{T_2}(\ell e_j)dt\\
		&=\int_{0}^{t_1} \sum_{j\in L_t}\sum_{\ell=1}^3\widehat{T_1}(\ell e_j)\widehat{T_2}(\ell e_j)dt
		+\int_{t_1}^{\infty} \sum_{j\in L_t}\sum_{\ell=1}^3\widehat{T_1}(\ell e_j)\widehat{T_2}(\ell e_j)dt\\
		&:=Y_{1,j}+Y_{2,j},
	\end{align*}
	where $t_1$ is chosen to satisfy
	\begin{equation}\label{t1}
		\mathcal{W}_{1,J}(T_1,T_2)=\exp\{-\frac{(d-1)(\frac{t_1}{12})^{\frac{1}{d-1}}}{4e}\}.
	\end{equation}
	Note that the assumption $\mathcal{W}_{1,J}(T_1,T_2)\leq \mathcal{W}_1(T_1,T_2)\leq e^{-2(d-1)}$ and  the above equation  imply that
	$$t_1\geq 12(8e)^{d-1}\quad\mbox{and \quad $t_1=12(\frac{4e}{d-1})^{d-1}(\log\frac{1}{\mathcal{W}_{1,J}(T_1,T_2)})^{d-1}$}.$$
	Then,
	\begin{equation}\label{y1j}
		Y_{1,j}\leq t_{1}\mathcal{W}_{1,J}(T_1,T_2)=12(\frac{4e}{d-1})^{d-1}\mathcal{W}_{1,J}(T_1,T_2)\left(\log\frac{1}{\mathcal{W}_{1,J}(T_1,T_2)}\right)^{d-1}.
	\end{equation}
	
	

	
	
	
	\textbf{Step II.} 
Since $T_1,T_2\in\mathbb{M}_{2^n}$ are  monotonically increasing and  satisfy the condition \eqref{weak-similar}, we can apply \eqref{lj} to get 
	\begin{equation}\label{lj-2}
		\sum_{j\in L_t}\sum_{\ell=1}^3\widehat{T_1}(\ell e_j)\widehat{T_2}(\ell e_j)\leq \sum_{j\in L_t^{T_1}}\sum_{\ell=1}^3\widehat{T_1}(\ell e_j)\widehat{T_2}(\ell e_j)+\sum_{j\in L_t^{T_2}}\sum_{\ell=1}^3\widehat{T_1}(\ell e_j)\widehat{T_2}(\ell e_j).
	\end{equation}
	By Cauchy-Schwarz inequality and \eqref{upper-bound-1},  we can get $\mathcal{W}_1(T_2)\leq1$ and
	\begin{align*}
		\sum_{j\in L_{t}^{T_1}}\sum_{\ell=1}^3\widehat{T_1}(\ell e_j)\widehat{T_2}(\ell e_j)&\leq\big(\sum_{j\in L_{t}^{T_1}}\sum_{\ell=1}^3\widehat{T_1}(\ell e_j)^2\big)^\frac{1}{2}\big(\sum_{j\in L_{t}^{T_1}}\sum_{\ell=1}^3\widehat{T_2}(\ell e_j)^2\big)^\frac{1}{2}\\
		&\leq \big(\sum_{j\in L_{t}^{T_1}}\sum_{\ell=1}^3\widehat{T_1}(\ell e_j)^2\big)^\frac{1}{2}.
	\end{align*}
Similarly,
	$$\sum_{j\in L_{t}^{T_2}}\sum_{\ell=1}^3\widehat{T_1}(\ell e_j)\widehat{T_2}(\ell e_j)\leq \big(\sum_{j\in L_{t}^{T_2}}\sum_{\ell=1}^3\widehat{T_2}(\ell e_j)^2\big)^\frac{1}{2}.$$
Hence, combining the above argument,   \eqref{lj-2} and \eqref{tala-correlation-lemma1}, for all $t\geq 12(4e)^{d-1}$, we have
	\begin{equation}\label{tala-correlation-lemma2}
		\begin{split}
			\sum_{j\in L_t}\sum_{\ell=1}^3\widehat{T_1}(\ell e_j)\widehat{T_2}(\ell e_j)&\leq \big(\sum_{j\in L_{t}^{T_1}}\sum_{\ell=1}^3\widehat{T_1}(\ell e_j)^2\big)^\frac{1}{2}+\big(\sum_{j\in L_{t}^{T_2}}\sum_{\ell=1}^3\widehat{T_2}(\ell e_j)^2\big)^\frac{1}{2}\\
			&\leq 2\cdot (\frac{5e}{3})^{\frac{1}{2}}(\frac{t}{12})^{-\frac{1}{2(d-1)}}\exp\left\{-\frac{(d-1)\left(\frac{t}{12}\right)^\frac{1}{d-1}}{4e}\right\}.
		\end{split}
	\end{equation}
	
	
	
	\textbf{Step III.} For $Y_{2,j}$ and all $t\geq t_1\geq 12(8e)^{d-1}\geq12(4e)^{d-1}$, by the above estimate \eqref{tala-correlation-lemma2}, we have
$$Y_{2,j}\leq \int_{t_1}^{\infty} 2\cdot \left(\frac{5e}{3}\right)^{\frac{1}{2}}\left(\frac{t}{12}\right)^{-\frac{1}{2(d-1)}}\exp\left\{-\frac{(d-1)\left(\frac{t}{12}\right)^\frac{1}{d-1}}{4e}\right\}dt.$$
To estimate the above inequality, we first apply a change of variables, setting 
\begin{equation}\label{s}
s=\frac{(d-1)\left(\frac{t}{12}\right)^\frac{1}{d-1}}{4e},
\end{equation}
obviously, $s\geq 2(d-1)$, and we obtain
$$Y_{2,j}=2\times 12\cdot \left(\frac{5e}{3}\right)^{\frac{1}{2}}\Big(\frac{4e}{d-1}\Big)^{d-\frac{3}{2}}(d-1)\int_{s=s_1}^{\infty} s^{d-\frac{5}{2}}e^{-s}ds,$$
where $s_1$ is derived by substituting  $t_1$ into \eqref{s}. If we denote $\varphi(s)=s^{d-\frac{5}{2}}e^{-s}$, then $\varphi(s)$ is decreasing and $\frac{\varphi(s+1)}{\varphi(s)}\leq e^{-\frac{1}{2}}$ on $s\geq 2(d-1)$ and $d\geq 2$, and we have
$$\int_{s=s_1}^{\infty} s^{d-\frac{5}{2}}e^{-s}ds\leq \frac{1}{1-e^{-1/2}} s_1^{d-\frac{5}{2}}e^{-s_1}.$$ 
It therefore follows that
\begin{align*}	
Y_{2,j}&\leq 
2\times 12\cdot \left(\frac{5e}{3}\right)^{\frac{1}{2}}\times 4e\cdot\frac{1}{1-e^{-1/2}}\cdot \left(\frac{t_1}{12}\right)^{1-\frac{3}{2(d-1)}}\cdot\exp\left\{-\frac{(d-1)(\frac{t_1}{12})^\frac{1}{d-1}}{4e}\right\}\\
		&\leq 2\times 12\cdot \left(\frac{5e}{3}\right)^{\frac{1}{2}}\times 4e\cdot\frac{1}{1-e^{-1/2}}\cdot (8e)^{-3/2} \left(\frac{t_1}{12}\right)\cdot\exp\left\{-\frac{(d-1)(\frac{t_1}{12})^\frac{1}{d-1}}{4e}\right\}\\
		&\leq 14 \left(\frac{4e}{d-1}\right)^{d-1}\mathcal{W}_{1,J}(T_1,T_2)\left(\log\frac{1}{\mathcal{W}_{1,J}(T_1,T_2)}\right)^{d-1},
\end{align*}
where the last inequality follows from the choice of $t_1$ in \eqref{t1}. Combining \eqref{y1j} with the above estimate, we finish the proof.
\end{proof}


We now prove Proposition \ref{tala-correlation-pro1}.


\begin{proof}[Proof of Proposition \ref{tala-correlation-pro1}]
For each $\s\in \{0,1,2,3\}^{n}$ with $|\s|=d$, we know the probability of $J$ such that $\s=\mathbf{u}\oplus \alpha e_{j}$ with $j\in J$ and $\supp(\mathbf{u})\subseteq J^{c}$ is $\binom{d}{1}\cdot\frac{1}{d}(1-\frac{1}{d})^{d-1}=(1-\frac{1}{d})^{d-1}\geq \frac{1}{e}$. Hence, 
\begin{equation}\label{wd-1}
\begin{aligned}
&\E_{J}\left(\sum_{\substack{\supp(\s)\subseteq J^{c}, |\s|=d-1\\j\in J,\alpha \in\{1,2,3\}}}\left|\widehat{T_1}\left(\s\oplus \alpha e_j\right)\widehat{T_2}\left(\s\oplus \alpha e_j\right)\right|\right)\\
=&\left(1-\frac{1}{d}\right)^{d-1}\sum_{\substack{\s\in\{0,1,2,3\}^{n},\\|\s|=d}}\left|\widehat{T_1}(\s )\widehat{T_2}(\s)\right|\\
\geq&\frac{1}{e}\sum_{\substack{\s\in\{0,1,2,3\}^{n},\\|\s|=d}}\left|\widehat{T_1}(\s)\widehat{T_2}(\s)\right|\geq \frac{1}{e}|\mathcal{W}_{d}(T_1,T_2)|.
\end{aligned}
\end{equation}
Secondly, Proposition \ref{tala-correlation-pro2}, together with the assumption $\mathcal{W}_1(T_1,T_2)\leq e^{-2(d-1)}$ and \eqref{wd-1}, yields that
  $$|\mathcal{W}_{d}(T_1,T_2)|\leq e\mathbb{E}_{J}\left[26\left(\frac{4e}{d-1}\right)^{d-1}\mathcal{W}_{1,J}(T_1,T_2)\left(\log\left(\frac{1}{\mathcal{W}_{1,J}(T_1,T_2)}\right)\right)^{d-1}\right].$$
Note that $x\mapsto x\big(\log(\frac{1}{x})\big)^{d-1}$ is a concave function on $(0,e^{-(d-1)})$. Then it follows from the above inequality and the Jensen inequality that
\begin{equation*}\label{wd-2}
\begin{aligned}
|\mathcal{W}_{d}(T_1,T_2)|\leq&26e\left(\frac{4e}{d-1}\right)^{d-1}\left(\mathbb{E}_{J}\left(\mathcal{W}_{1,J}(T_1,T_2)\right)\right)\left(\log\left(\frac{1}{\mathbb{E}_{J}(\mathcal{W}_{1,J}(T_1,T_2))}\right)\right)^{d-1}\\
=&26e\left(\frac{4e}{d-1}\right)^{d-1}\left(\frac{\mathcal{W}_1\left(T_1,T_2\right)}{d}\right)\left(\log\left(\frac{d}{\mathcal{W}_1\left(T_1,T_2\right)}\right)\right)^{d-1},
\end{aligned}
\end{equation*}
where we used $\mathbb{E}_{J}(\mathds{1}_{J}(j))=\frac{1}{d}$ in the last equality. Hence,  the desired inequality follows. 
\end{proof}




%\begin{proof} 
%	For each $j\in [n]$ and $T\in\mathbb{M}_{2^n}$, write
%	\begin{align*}
%		\|d_j(T)\|_{L_1(\mathbb{M}_{2^n})}&=\left\|\sum_{l=1}^3d_{j,l}(T)\right\|_{L_1}
%		\leq \sum_{l=1}^3\|d_{j,l}(T)\|_{L_1}\\
%		&=\sum_{l=1}^3\|\sigma_{le_j}d_{j,l}(T)\|_{L_1}
%		=\sum_{l=1}^3\mathrm{tr}(\sigma_{le_j}d_{j,l}(T))\\
%		&=\sum_{l=1}^3\sum_{\s:\s_j=l} \widehat{T}(\s) \mathrm{tr}(\sigma_{le_j}\sigma_{\s})
%		=\sum_{l=1}^3\widehat{T}(le_j).
%	\end{align*}
%By \eqref{weak-similar}, there exist $c_j>0$ such that
%$$\Big(\sum_{l=1}^3\widehat{T_1}(le_j)\Big)\Big(\sum_{l=1}^3\widehat{T_2}(le_j)\Big)=c_j\Big(\sum_{l=1}^3\widehat{T_2}(le_j)\Big)^2\leq 3c_j\sum_{l=1}^3\widehat{T_2}(le_j)^2=3\sum_{l=1}^3\widehat{T_1}(le_j)\widehat{T_2}(le_j)$$
%	Here we used the fact that $T$ is monotonically increasing. Thus, we have
%	\begin{align*}
%		\sum_{j=1}^n\|d_j(T_1)\|_{L_1}\|d_j(T_2)\|_{L_1}&\leq\sum_{j=1}^n\Big(\sum_{l=1}^3\widehat{T_1}(le_j)\Big)\Big(\sum_{l=1}^3\widehat{T_2}(le_j)\Big)\\
%		&\leq 3\sum_{j=1}^n\sum_{l=1}^3\widehat{T_1}(le_j)\widehat{T_2}(le_j)
%		\leq 3\mathrm{Cov}(T_1,T_2).
%	\end{align*}
%	Note that for each $T_1,T_2\in\mathbb{M}_{2^n}$, we have
%	$$\mathrm{Cov}(T_1,T_2)=\sum_{|\s|\neq0}\widehat{T_1}(\s)^*\widehat{T_2}(\s).$$
%	The proof is complete.
%\end{proof}
\section{Quantum KMS-correlation inequalities}
In this section, we aim to establish the quantum analogy  of the Keller-Mossel-Sen correlation inequality \eqref{kms-correlation}.  Recall the \emph{quantum reverse hypercontractivity inequality} as follows: Let $\{e^{-t\La}\}_{t\geq 0}$ be the quantum Ornstein-Uhlenbeck semigroup acting on $L_{p}(\M_{2^{n}})$, then the following reverse hypercontractivity of $\{e^{-t\La}\}_{t\geq 0}$ holds
\begin{equation}\label{reverse-hyper}
	\mathrm{tr}\left(|T|e^{-t\La}(|S|)\right)\geq \|T\|_{L_{p}}\|S\|_{L_{q}},
\end{equation}
whenever $e^{-2t}\leq (1-p)(1-q)$ with $(p,q)\in (0,1)\times (0,1)$. For further details, we refer the reader to \cite{BDR2020,CKMT2015}.


\begin{theorem}\label{quantum-kms}
Let monotonically increasing $T_1$, $T_2\in \M_{2^{n}}$ satisfy  
$$\max\{\|T_1\|_{L_{\infty}},\|T_2\|_{L_{\infty}}\}\leq 1.$$	There exists a universal constant $K>0$ such that the following inequality holds 
	\begin{equation}\label{correlation}
		\mathrm{Cov}(T_1,T_2)\geq K\sum_{j=1}^{n}\frac{\|\dc_{j}(T_1)\|_{L_{1}}\|\dc_{j}(T_2)\|_{L_{1}}}{\sqrt{\log(e/\|\dc_{j}(T_1)\|_{L_{1}})\log(e/\|\dc_{j}(T_2)\|_{L_{1}})}}.
	\end{equation}
\end{theorem}



The following lemma motivated by Keller, Mossel and Sen (\cite{KMS2014}).  
\begin{lemma}\label{technical-1}
Let $T\in \M_{2^{n}}$ with $\|T\|_{L_{\infty}}\leq 1$, and $0<p<\infty$. 
	\begin{enumerate}[{\rm(i)}]
		\item For each $p\geq 1$, we have 
		\[
		\left\|\partial_{j}\left(e^{-t\La}(T)\right)\right\|_{L_{p}}\leq 3^{\frac{p-1}{p}} t^{-(p-1)/2p}e^{-t/p}\|\partial_{j}(T)\|^{1/p}_{L_{1}},\quad t>0.
		\]
		\item If $T$ is monotonically increasing, then the following holds for every $0<p<1$,
		\[
		\left\|\partial_{j}\left(e^{-t\La}(T)\right)\right\|_{L_{p}}\geq 3^{\frac{p-1}{p}} t^{(1-p)/2p}e^{-t/p}\|\partial_{j}(T)\|^{1/p}_{L_{1}},\quad t>0.
		\]
	\end{enumerate}
\end{lemma}

\begin{proof} (i) Fix $p\geq 1$.  Note that
$$
\|\dc_{j}(e^{-t\La}(T))\|_{L_{\infty}}\leq t^{-1/2}\|T\|_{L_{\infty}}\leq t^{-1/2}
$$
by the Buser inequality \eqref{Q-Buser} with \(p=\infty\), the fact that \(e^{2t}-1\geq t\) for \(t>0\) and the assumption that \(\|T\|_{L_\infty }\leq 1\). Then by the above estimate and Lemma \ref{comparison} (i), we have
\begin{equation}\label{t-1-2}
\|\partial_{j}(e^{-t\La}(T))\|_{L_{\infty}}\leq3\|\dc_{j}(e^{-t\La}(T))\|_{L_{\infty}}\leq 3t^{-1/2}.
\end{equation}	
By this inequality and  Lemma \ref{comparison} (i) (ii),  we get
\begin{align*}
\left\|\partial_{j}\left(e^{-t\La}(T)\right)\right\|_{L_{p}}^p&=\mathrm{tr}(|\partial_{j}(e^{-t\La}(T))|^{p-1}|\partial_{j}(e^{-t\La}(T))|)\\
&\leq (3t^{-1/2})^{p-1}\|\partial_{j}(e^{-t\La}(T)\|_{L_{1}}\\
&=3^{p-1}t^{-(p-1)/2}e^{-t}\|e^{-t\La}(\partial_{j}(T))\|_{L_{1}}\\
&\leq 3^{p-1}t^{-(p-1)/2}e^{-t}\|\partial_{j}(T)\|_{L_{1}},\quad t>0.
\end{align*}
We finish the proof of (i).
	
	(ii) Fix $0<p<1$ and $t>0$. Since $T$ is monotonically increasing, it follows that 
	$$\|\partial_{j}(e^{-t\La}(T))\|_{L_{1}}=e^{-t}\mathrm{tr}\left(e^{-t\La}\left(\partial_{j}(T)\right)\right)=e^{-t}\|\partial_{j}(T)\|_{L_{1}},$$
where the first equality is due to Lemma \ref{comparison} (ii) and Lemma \ref{basic} (iii), and the second equality is due to that $e^{-t\mathcal{L}}$ is trace preserving.  As $0<p<1$,  by the elementary inequality and \eqref{t-1-2},
\[
\|\partial_{j}(e^{-t\La}(T))\|_{L_{1}}\leq  	\|\partial_{j}(e^{-t\La}(T))\|^{p}_{L_{p}}\|\partial_{j}(e^{-t\La}(T))\|_{L_{\infty}}^{1-p}\leq 3^{1-p}t^{-\frac{1-p}{2}} \|\partial_{j}(e^{-t\La}(T))\|^{p}_{L_{p}} .
\]
Thus, combining the above argument, we further have
\begin{align*}
	\|\partial_{j}(T)\|_{L_{1}}= e^{t} \|\partial_{j}(e^{-t\La}(T))\|_{L_{1}}\leq 3^{1-p}e^t t^{-\frac{1-p}{2}} \|\partial_{j}(e^{-t\La}(T))\|^{p}_{L_{p}} .
\end{align*}
This verifies item (ii).
\end{proof}


The lemma below was stated in \cite{KMS2014} without proof, we here provide its proof for the sake of completeness. 

\begin{lemma}\label{optimal}
	For each $t,s\in (0,1)$, we have
	\[
\frac{t^{2}+2st+1}{1-s^{2}}=	\inf_{\substack{s^{2}\leq (1-p)(1-q)\\(p,q)\in(0,1)^{2}}}\frac{t^{2}}{p}+\frac{1}{q}.
	\]
\end{lemma}
\begin{proof}
	On the one hand, for the range $s^{2}\leq (1-p)(1-q)$ with $(s,p,q)\in(0,1)^{3}$, since the function $s\mapsto \frac{t^{2}+2st+1}{1-s^{2}}$ is increasing, it follows that 
	\begin{equation}\label{main-optimal-1}
		\frac{t^{2}+2st+1}{1-s^{2}}\leq \frac{t^{2}+2\sqrt{(1-p)(1-q)}t+1}{1-(1-p)(1-q)}=\frac{t^{2}+2\sqrt{(1-p)(1-q)}t+1}{p+q-pq}.
	\end{equation}
	We claim here that
	\begin{equation}\label{main-optimal-2}
		\frac{t^{2}+2\sqrt{(1-p)(1-q)}t+1}{p+q-pq}\leq \frac{t^{2}}{p}+\frac{1}{q},\quad\mbox{for }(p,q)\in(0,1)^{2}.
	\end{equation}
Indeed,	rearranging \eqref{main-optimal-2} yields that
	\[
	q^{2}(1-p)t^{2}-2pq\sqrt{(1-p)(1-q)}t+p^{2}(1-q)=\left(q\sqrt{1-p}t-p\sqrt{1-q}\right)^{2}\geq 0,
	\]
	which verifies the claim. Combining \eqref{main-optimal-1} with \eqref{main-optimal-2} entails that
	\[
	\inf_{\substack{s^{2}\leq (1-p)(1-q)\\(p,q)\in(0,1)^{2}}}\frac{t^{2}}{p}+\frac{1}{q}\geq \frac{t^{2}+2st+1}{1-s^{2}}.
	\]
	
	On the other hand, for given $(s,t)\in (0,1)^{2}$, we define
	\[
	p_{0}\coloneqq \frac{t(1-s^{2})}{s+t}\quad\mbox{and}\quad q_{0}\coloneqq \frac{1-p_{0}-s^{2}}{1-p_{0}}
	\]
	It is clear that $(p_{0},q_{0})\in (0,1)^{2}$ with $s^{2}=(1-p_{0})(1-q_{0})$. Moreover, substituting the fact $s^{2}=(1-p_{0})(1-q_{0})$ into $F(s)\coloneqq\frac{t^{2}+2st+1}{1-s^{2}}$ yields that
	\[
	F([(1-p_{0})(1-q_{0})]^{\frac{1}{2}})=\frac{t^{2}}{p_{0}}+\frac{1}{q_{0}}\geq \inf_{\substack{s^{2}\leq (1-p)(1-q)\\(p,q)\in(0,1)^2}}\frac{t^{2}}{p}+\frac{1}{q}.
	\]
	Combining the two inequalities completes the proof.
\end{proof}


Now we are ready to prove Theorem \ref{quantum-kms}.  
\begin{proof}[Proof of Theorem \ref{quantum-kms}]
By Lemma \ref{jlzz-lem}, it follows from the assumption  that from  
\begin{equation}\label{dj T1 T2 1}
	\max\{\|\dc_{j}(T_1)\|_{L_{\infty}},\|\dc_{j}(T_2)\|_{L_{\infty}}\}\leq 1.
\end{equation}
for each $j\in [n]$. Applying the representation of covariance \eqref{key-identity} and Lemma \ref{comparison} (ii), we have
	\begin{equation*}\label{correlation-1}
		\begin{split}
			\mathrm{Cov}(T_1,T_2)&=\sum_{j=1}^{n}\int_{0}^{\infty}e^{-t}\mathrm{tr}\left(\partial_{j}(T_1)e^{-t\La}\left(\partial_{j}(T_2)\right)\right)dt\\
			&\geq \sum_{j=1}^{n}\int_{1}^{\infty}e^{-(t-1)}\mathrm{tr}\left[\partial_{j}\left(e^{-\La/2}(T_1)\right)\cdot e^{-(t-1)\La}\left(\partial_{j}\left(e^{-\La/2}(T_2)\right)\right)\right]dt\\
			&=\sum_{j=1}^{n}\int_{0}^{\infty}e^{-t}\mathrm{tr}\left[\partial_{j}\left(e^{-\La/2}(T_1)\right)\cdot e^{-t\La}\left(\partial_{j}\left(e^{-\La/2}(T_2)\right)\right)\right]dt.
		\end{split}
	\end{equation*}
	Changing the variable $e^{-t}\mapsto s$, we further have
	\begin{equation}\label{correlation-2}
	\begin{split}
\mathrm{Cov}(T_1,T_2)&\geq \sum_{j=1}^{n}\int_{0}^{1}\mathrm{tr}\left[\partial_{j}\left(e^{-\La/2}(T_1)\right)\cdot e^{\log(s)\La}\left(\partial_{j}\left(e^{-\La/2}(T_2)\right)\right)\right]ds\\
&=\sum_{j=1}^{n} \int_{0}^{1}A^{j}(s)ds.
\end{split}	
	\end{equation}
	
It suffices to estimate the integral $\int_{0}^{1}A^{j}(s)ds$. 	For each $j\in [n]$ and $s\in (0,1)$, by \eqref{reverse-hyper} and Lemma \ref{technical-1}, we have
	\begin{equation}\label{correlation-3}
\begin{split}
		A^{j}(s)&\geq \left\|\partial_{j}\left(e^{-\La/2}(T_1)\right)\right\|_{L_{p}}\cdot\left\|\partial_{j}\left(e^{-\La/2}(T_2)\right)\right\|_{L_{q}}\\
       &\geq2\cdot3^{\frac{p-1}{p}+\frac{q-1}{q}}\cdot (2e)^{-(\frac{1}{2p}+\frac{1}{2q})}\|\partial_{j}(T_1)\|^{1/p}_{L_{1}}\cdot\|\partial_{j}(T_2)\|^{1/q}_{L_{1}},
\end{split}
	\end{equation}
	for each $(p,q)\in (0,1)^{2}$ with $s^{2}\leq (1-p)(1-q)$. Applying Lemma \ref{basic} (i) to \eqref{correlation-3} yields that,
	\begin{equation}\label{correlation-4}
		A^{j}(s)\geq 18\cdot3^{-(\frac{1}{p}+\frac{1}{q})}\cdot(2e)^{-(\frac{1}{2p}+\frac{1}{2q})}\|\dc_{j}(T_1)\|^{1/p}_{L_{1}}\cdot\|\dc_{j}(T_2)\|^{1/q}_{L_{1}}.
	\end{equation}
	For each $j\in [n]$, take $a_{j}$, $b_{j}\in (0,\infty)$ such that
	\begin{equation}\label{parameter}
		\frac{1}{3\cdot\sqrt{2e}}\|\dc_{j}(T_1)\|_{L_{1}}=\exp\left\{-\frac{a^{2}_{j}}{2}\right\},\quad	\frac{1}{3\cdot\sqrt{2e}}\|\dc_{j}(T_2)\|_{L_{1}}=\exp\left\{-\frac{b^{2}_{j}}{2}\right\}.
	\end{equation}
	Then \eqref{correlation-4} turns into
	\begin{equation}\label{correlation-5}
	\begin{split}
A^{j}(s)&\geq 18\sup_{\substack{s^{2}\leq (1-p)(1-q)\\(p,q)\in(0,1)^{2}}}\exp\left\{-\frac{1}{2}\left(\frac{a^{2}_{j}}{p}+\frac{b^{2}_{j}}{q}\right)\right\}\\
&=18\exp\left\{-\frac{1}{2}\left(\inf_{\substack{s^{2}\leq (1-p)(1-q)\\(p,q)\in(0,1)^{2}}}\frac{a^{2}_{j}}{p}+\frac{b^{2}_{j}}{q}\right)\right\}.
\end{split}	
	\end{equation}
	Without loss os generality, by symmetry, we assume from now on that $a^{2}_{j}$ and $b^{2}_{j}$ are non-zeros and $\frac{a^{2}_{j}}{b^{2}_{j}}\in(0,1)$. Applying Lemma \ref{optimal} with $t=a_{j}/b_j$  to \eqref{correlation-5} entails that
	\begin{equation}\label{correlation-6}
		A^{j}(s)\geq18 \exp\left\{-\frac{1}{2}\left(\frac{a^{2}_{j}+2sa_{j}b_{j}+b^{2}_{j}}{1-s^{2}}\right)\right\},\mbox{for all }s\in(0,1).
	\end{equation}
	Since $F(s)\coloneqq \exp\left\{-\frac{1}{2}\left(\frac{a^{2}_{j}+2sa_{j}b_{j}+b^{2}_{j}}{1-s^{2}}\right)\right\}$ is decreasing on $s\in(0,1)$, we have
	\begin{equation}\label{correlation-7}
    \begin{split}
		\int_0^1A^{j}(s) ds&\geq 18\int_{0}^{1}\exp\left\{-\frac{1}{2}\left(\frac{a^{2}_{j}+2sa_{j}b_{j}+b^{2}_{j}}{1-s^{2}}\right)\right\}ds\\
       &\geq18 \epsilon \exp\left\{-\frac{1}{2}\left(\frac{a^{2}_{j}+2\epsilon a_{j}b_{j}+b^{2}_{j}}{1-\epsilon^{2}}\right)\right\},
    \end{split}
	\end{equation}
	for every $\epsilon\in(0,1)$ and $j\in [n]$. By \eqref{parameter} and \eqref{dj T1 T2 1}, it follows that
	\[
	\min\{a_{j},b_{j}\}\geq (\log(18e))^{1/2},\quad j\in[n],
	\]
and hence, \(\frac{1}{a_jb_j}\in(0,1)\) for each $j\in [n]$. We further have	\begin{equation}\label{correlation-8}
		\begin{split}
			\sup_{\epsilon\in(0,1)}\epsilon \exp\left\{-\frac{1}{2}\left(\frac{a^{2}_{j}+2\epsilon a_{j}b_{j}+b^{2}_{j}}{1-\epsilon^{2}}\right)\right\}&\geq \liminf_{\epsilon\to \frac{1}{a_{j}b_{j}}}\epsilon \exp\left\{-\frac{1}{2}\left(\frac{a^{2}_{j}+2\epsilon a_{j}b_{j}+b^{2}_{j}}{1-\epsilon^{2}}\right)\right\}\\
			&=\frac{1}{a_{j}b_{j}}\exp\left\{-\frac{1}{2}\left(\frac{a^{2}_{j}+2+b^{2}_{j}}{1-\frac{1}{a^{2}_{j}b^{2}_{j}}}\right)\right\}.
		\end{split}
	\end{equation}
	We now claim that there exists a universal constant $\delta>0$ such that
	\begin{equation}\label{estimate-1}
		\frac{a^{2}_{j}+2+b^{2}_{j}}{1-\frac{1}{a^{2}_{j}b^{2}_{j}}}\leq a^{2}_{j}+b^{2}_{j}+\delta,
	\end{equation}
provided that $a_{j}$, $b_{j}\geq (\log(18e))^{1/2}$. Rearranging \eqref{estimate-1}, it remains to show that there exists a universal constant $\delta>0$ such that
	\begin{equation*}
		\left(\frac{ab}{ab-1}\right)\left(2+\frac{1}{a}+\frac{1}{b}\right)\leq \delta,\quad\mbox{when }a=a_{j}^2,~b=b_{j}^2\geq \log(18e).
	\end{equation*}
	Since $ab\geq \frac{4}{3}$ for $a$, $b\geq \log(18 e)$, it follows that $ab-1\geq ab/4$. Hence, we can choose $\delta=12\geq 8\left(1+\frac{1}{\log(18e)}\right)$ such that the following holds
	\[
	\left(\frac{ab}{ab-1}\right)\left(2+\frac{1}{a}+\frac{1}{b}\right)\leq 4\left(2+\frac{1}{a}+\frac{1}{b}\right)\leq 8\left(1+\frac{1}{\log(18e)}\right)\leq 12.
	\]
	Combining $\delta=12$, \eqref{estimate-1}, \eqref{correlation-8} and \eqref{parameter} yields that
	\begin{equation}\label{correlation-9}
	\begin{split}
	\int_0^1 A^{j}(s)ds&\geq \frac{18}{a_{j}b_{j}e^{6}}\exp\left\{-\frac{a^{2}_{j}+b^{2}_{j}}{2}\right\}\\
&\geq 	\frac{1}{4e^7} \frac{\|\dc_{j}(T_1)\|_{L_{1}}\|\dc_{j}(T_2)\|_{L_{1}}}{\sqrt{\log(e/\|\dc_{j}(T_1)\|_{L_{1}})\log(e/\|\dc_{j}(T_2)\|_{L_{1}})}}.
\end{split}
	\end{equation}
Indeed,  for each $j\in[n]$, it follows from \eqref{parameter} that
	\[
	\exp\left\{-\frac{a_j^2+b_j^2}{2}\right\}
	=
	\frac{
		\|\dc_j(T_1)\|_{L_1}
		\|\dc_j(T_2)\|_{L_1}
	}{18e}.
	\]
	Moreover, since
	$\|\dc_j(T_i)\|_{L_1}\leq1$ and $3\sqrt{2e}<e^2$, we have
	\[
	a_j^2
	=
	2\log\left(
	\frac{3\sqrt{2e}}{\|\dc_j(T_1)\|_{L_1}}
	\right)
	\leq
	4\log\left(
	\frac{e}{\|\dc_j(T_1)\|_{L_1}}
	\right),
	\]
	and similarly,
	\[
	b_j^2
	\leq
	4\log\left(
	\frac{e}{\|\dc_j(T_2)\|_{L_{1}}}
	\right).
	\]
	which proves the last inequality in \eqref{correlation-9}. The proof is complete.

\end{proof}

\chapter{Quantum BKS inequality and its inverse}\label{sec-bks}
This chapter aims to investigate the quantum BKS-inequality in detail. In a seminal paper, Benjamini, Kalai and Schramm \cite{BKS1999} (classical BKS theorem) proved that if the sum of the squares of the influences of $f$ is close to zero, then $f$ must be noise sensitive. In \cite{KK2013}, Keller and Kindler obtained a quantitative version of the BKS theorem.  It was shown in \cite{KMS2014} that the quantitative BKS theorem also holds for the standard Gaussian measure on $\mathbb{R}^n$ with respect to geometric influences.

 Recently, a quantum quantitative BKS theorem has been obtained by Blecher, Gao and Xu \cite{BGX2025} via a \emph{random restriction technique} (see also \cite{JLLZ2025} for the random restriction technique used there to derive a quantum Eldan-Gross inequality). In this chapter, we will combine the semigroup tools with the noncommutative martingales to provide a simpler proof of the quantum BKS-inequality in Theorem \ref{geometric interpolations}. Our proof is actually motivated by \cite{JLZ2025}. We prefer to provide all the details for the reader's convenience.  
  Moreover, inspired by the semigroup approach by Keller, Mossel and Sen \cite{KMS2014}, we obtain an \emph{inverse} quantum BKS-inequality.
 
\section{The quantum BKS-inequality}\label{sec BKS 1}
We first recall some basic facts on noncommutative martingale theory. Recall that conditional expectation in the noncommutative setting could be defined as follows. Let $\mathcal{B}$ be a $\ast$-subalgebra of $\M_{2^{n}}$. There exists a linear operator $\mathcal{E}_{\mathcal{B}}: \M_{2^{n}}\to \mathcal{B}$ satisfying the following conditions: for every $T\in \M_{2^{n}}$
\begin{enumerate}[{\rm (i)}]
	\item $\mathcal{E}_{\mathcal{B}}$ is a normal contractive positive projection from $\M_{2^{n}}$ onto $\mathcal{B}$;
	\item $\mathcal{E}_{\mathcal{B}}(aTb)=a\mathcal{E}_{\mathcal{B}}(T)b$ provided $a,b\in\mathcal{B}$;
	\item $\tau(\mathcal{E}_{\mathcal{B}}(T))=\tau(T)$.
\end{enumerate}
The operator $\mathcal{E}_{\mathcal{B}}$ given as above is called the conditional expectation with respect to $\mathcal{B}$.  By (i) and (iii), one can deduce via interpolation that $\mathcal{E}_{B}$ is a contraction on $L_{p}(\M_{2^{n}})$ for all $1\leq p\leq\infty$.

For every $\mathbf{s}\in\{0,1,2,3\}^{n}$, we denote that $\max(\mathbf{s})\coloneqq\max\{k\in[n]:s_{k}\neq 0\}$ and let $\mathcal{M}_{j}$ be the $\ast$-subalgebra generated by $\{\sigma_{\mathbf{s}}\}_{\max(\mathbf{s})\leq j}$. Since $(\mathcal{M}_{j})_{j=1}^n$ forms an increasing sequence of $*$-subalgebras, it follows that, for each $T\in \M_{2^{n}}$, $(\mathcal{E}_{\mathcal{M}_{j}}(T))_{j=1}^{n}$ forms a noncommutative martingale (see e.g. \cite{PX1997} for its definition), and the related martingale difference is given as follows: $\mathcal{D}_{1}(T)\coloneqq\mathcal{E}_{\mathcal{M}_{1}}(T)$ and
\[
\mathcal{D}_{j}(T)\coloneqq\mathcal{E}_{\mathcal{M}_{j}}(T)-\mathcal{E}_{\mathcal{M}_{j-1}}(T),\quad 2\leq j\leq n.
\]


The next elementary lemma contains necessary properties of conditional expectations.

\begin{lemma}\label{BPce} 
	Let $T=\sum_{\mathbf{s}\in\{0,1,2,3\}^{n}}\widehat{T}(\mathbf{s})\sigma_{\mathbf{s}}\in \M_{2^{n}}$. 
	\begin{enumerate}[{\rm (i)}]
		\item For each $j\in[n]$, we have $\mathcal{E}_{\mathcal{M}_{j}}(T)=\sum_{\max(\mathbf{s})\leq j}\widehat{T}(\mathbf{s})\sigma_{\mathbf{s}}$ and
		\begin{equation}\label{martingale decomposition}
			\mathcal{D}_{j}(T)=\sum_{\max(\mathbf{s})=j}\widehat{T}(\mathbf{s})\sigma_{\mathbf{s}}.
		\end{equation}
		\item We have
		\begin{equation}\label{martingale version 2}
			\|T\|^{2}_{L_{2}}=\sum_{j=1}^n\|\mathcal{D}_{j}(T)\|^{2}_{L_{2}}.
		\end{equation}
	\end{enumerate}
\end{lemma}
\begin{proof}
	Since $\mathcal{M}_{j}$ is a $\ast$-subalgebra of $\M_{2^{n}}$ generated by $\{\sigma_{\mathbf{s}}\}_{\max(\mathbf{s})\leq j}$, by the orthonormality of stabilizer operators, we have
	\begin{equation*}
		\mathcal{E}_{\mathcal{M}_{j}}(\sigma_{\mathbf{s}})=
		\begin{cases}
			\sigma_{\mathbf{s}},~\mbox{if}~\max(\mathbf{s})\leq j\\
			0,~~~~~~\mbox{otherwise}.
		\end{cases}
	\end{equation*}
	Hence the claim follows easily from the orthnormality of the martingale differences $\{\mathcal{D}_{j}(T)\}_{j=1}^{n}$.
\end{proof}



\begin{lemma}\label{contraction}
	Let $\mathcal{M}_{[n]\setminus \{j\}}$ be the $\ast$-subalgebra of $\M_{2^{n}}$ generated by $\{\sigma_{\mathbf{s}}\}_{s_{j}=0}$. Then for $T=\sum_{\mathbf{s}\in\{0,1,2,3\}^{n}}\widehat{T}(\mathbf{s})\sigma_{\mathbf{s}}\in \M_{2^{n}}$ and $j\in [n]$, we have
	\[
	\mathcal{E}_{\mathcal{M}_{[n]\setminus \{j\}}}(T)=\sum_{s_{j}=0}\widehat{T}(\mathbf{s})\sigma_{\mathbf{s}}.
	\]
	Consequently, the following holds
	\begin{equation*}
		\dc_{j}(T)=T-\mathcal{E}_{\mathcal{M}_{[n]\setminus \{j\}}}(T).
	\end{equation*}
\end{lemma}
\begin{proof}
	The proof is trivial, which is left to the reader.
\end{proof}

\begin{corollary}\label{cor-bdd-p}
	For $T\in \M_{2^{n}}$ and each $j\in [n]$, we have
	\begin{equation*}\label{conditonal expectation 3}
		\mathcal{D}_{j}(T)=\mathcal{E}_{\mathcal{M}_{j}}(\dc_{j}(T)).
	\end{equation*}
	Consequently, $\|\mathcal{D}_{j}(T)\|_{L_{p}}\leq \|\dc_{j}(T)\|_{L_{p}}$ for every $1\leq p\leq \infty$.
\end{corollary}

\begin{proof}
	By  the last lemma,  we have
	\begin{align*}
		\mathcal{E}_{\mathcal{M}_{j}}(\dc_{j}(T))&= \mathcal{E}_{\mathcal{M}_{j}}(T)-\mathcal{E}_{\mathcal{M}_{j}} \left( \sum_{s_{j}=0}\widehat{T}(\mathbf{s})\sigma_{\mathbf{s}}\right) \\
		&=\mathcal{E}_{\mathcal{M}_{j}}(T)- \sum_{\max(\mathbf{s})\leq j-1}\widehat{T}(\mathbf{s})\sigma_{\mathbf{s}}\\
		&= \mathcal{E}_{\mathcal{M}_{j}}(T)-\mathcal{E}_{\mathcal{M}_{j-1}}(T)=\mathcal{D}_{j}(T),
	\end{align*}
	where we used  Lemma \ref{BPce} (i) in the third line. Since every conditional expectation is a contraction on $L_p(\M_{2^{n}})$ for $1\leq  p\leq \infty$, the desired inequality follows immediately. 
\end{proof}

\begin{lemma}\label{decomposition}
	For each $T\in\M_{2^{n}}$ with $\mathrm{tr}(T)=0$, we have
	\begin{equation*}
		\sum\limits_{j=1}^{n}\|\dc_{j}(T)\|^{2}_{L_{2}}=\sum\limits_{j=1}^{n}\sum\limits_{i=1}^{n}\|\dc_{j}(\mathcal{D}_{i}(T))\|^{2}_{L_{2}}.
	\end{equation*}
\end{lemma}
\begin{proof}
	Clearly, by the martingale decomposition of $T$ and $\mathrm{tr}(T)=0$, we have $T=\sum_{i=1}^{n}\mathcal{D}_{i}(T)$. Hence, the orthonormality of martingale differences yields that
	\[
	\|\dc_{j}(T)\|^{2}_{L_{2}}=\sum_{i=1}^{n}\|\dc_{j}(\mathcal{D}_{i}(T))\|^{2}_{L_{2}},~\forall~j\in[n].
	\]
	Summing over $j\in[n]$ entails the desired identity.
\end{proof}



We now arrive to state and prove the quantum BKS-inequality as follows.
\begin{theorem}\label{geometric interpolations}
	Suppose that $T\in \M_{2^{n}}$ with $\|T\|_{L_{\infty}}\leq 1$ and $T^*=T$. Then
	\[
	\mathrm{Var}\left(e^{-t\mathcal{L}}(T)\right)\leq \Big(\sum\limits_{j=1}^{n}\|\dc_{j}(T)\|^{2}_{L_{1}}\Big)^{\tanh(t)},\quad t\geq 0.
	\]
\end{theorem}

\begin{proof}
	We assume without loss of generality that $\mathrm{tr}(T)=0$. 
	It is not hard to see that $e^{-t\mathcal{L}}(\mathcal{D}_{j}(T))=\mathcal{D}_{j}(e^{-t\mathcal{L}}(T))$ for each $j\in [n]$. Indeed, by Lemma \ref{BPce} (i), we have
	\[
	e^{-t\mathcal{L}}\left(\mathcal{D}_{j}(T)\right)=e^{-t\mathcal{L}}\Big(\sum_{\max(\s)=j}\widehat{T}(\s)\sigma_{\s}\Big)=\sum_{\max(\s)=j}\widehat{T}(\s)e^{-t|\s|}\sigma_{\s}=\mathcal{D}_{j}(e^{-t\mathcal{L}}(T)).
	\]
	Thus, for each $t\geq 0$, by \eqref{martingale version 2} and the hypercontractivity (see Lemma \ref{hc-1}), we obtain
	\begin{align*}
		\mathrm{Var}\big(e^{-t\mathcal{L}}(T)\big)&\leq\sum_{j=1}^{n}\big\|e^{-t\mathcal{L}}\big(\mathcal{D}_{j}(T)\big)\big\|^{2}_{L_{2}}\leq \sum_{j=1}^{n}\left\|\mathcal{D}_{j}(T)\right\|^{2}_{L_{1+e^{-2t}}}\\
		&\leq \sum_{j=1}^{n}\left\|\mathcal{D}_{j}(T)\right\|^{2(1-\tanh(t))}_{L_{2}}\cdot\left\|\mathcal{D}_{j}(T)\right\|^{2\tanh(t)}_{L_{1}}\\
		&\leq \Big(\sum_{j=1}^{n}\left\|\mathcal{D}_{j}(T)\right\|^{2}_{L_{2}}\Big)^{1-\tanh(t)}\Big(\sum_{j=1}^{n}\left\|\mathcal{D}_{j}(T)\right\|^{2}_{L_{1}}\Big)^{\tanh(t)}\\
		&=\|T\|^{2(1-\tanh(t))}_{L_{2}}\Big(\sum_{j=1}^{n}\left\|\mathcal{D}_{j}(T)\right\|^{2}_{L_{1}}\Big)^{\tanh(t)}\leq\Big(\sum_{j=1}^{n}\|\dc_{j}(T)\|^{2}_{L_{1}}\Big)^{\tanh(t)}.
	\end{align*}
Here, in the second inequality, we used the H\"older inequality with parameters $\frac{1}{1+e^{-2t}}=\frac{\tanh(t)}{1}+\frac{1-\tanh(t)}{2}$, and used both the assumption $\|T\|_{L_{\infty}}\leq 1$ and Corollary \ref{cor-bdd-p} in the last inequality.
\end{proof}


\begin{remark}
It is not hard to see 
$$\frac{1}{2}(1-e^{-2t})\leq \tanh(t)\leq 1-e^{-2t},\quad \forall t>0.$$
Recall that the classical quantitative BKS inequality proved in   \cite[Theorem 4]{KK2013} states that for some $c>0.234$,
\[\mathrm{Var}\big(e^{-t\mathcal{L}}(f)\big)\leq 20 \Big(\sum\limits_{j=1}^{n}\|d_{j}(T)\|^{2}_{L_{1}(\{-1,1\}^n)}\Big)^{c (1-e^{-2t})},\quad t\geq 0.\]
It seems that our quantum BKS inequality is better, more or less, since we find that $c\in [1/2,1]$ is good enough for the purpose. Additionally, our proof is simpler than the one given in \cite{KK2013}. 
\end{remark}

\begin{remark}
 The classical quantitative BKS inequality  \cite[Theorem 4]{KK2013} was recently extended into the quantum setting via the some quantum random restriction method by Blecher et al. \cite{BGX2025}. Our proof of Theorem \ref{geometric interpolations} is simpler than \cite{BGX2025}. 
\end{remark}

\section{The inverse quantum BKS-inequality}
\begin{theorem}
	For each $t>0$ and monotone increasing $T\in \M_{2^{n}}$ defined as in Lemma \ref{mo-in-o}, we have
	\begin{equation}\label{inverse BKS}
		\mathrm{Var}\left(e^{-t\mathcal{L}}(T)\right)\geq e^{-2t}\sum_{j=1}^{n}\|\dc_{j}(T)\|^{2}_{L_{1}}.
	\end{equation}
\end{theorem}

\begin{proof}
	For a monotone increasing element $T\in \M_{2^{n}}$, we apply the covariance identity \eqref{key-identity} in Lemma \ref{basic} (iii) to get
	\[
	\mathrm{Var}(T)=\sum_{j=1}^{n}\int_{0}^{\infty}e^{-t}\mathrm{tr}\left(\partial_{j}(T)\cdot e^{-t\mathcal{L}}\left(\partial_{j}(T)\right)\right)dt.
	\]
	For each $j\in [n]$, we define
	\[
	h_{j}(t)\coloneqq \mathrm{tr}\left(\partial_{j}(T)\cdot e^{-t\mathcal{L}}\left(\partial_{j}(T)\right)\right),\quad\mbox{for } t>0.
	\]
	It is clear that
	\begin{align*}
		\frac{d}{dt}h_{j}(t)&=-\mathrm{tr}\left(\partial_{j}(T)\cdot \mathcal{L} e^{-t\mathcal{L}}\left(\partial_{j}(T)\right)\right)\\
		&=-\sum_{l=1}^{n}\mathrm{tr}\left(\dc_{l}\partial_{j}(T)\cdot \dc_{l}e^{-t\mathcal{L}}\left(\partial_{j}(T)\right)\right)\\
		&=-\sum_{l=1}^{n}\left\|e^{-t\mathcal{L}/2}\dc_{l}\partial_{j}(T)\right\|^{2}_{L_{2}}\leq 0.
	\end{align*}
	This means that $h_{j}(t)$ is a decreasing function on $t>0$. Hence,
	\begin{equation}\label{inverse BKS 1}
		\begin{split}
			\mathrm{Var}(T)&=\sum_{j=1}^{n}\int_{0}^{\infty}e^{-t}h_{j}(t)dt\\
			&\geq\sum_{j=1}^{n}\int_{0}^{\infty}e^{-t}\lim_{t\to \infty}h_{j}(t)dt\\
			&=\sum_{j=1}^{n}\mathrm{tr}\left(\partial_{j}(T)\right)^{2}\cdot\left(\int_{0}^{\infty}e^{-t}dt\right)=\sum_{j=1}^{n}\|\partial_{j}(T)\|^{2}_{L_{1}},
		\end{split}
	\end{equation}
	where we used $T$ is monotone increasing in the last line.
	
	Now, for each $t>0$ note here that $e^{-t\mathcal{L}}(T)$ is still monotone increasing (see Lemma \ref{basic} (iii)),   substituting $e^{-t\mathcal{L}}(f)$ to \eqref{inverse BKS 1} and Lemma \ref{comparison} (ii) entail that
	\begin{align*}
		\mathrm{Var}\left(e^{-t\mathcal{L}}(f)\right)&\geq\sum_{j=1}^{n}\left\|\partial_{j}\left(e^{-t\mathcal{L}}(f)\right)\right\|^{2}_{L_{1}}\\
		&=e^{-2t}\sum_{j=1}^{n}\mathrm{tr}\left(e^{-t\mathcal{L}}\left(\partial_{j}(f)\right)\right)^{2}\\
		&=e^{-2t}\sum_{j=1}^{n}\mathrm{tr}\left( \partial_{j}(f)\right)^{2}\\
		&=e^{-2t}\sum_{j=1}^{n}\|\partial_{j}(f)\|^{2}_{L_{1}}\geq e^{-2t}\sum_{j=1}^{n}\|\dc_{j}(f)\|^{2}_{L_{1}},
	\end{align*}
	where the second equality is due to that $e^{-t\mathcal{L}}$ is trace preserving, and we used Lemma \ref{basic} (i) in the last inequality.
\end{proof}


\begin{theorem}\label{Var-low-1} Let $T\in\mathbb{M}_{2^n}$ be monotonically increasing operator. Then 
	$$\mathrm{Var}(T)\geq\sum_{j=1}^n\|\dc_j(T)\|_{L_1(\mathbb{M}_{2^n})}^2.$$
\end{theorem}
\begin{proof} Combining \eqref{inverse BKS 1} with Lemma \ref{basic} (i) yields the expected result.
\end{proof}



\chapter[Quantum BH Inequality and Approximation]{Quantum Bohnenblust--Hille inequality and approximation of polynomials from tail spaces }\label{sec BH distance}
In Chapter \ref{sec BH}, we improve the constant of the quantum Bohnenblust--Hille inequality, which is the central technical result behind recent progress in learning bounded local Hamiltonians and observables \cite{HCP2023,VZ2024}. In Chapter \ref{sec distance Tail}, we study $L_1$ distance of polynomials on the quantum Boolean from tail spaces. This study relies on the quantum Bohnenblust--Hille inequality. The reduction method used in this chapter is also useful in next chapter concerning quantum learning problems.


\section{Quantum Bohnenblust--Hille inequality}\label{sec BH}
In \cite{Bl2001}, Blei established Bohnenblust--Hille inequalities for polynomials on Boolean cubes with dimension-free constants:
for any function $f \in \mathcal{P}^{\leq d}(\{-1,1\}^n)$, we have the dimension-free estimate
\begin{equation}\label{classical BH}
\left( \sum_{|A|\leq d} |\widehat{f}(A)|^{\frac{2d}{d+1}} \right)^{\frac{d+1}{2d}} \leq  C_d\|f\|_\infty.
\end{equation}
Denoting $\mathrm{BH}_{\{\pm1\}}^{\leq d} $  the best constant such that \eqref{classical BH} holds.  Defant et al. \cite{DMP2019} proved that the dimension-free Bohnenblust--Hille constants grow at most subexponentially in the degree, that is, $\mathrm{BH}_{\{\pm1\}}^{\leq d} \leq C^{\sqrt{d \log d}}$ with $C > 1$ being a universal constant. Recently, this
Bohnenblust--Hille inequality  has found great applications in
learning bounded low-degree functions on Boolean cubes \cite{EI2022}, and in the $L_p$ distance estimate of polynomials on the Boolean from
Walsh tail spaces \cite{EZ2024}. 




An analog of \eqref{classical BH} on the quantum Boolean $\mathbb{M}_{2^n}$  was proved by Huang-Chen-Preskill \cite{HCP2023}. 
For any operator $T\in \mathbb{M}_{2^n}$ of degree at most $d\in [n]$, one has
\[
 \left( \sum_{|\mathbf{s}|\leq d} |\widehat{T}(\mathbf{s})|^{\frac{2d}{d+1}} \right)^{\frac{d+1}{2d}} \leq \mathrm{BH}_{\mathbb{M}_2}^{\leq d} \|T\|,
\]
where $\mathrm{BH}_{\mathbb{M}_2}^{\leq d} < \infty$ denotes the best constant.
Clearly, $\mathrm{BH}_{\{\pm1\}}^{\leq d} \leq \mathrm{BH}_{\M_2}^{\leq d}$. Huang-Chen-Preskill \cite{HCP2023} showed that 
$$\mathrm{BH}_{\mathbb{M}_2}^{\leq d}\leq d^{O(d)}.$$
Later, Volberg and Zhang \cite[Theorem 1.2]{VZ2024} improved this constant to be
\[
\mathrm{BH}_{\mathbb{M}_2}^{\leq d} \leq 3^d \mathrm{BH}_{\{\pm1\}}^{\leq d}
\]
via a reduction method. In the very recent work \cite{BSVZ2025}, Becker et al. further improved this to be 
\[
\mathrm{BH}_{\mathbb{M}_2}^{\leq d} \leq \sqrt{3}^{d+1} \mathrm{BH}_{\{\pm1\}}^{\leq d}.
\]
Our main contribution here is to show that 
\[
\mathrm{BH}_{\mathbb{M}_2}^{\leq d} \leq \sqrt{3}^{d} \mathrm{BH}_{\{\pm1\}}^{\leq d},
\]
and we will also show that this is optimal in some sense (see Remark \ref{rem BH 1} below). 
Our method draws inspiration from the reduction  method developed by Volberg and Zhang \cite{VZ2024}.

\begin{theorem}\label{Quantum BH}
The following inequality holds for every $T\in \mathcal{P}(\mathbb{M}_{2^n})$,
\[
\left(\sum_{\substack{\mathbf{s}\in\{0,1,2,3\}^{n}\\|\mathbf{s}|\leq d}}|\widehat{T}(\s)|^{\frac{2d}{d+1}}\right)^{\frac{d+1}{2d}}\leq 3^{d/2}\mathrm{BH}^{\leq d}_{\{\pm1\}}\|T\|_{L_{\infty}(\mathbb{M}_{2^n})}.
\]
\end{theorem}

To begin with, we shall present a useful lemma.  Write
	\[\varepsilon=
	\big(\varepsilon_1^{(1)},\dots,\varepsilon_n^{(1)},
	\varepsilon_1^{(2)},\dots,\varepsilon_n^{(2)},
	\varepsilon_1^{(3)},\dots,\varepsilon_n^{(3)}\big)\in\{-1,1\}^{3n}.
	\]
	For each \(1\le j\le n\), define $\rho_j: \{-1,1\}^{3} \to \M_2(\mathbb C)$ by setting
	\begin{equation}\label{rho-j}
		\rho_j(\varepsilon):=
		\mathbf{1}_2+\frac{\varepsilon_j^{(1)}\sigma_1+\varepsilon_j^{(2)}\sigma_2+\varepsilon_j^{(3)}\sigma_3}{\sqrt3}\in \M_2(\mathbb C),
	\end{equation}
and define $\rho: \{-1,1\}^{3n} \to \M_{2^n}(\mathbb C)$ by setting
	\begin{equation}\label{rho}
	\rho:=\rho_1\otimes\cdots\otimes \rho_n.
	\end{equation}

\begin{lemma}\label{lem rho}
Let $\rho$ be as in \eqref{rho}. Given $d\in [n]$, let $T\in \mathcal{P}^{\leq d}(\M_{2^n})$ and define the function \(f_T:\{-1,1\}^{3n}\to\mathbb C\)  by setting
	\[
	f_T(\varepsilon):=\mathrm{tr}\big(T\,\rho(\varepsilon)\big),\quad \varepsilon \in \{-1,1\}^{3n}.
	\]
We have the following  properties.
\begin{enumerate}[{\rm (i)}]
\item The matrix  $\rho$ is a density matrix, that is,  \(\rho(\varepsilon)\geq 0\) and \(\mathrm{tr}(\rho(\varepsilon))=1\) for each $\varepsilon \in \{-1,1\}^{3n}$.
\item 
 The function $f_T$ is at most of degree $d$, that is, $f_T\in \mathcal{P}^{\leq d}(\{-1,1\}^{3n})$. 
\item We have $\{\widehat{f_T}(A)\}_{A\subset[3n]}=\{\frac{\widehat{T}(\s)}{3^{|\s|/2}}\}_{\mathbf{s}\in\{0,1,2,3\}^n}$ and 
\begin{equation}\label{fT T}
\left(\sum_{|A|\leq d}|\widehat{f_T}(A)|^{\frac{2d}{d+1}}\right)^{\frac{d+1}{2d}}=\left(\sum_{\substack{\s\in\{0,1,2,3\}^{n}\\|\s|\leq d}}\left|\frac{\widehat{T}(\s)}{3^{|\s|/2}}\right|^{\frac{2d}{d+1}}\right)^{\frac{d+1}{2d}}.
\end{equation}
\end{enumerate}
\end{lemma}
\begin{proof}
(i) 	It is easy to verify via \eqref{car-relation} that for each $1\leq j\leq n$,
	\[\frac{\varepsilon_j^{(1)}\sigma_1+\varepsilon_j^{(2)}\sigma_2+\varepsilon_j^{(3)}\sigma_3}{\sqrt3}\]
is hermitian and unitary.  Hence, for each $\varepsilon\in \{-1,1\}^{3n}$, the eigenvalues of $\rho_j(\varepsilon)$ are \(0\) and \(2\), and therefore \(\rho_j\ge 0\). Also note that 
	\[\mathrm{tr}_2(\rho_j(\varepsilon))
	=\mathrm{tr}_2\Big(
	\mathbf{1}_2+\frac{\varepsilon_j^{(1)}\sigma_1+\varepsilon_j^{(2)}\sigma_2+\varepsilon_j^{(3)}\sigma_3}{\sqrt3}\Big)=1,\quad 1\leq j \leq n.\]
	Consequently, for each $\varepsilon$, 
	\begin{equation}\label{rho-0-trace-1}
		\rho\ge 0,\qquad
		\mathrm{tr}(\rho(\varepsilon))=\prod_{j=1}^n \mathrm{tr}_2(\rho_j(\varepsilon))=1.
	\end{equation}

(ii)  It is not hard to see that for each $\varepsilon\in \{-1,1\}^{3n}$
\begin{equation}\label{weak version}
\begin{split}
f_T(\varepsilon)&=\sum_{\substack{\mathbf{s}\in\{0,1,2,3\}^{n}\\|\mathbf{s}|\leq d}}\widehat{T}(\s)\langle \rho(\varepsilon),\sigma_{\s}\rangle\\
&=\sum_{\substack{\mathbf{s}\in\{0,1,2,3\}^{n}\\|\mathbf{s}|\leq d}}\widehat{T}(\s)\prod_{j\in\mathrm{supp}(\s)}\langle \rho_{j}(\varepsilon),\sigma_{s_j}\rangle\\
&=\sum_{\substack{\mathbf{s}\in\{0,1,2,3\}^{n}\\|\mathbf{s}|\leq d}}\left(\frac{\widehat{T}(\s)}{3^{|\s|/2}}\right)\prod_{j\in \mathrm{supp}(\s)}\varepsilon^{(s_{j})}_{j},
\end{split}
\end{equation}
where the second equality is due to \eqref{rho-0-trace-1} and the last equality is due to the fact \(\mathrm{tr}(\sigma_j)=0\) for \(j\in\{1,2,3\}\). The desired assertion follows. 

(iii) The formula \eqref{fT T} follows from \eqref{weak version} directly. 
\end{proof}


Now we turn to prove Theorem \ref{Quantum BH}. 

\begin{proof}[Proof of Theorem \ref{Quantum BH}]
Let $T\in \mathcal{P}^{\leq d}(\M_{2^n})$, and let $f_{T}:\{-1,1\}^{3n}\to \mathbb{C}$ be as  in Lemma \ref{lem rho}. 	By \eqref{rho-0-trace-1} and H\"older's inequality, for each $\varepsilon\in \{-1,1\}^{3n}$, we  have
\begin{equation}
\begin{split}
	|f_T(\varepsilon)|&=	|\mathrm{tr}(T\rho(\varepsilon))|	\leq	\|T\|_{L_{\infty}(\mathbb{M}_{2^n})} \|\rho(\varepsilon)\|_{L_1(\M_{2^n})}	\\
&=	\|T\|_{L_{\infty}(\mathbb{M}_{2^n})} \mathrm{tr}(\rho(\varepsilon))	=	\|T\|_{L_{\infty}(\mathbb{M}_{2^n})},
\end{split}
\end{equation}
and	consequently,
	\begin{equation}\label{fa-a-infty}
	\|f_T\|_\infty\le \|T\|_{L_{\infty}(\mathbb{M}_{2^n})}.
	\end{equation}
Applying the classical Bohnenblust--Hille inequality \eqref{classical BH} to the function $f_{T}:\{-1,1\}^{3n}\to \mathbb{C}$ defined in Lemma \ref{lem rho}, we have
\begin{equation*}\label{transference to classical BH}
\begin{split}
\left(\sum_{|A|\leq d}|\widehat{f_T}(A)|^{\frac{2d}{d+1}}\right)^{\frac{d+1}{2d}}  
\leq \mathrm{BH}^{\leq d}_{\{\pm1\}}  \|f_T\|_{\infty}\leq \mathrm{BH}^{\leq d}_{\{\pm1\}} \|T\|_{L_{\infty}(\mathbb{M}_{2^n})} . 
\end{split}
\end{equation*}
Thus, by Lemma \ref{lem rho} (iii), we have
$$\left(\sum_{\substack{\s\in\{0,1,2,3\}^{n}\\|\s|\leq d}}\left|\frac{\widehat{T}(\s)}{3^{|\s|/2}}\right|^{\frac{2d}{d+1}}\right)^{\frac{d+1}{2d}}\leq  \mathrm{BH}^{\leq d}_{\{\pm1\}} \|T\|_{L_{\infty}(\mathbb{M}_{2^n})} .$$
We can deduce the  desired quantum Bohnenblust--Hille inequality whenever we  note that $3^{-d/2}\leq 3^{-|\s|/2}$ for all $\s\in\{0,1,2,3\}^{n}$ with $|\s|\leq d$. The proof is complete.
\end{proof}



\begin{remark}\label{rem BH 1}
We now explain that the factor $3^{d/2}$ should be optimal in the case $d=1$. In fact, we shall show that there exists $T$ such that inequality in Theorem \ref{Quantum BH} can be a equality.  Let 
\[
T\coloneqq \sum_{l=1}^{3}\sigma_{le_{1}}=
\begin{pmatrix}
1&1-i\\
1+i&-1
\end{pmatrix} \otimes \mathbf{1}_2\otimes \cdots \otimes\mathbf{1}_2.
\]
It is not hard to see that $\|T\|_{L_{\infty}(\mathbb{M}_{2^n})}=3^{1/2}$ and 
$$\sum_{\substack{\s\in\{0,1,2,3\}^{n}\\|\s|= 1}}|\widehat{T}(\s)|=3.$$
Recall that  $\mathrm{BH}^{\leq 1}_{\{\pm1\}}=1$. Combing above, we verify the assertion. 
\end{remark}

\begin{remark}\label{rem BH}
Our proof of the quantum Bohnenblust--Hille inequality Theorem \ref{Quantum BH} essentially builds upon the reduction method introduced by Volberg and Zhang \cite{VZ2024}. However,  by judiciously choosing   the density matrix $\rho$ and the reduction function $f_T$ here, we are able to obtain improved constants. 
\end{remark}

In \cite{VZ2024}, the authors applied the quantum  Bohnenblust--Hille inequality to study quantum junta theorem and learning problems of low degree quantum observables. We also study quantum learning problems in Chapter \ref{sec QLT} but with different arguments motivated by \cite{EIS2023}. 

In Chapter \ref{sec distance Tail} below, motivated by \cite{EZ2024}, we shall provide a new and interesting application of the quantum  Bohnenblust--Hille inequality. 

\section{Approximation of polynomials from quantum tail spaces}\label{sec distance Tail}
In \cite{Ol2017}, Oleszkiewicz estimated the $L_1$-distance of the Rademacher sum $f=\sum_{|A|=1}w_A$ from the $d$-th tail space $\mathcal{T}^{\geq d}(\{-1,1\}^n)$. 
\begin{theorem}[Oleszkiewicz]\label{Oleszkiewicz inequality}
There exist universal constants $C_{1}$ and $C_{2}$ such that for every  Rademacher sums $n\in \mathbb{N}$ and every $f\in \mathcal{P}^{=1}(\{-1,1\}^{n})$, the following inequality holds
\[
C_{1}K(f,m,\ell^{n}_{2},\ell^{n}_{\infty})\leq \inf_{g\in \mathcal{T}^{>m}(\{-1,1\}^{n})}\|f-g\|_{L_{1}}\leq C_{2}K(f,m,\ell^{n}_{2},\ell^{n}_{\infty}),
\]
where $1\leq m\leq n$.
\end{theorem}
The $K$-functional used here is given later. This is a restatement of initial Oleszkiewicz's theorem in terms of the $K$-functional, which was observed by Eskenazis  and Zhang \cite{EZ2024}.   As  pointed out by Oleszkiewicz, the method of \cite{Ol2017} does not appear to extend beyond Rademacher sums. Recently,  using a different method by employing the Bohnenblust--Hille inequality \eqref{classical BH}  and a class bound of
Figiel on the Rademacher projection of polynomials, Eskenazis  and Zhang \cite{EZ2024} established various bounds for the $L_p$ distance of polynomials on the Boolean from Walsh tail spaces, extending some of Oleszkiewicz's results for Rademacher sums. In this section, we aim to establish counterpart results in \cite{EZ2024} for the quantum Boolean.


We now recall basic concept like $K$-functional in interpolation theory. For two compatible Banach spaces $X$ and $Y$, define norms $\|\cdot\|_{X\cap Y}$ and $\|\cdot\|_{X+Y}$ on $X\cap Y$ and $X+Y$, respectively as follows
\[
\|f\|_{X\cap Y}\coloneqq \max\{\|f\|_{X},\|f\|_{Y}\},\quad \forall~ f\in X\cap Y,
\]
and
\[
\|f\|_{X+Y}\coloneqq \inf\{\|g\|_{X}+\|h\|_{Y}:f=g+h,~g\in X,~h\in Y\}.
\]
More generally, for each $t>0$, we define
\[
\|f\|_{X\cap tY}\coloneqq \max\{\|f\|_{X},t\|f\|_{Y}\},\quad \forall~f\in X\cap Y,
\]
and the $K$-functional as follows
\[
K(f,t,X,Y)=\|f\|_{X+tY}\coloneqq \inf\{\|g\|_{X}+t\|h\|_{Y}:f=g+h,~g\in X,~h\in Y\}.
\]
It is well-known that for two compatible Banach spaces $X$ and $Y$, we have
\[
(X\cap tY,\|\cdot\|_{X\cap tY})^{*}=(X^{*}+t^{-1}Y^{*},\|\cdot\|_{X^{*}+t^{-1}Y^{*}}),\quad \forall~t>0.
\]
We refer to the monograph \cite{BL1976} for more details. 


We first derive a quantum version of the Figiel inequality (see e.g. \cite{EI2022}).
Suppose that $T_{d}(x)=\sum_{\ell=0}^{d}c(d,\ell)x^{\ell}$ is the Chebyshev polynomial of the first kind which is characterized by $T_{d}(\cos(x))=\cos(dx)$ for all $x\in \mathbb{R}$. We now define  for simplicity
\[
\tilde{c}(d,\ell)\coloneqq
\begin{cases}
\frac{T^{(\ell)}_{d}(0)}{\ell!},&\mbox{if }d-\ell~\mbox{is even},\\
\frac{T^{(\ell)}_{d-1}(0)}{\ell!},&\mbox{if }d-\ell~\mbox{is odd}.
\end{cases}
\]
Recall the following well-known Markov inequality. Let $h$ be a polynomial on $[-1,1]$ with degree $d$. Then for each positive integer $\ell\leq d$ we have
\begin{equation}\label{Markov inequalities}
|h^{(\ell)}(0)|\leq
\begin{cases}
|T^{(\ell)}_{d}(0)|\|h\|_{L_{\infty}},&\mbox{if }d-\ell~\mbox{is even},\\
|T^{(\ell)}_{d-1}(0)|\|h\|_{L_{\infty}},&\mbox{if }d-\ell~\mbox{is odd}.
\end{cases}
\end{equation}




\begin{lemma}\label{Figiel inequality}
Suppose that $T\in \mathcal{P}^{\leq d}(\M_{2^n})$. Then, for each $\ell\in \mathbb{N}$ with $\ell \leq d$, we have
\[
\big\|\mathrm{Rad}^{=\ell}(T)\big\|_{L_{\infty}}\leq \tilde{c}(d,\ell)\|T\|_{L_\infty}.
\]
\end{lemma}
\begin{proof}
Without loss of generality, we write $T=\sum_{\s\in\{0,1,2,3\}^{n},|\s|\leq d}\widehat{T}(\s)\sigma_{\s}$. It follows that the function defined by
\[
h_{T}(t)\coloneqq\sum_{\substack{\s\in \{0,1,2,3\}^{n}\\|\s|\leq d}}\widehat{T}(\s)\sigma_{\s}t^{|\s|},\quad\forall~t\in [-1,1],
\]
satisfies $\sup_{t\in[-1,1]}\|h_{T}(t)\|_{L_{\infty}}\leq \|T\|_{L_{\infty}}$. For each $\xi$, $\eta\in L_{2}(\M_{2^n})$ with $\|\xi\|_{L_{2}}=\|\eta\|_{L_{2}}=1$, we define
\[
h_{T}(t)[\xi,\eta]\coloneqq\langle \xi,h_{T}(t)\eta\rangle.
\]
It is clear that
\[
\sup\{h_{T}(t)[\xi,\eta]:t\in [-1,1],\|\xi\|_{L_{2}}=\|\eta\|_{L_{2}}=1\}\leq \|T\|_{L_{\infty}},
\]
and $h_{T}(\cdot)[\xi,\eta]$ is of degree $d$ on $[-1,1]$. Hence, by \eqref{Markov inequalities}, we get that
\begin{equation}\label{differential Markov inequalities}
|h^{(\ell)}_{T}(0)[\xi,\eta]|\leq
\begin{cases}
|T^{(\ell)}_{d}(0)|\|T\|_{L_{\infty}},&\mbox{if }d-\ell~\mbox{is even}\\
|T^{(\ell)}_{d-1}(0)|\|T\|_{L_{\infty}},&\mbox{if }d-\ell~\mbox{is odd},
\end{cases}
\end{equation}
for every $\xi$, $\eta\in L_{2}(\M_{2^n})$. By the definition of Rademacher projection $\mathrm{Rad}^{=\ell}$, it follows that for each $\xi$, $\eta\in L_{2}(\M_{2^n})$ with $\|\xi\|_{L_{2}}=\|\eta\|_{L_{2}}=1$ we have
\begin{equation}\label{pointwise version}
|\langle \xi, \mathrm{Rad}^{=\ell}(T)\eta\rangle|=\Big|\sum_{\substack{\s\in \{0,1,2,3\}^{n}\\|\s|=\ell}}\widehat{T}(\s)\langle \xi,\sigma_{\s}\eta\rangle\Big|\\
=\big|\frac{h^{(\ell)}(0)[\xi,\eta]}{\ell!}\big|.
\end{equation}
Combining with \eqref{differential Markov inequalities} and \eqref{pointwise version} and taking the supremum over $\xi$ and $\eta$ in the unit sphere of $L_{2}(\M_{2^n})$ yields the desired result.
\end{proof}





For each $k$, $d\in [n]$ with $d<k$, we define
\[
\mathcal{P}_{(d,k]}(\M_{2^n})\coloneqq\left\{T\in \M_{2^n}:T=\sum_{\substack{\s\in\{0,1,2,3\}^{n}\\d<|\s|\leq k}}\widehat{T}(\s)\sigma_{\s}\right\}.
\]
In what follows, we shall write $\mathrm{Rad}^{\leq d}=\sum_{\ell \leq d}\mathrm{Rad}^{= \ell}$, $d\in [n]$. Given $T\in \M_{2^n}$, denote by $\widehat{T}$ the vector of Fourier coefficients of $T$. 




\begin{lemma}\label{upper bound for Rademacher projection 1}
For each $d<k\leq n$, let $N_{d}\coloneqq \sum_{l=0}^{d}3^{l}\binom{n}{l}$. Then the following inequality holds for every $T\in \mathcal{P}^{\leq k}(\M_{2^n})$,
\[
\max\left\{\left\|\widehat{\mathrm{Rad}^{\leq d}(T)}\right\|_{\ell^{N_{d}}_{2}},z(k,d)^{-1}\left\|\widehat{\mathrm{Rad}^{\leq d}(T)}\right\|_{\ell^{N_{d}}_{2d/(d+1)}}\right\}\leq\inf_{G\in \mathcal{P}_{(d,k]}(\M_{2^n})}\|T-G\|_{L_{\infty}},
\]
where $z(k,d)\coloneqq \mathrm{BH}_{\mathbb{M}_2}^{\leq d} \sum_{l=0}^{d}|\tilde{c}(k,l)|$ and $\mathrm{BH}_{\mathbb{M}_2}^{\leq d} \leq 3^{d/2}\mathrm{BH}^{\leq d}_{\{\pm1\}}$ is the constant coming from the quantum Bohnenblust--Hille inequality.
\end{lemma}
\begin{proof}
Take $T\in \mathcal{P}^{\leq k}(\M_{2^n})$. It follows from the Parseval identity that 
\begin{align*}
\left\|\widehat{\mathrm{Rad}^{\leq d}(T)}\right\|_{\ell_{2}^{N_{d}}}&=\left\|\mathrm{Rad}^{\leq d}(T)\right\|_{L_{2}}
\leq \|T\|_{L_{2}}\leq \|T\|_{L_{\infty}}.
\end{align*}
On the other hand, by the quantum Bohnenblust--Hille inequality (Theorem \ref{Quantum BH}), we have
\begin{align*}
\left\|\widehat{\mathrm{Rad}^{\leq d}(T)}\right\|_{\ell^{N_{d}}_{2d/(d+1)}}&\leq \mathrm{BH}_{\mathbb{M}_2}^{\leq d} \left\|\mathrm{Rad}^{\leq d}(T)\right\|_{L_{\infty}}\\
&\leq \mathrm{BH}_{\mathbb{M}_2}^{\leq d} \sum_{l=0}^{d}\left\|\mathrm{Rad}^{=l}(T)\right\|_{L_{\infty}}\leq \mathrm{BH}_{\mathbb{M}_2}^{\leq d} \sum_{l=0}^{d}|\tilde{c}(k,l)|\|T\|_{L_{\infty}},
\end{align*}
where we used Lemma \ref{Figiel inequality} in the last inequality.

Therefore, the Rademacher projection $\widehat{\mathrm{Rad}^{\leq d}}$ is a linear operator which is bounded from $(\mathcal{P}^{\leq k}(\M_{2^n}),\|\cdot\|_{L_{\infty}})$ to $\ell^{N_{d}}_{2}\cap z(k,d)^{-1}\ell^{N_{d}}_{2d/(d+1)}$. Note here that $\mathrm{ker}(\mathrm{Rad}^{\leq d})=\mathcal{P}_{(d,k]}(\M_{2^n})$. Hence, the operator $\widehat{\mathrm{Rad}^{\leq d}}$ induces an bounded operator (which is still denoted by $\widehat{\mathrm{Rad}^{\leq d}}$) from the quotient space $\mathcal{P}^{\leq k}(\M_{2^n})/\mathcal{P}_{(d,k]}(\M_{2^n})$ to $\ell^{N_{d}}_{2}\cap z(k,d)^{-1}\ell^{N_{d}}_{2d/(d+1)}$ with the norm less than $1$.  and the desired
inequality follows directly from the definition of the quotient norm.  The proof is complete.
%Since the quotient norm of $T+\mathcal{P}_{(d,k]}(\M_{2^n})$ is
%$\inf_{G\in\mathcal{P}_{(d,k]}(\M_{2^n})}\|T-G\|_{L_\infty}$, this completes the proof.
\end{proof}


We now apply Lemma \ref{upper bound for Rademacher projection 1} to study the $L_1$-distance of the at most $d$-degree polynomials from the quantum tail spaces, which is 
 a quantum analogy of the approximate theorem of Eskenazis and Zhang \cite[Theorem 3]{EZ2024}.
\begin{theorem}\label{EZ theorem 1}
For any integers $n\geq k\geq d$, there exists a constant $c_{d}>0$ depending only on $d$ such that, for each $T\in \mathcal{P}^{\leq d}(\M_{2^n})$, we have
\[
\|T\|_{L_{1}}\leq K\left(\widehat{T},c_{d}k^{d},\ell^{N_{d}}_{2},\ell^{N_{d}}_{2d/(d-1)}\right),
\]
where $N_{d}\coloneqq \sum_{l=0}^{d}3^{l}\binom{n}{l}$. Consequently, we have
\begin{equation}\label{L1 distance upper}
\inf_{G\in \mathcal{T}^{>k}(\M_{2^n})}\|T-G\|_{L_{1}}\leq K(\widehat{T},c_{d}k^{d},\ell^{N_{d}}_{2},\ell^{N_{d}}_{2d/(d-1)}).
\end{equation}
\end{theorem}
\begin{proof}
The proof follows from an easy duality argument of Lemma \ref{upper bound for Rademacher projection 1} as follows. Taking the adjoint of the operator $\widehat{\mathrm{Rad}^{\leq d}}$, we get that $\widehat{\mathrm{Rad}^{\leq d}}^{*}$ is an operator from $\ell^{N_{d}}_{2}+z(k,d)\ell^{N_{d}}_{2d/(d-1)}$ to $\mathrm{ker}(\widehat{\mathrm{Rad}^{\leq d}}(\M_{2^n}))^{\bot}$ with norm less than $1$.

On the other hand, it is easy to see that


\[
\mathrm{ker}(\widehat{\mathrm{Rad}^{\leq d}}(\M_{2^n}))^{\bot}=\mathcal{P}^{\leq d}(\M_{2^n})\oplus\mathcal P^{>k}(\M_{2^n}), 
\]
If this is the case, the second desired  inequality follows directly.
and
\[
\widehat{\mathrm{Rad}^{\leq d}}^{*}\left(\{x_{\s}\}_{\substack{\s\in \{0,1,2,3\}^{n}\\|\s|\leq d}}\right)=\sum_{\substack{\s\in\{0,1,2,3\}^{n}\\|\s|\leq d}}x_{\s}\sigma_{\s}.
\]
Hence,
\[
\|T\|_{L_{1}}\leq K(\widehat{T},z(k,d),\ell^{N_{d}}_{2},\ell^{N_{d}}_{2d/(d-1)}).
\]
To get the desired bound, it suffices to note that
\[
z(k,d)=\mathrm{BH}_{\mathbb{M}_2}^{\leq d}\sum_{l=0}^{d}|\tilde{c}(k,l)|\leq \mathrm{BH}_{\mathbb{M}_2}^{\leq d}\sum_{l=0}^{d}\frac{k^{l}}{l!}\leq e\mathrm{BH}_{\mathbb{M}_2}^{\leq d}k^{d}.
\]

Now we turn to show \eqref{L1 distance upper}. We assume this inequality does not hold. Then, there exist $T\in \mathcal{P}^{\leq d}(\M_{2^n})$ and $\alpha\in (0,\infty)$ such that for some $d\leq k \in [n]$,
\[A:=\inf_{G\in \mathcal{T}^{>k}(\M_{2^n})}\|T-G\|_{L_{1}}> \alpha > K(\widehat{T},c_{d}k^{d},\ell^{N_{d}}_{2},\ell^{N_{d}}_{2d/(d-1)}).\]
Note that we can always take $G\in \mathcal{T}^{>k}(\M_{2^n})$ such that $\|G\|_{L_1}<\frac{A-\alpha}{2}$. Then
\begin{align*}
\|T\|_{L_1}&\geq \|T-G\|_{L_1}-\|G\|_{L_1}\\
&\geq A-\frac{A-\alpha}{2}\\
&>A-(A-\alpha)=\alpha,
\end{align*}
which contradicts the first assertion of the theorem. 
\end{proof}

Even in the commutative setting, Theorem \ref{EZ theorem 1} is far from optimal; see \cite[Theorem 3]{EZ2024} and its related discussion there. 
However, in the case $d=1$, Theorem \ref{EZ theorem 1} is optimal. We show this by establishing 
 the quantum version of the Oleszkiewicz inequality (Theorem \ref{Oleszkiewicz inequality} above). 

\begin{theorem}\label{thm Ol}
There exist universal constant $C_1,C_2>0$ that 
\[
C_1^{-1} K\left(\widehat{T},m,\ell^{3n}_{2},\ell^{3n}_{\infty}\right)\leq \inf_{G\in \mathcal{T}^{>m}(\M_{2^n})}\|G-T\|_{L_{1}} \leq C_2 K\left(\widehat{T},m,\ell^{3n}_{2},\ell^{3n}_{\infty}\right),
\]
for $T\in \mathcal{P}^{=1}(\M_{2^n})$ and $1\leq m\leq 3n$.
\end{theorem}
\begin{proof}
The right-hand side inequality follows from Theorem \ref{EZ theorem 1}. It suffices to show the left-hand side inequality. 
We here apply the reduction method used in the proof of Theorem \ref{Quantum BH}. Let $\rho_j$ be as in \eqref{rho-j} and define
\[
\gamma:=\frac{\rho_1}{2}\otimes \frac{\rho_2}{2}\cdots \otimes \frac{\rho_n}{2}.
\]
It was already shown in the proof of Lemma \ref{lem rho} that for each $j\in [n]$ and $\varepsilon\in \{-1,1\}^{3n}$, the eigenvalues of $\rho_j(\varepsilon)$  are \(0\) and \(2\). Thus, it follows that for each  $\varepsilon\in \{-1,1\}^{3n}$, $\|\gamma(\varepsilon)\|_{L_{\infty}}=1$.


For $T\in \mathcal{P}^{=1}(\M_{n})$ and $G\in \mathcal{T}^{>m}(\M_{n})$, we define two functions on $\{-1,1\}^{3n}$ as
\[
f(\cdot)\coloneqq \langle \gamma(\cdot),T\rangle,\quad \mbox{and}\quad  g(\cdot)\coloneqq \langle \gamma(\cdot),G\rangle.
\]
The same reason as to \eqref{weak version}, we have $f\in \mathcal{P}^{=1}(\{-1,1\}^{3n})$ and $g\in \mathcal{T}^{>m}(\{-1,1\}^{3n})$. Moreover, we have 
$$\{\widehat{T}(\mathbf{s})\}_{\mathbf{s}}= 2\sqrt{3}\{\widehat{f}(A)\}_{A}.$$
 Applying the Oleszkiewicz inequality (i.e., Theorem \ref{Oleszkiewicz inequality}) to the functions $f$ and $g$, it gives that
\begin{align*}
K\left(\widehat{T},m,\ell^{3n}_{2},\ell^{3n}_{\infty}\right)&=2\sqrt{3} K\left(f,m,\ell^{3n}_{2},\ell^{3n}_{\infty}\right)\\
&\leq C_1^{-1}\mathbb{E}|f-g|\\
&= C_1^{-1}\mathbb{E}\left(\left|\langle\gamma(\varepsilon),T-G \rangle\right|\right)\\
&\leq C_1^{-1}\mathbb{E}\left(\|\gamma(\varepsilon)\|_{L_{\infty}}\right)\|T-G\|_{L_{1}}\\
&=C_1^{-1}\|T-G\|_{L_{1}},
\end{align*}
where the expectation is taken over $\{-1,1\}^{3n}$. Since $G\in \mathcal{T}^{>m}(\M_{2^n})$ is arbitrary, the desired result follows.
\end{proof}


Particularly, they obtained a more accurate estimate for \emph{symmetric polynomials} in \cite[Theorem 4]{EZ2024}. 
Inspired by this,  we introduce the quantum analogies of \emph{symmetric polynomials} as follows.
\begin{definition}
 For each $l\in [n]$, we say
\[
E_{l}\coloneqq \sum_{\substack{\s\in \{0,1,2,3\}^{n}\\|\s|=l}}\sigma_{\s}
\]
is  the \emph{quantum elementary symmetric polynomial of order $l$}. We say $T\in \M_{2^n}$ is a \emph{quantum symmetric polynomial of degree at most $d\in [n]$}  if $T$ is  given by $T\coloneqq\sum_{l=0}^{d}\alpha_{l}E_{l}$ for some sequence $\{\alpha_{l}\}_l\subset \mathbb{C}$.
\end{definition}

We first collect a  basic property for each quantum elementary symmetric polynomial of order $l$. To this end,  for each $j\in [n]$ and $A\subseteq [n]$, we define
\begin{equation}\label{transference element}
u_{j}\coloneqq \frac{1}{\sqrt{3}}\sum_{l=1}^{3}\sigma_{le_{j}}\in \M_{2^n}
\end{equation}
and
\begin{equation}\label{transference Walsh}
u_{A}\coloneqq \prod_{j\in A}u_{j},\quad u_{\emptyset}\coloneqq \mathbf{1}.
\end{equation}
Note here that $u_{i}u_{j}=u_{j}u_{i}$ for each $i$, $j\in [n]$ and $\{u_{A}\}_{A\subseteq [n]}$ forms an orthonormal sequence in $L_{2}(\M_{2^n})$. Indeed, for each $A$, $B\subseteq [n]$ with $A\neq B$, we have
\begin{equation}\label{normal}
\|u_{A}\|^{2}_{L_{2}}=\prod_{j\in A}\|u_{j}\|^{2}_{L_{2}}=1,
\end{equation}
and we assume without loss of generality that $i\in A\setminus B$, then 
\begin{equation}\label{orthogonal}
\langle u_{A},u_{B}\rangle=\mathrm{tr}_{2}(u_{i})\cdot\langle u_{A\setminus \{i\}},u_{B}\rangle=0.
\end{equation}
Let $p\coloneqq \prod_{j=1}^{n}(\mathbf{1}+u_{j})$ be the \emph{Riesz product}. Then it is clear that \begin{equation}\label{p uA}
p=\sum_{A\subseteq [n]}u_{A}. 
\end{equation} 
Since each $\mathbf{1}+u_{j}$ is nonnegative and $\{1+u_j\}_{j\in [n]}$ are commutative, it follows from that $p$ is positive. Thus
\begin{equation}\label{p L1}
\|p\|_{L_{1}}=\mathrm{tr}\left(\sum_{A\subseteq [n]}u_{A}\right)=\mathrm{tr}(u_{\emptyset})=1.
\end{equation}


\begin{lemma}\label{El}
Let $E_l$ be a  elementary symmetric polynomials for some $l\in [n].$ Then
\begin{align*}
E_{l}=3^{l/2}\sum_{\substack{A\subseteq[n]\\|A|=l}}u_{A}.
\end{align*}
\end{lemma}
\begin{proof}
It is clear that
\begin{align*}
E_{l}=\sum_{\substack{\s\in\{0,1,2,3\}^{n}\\|\s|=l}}\sigma_{\s}=\sum_{\substack{A\subseteq [n]\\|A|=l}}\sum_{\substack{\s\in\{0,1,2,3\}^{n}\\\mathrm{supp}(\s)=A}}\sigma_{\s}=3^{l/2}\sum_{\substack{A\subseteq[n]\\|A|=l}}u_{A}.
\end{align*}
\end{proof}


\begin{theorem}\label{upper bound for symmetric polynomials}
Let $n\geq k\geq d$ be integers. Then for each symmetric polynomial $T$ of degree at most $d$, we have
\begin{equation}\label{fundamental L1 bound}
\|T\|_{L_{1}}\leq \sum_{l=0}^{d}3^{l/2}|\alpha_{l}|\cdot|\tilde{c}(k,l)|,\quad T=\sum_{l=0}^{d}\alpha_{l}E_{l}.
\end{equation}
\end{theorem}
\begin{proof}
Since every symmetric polynomial is a linear combination of elementary symmetric polynomials, by the triangle inequality it suffices to bound $\|E_{l}\|_{L_{1}}$ for each $l\in [d]$. 
Take $G\in L_{\infty}(\M_{2^n})$ with $\|G\|_{L_{\infty}}\leq 1$. By Lemma \ref{El}, we have
\begin{equation*}\label{elementary terms}
\mathrm{tr}(E_{l}\cdot G)=3^{l/2}\sum_{\substack{A\subseteq [n]\\|A|=l}}\langle u_{A},G\rangle=3^{l/2}\left\langle \sum_{A\subseteq [n]}u_{A},\mathrm{Rad}^{=l}(G)\right\rangle=3^{l/2}\left\langle p,\mathrm{Rad}^{=l}(G)\right\rangle,
\end{equation*}
where  $p$ is given as in \eqref{p uA}. 
Hence, for arbitrary $G\in L_{\infty}(\M_{n})$ with $\|G\|_{L_{\infty}}\leq 1$, we have
\[
|\mathrm{tr}(E_{l}\cdot G)|\leq 3^{l/2}\|p\|_{L_1}\|\mathrm{Rad}^{=l}(G)\|_{L_{\infty}}\leq 3^{l/2}|\tilde{c}(k,l)|,
\]
where we used the Figiel inequality (i.e., Lemma \ref{Figiel inequality}) and \eqref{p L1} in the last line. This implies $$\|E_{l}\|_{L_{1}}\leq 3^{l/2}|\tilde{c}(k,l)|$$ for each $l\in [d]$.
It now remains to apply the triangle inequality to conclude that
\begin{align*}
\|T\|_{L_{1}}\leq \sum_{l=0}^{d}3^{l/2}|\alpha_{l}|\cdot\|E_{l}\|_{L_{1}}\leq \sum_{l=0}^{d}3^{l/2}|\alpha_{l}|\cdot|\tilde{c}(k,l)|.
\end{align*}
The proof of this theorem is complete.
\end{proof}


\begin{remark}
In the classical setting of Boolean cubes, let $f_{l}$ be the $l$-th elementary symmetric polynomial, defined by
\[
f_{l}(\varepsilon)\coloneqq \sum_{\substack{A\subseteq [n]\\|A|=l}}w_{A}(\varepsilon),\quad \varepsilon\in\{-1,1\}^{n},~\mbox{with }l\in[n].
\]
The upper bound for symmetric polynomials with respect to $L_{1}$ norm is due to Esekenazis and Zhang (see \cite[Theorem 4]{EZ2024}), which states that $\|f_{l}\|_{L_{1}}\leq |\tilde{c}(k,l)|$ for each $n\geq k\geq l$. By contrast, in the quantum counterpart, the factor $3^{l/2}$ appearing on the right hand side of \eqref{fundamental L1 bound} is necessary, as shown in the elementary example below. For the case $n=l=1$, we have
\[
e_{1}=\sigma_{1}+\sigma_{2}+\sigma_{3}=
\begin{pmatrix}
1&1-i\\
1+i&-1
\end{pmatrix}.
\]
Hence, $\|e_{1}\|_{L_{1}}=3^{1/2}$. For this same case, the Chebyshev polynomial of the first kind of degree $1$ is $T_{1}(x)=x$ for all $x\in \mathbb{R}$, which gives $|\tilde{c}(1,1)|=1$. This show the necessity of the desired factor in \eqref{fundamental L1 bound}.
\end{remark}


We now conclude this section with the following lower estimate of quantum analogy of the $L_{1}$-distance of elementary symmetric polynomials from tail spaces.
\begin{theorem}
Let $E_d$ be the quantum elementary symmetric polynomial of order $d$.  For integers $n\geq k\geq d$, there exists a positive function $\varepsilon(n,k,d)$ with $\lim_{n\to \infty}\varepsilon(n,k,d)=0$ such that
\[
3^{d/2}\left(|\tilde{c}(k,d)|-\varepsilon(n,k,d)\right)\leq \inf_{G\in \mathcal{T}^{>k}(\M_{2^n})}\|E_{d}-G\|_{L_{1}}\leq 3^{d/2} |\tilde{c}(k,d)|.
\]
\end{theorem}
\begin{proof}
Let $T_{k}$ be the $k$-th Chebyshev polynomial of first kind, and we define
\[
h_{n,k}(\varepsilon)\coloneqq T_{k}\left(\frac{1}{n}\sum_{j=1}^{n}\varepsilon_{j}\right),\quad\forall~\varepsilon=\{\varepsilon_{j}\}_{j=1}^{n}\in\{-1,1\}^{n}.
\]
It is clear that $\|h_{n,k}\|_{L_{\infty}}\leq \sup_{t\in[-1,1]}|T_{k}(t)|\leq 1$.

Since $h_{n,k}$ is a symmetric polynomial, it follows that
\begin{equation}\label{auxiliary function 1}
h_{n,k}=\sum_{j=0}^{k}\beta_{j,k,n}f_{j},
\end{equation}
where $\beta_{j,k,n}\in \mathbb{R}$ and $f_{l}$ is the $l$-th elementary symmetric polynomial on $\{-1,1\}^{n}$. We now define $\xi_{l}\coloneqq \sum_{\substack{A\subseteq [n]\\|A|=l}}u_{A}$, where $u_{A}$ is given as in \eqref{transference Walsh} for each $A\subseteq [n]$. Furthermore, let
\[
H_{n,k}\coloneqq \sum_{j=0}^{k}\beta_{j,k,n}\xi_{j}=\sum_{j=0}^{k}\frac{\beta_{j,k,n}}{3^{j/2}}E_{j},
\]
where we used Lemma \ref{El}. 
We denote that $\mathcal{F}\coloneqq \mathrm{span}\{u_{A}:A\subseteq [n]\}\subseteq \mathbb{M}_{2^n}$, and define $\Gamma: L_{\infty}(\{-1,1\}^{n})\to \mathcal{F}$ by
\[
f=\sum_{A\subseteq [n]}\widehat{f}(A)w_{A}\mapsto \sum_{A\subseteq [n]}\widehat{f}(A)u_{A}.
\]
From the fact that $\{u_{A}\}_{A\subseteq [n]}$ commute with each other, it is easy to see that $\Gamma$ is a trace preserving $*$-homomorphism. Thus, we have
\begin{equation}\label{isometry on Lp}
\|f\|_{L_{2m}}=\left\|\Gamma(f)\right\|_{L_{2m}},\quad m\in \mathbb{N}.
\end{equation}
Letting $m\to \infty$ on both sides of \eqref{isometry on Lp} yields that
\[
\|f\|_{L_{\infty}}=\left\|\Gamma(f)\right\|_{L_{\infty}},\quad f\in L_{\infty}(\{-1,1\}^{n}).
\]
Due to $\Gamma(h_{n,k})=H_{n,k}$, we have $\|h_{n,k}\|_{L_{\infty}}=\|H_{n,k}\|_{L_{\infty}}=1$.



{\bf Case I.} If $k-d$ is even, we define
\[
\varphi_{d,k,n}\coloneqq \sum_{\substack{0\leq j\leq d\\d-j~\mbox{is even}}}\mathrm{sign}(\beta_{j,k,n})E_{j}.
\]
By Theorem \ref{upper bound for symmetric polynomials}, it follows that
\begin{equation}\label{upper bound}
\|\varphi_{d,k,n}\|_{L_{1}}\leq \sum_{\substack{0\leq j\leq d\\d-j~\mbox{is even}}}3^{j/2}|\tilde{c}(k,j)|.
\end{equation}

For any $G\in \mathcal{T}^{>k}(\M_{2^n})$, it follows from $\|H_{k,n}\|_{L_{\infty}}\leq 1$ that
\begin{equation}\label{test function}
\begin{split}
\|\varphi_{d,k,n}-G\|_{L_{1}}&\geq\mathrm{tr}\left(\varphi_{d,k,n} H_{k,n}\right)\\
&=\sum_{\substack{0\leq d\leq j\\d-j~\mbox{is even}}}\frac{|\beta_{j,k,n}|}{3^{j/2}}\|E_{j}\|^{2}_{L_{2}}\\
&=\sum_{\substack{0\leq j\leq d\\ d-j~\mbox{is even}}}3^{j/2}|\beta_{j,k,n}|\binom{n}{j}\\
&=\sum_{\substack{0\leq j\leq d\\d-j~\mbox{is even}}}3^{j/2}|\beta_{j,k,n}|f_{j}(1,\dots,1).
\end{split}
\end{equation}
Since \eqref{test function} holds for arbitrary $G\in \mathcal{T}^{>k}(\mathbb{M}_{2^n})$, it yields
\begin{equation}\label{inf-version}
\inf_{G\in \mathcal{T}^{>k}(\M_{2^n})}\|\varphi_{d,k,n}-G\|_{L_{1}}\geq \sum_{\substack{0\leq j\leq d\\d-j~\mbox{is even}}}3^{j/2}|\beta_{j,k,n}|f_{j}(1,\dots,1).
\end{equation}
It is clear that $E_{d}=\mathrm{sign}(\beta_{d,k,n})(\varphi_{d,k,n}-\varphi_{d-2,k,n})$ with $\varphi_{-1,k,n}=\varphi_{0,k,n}\equiv 0$. By the triangle inequality and \eqref{upper bound} we get that
\begin{align*}
\inf_{G\in \mathcal{T}^{>k}(\M_{2^n})}&\left\|E_{d}-G\right\|_{L_{1}}\\
&\geq \inf_{G\in \mathcal{T}^{>k}(\mathbb{M}_{2^n})}\left\|\varphi_{d,k,n}-G\right\|_{L_{1}}-\inf_{G\in \mathcal{T}^{>k}(\mathbb{M}_{2^n})}\left\|\varphi_{d-2,k,n}-G\right\|_{L_{1}}\\
&\geq \sum_{\substack{0\leq j\leq d\\d-j~\mbox{is even}}}3^{j/2}|\beta_{j,k,n}|f_{j}(1,\dots,1)-\sum_{\substack{0\leq j\leq d-2\\d-2-j~\mbox{is even}}}3^{j/2}|\tilde{c}(k,j)|\\
&\geq \sum_{\substack{0\leq j\leq d,\\d-j~\mbox{is even}}}3^{j/2}|\tilde{c}(k,j)|-\sum_{\substack{0\leq j\leq d\\d-j~\mbox{is even}}}\frac{2(k+1)!3^{j/2}\max_{l\geq j}|\tilde{c}(k,l)|}{n}\\
&\;\;\;-\sum_{\substack{0\leq j\leq d-2\\d-2-j~\mbox{is even}}}3^{j/2}|\tilde{c}(k,j)|\\
&\geq 3^{d/2}\left(|\tilde{c}(k,d)|-\frac{2d(k+1)!\max_{k\geq l\geq 0}|\tilde{c}(k,l)|}{n}\right),
\end{align*}
where we used \cite[Lemma 27]{IRRRY2021} in the third inequality. This completes the proof for the case $k-d$ is even.

{\bf Case II.} If $k-d$ is odd, we need to reduce the approach to the even case by considering $k-d-1$. Indeed, we define $\varphi_{-1,k-1,n}=\varphi_{0,k-1,n}\equiv 0$ and
\[
E_{d}=\mathrm{sign}(\beta_{d,k-1,n})(\varphi_{d,k-1,n}-\varphi_{d-2,k-1,n})\quad 0\leq d<k.
\]
Recall that
\begin{equation}\label{key parameter}
|\tilde{c}(k,d)|=
\begin{cases}
\frac{|T_{k}^{(d)}(0)|}{d!},&\mbox{if }k-d~\mbox{is even},\\
\frac{|T_{k-1}^{(d)}(0)|}{d!},&\mbox{if }k-d~\mbox{is odd}.
\end{cases}
\end{equation}
Hence,
\begin{align*}
\inf_{G\in \mathcal{T}^{>k}(\mathbb{M}_{2^n})}\|E_{d}-G\|_{L_{1}}&\geq \inf_{G\in \mathcal{T}^{>k}(\mathbb{M}_{2^n})}\|\varphi_{d,k-1,n}-G\|_{L_{1}}-\inf_{G\in \mathcal{T}^{>k}(\mathbb{M}_{2^n})}\|\varphi_{d-2,k-1,n}-G\|_{L_{1}}\\
&\geq 3^{d/2}\left(|\tilde{c}(k,d)|-\frac{2d(k+1)!\max_{k\geq l\geq 0}|\tilde{c}(k,l)|}{n}\right),
\end{align*}
where we used \cite[Lemma 27]{IRRRY2021} and \eqref{key parameter} in the second inequality. The proof is complete.
\end{proof}




\chapter{Learning quantum observables}\label{sec QLT}
In this chapter, we have two main objectives. The first is to investigate the optimal asymptotics of the quantum random query model in quantum learning theory, and the second is to study quantum learning problems of related types.
To describe the quantum learning problem, let us first review the classical one. 
Let $\mathcal{F}$ be a class of functions on $\{-1,1\}^n$.  Given $f\in \mathcal{F}$ and access to examples $(X_1, f(X_1))$, $\dots$, $(X_Q, f(X_Q))$, the goal is to algorithmically construct a hypothesis function $h : \{-1,1\}^n \to \mathbb{R}$ which effectively approximates $f$. Different access models to examples give rise to concrete versions of the learning problem. The two most standard  models are the query model and random query model which are stated as follows
\begin{problem}[Query model] Let $\epsilon\in (0,1)$.  Determine  the least number $Q=Q_{\{\pm1\}}(\mathcal{F},\epsilon)$ such that for any $f\in \mathcal{F}$ together with $Q$ queries
$$(X_1, f(X_1)), \dots, (X_Q, f(X_Q)),$$
 one can construct a function $h: \{-1,1\}^n \to \mathbb{R}$ such that 
$$\|h - f\|_{L_2}^2 \leq \varepsilon.$$
\end{problem}

\begin{problem}[Random query model]
 Let $\epsilon,\delta\in (0,1)$.  Determine  the least number $Q_r=Q_{\{\pm1\},r}(\mathcal{F},\epsilon,\delta)$ such that for any $f\in \mathcal{F}$ together with $Q$ random queries
$$(X_1, f(X_1)), \dots, (X_Q, f(X_Q)),$$
 one can construct a random function $h: \{-1,1\}^n \to \mathbb{R}$ such that 
$$\|h - f\|_{L_2}^2 \leq \varepsilon$$
with probability at least $1-\delta$. 
\end{problem}
The least number $Q$ (resp. $Q_r$) of examples required to solve the learning problem in each case is called the query complexity of the  model. The query complexity of learning problems has been studied extensively for various classes $\mathcal{F}$ of functions on the discrete Boolean $\{-1,1\}^n$.  Set
\begin{equation}\label{Fnd}
\mathcal{F}_{n,d} \stackrel{\text{def}}{=} \left\{ f : \{-1,1\}^n \to [-1,1] : f \text{ has degree at most } d \right\} .
\end{equation}
In \cite{EI2022}, by  applying the  Bohnenblust--Hille inequality \eqref{classical BH}, Eskenazis and Ivanisvili,  proved that 
\begin{equation}\label{EI 2022}
Q_{\{\pm1\},r}(\mathcal{F}_{n,d},\epsilon,\delta)\leq c_d \frac{1}{\epsilon^{d+1}} \log\left(\frac{n}{\delta}\right),
\end{equation}
where $c_d$ depends on the Bohnenblust--Hille constant $\text{BH}_{\{\pm1\}}^{\leq d}$. This improves significantly the previous work studied in \cite{LMN1993} and \cite{IRRRY2021}. Recently, Eskenazis, Ivanisvili, and Streck \cite{EIS2023} established the sharpness as of the above upper bound \eqref{EI 2022} as $n\to \infty$. More precisely, they showed that there exists constant $c_{d,\epsilon,\delta} $ such that 
\begin{equation}\label{EIS 2023}
Q_{\{\pm1\},r}(\mathcal{F}_{n,d},\epsilon,\delta)\geq c_{d,\epsilon,\delta}\log(n).
\end{equation}
In the same paper, Eskenazis, Ivanisvili, and Streck \cite{EIS2023} also established a  more robust version of the upper bound \eqref{EI 2022} that applies to different concept classes $\mathcal{F}$ which are not necessarily of bound degree.

Now suppose that we want to learn quantum observable $T$ over $n$-qubits, i.e. $T \in \M_{2^n}$. Recall that a matrix $\rho\in \M_{2^n}$ is a density matrix whenever $\rho\geq0$  and $\mathrm{tr}(\rho)=1$. Let $\mathcal{S}$ be some set of density matrices in $\M_{2^n}$. The quantum learning problem is similar to the classical setting, however the queries are now replaced with
\[
\operatorname{tr}[T\rho], \quad \rho \sim \mathcal{S},
\]
where $\rho$ are density matrices in $\M_{2^n}$.    We elaborate on the concrete matters of quantum learning problems as follows. 
\begin{problem}[Quantum query model] Let $\mathcal{Q}\subset \M_{2^n}$ and let $\epsilon\in (0,1)$.  Determine  the least number $Q_{\M_2}=Q_{\M_2}(\mathcal{Q},\epsilon)$ such that for any $T\in \mathcal{Q}$ together with $Q_{\M_2}$ queries
$$(\rho^{(1)}, \tr(T\rho^{(1)})), \dots, (\rho^{(Q_{\M_2})}, \tr(T\rho^{(Q_{\M_2})})), \quad \rho^{(k)}\in \mathcal{S},$$
 one can construct another  observable $\widetilde{T}\in \M_{2^n}$ such that 
$$\|T-\widetilde{T}\|_{L_2}^2 \leq \epsilon.$$
\end{problem}

\begin{problem}[Quantum random query model]
 Let $\mathcal{Q}\subset \M_{2^n}$ and let $\epsilon,\delta\in (0,1)$.  Determine  the least number $Q_{\M_2,r}=Q_{\M_2,r}(\mathcal{Q},\epsilon,\delta)$ such that for any $T\in \mathcal{Q}$ together with $Q_{\M_2,r}$ queries
$$(\rho^{(1)}, \tr(T\rho^{(1)})), \dots, (\rho^{(Q_{\M_2,r})}, \tr(T\rho^{(Q_{\M_2,r})})), \quad \rho^{(k)}\in \mathcal{S},$$
 one can construct another `random' observable $\widetilde{T}\in \M_{2^n}$ such that 
$$\|T-\widetilde{T}\|_{L_2}^2 \leq \epsilon$$
with probability at least $1-\delta$. 
\end{problem}

Recently, the learning problems similar to the upper bound \eqref{EI 2022} for low degree
quantum observables were investigated by Volberg and Zhang \cite{VZ2024} via the newly constructed quantum Bohnenblust--Hille inequality (see e.g. Theorem \ref{Quantum BH}). For some $d\in [n]$, let
\[
\mathcal Q_{n,d}
:=
\{T=T^*\in \M_{2^n}: \|T\|_\infty\leq 1,\, T\in \mathcal{P}^{\leq d}(\M_{2^n})\}. 
\]
Volberg and Zhang's result states that for $\epsilon,\delta\in (0,1)$, 
\begin{equation}\label{VZ logn}
Q_{\M_2,r}(\mathcal{Q}_{n,d},\epsilon,\delta)\leq \frac{C^{d\sqrt{d\log d}}}{\epsilon^{d+1}}\log(\frac{n}{\delta})
\end{equation}
for some large enough $C>0$. Our first goal of the chapter is to show that \eqref{VZ logn} is asymptotically optimal as $n\to \infty$. 

\begin{theorem}\label{thm sharpness}
For $\epsilon,\delta\in (0,1)$, we have 
$$Q_{\M_2,r}(\mathcal{Q}_{n,d},\epsilon,\delta)\simeq_{d,\epsilon,\delta}\log(n). $$
\end{theorem}

This is the quantum analogue of \eqref{EIS 2023}. In Chapter \ref{sec sharpness}, we provide its proof. 

\section{Asymptotic sharpness}\label{sec sharpness}
The main objective here is to prove Theorem \ref{thm sharpness}. 
Since we already know \eqref{VZ logn}, it remains to show 
\begin{equation}\label{desired random}
Q_{\M_2,r}(\mathcal{Q}_{n,d},\epsilon,\delta)\geq c_{d,\epsilon,\delta}\log(n).
\end{equation}
 Let
\[
\mathcal B_n:=\{T\in \mathbb M_{2^n}: T=T^*=T^{-1}\}
\]
and let
\[
\mathcal B_{n,d}:=\{T\in \mathcal B_n:T\in \mathcal{P}^{\leq d}(\M_{2^n})\},\quad d\in [n].
\]
As it is clear that $Q_{\M_2,r}(\mathcal{Q}_{n,d},\epsilon,\delta)\geq Q_{\M_2,r}(\mathcal{B}_{n,d},\epsilon,\delta)$, then \eqref{desired random} follows by  combining Theorem \ref{boolean-Proposition-8-packing} and Theorem \ref{low bdd} (or Theorem \ref{quantum packing}). 

\begin{proof}[Proof of Theorem \ref{thm sharpness}]
The assertion follows by \eqref{VZ logn}, Theorem \ref{boolean-Proposition-8-packing} and Theorem \ref{low bdd} or Theorem \ref{quantum packing}. 
\end{proof}

Now our primary remaining task is to establish Theorem \ref{boolean-Proposition-8-packing} and Theorem \ref{low bdd}. 
To proceed, let us recall the \emph{packing number} of a metric space $(\mathcal{M},\wp)$ as follows. For any $\epsilon>0$, the subset $\mathcal{N}\subseteq \mathcal{M}$ is said to be $\epsilon$-packing, if for each $\xi$, $\xi^{\prime}\in \mathcal{N}$, we have $\wp(\xi,\xi^{\prime})>\epsilon$. For arbitrary $\epsilon>0$, the packing number of the metric space $(\mathcal{M},\wp)$, denoted by $M(\mathcal{M},\wp,\epsilon)$, is defined by
\[
M(\mathcal{M},\wp,\epsilon)\coloneqq \sup\left\{|\mathcal{N}|:\mathcal{N}\subseteq \mathcal{M},~\mbox{and }\mathcal{N}~\mbox{is an }\mbox{$\epsilon$-packing subset}\right\}.
\]

 

The next result provides a lower bound for the (random) query complexity  	$Q_{\M_2,r}(\mathcal Q_{n,d},\epsilon,\delta)$ via the packing number. 
\begin{theorem}\label{boolean-Proposition-8-packing}
Let $d\in [n]$.	For every $\epsilon>0$ and every $\delta\in(0,1)$,
	\begin{equation}\label{Qbr-20}
		Q_{\M_2,r}(\mathcal Q_{n,d},\epsilon,\delta)
		\geq \log_2 M(\mathcal Q_{n,d},\|\cdot\|_2,2\sqrt{\epsilon})+\log_2(1-\delta).
	\end{equation}
\end{theorem}

\begin{proof}
 For simplicity,  let
	\[
M:=M(\mathcal Q_{n,d},\|\cdot\|_2,2\sqrt{\epsilon}).
	\]
By the definition of the packing number, we can take \(A_1,\cdots,A_M\in \mathcal Q_{n,d}\) such that 
\begin{equation}\label{geq epsilon}
\|A_i-A_j\|_{L_2}>2\sqrt{\epsilon},\quad i\neq j. 
\end{equation} 


A random-example learning algorithm of length $Q$ is a deterministic map
\[H:(\mathcal S\times\{-1,1\})^Q\to L_2(\mathbb M_{2^n},\mathrm{tr}).\]
According to the proof of \eqref{VZ logn} provided in \cite{VZ2024}, we can assume that the samples are generated by
\[\varepsilon^{(1)},\dots,\varepsilon^{(Q)}\in \{-1,1\}^{3n}.\]
Define a probability  measure on \(\{-1,1\}\) with respect to some $A\in \mathcal{Q}_{n,d}$ and $\varepsilon\in \{-1,1\}^{3n}$ by setting
		\begin{equation}\label{P-YA}
			\mathbb{P}\bigl(Y=\pm1\mid \varepsilon,A\bigr)
			=\mathrm{tr} \left(\frac{1\pm A}{2}\rho\bigl(\varepsilon\bigr)\right)
			=\frac{1\pm \mathrm{tr}\bigl(A\rho(\varepsilon)\bigr)}{2},
		\end{equation}
		where $\rho$ is defined as in \eqref{rho}. 
		
In the sequel, we denote by
		\[
		Q:=Q_{\M_2,r}(\mathcal Q_{n,d},\epsilon,\delta)
		\]
		the least $Q\in\mathbb N$ such that there exists such an algorithm satisfying for each $A\in\mathcal Q_{n,d}$
		\begin{equation}\label{Hr-de}
			\mathbb{P}_{\varepsilon^{(1)},\dots,\varepsilon^{(Q)}\in\{-1,1\}^{3n}}\Bigl(
			\bigl\|H\bigl((\rho(\varepsilon^{(1)}),Y_1),\dots,(\rho(\varepsilon^{(Q)}),Y_Q)\bigr)-A\bigr\|_2^2\le \epsilon
			\Bigr)
			\ge 1-\delta.
		\end{equation}
Let
	\[\widetilde{\varepsilon}=(\varepsilon^{(1)},\dots,\varepsilon^{(Q)})\]
be  independent random vectors such that there exists such an algorithm satisfying  $ \rho(\varepsilon^{(t)})\in\mathcal S, t=1,\dots,Q$ and \eqref{Hr-de} for some
learning algorithm $H$. For each $i\in\{1,\dots,M\}$ and each fixed $\widetilde{\varepsilon}$, let $\mathbb P_{i,\widetilde{\varepsilon}}$
denote the conditional distribution on $\{-1,1\}^Q$ of the measurement outcome string when the target observable is $A_i$. That is,  by \eqref{P-YA},  we have
\[\mathbb P_{i,\widetilde{\varepsilon}}(Y)=\prod_{t=1}^Q\frac{1+Y_t\mathrm{tr}(A_i\rho(\varepsilon^{(t)}))}{2},\quad Y=(Y_1,\cdots,Y_Q)\in \{-1,1\}^Q.\]
It is clear that $\mathbb P_{i,\widetilde{\varepsilon}}$ is a probability measure on \(\{-1,1\}^Q\).
 For each $i\in\{1,\dots,M\}$, define
	\[
	E_i(\widetilde{\varepsilon}):=\{Y\in\{-1,1\}^Q:\|H\bigl((\rho(\varepsilon^{(1)}),Y_1),\dots,(\rho(\varepsilon^{(Q)}),Y_Q)\bigr)-A_i\|_2^2\leq \epsilon\}.
	\]
	For every fixed $\widetilde{\varepsilon}$,  the sets $	E_i(\widetilde{\varepsilon}),\dots,	E_M(\widetilde{\varepsilon})$,
	are pairwise disjoint. Indeed, if $Y\in E_i(\widetilde{\varepsilon})\cap E_j(\widetilde{\varepsilon})$ with
	$i\neq j$, then
	\begin{align*}
	\|A_i-A_j\|_2
	&\le \|A_i-H((\rho(\varepsilon^{(1)}),Y_1),\cdots,(\rho(\varepsilon^{(Q)}),Y_Q))\|_2\\
	&\qquad+\|H((\rho(\varepsilon^{(1)}),Y_1),\cdots,(\rho(\varepsilon^{(Q)}),Y_Q))-A_j\|_2
	\le 2\sqrt{\varepsilon},
	\end{align*}
	contradicting the packing property \eqref{geq epsilon}. Hence,  by \eqref{Hr-de}, for each \(i=1,\dots,M\), we have
	\begin{align*}
		\mathbb E_{\widetilde{\varepsilon}}&\big[\mathbb{P}_{i,\widetilde{\varepsilon}}(E_i(\widetilde{\varepsilon}))\big]=\mathbb E_{\widetilde{\varepsilon}}\big[\mathbb P(Y\in E_i(\widetilde{\varepsilon})|\widetilde{\varepsilon},A_i)\big]\\
		&=\mathbb{P}_{\varepsilon^{(1)},\dots,\varepsilon^{(Q)}\in\{-1,1\}^{3n}}\Bigl(
		\bigl\|H\bigl((\rho(\varepsilon^{(1)}),Y_1),\dots,(\rho(\varepsilon^{(Q)}),Y_Q)\bigr)-A_i\bigr\|_2^2\le \epsilon
		\Bigr)\\
		&\geq 1-\delta.
	\end{align*}
	Summing over $i$,  by the disjointness of  $\{E_i(\widetilde{\varepsilon})\}_i$, we obtain
	\begin{align*}
		M(1-\delta)
		&\le \sum_{i=1}^M \mathbb E_{\widetilde{\varepsilon}}
		\big[\mathbb{P}_{i,\widetilde{\varepsilon}}(E_i(\widetilde{\varepsilon}))\big] \\
		&= \mathbb E_{\widetilde{\varepsilon}}
		\Big[\sum_{Y\in\{-1,1\}^Q}\sum_{i=1}^M
		\chi_{E_i(\widetilde{\varepsilon})}(Y)\mathbb P_{i,\widetilde{\varepsilon}}(Y)
		\Big] \\
		&\le \mathbb E_{\widetilde{\varepsilon}}
		\Big[\sum_{Y\in\{-1,1\}^Q}1\Big]
		= 2^Q.
	\end{align*}
 Taking logarithms gives \eqref{Qbr-20}. The proof is complete.
\end{proof}




\begin{remark}
Theorem \ref{boolean-Proposition-8-packing} does not require $A^2=I$, self-adjointness and contraction $$\|A\|_{L_\infty}\leq 1$$ suffices. Indeed, $\{-1,1\}^Q$ is not the range of $A$ but of the outcomes $Y$ produced by the two-outcome $\{\tfrac{1+A}{2},\tfrac{1-A}{2}\}$ in \eqref{P-YA}, which is valid for any self-adjoint $A$ with $\|A\|_{L_\infty}\leq 1$, this is exactly the condition $\tfrac{1\pm A}{2}\geq0$, and satisfies $\mathbb{E}[Y_t]=\mathrm{tr}(A\rho(\varepsilon^{(t)}))$ regardless of whether $A$ is Boolean. Consequently the proof gives \eqref{Qbr-20} for all of $\mathcal{Q}_{n,d}$, not merely its Boolean subclass.
		
The classical counterpart of this observation is that the randomized bound in \cite[(20) of Proposition 8]{EIS2023} already holds for the full class $\mathcal{F}_{n,d}$ of bounded (not necessarily $\{-1,1\}$-valued) functions, with no Boolean assumption: querying an oracle that, on input $x$, outputs $+1$ with probability $\tfrac{1+f(x)}{2}$ and $-1$ with probability $\tfrac{1-f(x)}{2}$ gives an unbiased $\{-1,1\}$-valued estimate of $f(x)$, and running the same counting argument with this randomized oracle in place of $f$ itself recovers Proposition 8 (20) unchanged. 
\end{remark}


%{\color{red} in the following, \cite[Theorem 9 (22)]{EIS2023}. Luo prove it in another draft.}
%\begin{lemma}[{\cite[Lemma 12]{EIS2023}}]\label{m-1t-omega}
%	Let \(\Omega\) be a finite set with \(|\Omega|=m\), let \(k<m\), and let \(\epsilon\in(0,1)\). Then there exist
%	\(t\ge (2k)^{-k/2}m^{(1-\epsilon)k/2}\)	 and subsets \(\Sigma_1,\dots,\Sigma_t\subseteq\Omega\)  of size $k$ satisfying
%	\[|\Sigma_r\cap\Sigma_s|<(1-\epsilon)k,\quad r\neq s\in\{1,\cdots,t\}.\]
%\end{lemma}
%
%
%\begin{theorem}\label{low-le-Theorem9(22)nc-packing}
%	Fix \(n\in\mathbb N\), \(d\in\mathbb N\) with \(d\le \log_2 n\), and let \(\epsilon\in(0,1)\). Then
%	\begin{equation}\label{nc-packing-sqrt2e}
%		\log_2 M\bigl(\mathcal B_{n,d},\|\cdot\|_2,\sqrt{2\epsilon}\bigr)
%		\ge
%		(1-\epsilon)2^{d-2}\log_2  3(n-d+1)-d2^{d-2}.
%	\end{equation}
%\end{theorem}
%\begin{proof} Set $k:=2^{d-1}$, and 
%\[\mathcal U:=\{\sigma_\kappa^{(j)}: j=d,\dots,n,\ \kappa=1,2,3\},
%	\qquad
%	m:=|\mathcal U|=3(n-d+1).
%	\]
%	We may order the elements of set \(\mathcal U\)
%	according to the lexicographical order with double indices, that is,
%	\[\sigma_1^{(d)}<\sigma_2^{(d)}<\sigma_3^{(d)}<\sigma_1^{(d+1)}<\cdots<\sigma_3^{(n)}.\]
%	Since \(d\le \log_2 n\), we have \(k=2^{d-1}\le n/2<m\), so Lemma~\ref{m-1t-omega} applies to \(\Omega=\mathcal U\).
%	For \(\epsilon=(\epsilon_1,\cdots,\epsilon_{d-1})\in\{-1,1\}^{d-1}\), let
%	\[P_\epsilon=\frac{1}{2^{d-1}}\bigotimes_{r=1}^{d-1}(I_1+\epsilon_r\sigma_1)\otimes I^{\otimes(n-d+1)},\]
%	 and we can enumerate the address set
%	\(\{\pm1\}^{d-1}\) in lexicographic order as
%	\[
%	\epsilon^{(1)},\dots,\epsilon^{(k)}.
%	\]
%	For each \(\ell\in\{1,\dots,k\}\), let
%	\begin{equation}\label{pl-def}
%	P_\ell:=P_{\epsilon^{(\ell)}}
%	=\frac{1}{2^{d-1}}\bigotimes_{r=1}^{d-1}(I_2+\epsilon_r^{(\ell)}\sigma_1)\otimes I^{\otimes(n-d+1)}.
%	\end{equation}
%	Then the family \(\{P_\ell\}_{\ell=1}^k\) consists of pairwise orthogonal projections,
%	\[
%	P_\ell P_j=0 \quad \ell\neq j,
%	\qquad
%	\sum_{\ell=1}^k P_\ell=I,
%	\qquad
%	\mathrm{tr}(P_\ell)=2^{-(d-1)}.
%	\]
%	
%	By Lemma~\ref{m-1t-omega}, there exist subsets \(\Sigma_1,\dots,\Sigma_t\subseteq\mathcal U\), each of size
%	\(k\), such that
%	\[
%	|\Sigma_r\cap\Sigma_s|<(1-\epsilon)k,
%	\qquad r\neq s,
%	\]
%	and
%	\[
%	t\ge (2k)^{-k/2}m^{(1-\epsilon)k/2}.
%	\]
%	For each \(r\in\{1,\dots,t\}\), write the ordered set \(\Sigma_r\) as
%	\[
%	\Sigma_r=\{U_{1}<\cdots<U_{k}\},
%	\]
%	and define the controlled observable tree
%	\[
%	A_{\Sigma_r}:=\sum_{\ell=1}^k P_\ell U_{\ell}.
%	\]
%	We claim that \(A_{\Sigma_r}\in B_{n,d}\) for every \(\Sigma_r\).
%	
%	Indeed, for each \(U_{\ell}\), \(U_{\ell}^*=U_{\ell}\), \(\mathrm{deg}(U_{\ell})=1\), and it is supported on one of the leaf qubits \(d,\dots,n\), by the definition of \eqref{pl-def}, we have
%	\begin{equation}\label{pl-com}
%	P_\ell U_{j}=U_{j}P_\ell
%	\end{equation} 
%	\[A_{\Sigma_r}^*=\sum_{\ell=1}^k U_{\ell}P_\ell=A_{\Sigma_r}.
%	\]
%	Using pairwise orthogonality of the \(P_\ell\)'s, we get
%	\[
%	A_{\Sigma_r}^2
%	=\sum_{\ell,j=1}^k P_\ell P_j U_{\ell}U_{j}
%	=\sum_{\ell=1}^k P_\ell U_{\ell}^2
%	=\sum_{\ell=1}^k P_\ell
%	=I.
%	\]
%	Therefore \(A_{\Sigma_r}=A_{\Sigma_r}^*=A_{\Sigma_r}^{-1}\), i.e. \(A_{\Sigma_r}\in \mathcal B_n\). Finally, \(\mathrm{deg}(P_\ell)=d-1\), so every  \(P_\ell U_{\ell}\) has degree at most
%	\(d\), hence \(A_{\Sigma_r}\in \mathcal B_{n,d}\).
%	
%	We now estimate the pairwise Hilbert--Schmidt distances. Fix \(r,s\in\{1,2,\cdots,k\}\), and 
%	\[\Sigma_s=\{V_1<V_2<\cdots<V_k\}\]
%	Since the projections \(P_\ell\) are pairwise orthogonal and \eqref{pl-com},
%	\begin{align*}
%		\|A_{\Sigma_r}-A_{\Sigma_s}\|_2^2
%		&= \mathrm{tr}\Big(\big(\sum_{\ell=1}^k P_\ell(U_\ell-V_\ell)\big)^2\Big) \\
%		&= \sum_{\ell=1}^k \mathrm{tr}\bigl(P_\ell(U_\ell-V_\ell)^2\bigr).
%	\end{align*}
%	Now each term can be computed explicitly. Since \(P_\ell\) acts on the address tensor factor and
%	\(U_\ell-V_\ell\) acts on the leaf tensor factor, and since \(\widetilde P_\ell\) is rank one,
%	\[
%	\mathrm{tr}\bigl(P_\ell(U_\ell-V_\ell)^2\bigr)
%	= \mathrm{tr}_{d-1}(\widetilde P_\ell) \mathrm{tr}_{n-d+1}\bigl((U_\ell-V_\ell)^2\bigr)=2^{-(d-1)}\mathrm{tr}_{n-d+1}\bigl((U_\ell-V_\ell)^2\bigr).\]
%	If \(U_\ell=V_\ell\), the term is zero. If \(U_\ell\neq V_\ell\), 
%	\[  \mathrm{tr}_{n-d+1}\bigl((U_\ell-V_\ell)^2\bigr)= 2.
%	\]
%	Consequently,
%	\[
%	 \mathrm{tr}\bigl(P_\ell(U_\ell-V_\ell)^2\bigr)
%	= 2^{-(d-2)}\mathbf \chi_{\{U_\ell\neq V_\ell\}},
%	\]
%	and hence
%	\begin{equation}\label{Ars-L2}
%		\|A_{\Sigma_r}-A_{\Sigma_s}\|_2^2
%		= 2^{-(d-2)}\#\{\ell\in\{1,\dots,k\}: U_\ell\neq V_\ell\}.
%	\end{equation}
%	
%	We next compare this quantity with the intersection size \(|\Sigma_r\cap\Sigma_s|\). If \(U_\ell=V_\ell\), then  \(U_\ell=V_\ell\in\Sigma_r\cap\Sigma_s\). Hence
%	\[
%	\#\{\ell:U_\ell=V_\ell\}\le |\Sigma_r\cap\Sigma_s|,
%	\]
%	so
%	\[
%	\#\{\ell:U_\ell\neq V_\ell\}
%	= k-\#\{\ell:U_\ell=V_\ell\}
%	\ge k-|\Sigma_r\cap\Sigma_s|.
%	\]
%	Using Lemma \ref{m-1t-omega}, we obtain \(|\Sigma_r\cap\Sigma_s|<(1-\epsilon)k\) and
%	\[
%	\#\{\ell:U_\ell\neq V_\ell\}>\epsilon k=\epsilon 2^{d-1}.
%	\]
%	Combining this with \eqref{Ars-L2} yields
%	\[
%	\|A_{\Sigma_r}-A_{\Sigma_s}\|_2^2>2^{-(d-2)}\epsilon 2^{d-1}=2\epsilon,
%	\]
%	that is,
%	\[
%	\|A_{\Sigma_r}-A_{\Sigma_s}\|_2>\sqrt{2\epsilon}.
%	\]
%	Therefore \(\{A_1,\dots,A_t\}\subseteq \mathcal B_{n,d}\) is a \(\sqrt{2\epsilon}\)-packing of \(\mathcal B_{n,d}\), and so
%	\[
%	M\bigl(\mathcal B_{n,d},\|\cdot\|_2,\sqrt{2\epsilon}\bigr)\ge t\geq (2k)^{-k/2}m^{(1-\epsilon)k/2}.
%	\]
%	Taking binary logarithms and recalling that \(2k=2^d\) and \(k=2^{d-1}\), we obtain
%	\begin{align*}
%		\log_2 M\bigl(B_{n,d},\|\cdot\|_2,\sqrt{2\epsilon}\bigr)
%		&\ge \log_2 t \\
%		&\ge -\frac{k}{2}\log_2(2k)+\frac{(1-\epsilon)k}{2}\log_2 m \\
%		&= -d2^{d-2}+(1-\epsilon)2^{d-2}\log_2\!\bigl(3(n-d+1)\bigr).
%	\end{align*}
%The proof is complete.
%\end{proof}



Now we turn to derive the packing bound for quantum  functions. We need the following combinatorial lemma, which is analogous to \cite[Lemma 12]{EIS2023}. For each $m$, $k\in \mathbb{N}$ with $k\leq m$, let $\{0,1,2,3\}^{m}_{=k}$ be the set of all sequences with length $k$, more precisely,
\[
\{0,1,2,3\}^{m}_{=k}\coloneqq \left\{\eta\in \{0,1,2,3\}^{m}:\sum_{j=1}^{m}\mathrm{sign}(\eta(j))=k\right\}.
\]
We equip $\{0,1,2,3\}^{m}_{=k}$ with the uniform distribution $\mu$, that is, $\mu(\eta)=\frac{1}{\binom{m}{k}3^{k}}$ for each $\eta\in\{0,1,2,3\}^{m}_{=k}$. For $\eta_{1}$ and $\eta_{2}\in\{0,1,2,3\}^{m}_{=k}$, we define $$|\eta_{1}\cap\eta_{2}|\coloneqq |\{j\in [m]:\eta_{1}(j)= \eta_{2}(j)\neq 0\}|.$$
\begin{lemma}\label{packing number lemma}
For any $m$, $k\in \mathbb{N}$ with $k\leq m$ and every $0<\epsilon<\frac{1}{2}$, there exist at least $\{\eta_{j}\}_{j=1}^{t}\subseteq \{0,1,2,3\}^{m}_{=k}$ with $t\geq \left(\frac{3}{2k}\right)^{k}m^{(1-\epsilon)k}\left(\frac{m}{3m+2}\right)^{\epsilon k}$ such that
\[
|\eta_{i}\cap \eta_{j}|<(1-\epsilon)k,\quad i\neq j,~i,~j\in [t].
\]
\end{lemma}
\begin{proof}
We now provide a direct proof invoking the greedy algorithm. For a fixed $\eta\in \{0,1,2,3\}^{m}_{=k}$, we denote
\[
D(\eta,(1-\epsilon)k)\coloneqq \{\zeta\in \{0,1,2,3\}^{m}_{=k}:|\eta\cap \zeta|\geq (1-\epsilon)k\}.
\]
It is clear that
\begin{equation}\label{combinatorics 1}
\begin{split}
|D(\eta,(1-\epsilon)k)|&\leq\sum_{j=\lceil (1-\varepsilon)k\rceil}^{k}\sum_{l=0}^{k-j}\binom{k}{j}\binom{k-j}{l}\binom{m-k}{k-j-l}2^{l}3^{k-j-l}\\
&=\sum_{r=0}^{\lfloor \epsilon k\rfloor}\sum_{l=0}^{r}\binom{k}{r}\binom{r}{l}\binom{m-k}{r-l}2^{l}3^{r-l}.
\end{split}
\end{equation}
By the fact $\binom{m-k}{r-l}\leq m^{r-l}$, it follows that
\begin{equation}\label{combinatorics 2}
\sum_{l=0}^{r}\binom{r}{l}\binom{m-k}{r-l}2^{l}3^{r-l}\leq \sum_{l=0}^{r}\binom{r}{l}2^{l}(3m)^{r-l}=(2+3m)^{r}.
\end{equation}
Substituting \eqref{combinatorics 2} into \eqref{combinatorics 1} yields that
\begin{equation}\label{combinatorics 3}
\begin{split}
|D(\eta,(1-\epsilon)k))|&\leq \sum_{r=0}^{\lfloor \epsilon k\rfloor}\binom{k}{r}(2+3m)^{r}\\
&\leq(2+3m)^{\lfloor \epsilon k\rfloor}\sum_{r=0}^{k}\binom{k}{r}=2^{k}(2+3m)^{\lfloor \epsilon k\rfloor}.
\end{split}
\end{equation}
We apply the greedy algorithm as follows. Pick an $\eta_{1}\in \{0,1,2,3\}^{m}_{=k}$, and let $R_1:=\{0,1,2,3\}^m_{=k}\setminus D(\eta_1,(1-\epsilon)k)$. Having chosen $\eta_1,\dots,\eta_i$ and formed $R_{i}$, pick $\eta_{i+1}\in R_i$ and let $R_{i+1}:=R_{i}\setminus D(\eta_{i+1},(1-\epsilon)k)$.
%We apply the greedy algorithm as follows. Pick an $\eta_{1}\in \{0,1,2,3\}^{m}_{=k}$, then remove the set $D(\eta_{1},(1-\epsilon)k)$ from $\{0,1,2,3\}^{m}_{=k}$, which is denoted by $R_{1}$. We now proceed to pick an $\eta_{2}\in R_{1}$ and remove the set $D(\eta_{2},(1-\epsilon)k)$ from $R_{1}$, which is denoted by $R_{2}$. 
Since $\{0,1,2,3\}^{m}_{=k}$ is a finite set, repeating the previous processes $t$-times turns the set $\{0,1,2,3\}^{m}_{=k}$ into an empty set, that means
\[\{0,1,2,3\}^{m}_{=k}=\bigcup_{i=1}^tD(\eta_{i},(1-\epsilon)k).\]
By the fact
\[
|\{0,1,2,3\}^{m}_{=k}|=\binom{m}{k}3^{k},
\]
we obtain that
\begin{equation}\label{greedy estimate}
2^{k}(2+3m)^{\lfloor \epsilon k\rfloor}t\geq \binom{m}{k}3^{k}\geq \left(\frac{3m}{k}\right)^{k}.
\end{equation}
Rearranging \eqref{greedy estimate} yields the desired result.
\end{proof}




We now turn to show a low bound for the packing number. Given $K\leq L\leq n$, we introduce a new symbol defined by setting
\begin{align*}
&\{0,1,2,3\}^{\geq L}_{=K}\\
\coloneqq&\left\{\eta\in \{0,1,2,3\}^{n}:\eta(j)=0,~\mbox{for }1\leq j\leq L-1,~\mbox{and }\sum_{l=1}^{n}\mathrm{sign}(\eta(j))=K\right\},
\end{align*}
where we order the set $\{0,1,2,3\}^{\geq L}_{=K}$ with the lexicography order.
Hence, for each $\eta\in \{0,1,2,3\}^{\geq L}_{=K}$, there are $\{i_k\}_{1\leq k\leq K}$ such that $L\leq i_k\leq K$ and $\eta(i_k)\in \{1,2,3\}$ for each $1\leq k\leq K$. 


\begin{theorem}\label{low bdd}
For any $\epsilon\in (0,\frac{1}{2})$ and $n$, $d\in \mathbb{N}$ with $2^{d}\leq n$, we have
\[
\log M(\mathcal{Q}_{n,d},\|\cdot\|_{L_{2}},\sqrt{2\epsilon})\geq (1-\epsilon)2^{d-1}\log(n)-d2^{d-1}.
\]
\end{theorem}
\begin{proof}
The proof relies on an quantum analogy of decision tree.  For each $j\in [n]$, we let
\begin{equation}\label{uj}
u_{j}\coloneqq \sum_{l=1}^{3}\frac{\sigma_{le_j}}{\sqrt{3}}.
\end{equation}
Observe that for each $w=(w_1,\cdots,w_{d-1})\in\{0,1\}^{d-1}$, $k(w):=1+\sum_{k=1}^{d-1}2^{k-1}w_{k}$ runs over $\{1,2,\cdots,2^{d-1}\}$. 
For each $\eta\in \{0,1,2,3\}^{\geq d}_{=2^{d-1}}$, we define
\[
\Gamma_{\eta}\coloneqq\sum_{w\in\{0,1\}^{d-1}}\prod_{j=1}^{d-1}\left(\frac{\mathbf 1+(2w_{j}-1)u_{j}}{2}\right)\sigma_{\eta(i_{k(w)})e_{i_{k(w)}}}\in \M_{2^n}.
\]
We shall verify that $\Gamma_{\eta}\in \mathcal{Q}_{n,d}$. It is clear that $\Gamma_{\eta}$ is  of degree $d$. We can also verify that $\Gamma_{\eta}$ is self-adjoint and unitary. For $w\in \{0,1\}^{d-1}$, we denote  for simplicity
$$P_j(w)\coloneqq \frac{\mathbf{1}+(2w_{j}-1)u_{j}}{2},\quad 1\leq j\leq d-1$$
 and
 $$P(w)\coloneqq \prod_{j=1}^{d-1}P_j(w).$$
 Note that $\{u_{j}\}_{j=1}^{n}$ are unitaries. Then each $P_j(w)$ is a projection in $\mathbb{M}_{2^n}$. 
Moreover, since $d\leq i_{k(w)}\leq n$ for each $w$, it follows from the definition of $P_j$ that $P_j$ commutes with  $\sigma_{\eta(i_{k(w)})e_{i_{k(w)}}}$ for all $1\leq j\leq d-1$. Thus, it is easy to see that $\Gamma_{\eta}$ is self-adjoint. 
We now show that
\[
P(w)P(w^{\prime})=\delta_{w,w^{\prime}}P(w),\quad w,~w^{\prime}\in \{0,1\}^{d-1},
\]
and 
$$\sum_{w\in \{0,1\}^{d-1}}P(w)=\mathbf{1}.$$ Indeed, by the commuting of $\{u_{j}\}_{j=1}^{n}$, it follows that for arbitrary $w$, $w^{\prime}\in \{0,1\}^{d-1}$, we have
\begin{align*}
P(w)P(w^{\prime})&=\prod_{j=1}^{d-1}\frac{(\mathbf{1}+(2w_{j}-1)u_{j})(\mathbf{1}+(2w^{\prime}_{j}-1)u_{j})}{4}\\
&=\prod_{j=1}^{d-1}\frac{\mathbf{1}+(2w_{j}-1)(2w^{\prime}_{j}-1)+2(w_{j}+w^{\prime}_{j}-1)u_{j}}{4}\\
&=
\begin{cases}
\prod_{j=1}^{d-1}\frac{1+(2w_{j}-1)u_{j}}{2},&\mbox{if }w=w^{\prime}\\
0,&\mbox{otherwise}.
\end{cases}
\end{align*}
Moreover, it is clear that
\begin{align*}
\sum_{w\in \{0,1\}^{d-1}}\prod_{j=1}^{d-1}P_j(w)&=\left(\frac{1-u_{1}}{2}+\frac{1+u_{1}}{2}\right)\sum_{v\in \{0,1\}^{d-2}}\prod_{j=2}^{d-1}\frac{1+(2v_{j}-1)u_{j}}{2}\\
&=\sum_{v\in \{0,1\}^{d-2}}\prod_{j=2}^{d-1}\frac{1+(2v_{j}-1)u_{j}}{2}=\mathbf{1},
\end{align*}
where we applied the induction argument. Therefore, $\Gamma_{\eta}$ is unitary by seeing that
\[
\Gamma_{\eta}^{2}=\left(\sum_{w\in \{0,1\}^{d-1}}P(w)\sigma_{\eta(i_{k(w)})e_{i_{k(w)}}}\right)^{2}=\sum_{w\in \{0,1\}^{d-1}}P(w)=\mathbf{1}.
\]

Applying Lemma \ref{packing number lemma} with parameters $k=2^{d-1}$ and $m=n-d+1$, it follows that there exist $\{\eta_{j}\}_{j=1}^{t}$ with
\[
t\geq \left(\frac{3}{2^{d}}\right)^{2^{d-1}}n^{(1-\epsilon)2^{d-1}}\left(\frac{n}{3n+2}\right)^{\epsilon 2^{d-1}}
\]
such that $|\eta_{i}\cap \eta_{j}|<(1-\epsilon)2^{d-1}$. 
 For any $\eta$, $\eta^{\prime}\in \{\eta_{j}\}_{j=1}^{t}$, we have
\begin{equation}\label{lower bound for decision tree 1}
\begin{split}
\|\Gamma_{\eta}-\Gamma_{\eta^{\prime}}\|^{2}_{L_{2}}&=\sum_{w\in \{0,1\}^{d-1}}\|P(w)\|^{2}_{L_{2}}\left\|\sigma_{\eta(i_{k(w)})e_{i_{k(w)}}}-\sigma_{\eta^{\prime}(i_{k(w)})e_{i_{k(w)}}}\right\|^{2}_{L_{2}}\\
&\geq 2^{2-d}|\{j\in \mathbb{N}:\eta_{i(j)}\neq \eta^{\prime}_{i(j)}\}|\\
&\geq 2^{2-d}\left(2^{d-1}-|\eta\cap\eta^{\prime}|\right)\geq 2\epsilon,
\end{split}
\end{equation}
where we used the fact
\[
\|P(w)\|^{2}_{L_{2}}=\prod_{j=1}^{d-1}\|P_j(w)\|^{2}_{L_{2}}=\prod_{j=1}^{d-1}\mathrm{tr}\left(\frac{1+(2w_{j}-1)u_{j}}{2}\right)=\frac{1}{2^{d-1}}.
\]
It follows from \eqref{lower bound for decision tree 1} that $\|\Gamma_{\eta}-\Gamma_{\eta^{\prime}}\|^{2}_{L_{2}}\geq 2\epsilon$ for each $i\neq j$. Therefore, we have
\[
\log M(\mathcal{Q}_{n,d},\|\cdot\|_{L_{2}},\sqrt{2\epsilon})\geq (1-\epsilon)2^{d-1}\log(n)-2^{d-1}\log(2^{d}/3)-\epsilon 2^{d-1}\log\left(\frac{3n+2}{n}\right).
\]
Since $2^{d-1}\log(3)\geq \epsilon 2^{d-1}\log(3+\frac{2}{n})$ for $d$, $n\in \mathbb{N}$ as $\epsilon\in (0,1/2)$, it follows that
\[
\log M(\mathcal{Q}_{n,d},\|\cdot\|_{L_{2}},\sqrt{2\epsilon})\geq (1-\epsilon)2^{d-1}\log(n)-d2^{d-1}.
\]
This completes the proof.
\end{proof}

Modifying the proof of Theorem \ref{low bdd}, we can provide an improved version of the lower bound for $\log M(\mathcal{Q}_{n,d},\|\cdot\|_{L_{2}},\sqrt{2\epsilon})$, whose construction is rather invoked.
\begin{theorem}\label{quantum packing}
There exists a universal constant $C>0$ such that the following holds. For every $0<\epsilon<1$ and every $n,d\in\mathbb{N}$ satisfying $n\geq\frac{3^{d}-1}{2}$ we have
\[
\log M\left(\mathcal{Q}_{n,d},\sqrt{2\epsilon},\left\|\cdot\right\|_{L_{2}}\right)\geq(1-\epsilon)3^{d-1}\log n-Cd3^{d-1}.
\]
\end{theorem}
\begin{proof}
Let  $n,d\in\mathbb{N}$ be such that $n\geq\frac{3^{d}-1}{2}$. Set
\[
K\coloneqq3^{d-1},\quad q\coloneqq1+3+\cdots+3^{d-2}=\frac{3^{d-1}-1}{2}.
\]
The integer $q$ is the number of internal vertices of the tree of depth $d-1$, and $K$ is the number of its leaves with
\[
n\geq q+K=\frac{3^{d}-1}{2}.
\]

Let
\[
\mathcal{V}\coloneqq\bigcup_{r=0}^{d-2}\{1,2,3\}^{r},
\]
where $[3]^{0}=\left\{\emptyset\right\}$ for convenience. Let $\mathcal{B}: \mathcal{V}\to[q]\coloneqq \{1,2,3,\dots, q\}$ be a fixed bijection. The set of leaves $\mathcal{L}\coloneqq \{1,2,3\}^{d-1}$ are indexed by the lexicographic order.


For $\alpha=(\alpha_{1},\dots,\alpha_{d-1})\in\mathcal{L}$, let $r(\alpha)\in\left\{1,\dots,K\right\}$ be the position of $\alpha$ in this lexicographic order.
For $0\leq s\leq d-2$, write
\[
\alpha|_{s}\coloneqq(\alpha_{1},\dots,\alpha_{s})\in[3]^{s},
\]
where $\alpha|_{0}=\emptyset$. Define the quantum path corresponding to $\alpha\in \mathcal{L}$ by setting
\[
\mathcal{P}_{\alpha}\coloneqq\prod_{s=0}^{d-2}\sigma_{\alpha_{s+1}e_{\mathcal{B}(\alpha|_{s})}}.
\]
Since \(\mathcal{B}\) is a bijection,  \(\alpha|_{0},\alpha|_{1},\cdots,\alpha|_{d-2}\)  are \(d-1\) distinct elements of \(\mathcal{V}\).
For each $\eta\in\left\{0,1,2,3\right\}^{\geq q+1}_{=K}$, let
\[
q<i_{\eta}(1)<i_{\eta}(2)<\cdots<i_{\eta}(K)\leq n
\]
be the nonzero coordinates of $\eta$.
%For each leaf $\alpha\in\mathcal{L}$, we define
%\[
%V_{\eta,\alpha}\coloneqq\sigma_{\eta(i_{\eta}(r(\alpha)))e_{i_{\eta}(r(\alpha))}}.
%\]
We define the quantum decision tree associated with $\left\{0,1,2,3\right\}^{\geq q+1}_{=K}$ by the Walsh expansion
\begin{equation}\label{Walsh expansion tree}
\Gamma_{\eta}\coloneqq3^{-\frac{d-1}{2}}\sum_{\alpha\in[3]^{d-1}}\mathcal{P}_{\alpha}\cdot \sigma_{\eta(i_{\eta}(r(\alpha)))e_{i_{\eta}(r(\alpha))}}.
\end{equation}
We first prove that $\Gamma_{\eta}\in\mathcal{Q}_{n,d}$. The fact that $\deg(\Gamma_{\eta})\leq d$ follows from \eqref{Walsh expansion tree}. Moreover, every $\mathcal{P}_{\alpha}$ and every $\sigma_{\eta(i_{\eta}(r(\alpha)))e_{i_{\eta}(r(\alpha))}}$ is self-adjoint, and their supports are disjoint. Thus, $\Gamma^{*}_{\eta}=\Gamma_{\eta}$.

It remains to prove that $\Gamma_{\eta}$ is unitary, that is $\Gamma_{\eta}^{2}=\mathbf{1}$. By the definition of $\Gamma_{\eta}$, we have
\begin{equation}\label{expansion of quantum tree 1}
\begin{split}
\Gamma_{\eta}^{2}&=3^{-(d-1)}\sum_{\alpha,\beta\in [3]^{d-1}}\mathcal{P}_{\alpha}\mathcal{P}_{\beta}\sigma_{\eta(i_{\eta}(r(\alpha)))e_{i_{\eta}(r(\alpha))}}\sigma_{\eta(i_{\eta}(r(\beta)))e_{i_{\eta}(r(\beta))}}\\
&=3^{-(d-1)}\sum_{\substack{\alpha, \beta\in [3]^{d-1}\\\alpha=\beta}}\mathcal{P}^{2}_{\alpha}\sigma^{2}_{\eta(i_{\eta}(r(\alpha)))e_{i_{\eta}(r(\alpha))}}\\
&\quad+3^{-(d-1)}\sum_{\substack{\alpha, \beta\in [3]^{d-1}\\\alpha\neq \beta}}\mathcal{P}_{\alpha}\mathcal{P}_{\beta}\sigma_{\eta(i_{\eta}(r(\alpha)))e_{i_{\eta}(r(\alpha))}}\sigma_{\eta(i_{\eta}(r(\beta)))e_{i_{\eta}(r(\beta))}}\\
&=\mathbf{1}+3^{-(d-1)}\sum_{\substack{\alpha, \beta\in [3]^{d-1}\\\alpha\neq \beta}}\mathcal{P}_{\alpha}\mathcal{P}_{\beta}\sigma_{\eta(i_{\eta}(r(\alpha)))e_{i_{\eta}(r(\alpha))}}\sigma_{\eta(i_{\eta}(r(\beta)))e_{i_{\eta}(r(\beta))}}.
\end{split}
\end{equation}
%It is not hard to check that
For $\alpha,\beta\in[3]^{d-1}$ with $\alpha\neq\beta$, let
\[
\ell_{0}\coloneqq\min\{\ell\in[d-1]:\alpha_{\ell}\neq\beta_{\ell}\},
\qquad
s_{0}\coloneqq\ell_{0}-1.
\]
For $s<s_{0}$, the corresponding factors in $\mathcal P_{\alpha}$ and
$\mathcal P_{\beta}$ coincide and cancel. At the first branching vertex,
\[
\alpha|_{s_{0}}=\beta|_{s_{0}},
\qquad
\alpha_{s_{0}+1}\neq\beta_{s_{0}+1},
\]
so the two distinct Pauli matrices supported on the same coordinate
$c(\alpha|_{s_{0}})=c(\beta|_{s_{0}})$ anticommute. All the remaining
factors commute, since they are supported on distinct coordinates.
Therefore,
\[
\mathcal{P}_{\alpha}\mathcal{P}_{\beta}\sigma_{\eta(i_{\eta}(r(\alpha)))e_{i_{\eta}(r(\alpha))}}\sigma_{\eta(i_{\eta}(r(\beta)))e_{i_{\eta}(r(\beta))}}=-\mathcal{P}_{\beta}\mathcal{P}_{\alpha}\sigma_{\eta(i_{\eta}(r(\beta)))e_{i_{\eta}(r(\beta))}}\sigma_{\eta(i_{\eta}(r(\alpha)))e_{i_{\eta}(r(\alpha))}}.
\]
Pairing the terms indexed by $(\alpha,\beta)$ and $(\beta,\alpha)$ shows
that all the off-diagonal terms in \eqref{expansion of quantum tree 1}
cancel. Consequently,
\[\sum_{\substack{\alpha, \beta\in [3]^{d-1}\\\alpha\neq \beta}}\mathcal{P}_{\alpha}\mathcal{P}_{\beta}\sigma_{\eta(i_{\eta}(r(\alpha)))e_{i_{\eta}(r(\alpha))}}\sigma_{\eta(i_{\eta}(r(\beta)))e_{i_{\eta}(r(\beta))}}=0.\]
Thus, we deduce from \eqref{expansion of quantum tree 1} that $\Gamma_{\eta}$ is unitary. 
%{\color{red}For each $\alpha$, $\beta\in [3]^{d-1}$ with $\alpha\neq \beta$, we define
%\[
%\ell_{0}\coloneqq \min\{\ell\in [d-1]:\alpha_{\ell}\neq \beta_{\ell}\},\quad s_{0}\coloneqq \ell_{0}-1.
%\]
%Since $\mathcal{B}:\mathcal{V}\to [q]$ is a bijection, it follows that $\{\sigma_{\alpha_{s+1}e(\mathcal{B}(\alpha|s))}\}_{s=0}^{d-2}$ commute with each other. Moreover, we have
%{\small
%\begin{equation}\label{anti-commuting property}
%\begin{split}
%&\left(\prod_{s=0}^{d-2}\sigma_{\alpha_{s+1}e(c(\alpha|s))}\right)\left(\prod_{s^{\prime}=0}^{d-2}\sigma_{\beta_{s^{\prime}+1}e(c(\beta|s^{\prime}))}\right)\sigma_{\eta(i_{\eta}(r(\alpha)))e_{i_{\eta}(r(\alpha))}}\sigma_{\eta(i_{\eta}(r(\beta)))e_{i_{\eta}(r(\beta))}}\\
%=&\left(\sigma_{\alpha_{s_{0}+1}e(c(\alpha|s_{0}))}\sigma_{\beta_{s_{0}+1}e(c(\beta|s_{0}))}\right)\!\!\left(\prod_{s=s_{0}+1}^{d-2}\sigma_{\alpha_{s+1}e(c(\alpha|s))}\right)\!\!\left(\prod_{s^{\prime}=s_{0}+1}^{d-2}\sigma_{\beta_{s^{\prime}+1}e(c(\beta|s^{\prime}))}\right)\!\!\sigma_{\eta(i_{\eta}(r(\alpha)))e_{i_{\eta}(r(\alpha))}}\sigma_{\eta(i_{\eta}(r(\beta)))e_{i_{\eta}(r(\beta))}}\\
%=&-\!\!\left(\sigma_{\beta_{s_{0}+1}e(c(\beta|s_{0}))}\sigma_{\alpha_{s_{0}+1}e(c(\alpha|s_{0}))}\right)\!\!\left(\prod_{s=s_{0}+1}^{d-2}\sigma_{\alpha_{s+1}e(c(\alpha|s))}\right)\!\!\left(\prod_{s^{\prime}=s_{0}+1}^{d-2}\sigma_{\beta_{s^{\prime}+1}e(c(\beta|s^{\prime}))}\right)\!\!\sigma_{\eta(i_{\eta}(r(\alpha)))e_{i_{\eta}(r(\alpha))}}\sigma_{\eta(i_{\eta}(r(\beta)))e_{i_{\eta}(r(\beta))}},
%\end{split}
%\end{equation}}
%where the last equality follows from the anticommutativity of \(\{\sigma_{\s}\}_{\s\in{\{0,1,2,3\}^n}}\).
%Since the role of $\alpha$ and $\beta\in [3]^{d-1}$ is symmetric, substituting \eqref{anti-commuting property} into \eqref{expansion of quantum tree 1} yields that
%\[
%\sum_{\substack{\alpha,\beta\in [3]^{d-1}\\\alpha\neq \beta}}\mathcal{P}_{\alpha}\mathcal{P}_{\beta}\sigma_{\eta(i_{\eta}(r(\alpha)))e_{i_{\eta}(r(\alpha))}}\sigma_{\eta(i_{\eta}(r(\beta)))e_{i_{\eta}(r(\beta))}}=0,
%\]
%which yields that $\Gamma^{2}_{\eta}=\mathbf{1}$.}

For $\eta,\zeta\in\left\{0,1,2,3\right\}^{\geq q+1}_{=K}$, we now turn to estimate $\|\Gamma_\eta-\Gamma_\zeta\|_{L_{2}}$. Recall here that
\[
|\eta\cap\zeta|\coloneqq\#\left\{j\in\left\{q+1,\dots,n\right\}:\eta(j)=\zeta(j)\neq0\right\},
\]
then by $\|\Gamma_{\eta}-\Gamma_{\zeta}\|^{2}_{L_{2}}=2-2\langle \Gamma_{\eta},\Gamma_{\zeta}\rangle_{L_2}$, it remains to estimate $\langle \Gamma_{\eta},\Gamma_{\zeta}\rangle_{L_2}$ from below. We claim that
\[
\left\langle\Gamma_{\eta},\Gamma_{\zeta}\right\rangle_{L_{2}}
\leq\frac{|\eta\cap\zeta|}{K}.
\]
Indeed, by the Walsh expansion,
\[
\begin{split}
\left\langle\Gamma_{\eta},\Gamma_{\zeta}\right\rangle_{L_{2}}&=3^{-(d-1)}\sum_{\alpha,\beta\in[3]^{d-1}}\langle \sigma_{\eta(i_{\eta}(r(\alpha)))e_{i_{\eta}(r(\alpha))}},\sigma_{\zeta(i_{\zeta}(r(\beta)))e_{i_{\zeta}(r(\beta))}}\rangle\cdot\langle \mathcal{P}_{\alpha},\mathcal{P}_{\beta}\rangle.\\
\end{split}
\]

If $\alpha\neq\beta$, then by the construction of quantum path, it follows that at least one position of $\mathcal{P}_{\alpha}$ and $\mathcal{P}_{\beta}$ are different. Hence, by the orthogonality of $\{\sigma_{\mathbf{s}}\}_{\mathbf{s}\in\{0,1,2,3\}^{n}}$, it follows that $\langle \mathcal{P}_{\alpha},\mathcal{P}_{\beta}\rangle=0$.

If $\alpha=\beta$, then by $\|\mathcal{P}_{\alpha}\|_{L_{2}}=1$ for each $\alpha\in \{1,2,3\}^{d-1}$, we have
\[
\left\langle\Gamma_{\eta},\Gamma_{\zeta}\right\rangle_{L_{2}}=\frac{1}{K}\sum_{\alpha\in[3]^{d-1}}\langle \sigma_{\eta(i_{\eta}(r(\alpha)))e_{i_{\eta}(r(\alpha))}},\sigma_{\zeta(i_{\zeta}(r(\alpha)))e_{i_{\zeta}(r(\alpha))}}\rangle.
\]
By the orthogonality of $\{\sigma_{\mathbf{s}}\}_{\mathbf{s}\in\{0,1,2,3\}^{n}}$ again, we have
\[
\langle\sigma_{\eta(i_{\eta}(r(\alpha)))e_{i_{\eta}(r(\alpha))}},\sigma_{\zeta(i_{\zeta}(r(\alpha)))e_{i_{\zeta}(r(\alpha))}}\rangle=
\begin{cases}
1,&\mbox{if two index coincide},\\
0,&\mbox{otherwise}.
\end{cases}
\]
Since the number of leaves $\alpha$ for which $\sigma_{\eta(i_{\eta}(r(\alpha)))e_{i_{\eta}(r(\alpha))}}=\sigma_{\zeta(i_{\zeta}(r(\alpha)))e_{i_{\zeta}(r(\alpha))}}$ is at most $|\eta\cap\zeta|$.
Hence
\[
\left\langle\Gamma_{\eta},\Gamma_{\zeta}\right\rangle_{L_{2}}\leq\frac{|\eta\cap\zeta|}{K},
\]
which completes the proof of our claim.

Therefore, if
\[
|\eta\cap\zeta|<(1-\epsilon)K,
\]
then
\[
\left\|\Gamma_{\eta}-\Gamma_{\zeta}\right\|_{L_{2}}^{2}>2-2(1-\epsilon)=2\epsilon.
\]

It remains to choose many such $\eta$. We shall apply Lemma \ref{packing number lemma} with \(m=n-q\) and \(k=K=3^{d-1}.\)
Since $n\geq q+K$, we have $K\leq m$. Thus, there exists a family
\[
\mathcal{E}\subseteq\left\{0,1,2,3\right\}^{\geq q+1}_{=K}
\]
such that
\[
|\eta\cap\zeta|<(1-\epsilon)K,\quad\eta\neq\zeta,\quad\eta,\zeta\in\mathcal{E},
\]
and
\[
\log\left|\mathcal{E}\right|
\geq
(1-\epsilon)K\log m-K\log K-CK.
\]
As $n\geq q+K$ and $q=(K-1)/2$, we have
\[
m=n-q\geq\frac{2}{3}n.
\]
Hence,
\[
\log m\geq\log n-\log\frac{3}{2}.
\]
Consequently,
\[
\begin{split}
\log\left|\mathcal{E}\right|
&\geq(1-\epsilon)K\log n-K\log K-CK\\
&\geq(1-\epsilon)3^{d-1}\log n-Cd3^{d-1}.
\end{split}
\]
For every distinct $\eta,\zeta\in\mathcal{E}$, the preceding distance estimate gives
\[
\left\|\Gamma_{\eta}-\Gamma_{\zeta}\right\|_{L_{2}}>\sqrt{2\epsilon}.
\]
Therefore $\left\{\Gamma_{\eta}:\eta\in\mathcal{E}\right\}$ is a $\sqrt{2\epsilon}$-packing of $\mathcal{Q}_{n,d}$ in $L_{2}$, and hence
\[
\log M\left(\mathcal{Q}_{n,d},\sqrt{2\epsilon},\|\cdot\|_{L_{2}}\right)\geq\log\left|\mathcal{E}\right|\geq(1-\epsilon)3^{d-1}\log n-Cd3^{d-1}.
\]
The proof is complete.
\end{proof}


In Theorem \ref{boolean-Proposition-8-packing}, we provide the lower bound for random query complexity. We also have similar estimate for query model. However, the upper estimate for this remains open even in the classical setting; see \cite[page 12]{EIS2023}. 


A deterministic learning algorithm of length $Q$ consists of a map
\[
H : (\mathcal{S} \times \{-1,1\})^Q \to L_2(\mathbb{M}_{2^n}, \mathrm{tr}),
\]
together with an initial query state $\rho_1 \in \mathcal{D}_n$ and deterministic maps
\[
\Phi_t : (\mathcal{S}\times \{-1,1\})^{t-1} \to \mathcal{S}, \qquad t = 2, \dots, Q,
\]
where $\mathcal{S}$ denotes the set of density matrices on $\mathbb M_{2^n}$. We emphasize that ``deterministic'' here refers only to the maps $H, \Phi_2, \dots, \Phi_Q$ themselves: given an unknown observable $A \in \mathcal{Q}_{n,d}$, the sequence of query states $\rho_1, \rho_2, \dots, \rho_Q$ and the corresponding measurement outcomes $Y_1, \dots, Y_Q$ are still random, since each $Y_t$ is the outcome of a genuine quantum measurement on $\rho_t$. Given a realized outcome string $y = (y_1, \dots, y_Q) \in \{-1,1\}^Q$, we define recursively
\[
\rho_1^y := \rho_1, \qquad \rho_t^y := \Phi_t\big((\rho_1^y, y_1), \dots, (\rho_{t-1}^y, y_{t-1})\big), \qquad t = 2, \dots, Q;
\]
note that $\rho_t^y$ depends only on the prefix $(y_1, \dots, y_{t-1})$, and not on $y_t, \dots, y_Q$. At stage $t$, given the query state $\rho_t^y$ and the true observable $A$, the oracle returns a bit $Y_t \in \{-1,1\}$ via the two-outcome measurement
\[
\mathbb{P}(Y_t = \pm 1 \mid \rho_t^y, A) = \mathrm{tr}\left(\frac{1 \pm A}{2} \rho_t^y\right) = \frac{1 \pm \mathrm{tr}(A \rho_t^y)}{2}.
\]
For $A \in \mathcal{B}_n$, we say that an outcome string $y \in \{-1,1\}^Q$ is $A$-admissible if the transcript $(\rho_1^y, y_1), \dots, (\rho_Q^y, y_Q)$ occurs with positive probability when the unknown observable is $A$, that is,
\[
\prod_{t=1}^{Q} \frac{1 + y_t \, \mathrm{tr}(A \rho_t^y)}{2} > 0.
\]
Because the outcome of each measurement is genuinely random, a fixed $A$ may admit several distinct admissible outcome strings; the algorithm is required to succeed simultaneously on all of them. We denote by
\[
Q_{\M_2}(\mathcal{Q}_{n,d},\epsilon)
\]
the least $Q \in \mathbb{N}$ for which there exists such a deterministic learning algorithm satisfying
\begin{equation}\label{H-de}
	\big\| H\big((\rho_1^y, y_1), \dots, (\rho_Q^y, y_Q)\big) - A \big\|_{L_2}^2 \le \varepsilon
\end{equation}
for every $A \in \mathcal{Q}_{n,d}$ and every $A$-admissible outcome string $y \in \{-1,1\}^Q$.
\begin{remark}
Let $d\in [n]$.	For every $\epsilon>0$, we have
	\begin{equation}\label{Qb-19}
		Q_{\M_2}(\mathcal Q_{n,d},\epsilon)
		\ge \log_2 M(\mathcal Q_{n,d},\|\cdot\|_{L_2},2\sqrt{\epsilon}).
	\end{equation}
\end{remark}

\begin{proof}
For simplicity, let
	\[
	M:=M(\mathcal Q_{n,d},\|\cdot\|_2,2\sqrt{\epsilon}),
	\]
	and choose $A_1,\dots,A_M\in\mathcal Q_{n,d}$ such that
	\[
	\|A_i-A_j\|_2>2\sqrt{\epsilon}
	\qquad (i\neq j).
	\]
	 
	 Denote $Q:=Q_{\M_2}(\mathcal Q_{n,d},\epsilon),$
	and let $H,\rho_1,$ and $\Phi_2,\dots,\Phi_Q$ be a deterministic    learning algorithm of length
	$Q$ satisfying the definition \eqref{H-de}. For each $i \in \{1, \dots, M\}$, since $A_i$ admits at least one admissible outcome string, fix one such string $y^{(i)} \in \{-1,1\}^Q$, and consider the set
	\[
	\Sigma:=\{y^{(i)}:i=1,\dots,M\}\subseteq\{-1,1\}^Q.
	\]
	We claim that $|\Sigma|=M$.
	Indeed, if $|\Sigma|<M$, then there exist $i\neq j$ such that
	\[
	y^{(i)}=y^{(j)}.
	\]
	Since $\rho_t^y$ depends only on  $(y_1, \dots, y_{t-1})$, the coincidence of the bit strings implies
	\[
	\rho_t^{y^{(i)}}=\rho_t^{y^{(j)}},
	\qquad t=1,\dots,Q.
	\]
	Hence the full data transcripts coincide, and the common output:
	\[
	\widetilde A:=H\big((\rho_1^{y^{(i)}},y^{(i)}_1),\dots,(\rho_Q^{y^{(i)}},y^{(i)}_Q)\big)=H\big((\rho_1^{y^{(j)}},y^{(j)}_1),\dots,(\rho_Q^{y^{(j)}},y^{(j)}_Q)\big)
	\]
	satisfies
	\[
	\|\widetilde A-A_i\|_2^2\le \varepsilon
	\qquad\text{and}\qquad
	\|\widetilde A-A_j\|_2^2\le \varepsilon,
	\]
	since $y^{(i)}$ is $A_i$-admissible and $y^{(j)}$ is $A_j$-admissible. This contradicts the packing condition, since
	\[
	\|A_i-A_j\|_2
	\le \|A_i-\widetilde A\|_2+\|\widetilde A-A_j\|_2
	\le 2\sqrt{\varepsilon}.
	\]
	This proves the claim. Since $\Sigma\subseteq\{-1,1\}^Q$, we conclude that
	\[
	M=|\Sigma|\le 2^Q,
	\]
	which yields \eqref{Qb-19}.
\end{proof}



We end this section by the following exact learning result.  
\begin{theorem}\label{exact-lea-thm10}
	Let $n\in\mathbb N$, $d\in\{1,\dots,n\}$ and $\delta\in(0,1)$.
	Then
	\begin{equation}\label{exac-Q}
	Q_{\M_2}(\mathcal Q_{n,d},0)
	=\sum_{\ell=0}^d 3^\ell \binom{n}{\ell}.
	\end{equation}
	Moreover, there exists some universal constant $C>0$ such that
\begin{equation}\label{exac-Qr}
	Q_{\M_2,r}(\mathcal Q_{n,d},0,\delta)
	\le
	C\,d^2\,6^d\,n^d\,\log \Bigl(\frac{n}{\delta}\Bigr).
\end{equation}	
\end{theorem}

\begin{proof}
We now claim that $Q_{\M_2}(\mathcal Q_{n,d},0)
\leq \sum_{\ell=0}^d 3^\ell \binom{n}{\ell}$. Indeed, let
	\[
	A=\sum_{|\s|\le d}\widehat A(\s)\sigma_{\s}\in \mathcal Q_{n,d}.
	\]	For $\s\neq 0$ with $|\s|\le d$, define
	\[
	\rho_{\s}:=\mathbf{1}+\sigma_{\s},
	\qquad
	\rho_0:=\mathbf{1}.
	\]
	Since each $\sigma_{\s}$ is  self-adjoint unitary, $\mathbf{1}+\sigma_{\s}\ge 0$, and therefore every
	$\rho_{\s}$ is a density matrix.
	Then
	\[
	\mathrm{tr}(A\rho_0)=\mathrm{tr}(A)=\widehat A(0),
	\]
	while for every $\s\neq 0$ with $|\s|\le d$,
	\[
	\mathrm{tr}(A\rho_{\s})
	=
	\mathrm{tr}(A)+\mathrm{tr}(A\sigma_{\s})
	=
	\widehat A(0)+\widehat A(\s).
	\]
It is clear that $A$ has $\sum_{\ell=0}^d 3^\ell \binom{n}{\ell}$ Fourier coefficients as $|\s|\leq d$. Thus, according to the above correspondence, it suffices to take $\sum_{\ell=0}^d 3^\ell \binom{n}{\ell}$ query values
	\[
(\rho_0,	\mathrm{tr}(A\rho_0)), \cdots
(\rho_{\s},	\mathrm{tr}(A\rho_{\s})) \quad (0< |\s|\le d),
	\]
to  recover $A$. 
	Therefore $A$ is exactly determined by $\sum_{\ell=0}^d 3^\ell \binom{n}{\ell}$ deterministic queries, so
	\[
	Q_{\M_2}(\mathcal Q_{n,d},0)\le \sum_{\ell=0}^d 3^\ell \binom{n}{\ell}.
	\]

It remains to show $Q_{\M_2}(\mathcal Q_{n,d},0)\geq \sum_{\ell=0}^d 3^\ell \binom{n}{\ell}$. We proceed by contradiction. Denote $Q\coloneqq Q_{\M_2}(\mathcal Q_{n,d},0)$and assume that $Q<\sum_{\ell=0}^d 3^\ell \binom{n}{\ell}$.Let $\{\rho_1,\dots,\rho_Q\}$ be the adaptive density matrix queries. Now run the algorithm on the zero matrix $0$. Consider the linear map
	\[
	\Psi_0:\mathcal P^{\leq d}\to \mathbb R^Q\]
	defined by
	\[
	\Psi_0(B)=
	\bigl(\mathrm{tr}(B\rho_1),\dots,\mathrm{tr}(B\rho_Q)\bigr).
	\]
	Since
	\[
\mathrm{dim}(\mathcal P^{\leq d})=\sum_{\ell=0}^d 3^\ell \binom{n}{\ell}>Q=\mathrm{dim}(\mathbb R^Q),
	\]
	then, $\mathrm{ker}\Psi_0\neq\emptyset$. Choose $0\neq B\in\ker\Psi_0$ and set
	\[
	A:=\frac{B}{\|B\|_\infty}\in \mathcal Q_{n,d}.
	\]
	Then
	\[
	\mathrm{tr}(A\rho_j)=0,\qquad \mbox{for each}\ j=1,\dots,Q.
	\]
	Since all answers received by the algorithm on $A$ are identical to those on \(0\).
	This is impossible, because exact learning requires the output to equal $0$ on the \(0\) matrix
	and to equal $A\neq 0$ on the input $A$.
	Thus
	\[
	Q\ge \sum_{\ell=0}^d 3^\ell \binom{n}{\ell},
	\]
which proves \eqref{exac-Q}.
	
	For the  proof of the random example model \eqref{exac-Qr}.  Set $\varepsilon^{(1)},\dots,\varepsilon^{(Q)}\in\{-1,1\}^{3n}$,  each \(\varepsilon^{(t)}\in\{-1,1\}^{3n}\) determines a density matrix  \(\rho(\varepsilon^{(t)})\) by \eqref{rho-j} and \eqref{rho}. Define the operator
	\[
	\Psi_{\varepsilon^{(1)},\dots,\varepsilon^{(Q)}}:
	\mathcal Q_{n,d}\to \mathbb R^Q
	\]
	by
	\[
	\Psi_{\varepsilon^{(1)},\dots,\varepsilon^{(Q)}}(A)
	:=
	\bigl(
	\mathrm{tr}(A\rho(\varepsilon^{(1)})),
	\dots,
	\mathrm{tr}(A\rho(\varepsilon^{(Q)}))
	\bigr).
	\]
	If this operator is injective, then the $Q$ sample values uniquely determine the unknown
	$A\in \mathcal Q_{n,d}$, and one can reconstruct $A$ exactly by solving a linear system
	for the Pauli coefficients $\{\widehat A(\s): |\s|\le d\}$.
	In order to prove the upper bound on the query complexity of the random example model without error, it suffices to show
	that if \(Q\) is large enough and $\varepsilon^{(1)},\dots,\varepsilon^{(Q)}\in\{-1,1\}^{3n}$ are independent and uniformly distributed random vectors, then $\Psi_{\varepsilon^{(1)},\dots,\varepsilon^{(Q)}}$ is injective with high
	probability.
	
	Fix points $\eta_1,\dots,\eta_q\in\{-1,1\}^{3n}$ and suppose that
	\[
	\Psi_{\eta_1,\dots,\eta_q}
	\]
	is not injective. Choose a nonzero matrix
	\[
	B\in \ker \Psi_{\eta_1,\dots,\eta_q}\subseteq \mathcal Q_{n,d}.
	\]
	Let $\varepsilon$ be an independent uniform random vector on $\{-1,1\}^{3n}$. We have
	\begin{equation}\label{p-2d}
	\mathbb{P}\Bigl(
	\dim \ker \Psi_{\eta_1,\dots,\eta_q,\varepsilon}
	<
	\dim \ker \Psi_{\eta_1,\dots,\eta_q}
	\Bigr)\geq
	\mathbb{P}\bigl(\mathrm{tr}(B\rho(\varepsilon))\neq 0\bigr)\ge 2^{-d},
	\end{equation}
	where the last inequality is due to \cite[Lemma 3.5]{ODo2014}.
	
	
	Consider $Q>\sum_{\ell=0}^d 3^\ell \binom{n}{\ell}=:K$.
	Suppose that
	\[
	\Psi_{\varepsilon^{(1)},\dots,\varepsilon^{(Q)}}
	\]
	is not injective. Then its kernel has dimension at least $1$, so along the chain
	\[
	K\ge \dim\ker \Psi_{\varepsilon^{(1)}}
	\ge
	\dim\ker \Psi_{\varepsilon^{(1)},\varepsilon^{(2)}}
	\ge
	\cdots
	\ge
	\dim\ker \Psi_{\varepsilon^{(1)},\dots,\varepsilon^{(Q)}}
	\]
	there can be at most $K-1$ strict drops. Hence at least $Q-K+1$ steps are equalities.
	By \eqref{p-2d} and independence, this yields
	\[
	\mathbb{P}\Bigl(
	\Psi_{\varepsilon^{(1)},\dots,\varepsilon^{(Q)}}
	\text{ is not injective}
	\Bigr)
	\le
	\binom{Q}{Q-K+1}\bigl(1-2^{-d}\bigr)^{Q-K+1}.
	\]
	Using
	\[
	\binom{Q}{Q-K+1}=\binom{Q}{K-1}\le Q^{K-1},
	\qquad
	1-2^{-d}\ge 1/2 \quad \text{for } d\ge 1,
	\]
	we obtain
	\[
	\mathbb P\Bigl(
	\Psi_{\varepsilon^{(1)},\dots,\varepsilon^{(Q)}}
	\text{ is not injective}
	\Bigr)
	\le
	(2Q)^{K-1}\bigl(1-2^{-d}\bigr)^Q\leq \exp\Big((K-1)\log(2Q)-Q\,2^{-d}\Big),
	\]
	the last inequality is due to the fact
	\[
	1-2^{-d}\le e^{-2^{-d}}.
	\]
	Therefore, if
	\[
	Q\ge C\,2^d K \log \Big(\frac{K}{\delta}\Big)
	\]
	for a large enough universal constant $C>0$, then
	\[
	\mathbb{P}\Bigl(
	\Psi_{\varepsilon^{(1)},\dots,\varepsilon^{(Q)}}
	\text{ is injective}
	\Bigr)\ge 1-\delta.
	\]
Note that
	\[K=
	\sum_{\ell=0}^d 3^\ell \binom n\ell
	\le
	(d+1)3^d n^d.
	\]
	Hence,
	\[
	Q_{\M_2,r}(\mathcal Q_{n,d},0,\delta)
	\le
	C\,2^d(d+1)3^d n^d
	\log \Bigl(\frac{(d+1)3^d n^d}{\delta}\Bigr).
	\]
	Since $d\le n$, one has
	\[
	\log \Bigl(\frac{(d+1)3^d n^d}{\delta}\Bigr)
	\le
	C\,d\,\log \Bigl(\frac{n}{\delta}\Bigr),
	\]
	possibly after increasing $C$.
	Therefore
	\[
	Q_{\M_2,r}(\mathcal Q_{n,d},0,\delta)
	\le
	C\,d^2\,6^d\,n^d\,\log \Big(\frac{n}{\delta}\Big),
	\]
	as claimed. The proof is complete.
\end{proof}




\begin{remark}\label{rho-re}
Note that the number in \eqref{exac-Q} is different from its classical counterpart
$$Q_{\{\pm1\}}(\mathcal{F}_{n,d},0)=\sum_{\ell=1}^d \binom{n}{\ell} $$
 provided in \cite[Theorem 10]{EIS2023}.
 The proof of deterministic identity \eqref{exac-Q} is actually simpler. 
In the classical case, one recovers the top Walsh layer from values on a Hamming ball
	via discrete derivatives. In the quantum setting, one can read off every low-degree Pauli coefficient
	directly by querying the density matrices
	\[
	\rho_0=\mathbf{1},
	\qquad
	\rho_{\s}=\mathbf{1}+\sigma_{\s}\quad (1\leq |\s|\le d).
	\]
\end{remark}

\section{More interesting estimates}
This section is devoted to investigating the following two fundamental objects:
\begin{enumerate}[(i)]
\item Deriving a broader upper bound for quantum learning without requiring the target quantum function to be of degree $d$;
\item Deriving the upper bound  for the metric entropy of the class $(\mathcal Q_{n,d}, \|\cdot\|_{L_2})$.
\end{enumerate}
The first main result in this section is now stated as follows.
\begin{theorem}\label{low-lea-thm2}
Let $\epsilon,\delta\in (0,1)$, \(d\in[n]\) and \(t,\eta\in[0,1)\), and let $A^*=A\in \M_{2^n}$  with  $\|A\|_\infty\le 1$. Assume that there exists a set
	$\mathcal T(A)\subseteq\{0,1,2,3\}^n$ with $|\mathcal{T}(A)|\le m$, $m\in\mathbb{N}$, such that
	\[
	\sum_{|\s|>d}|\widehat A(\s)|^2\le t,
	\qquad
	\sum_{s\notin \mathcal T(A)}|\widehat A(\s)|^2\le \eta.
	\]
If
	\[
	N\ge 	\Big\lceil
	\frac{18\cdot 3^dm}{\varepsilon}
	\log \Big(\frac{2\sum_{\ell=0}^d 3^\ell\binom{n}{\ell}}{\delta}\Big)
	\Big\rceil,
	\]
then, from $N$ independent random queries
	\[
	\bigl(\rho(x^{(1)}),\mathrm{tr}(A\rho(x^{(1)}))\bigr),\dots,
	\bigl(\rho(x^{(N)}),\mathrm{tr}(A\rho(x^{(N)}))\bigr),
	\qquad x^{(1)},\dots,x^{(N)} \in \{-1,1\}^{3n},
	\]
	one can construct a random observable $\widetilde A$ such that
	\[
	\|A-\widetilde A\|_2^2\le \eta+t+\varepsilon
	\]
	with probability at least $1-\delta$.
\end{theorem}

\begin{proof}
	For $\s\in\{0,1,2,3\}^n$ with $|\s|\le d$, define the empirical coefficients
	\[
	\alpha_{\s}:=
	\frac1N\sum_{r=1}^N \mathrm{tr}(A\rho(x^{(r)}))\omega_{S}(x^{(r)})=\frac1N\sum_{r=1}^N 3^{-|\s|/2}\widehat A(\s)\omega_{S}(x^{(r)}), \quad x^{(r)}\in\{-1,1\}^{3n},
	\]
where \(S\) is determined by  $\supp(\s)$, as defined in \eqref{weak version}  within the proof of Lemma \ref{lem rho}.
	Then, $
	\mathbb E[\alpha_{\s}]=3^{-|\s|/2}\widehat A(\s),$ where  $\mathbb E$ is with respect to the uniform distribution. Since $\alpha_{\s}$ is an average of bounded i.i.d. random variables and \(\|f_A\|_{L_\infty(\{-1,1\}^{3n})}\le \|A\|_{L_\infty}\leq 1\), Hoeffding's inequality gives
	\[
	\mathbb P\Big(
	|\alpha_{\s}-3^{-|\s|/2}\widehat A(\s)|>b
	\Big)\le 2e^{-Nb^2/2},
	\qquad |\s|\le d.
	\]
	Denote
	\[G_b:=\Big\{
	|\alpha_{\s}-3^{-|\s|/2}\widehat A(\s)|\le b
	\ \text{for all }|\s|\le d
	\Big\}.\]
	Using the union bound over $|\s|\leq d$, we obtain that
	\[\mathbb P(G_b)\geq 1-2\sum_{\ell=0}^d 3^\ell\binom{n}{\ell}e^{-Nb^2/2}\geq 1-\delta,\]
	provided
	\[
	N\ge \frac{2}{b^2}\log \Big(\frac{2\sum_{\ell=0}^d 3^\ell\binom{n}{\ell}}{\delta}\Big).
	\]
	Define
	\[
	T_b:=\{\s\in\{0,1,2,3\}^n:\ |\s|\le d,\ |\alpha_{\s}|\ge 2b\},
	\qquad
	\widetilde A_b:=\sum_{s\in T_b}3^{|\s|/2}\alpha_{\s}\,\sigma_{\s}.
	\]
	On the event $G_b$ we have
	\begin{equation}\label{s-noin-Tb}
		\s\notin T_b \Longrightarrow |3^{-|\s|/2}\widehat A(\s)|\le |\alpha_{\s}|+|\alpha_{\s}-3^{-|\s|/2}\widehat A(\s)|<2b+b=3b,
	\end{equation}
	and
	\begin{equation}\label{s-in-Tb}
		\s\in T_b\Longrightarrow |3^{-|\s|/2}\widehat A(\s)|\ge |\alpha_{\s}|-|\alpha_{\s}-3^{-|\s|/2}\widehat A(\s)|\ge 2b-b=b.
	\end{equation}
	Therefore, 
	\begin{align*}
		\|A-\widetilde A_b\|_{2}^2
		&=
		\sum_{\s\in T_b}|\widehat{A}(\s)-3^{|\s|/2}\alpha_{\s}|^2
		+\sum_{\substack{|\s|\le d\\ \s\notin T_b}}|\widehat{A}(\s)|^2
		+\sum_{|\s|>d}|\widehat{A}(\s)|^2 \\
		&= \sum_{\s\in T_b}3^{|\s|}|3^{-|\s|/2}\widehat{A}(\s)-\alpha_{\s}|^2
		+\sum_{\substack{|\s|\le d\\ \s\notin T_b}}|\widehat{A}(\s)|^2
		+\sum_{|\s|>d}|\widehat{A}(\s)|^2 \\
		&=\sum_{\s\in T_b\cap \mathcal T(A)}3^{|\s|}|\alpha_{\s}-3^{-|\s|/2}\widehat{A}(\s)|^2
		+\sum_{\substack{|\s|\le d\\ \s\notin T_b,\ \s\in \mathcal T(A)}}|\widehat{A}(\s)|^2 \\
		&\qquad
		+\sum_{\s\in T_b,\s\notin\mathcal T(A)}3^{|\s|}|\alpha_{\s}-3^{-|\s|/2}\widehat{A}(\s)|^2
		+\sum_{\substack{|\s|\le d\\ \s\notin T_b,\ \s\notin \mathcal T(A)}}|\widehat{A}(\s)|^2
		+\sum_{|\s|>d}|\widehat{A}(\s)|^2.
	\end{align*}
	On $G_b$, by \eqref{s-noin-Tb}, we have
\begin{align*}
&\sum_{\s\in T_b\cap \mathcal T(A)}3^{|\s|}|\alpha_{\s}-3^{-|\s|/2}\widehat{A}(\s)|^2
+\sum_{\substack{|\s|\le d\\ \s\notin T_b,\ \s\in \mathcal T(A)}}|\widehat{A}(\s)|^2\\
&=\sum_{s\in T_b\cap \mathcal T(A)}3^{|\s|}|\alpha_{\s}-3^{-|\s|/2}\widehat{A}(\s)|^2
+\sum_{\substack{|\s|\le d\\ s\notin T_b,\ s\in \mathcal T(A)}}3^{|\s|}|3^{-|\s|/2}\widehat{A}(\s)|^2\\
&\leq 3^d9b^2\#\mathcal T(A)\leq 3^d9b^2m.
\end{align*}
	And  on $G_b$ for $\s\in T_b$, by \eqref{s-in-Tb}, 
	\[3^{|\s|}|\alpha_{\s}-3^{-|\s|/2}\widehat{A}(\s)|^2\le 3^{|\s|}b^2\le 3^{|\s|}|3^{-|\s|/2}\widehat{A}(\s)|^2=|\widehat{A}(\s)|^2.
	\]
	Therefore, the third and fourth terms are bounded by
	\[
	\sum_{\s\notin \mathcal T(A)}|\widehat A(\s)|^2
	\le \eta,
	\]
	while the last term is bounded by
	\[
	\sum_{|\s|>d}|\widehat A(\s)|^2\le t.
	\]
	Thus
	\[
	\|A-\widetilde A_b\|_{2}^2\le 3^d9b^2m+\eta+t=\varepsilon+\eta+t,
	\]
	the last equality follows by choosing $b^2=\varepsilon/(9\cdot3^dm)$, and this choice gives exactly
	\[
	N\ge
	\Big\lceil
	\frac{18\cdot 3^dm}{\varepsilon}
	\log \Big(\frac{2\sum_{\ell=0}^d 3^\ell\binom{n}{\ell}}{\delta}\Big)
	\Big\rceil.
	\]
	The proof is complete.
\end{proof}


We observe in \cite{EIS2023} that  the classical counterpart of Theorem \ref{low-lea-thm2}  yields numerous applications concerning Fourier concentration. We skip  applications of Theorem \ref{low-lea-thm2} herein, as they fall beyond the scope of the present paper to some extent.


Finally, we apply the reduction method used in the last section to prove the following upper bound for the $\epsilon$-metric entropy of $(\mathcal{Q}_{n,d}, \|\cdot\|_{L_2})$. To proceed, let us recall the \emph{covering number} of a metric space $(\mathcal{M},d)$ as follows. For any $\epsilon>0$, the subset $\mathcal{N}\subseteq \mathcal{M}$ is said to be $\epsilon$-covering, if for each $\xi\in \mathcal{M}$, there exists some $\xi^{\prime}\in \mathcal{N}$ with $d(\xi,\xi^{\prime})\leq\epsilon$, and the covering number of the metric space $(\mathcal{M},d)$, denoted by $N(\mathcal{M},d,\epsilon)$, is defined by
\[
N(\mathcal{M},d,\epsilon)\coloneqq \inf\left\{|\mathcal{N}|:\mathcal{N}\subseteq \mathcal{M},~\mbox{and }\mathcal{N}~\mbox{is an }\mbox{$\epsilon$-covering subset}\right\}.
\]
In the classical setting, the corresponding result was proved (see \cite{EIS2023}) by using a deep junta theorem due to Dinur et al. \cite{DFKO2007}. Volberg and Zhang \cite{VZ2024} established a noncommutative junta theorem only for operators that are exactly of degree at most $d$, corresponding to the special case in the classical framework.  However,  a noncommutative analogue of the more general junta theorem in  \cite{DFKO2007} has not yet been proved in the quantum setting.
%which applies to functions whose spectra are merely approximately concentrated up to degree $d$.

\begin{theorem}\label{mn-thm9(23)}
	For every $n,d\in\mathbb N$ and every $\epsilon>0$, there exist a universal constant $C>0$ and 
	$\widetilde{\kappa}(d,\varepsilon)>0$ depending only on $d$ and $\epsilon$ such that
	\[
	\log_2 N \left(\mathcal Q_{n,d},\|\cdot\|_{L_2},\epsilon\right)
	\le
	\frac{2^{Cd}3^{2d}}{\epsilon^4}\log n+\widetilde{\kappa}(d,\epsilon).
	\]
\end{theorem}

\begin{proof}
Recall that it was proved in \cite[Theorem 9]{EIS2023} that
there exists a universal constant
	$C_0>0$ and a function $\kappa(\cdot,\cdot)$ such that
	\[
	\log_2 N(\mathcal F_{m,d},\|\cdot\|_{L_2},\eta)
	\le
	\frac{2^{C_0 d}}{\eta^4}\log m+\kappa(d,\eta)
	\qquad m,d\in\mathbb N,\ \eta\in(0,1),
	\]
we claim that the following reduction holds
	\[
	N \left(\mathcal Q_{n,d},\|\cdot\|_{L_2},\varepsilon\right)
	\le
	N \left(\mathcal F_{3n,d},\|\cdot\|_{L_2},\frac{\epsilon}{2\cdot 3^{d/2}}\right)
	\]
where $\mathcal F_{m,d},\mathcal F_{3n,d}$ are defined as in \eqref{Fnd}.
Once this claim has been proved, taking $m=3n$ and $\eta=\epsilon/(2\cdot 3^{d/2})$, combining  the last two inequalities, we have
	\[
	\log_2 N\ \left(\mathcal Q_{n,d},\|\cdot\|_{L_2},\epsilon\right)
	\le
	16\cdot \frac{2^{C_0 d}3^{2d}}{\epsilon^4}\log(3n)
	+
	\kappa \left(d,\frac{\epsilon}{2\cdot 3^{d/2}}\right).
	\]
	Since $16\le 2^{4d}$ for every $d\ge 1$, we may increase the universal constant and absorb
	the harmless $\log 3$ term into the function depending only on $(d,\epsilon)$. Thus,
	for a possibly larger universal constant $C>0$,
	\[
	\log_2 N\ \left(\mathcal Q_{n,d},\|\cdot\|_{L_2},\epsilon\right)
	\le
	\frac{2^{Cd}3^{2d}}{\epsilon^4}\log n+\widetilde{\kappa}(d,\epsilon),
	\]
	where one may take for instance
	\[
	\widetilde{\kappa}(d,\epsilon)
	:=
	\kappa\left(d,\frac{\epsilon}{2\cdot 3^{d/2}}\right)
	+
	16\cdot \frac{2^{C_0 d}3^{2d}}{\epsilon^4}\log 3.
	\]


Now we turn to verify the claim. Let
	\[
	\mathcal G_{n,d}:=\{f_A:\ A\in \mathcal Q_{n,d}\}\subseteq \mathcal F_{3n,d},
	\]
	where $f_A$ is defined as in Lemma \ref{lem rho}. We first compare the \(L_2\) distance on $\mathcal Q_{n,d}$ with the
	$L_2$-distance on $\mathcal G_{n,d}$. If $A,B\in \mathcal Q_{n,d}$, then
	$\deg(A-B)\le d$, hence
	\begin{align*}
		\|A-B\|_2^2
		&=
		\sum_{|\s|\le d} |\widehat{A-B}(s)|^2 \\
		&=
		\sum_{|S|\le d} 3^{|\s|}
		\bigl|\widehat{f_A-f_B}(S)\bigr|^2 \\
		&\le
		3^d \sum_{|S|\le d}
		\bigl|\widehat{f_A-f_B}(S)\bigr|^2 \\
		&\le
		3^d \|f_A-f_B\|_{L_2}^2.
	\end{align*}
	Therefore
	\begin{equation}\label{AB-fafb}
	\|A-B\|_2\le 3^{d/2}\|f_A-f_B\|_{L_2}.
	\end{equation}
Now set $\{h_1,\dots,h_M\}\subseteq \mathcal F_{3n,d}$ be an $\epsilon/(2\cdot 3^{d/2})$-cover of $\mathcal F_{3n,d}$ with
	\[
	M=N \big(\mathcal F_{3n,d},\|\cdot\|_{L_2},\frac{\epsilon}{2\cdot 3^{d/2}}\big).
	\]
	Fix $A \in \mathcal{Q}_{n,d}$. Since $f_A \in \mathcal{G}_{n,d}
	\subseteq \mathcal{F}_{3n,d}$ and $\{h_{j}\}_{j=1}^{M}$ covers $\mathcal{F}_{3n,d}$,
	there exists $j$ with
\[
\|f_A - h_j\|_{L_2} \le   \frac{\epsilon}{2\cdot 3^{d/2}}.
\]
In particular, we have $f_A \in \mathrm{Ball}_{L_2}(h_j,\epsilon/(2\cdot 3^{d/2})) \cap \mathcal{G}_{n,d}$, so this
	intersection is automatically non-empty. Choose any point $g_j$ from
	this intersection, and write $g_j = f_{A_j}$ for some $A_j \in
	\mathcal{Q}_{n,d}$. Then
	$$
	\|g_j - h_j\|_{L_2} \le \frac{\epsilon}{2\cdot 3^{d/2}}.
	$$
	We claim that $\{A_j\}$ is an $\varepsilon$-cover of $(\mathcal Q_{n,d},\|\cdot\|_{L_2})$. Indeed, let $A\in \mathcal Q_{n,d}$.  By the triangle inequality,
\[
	\|f_A - g_j\|_{L_2} \le \|f_A - h_j\|_{L_2} + \|h_j - g_j\|_{L_2}
	\le 2\frac{\epsilon}{2\cdot 3^{d/2}} = \frac{\varepsilon}{3^{d/2}}.
\]
	Applying \eqref{AB-fafb},
	\[
	\|A-A_j\|_2\le 3^{d/2}\|f_A-g_j\|_{L_2}\le \varepsilon.
	\]
	This proves
	\[
	N\ \left(\mathcal Q_{n,d},\|\cdot\|_2,\varepsilon\right)
	\le
	N \big(F_{3n,d},\|\cdot\|_{L_2},\frac{\varepsilon}{2\cdot 3^{d/2}}\big).
	\]
	The proof is complete.
\end{proof}

%\section*{Acknowledgment}
% This work was supported by the National Natural Science Foundation of China (Grant Nos. 12125109, W2411005 \& 12671165); the Natural Science Foundation of Hunan Province (Grant Nos: 2025ZYJ002, 2024JJ1010 \& 2024RC3040); the Scientific Research Fund of Hunan Provincial Education Department (Grant Nos. 25A0009 \& 25B0008).



%\backmatter
%\bibliographystyle{amsalpha}
%\bibliography{References-memoirs}


\newcommand{\etalchar}[1]{$^{#1}$}
\providecommand{\bysame}{\leavevmode\hbox to3em{\hrulefill}\thinspace}
\providecommand{\MR}{\relax\ifhmode\unskip\space\fi MR }
% \MRhref is called by the amsart/book/proc definition of \MR.
\providecommand{\MRhref}[2]{%
	\href{http://www.ams.org/mathscinet-getitem?mr=#1}{#2}
}
\providecommand{\href}[2]{#2}
\begin{thebibliography}{ADEGP25}
	
	\bibitem[AA14]{A-A2014}
	S.~Aaronson and A.~Ambainis, \emph{The need for structure in quantum speedups},
	Theory Comput. \textbf{10} (2014), 133--166. \MR{3249097}
	
	\bibitem[Ada79]{Ad1979}
	R.~A. Adams, \emph{General logarithmic {S}obolev inequalities and {O}rlicz
		imbeddings}, J. Functional Analysis \textbf{34} (1979), no.~2, 292--303.
	\MR{552707}
	
	\bibitem[ADEGP25]{A-S-E-P2025}
	S.~Arunachalam, A.~Dutt, F.~Escudero~Guti\'errez, and C.~Palazuelos, \emph{A
		cb-{B}ohnenblust-{H}ille inequality with constant one and its applications in
		learning theory}, Math. Ann. \textbf{392} (2025), no.~3, 3367--3396.
	\MR{4939666}
	
	\bibitem[BDR20]{BDR2020}
	S.~Beigi, N.~Datta, and C.~Rouz\'{e}, \emph{Quantum reverse hypercontractivity:
		its tensorization and application to strong converses}, Comm. Math. Phys.
	\textbf{376} (2020), no.~2, 753--794. \MR{4103958}
	
	\bibitem[BE95]{BE1995}
	P.~Borwein and T.~Erd\'{e}lyi, \emph{Polynomials and polynomial inequalities},
	Graduate Texts in Mathematics, vol. 161, Springer-Verlag, New York, 1995.
	\MR{1367960}
	
	\bibitem[Bec75]{Bec1975}
	W.~Beckner, \emph{Inequalities in {F}ourier analysis}, Ann. of Math.
	\textbf{102} (1975), no.~1, 159--182.
	
	\bibitem[Beu70]{B1970}
	A.~Beurling, \emph{On analytic extension of semigroups of operators}, J.
	Functional Analysis \textbf{6} (1970), 387--400. \MR{282248}
	
	\bibitem[BGX25]{BGX2025}
	David~P. Blecher, Li~Gao, and Bang Xu, \emph{Geometric influences on quantum
		{B}oolean cubes}, J. Funct. Anal. \textbf{289} (2025), no.~11, Paper No.
	111132, 47. \MR{4933761}
	
	\bibitem[BH31]{B-H1931}
	H.~F. Bohnenblust and E.~Hille, \emph{On the absolute convergence of
		{D}irichlet series}, Ann. of Math. (2) \textbf{32} (1931), no.~3, 600--622.
	\MR{1503020}
	
	\bibitem[Bha26]{B-S2026}
	Sreejata~K. Bhattacharya, \emph{{A}aronson-{A}mbainis conjecture is true for
		random restrictions}, arXiv preprint arXiv:2402.13952 (2026).
	
	\bibitem[BIM23]{BIM2023}
	D.~Beltran, P.~Ivanisvili, and J.~Madrid, \emph{On sharp isoperimetric
		inequalities on the hypercube}, arXiv preprint arXiv:2303.06738 (2023).
	
	\bibitem[BKS99]{BKS1999}
	I.~Benjamini, G.~Kalai, and O.~Schramm, \emph{Noise sensitivity of {B}oolean
		functions and applications to percolation}, Inst. Hautes \'{E}tudes Sci.
	Publ. Math. (1999), no.~90, 5--43 (2001). \MR{1813223}
	
	\bibitem[BL76]{BL1976}
	J.~Bergh and J.~L\"{o}fstr\"{o}m, \emph{Interpolation spaces. {A}n
		introduction}, Grundlehren der Mathematischen Wissenschaften, No. 223,
	Springer-Verlag, Berlin-New York, 1976. \MR{0482275}
	
	\bibitem[Ble01]{Bl2001}
	R.~Blei, \emph{Analysis in integer and fractional dimensions}, Cambridge
	Studies in Advanced Mathematics, vol.~71, Cambridge University Press,
	Cambridge, 2001. \MR{1853423}
	
	\bibitem[Bon70]{Bo1970}
	A.~Bonami, \emph{\'{E}tude des coefficients de {F}ourier des fonctions de
		{$L^{p}(G)$}}, Ann. Inst. Fourier (Grenoble) \textbf{20} (1970), no.~fasc. 2,
	335--402 (1971). \MR{283496}
	
	\bibitem[BSVZ25]{BSVZ2025}
	L.~Becker, J.~Slote, A.~Volberg, and H.~Zhang, \emph{Approximating the operator
		norm of local hamiltonians via few quantum states}, arXiv preprint
	arXiv:2509.11979 (2025).
	
	\bibitem[CC25]{CC2025}
	F.~Chang and Y.~Chen, \emph{Talagrand-type correlation inequalities for
		supermodular and submodular functions on the hypercube}, arXiv preprint
	arXiv:2510.22307 (2025).
	
	\bibitem[CEE23a]{CEE2020}
	D.~Cordero-Erausquin and A.~Eskenazis, \emph{Talagrand's influence inequality
		revisited}, Anal. PDE \textbf{16} (2023), no.~2, 571--612. \MR{4593775}
	
	\bibitem[CEE23b]{CE2023}
	\bysame, \emph{Talagrand's influence inequality revisited}, Anal. PDE
	\textbf{16} (2023), no.~2, 571--612. \MR{4593775}
	
	\bibitem[CEE24]{CE2024}
	\bysame, \emph{Discrete logarithmic {S}obolev inequalities in {B}anach spaces},
	J. Lond. Math. Soc. (2) \textbf{109} (2024), no.~2, Paper No. e12873, 24.
	\MR{4704160}
	
	\bibitem[Cha14]{Cha2014}
	S.~Chatterjee, \emph{Superconcentration and related topics}, Springer
	Monographs in Mathematics, Springer, Cham, 2014.
	
	\bibitem[CKMT15]{CKMT2015}
	T.~Cubitt, M.~Kastoryano, A.~Montanaro, and K.~Temme, \emph{Quantum reverse
		hypercontractivity}, J. Math. Phys. \textbf{56} (2015), no.~10, 102204, 11.
	\MR{3412699}
	
	\bibitem[CL93]{CL1993}
	A.~Carlen and H.~Lieb, \emph{Optimal hypercontractivity for {F}ermi fields and
		related noncommutative integration inequalities}, Comm. Math. Phys.
	\textbf{155} (1993), no.~1, 27--46. \MR{1228524}
	
	\bibitem[DDF{\etalchar{+}}25]{DDFKOSW2025}
	K.~Domelevo, P.~Durcik, V.~Fragkiadaki, O.~Klein, D.~Oliveira~e Silva,
	L.~Slav\'{\i}kov\'{a}, and B.~Wr\'{o}bel, \emph{Dimension-free estimates for
		low degree functions on the {H}amming cube}, Studia Math. \textbf{280}
	(2025), no.~2, 161--173. \MR{4864451}
	
	\bibitem[DFKO07]{DFKO2007}
	I.~Dinur, E.~Friedgut, G.~Kindler, and R.~O'Donnell, \emph{On the {F}ourier
		tails of bounded functions over the discrete cube}, Israel J. Math.
	\textbf{160} (2007), 389--412. \MR{2342503}
	
	\bibitem[DMoP19]{D-M-P2019}
	A.~Defant, M.~Masty\l~o, and A.~P\'erez, \emph{On the {F}ourier spectrum of
		functions on {B}oolean cubes}, Math. Ann. \textbf{374} (2019), no.~1-2,
	653--680. \MR{3961323}
	
	\bibitem[DMP19]{DMP2019}
	A.~Defant, M.~Mastylo, and A.~P\'{e}rez, \emph{On the {F}ourier spectrum of
		functions on {B}oolean cubes}, Math. Ann. \textbf{374} (2019), no.~1-2,
	653--680. \MR{3961323}
	
	\bibitem[EG20]{EG2020}
	R.~Eldan and R.~Gross, \emph{Concentration on the {B}oolean hypercube via
		pathwise stochastic analysis}, S{TOC} '20---{P}roceedings of the 52nd
	{A}nnual {ACM} {SIGACT} {S}ymposium on {T}heory of {C}omputing, ACM, New
	York, 2020, pp.~208--221.
	
	\bibitem[EG22]{EG2022}
	\bysame, \emph{Concentration on the {B}oolean hypercube via pathwise stochastic
		analysis}, Invent. Math. \textbf{230} (2022), no.~3, 935--994. \MR{4506770}
	
	\bibitem[EI20a]{EI2020-JAT}
	A.~Eskenazis and P.~Ivanisvili, \emph{Dimension independent
		{B}ernstein-{M}arkov inequalities in {G}auss space}, J. Approx. Theory
	\textbf{253} (2020), 105377, 20. \MR{4064847}
	
	\bibitem[EI20b]{EI2020}
	\bysame, \emph{Polynomial inequalities on the {H}amming cube}, Probab. Theory
	Related Fields \textbf{178} (2020), no.~1-2, 235--287. \MR{4146538}
	
	\bibitem[EI22]{EI2022}
	\bysame, \emph{Learning low-degree functions from a logarithmic number of
		random queries}, S{TOC} '22---{P}roceedings of the 54th {A}nnual {ACM}
	{SIGACT} {S}ymposium on {T}heory of {C}omputing, ACM, New York, [2022]
	\copyright 2022, pp.~203--207. \MR{4489995}
	
	\bibitem[EIS23]{EIS2023}
	A.~Eskenazis, P.~Ivanisvili, and L.~Streck, \emph{Low-degree learning and the
		metric entropy of polynomials}, Discrete Anal. (2023), Paper No. 17, 23.
	\MR{4684859}
	
	\bibitem[EKLM22]{EKLM2022}
	R.~Eldan, G.~Kindler, N.~Lifshitz, and D.~Minzer, \emph{Isoperimetric
		inequalities made simpler}, arXiv preprint arXiv:2204.06686 (2022).
	
	\bibitem[Eld22]{El2022}
	R.~Eldan, \emph{Second-order bounds on correlations between increasing
		families}, Combinatorica \textbf{42} (2022), no.~suppl. 1, 1099--1118.
	\MR{4543725}
	
	\bibitem[ELP08a]{BL2008}
	L.~Ben E. and F.~Lust-Piquard, \emph{Poincar\'{e} type inequalities on the
		discrete cube and in the {CAR} algebra}, Probab. Theory Related Fields
	\textbf{141} (2008), no.~3-4, 569--602. \MR{2391165}
	
	\bibitem[ELP08b]{EL2008}
	L.~Efraim and F.~Lust-Piquard, \emph{Poincar\'{e} type inequalities on the
		discrete cube and in the {CAR} algebra}, Probab. Theory Related Fields
	\textbf{141} (2008), no.~3-4, 569--602. \MR{2391165}
	
	\bibitem[ELP08c]{ELP2008}
	L.~B. Efraim and F.~Lust-Piquard, \emph{Poincar\'{e} type inequalities on the
		discrete cube and in the {CAR} algebra}, Probab. Theory Related Fields
	\textbf{141} (2008), no.~3-4, 569--602.
	
	\bibitem[Erd20]{E2020}
	T.~Erd\'{e}lyi, \emph{Reverse {M}arkov- and {B}ernstein-type inequalities for
		incomplete polynomials}, J. Approx. Theory \textbf{251} (2020), 105341, 10.
	\MR{4038647}
	
	\bibitem[Esk23]{EI2023}
	P.~Eskenazis, A.and~Ivanisvili, \emph{Sharp growth of the
		{O}rnstein-{U}hlenbeck operator on {G}aussian tail spaces}, Israel J. Math.
	\textbf{253} (2023), no.~1, 469--485. \MR{4575416}
	
	\bibitem[EZ24]{EZ2024}
	A.~Eskenazis and H.~Zhang, \emph{Approximation of polynomials from {W}alsh tail
		spaces}, Studia Math. \textbf{275} (2024), no.~3, 285--296. \MR{4754579}
	
	\bibitem[Fei75]{Fe1975}
	G.~F. Feissner, \emph{Hypercontractive semigroups and {S}obolev's inequality},
	Trans. Amer. Math. Soc. \textbf{210} (1975), 51--62. \MR{511867}
	
	\bibitem[FHKL16]{FHKL2016}
	Y.~Filmus, H.~Hatami, N.~Keller, and N.~Lifshitz, \emph{On the sum of the
		{$L_1$} influences of bounded functions}, Israel J. Math. \textbf{214}
	(2016), no.~1, 167--192. \MR{3540610}
	
	\bibitem[FKKK18]{FKKK2018}
	E.~Friedgut, J.~Kahn, G.~Kalai, and N.~Keller, \emph{Chv\'{a}tal's conjecture
		and correlation inequalities}, J. Combin. Theory Ser. A \textbf{156} (2018),
	22--43. \MR{3762101}
	
	\bibitem[FKN02]{FKN2002}
	E.~Friedgut, G.~Kalai, and A.~Naor, \emph{Boolean functions whose {F}ourier
		transform is concentrated on the first two levels}, Adv. in Appl. Math.
	\textbf{29} (2002), no.~3, 427--437. \MR{1942632}
	
	\bibitem[Gor26]{Go2026}
	M.~Gordin, \emph{Vector-valued concentration inequalities on the biased
		discrete cube}, J. Funct. Anal. \textbf{290} (2026), no.~7, Paper No. 111334.
	\MR{5013120}
	
	\bibitem[Gra14]{Gr2014}
	L.~Grafakos, \emph{Classical {F}ourier analysis}, third ed., Graduate Texts in
	Mathematics, vol. 249, Springer, New York, 2014. \MR{3243734}
	
	\bibitem[Gro75a]{Gr1975}
	L.~Gross, \emph{Logarithmic {S}obolev inequalities}, Amer. J. Math. \textbf{97}
	(1975), no.~4, 1061--1083.
	
	\bibitem[Gro75b]{Gr1975-Duke}
	Leonard Gross, \emph{Hypercontractivity and logarithmic {S}obolev inequalities
		for the {C}lifford {D}irichlet form}, Duke Math. J. \textbf{42} (1975),
	no.~3, 383--396. \MR{372613}
	
	\bibitem[Har60]{Ha1960}
	T.~E. Harris, \emph{A lower bound for the critical probability in a certain
		percolation process}, Proc. Cambridge Philos. Soc. \textbf{56} (1960),
	13--20. \MR{115221}
	
	\bibitem[HCP23]{HCP2023}
	H.~Huang, S.~Chen, and J.~Preskill, \emph{Learning to predict arbitrary quantum
		processes}, PRX Quantum \textbf{4} (2023), no.~2, Paper No. 040337.
	
	\bibitem[HMO17]{HMO2017}
	S.~Heilman, E.~Mossel, and K.~Oleszkiewicz, \emph{Strong contraction and
		influences in tail spaces}, Trans. Amer. Math. Soc. \textbf{369} (2017),
	no.~7, 4843--4863. \MR{3632552}
	
	\bibitem[IN22]{IN2022}
	P.~Ivanisvili and F.~Nazarov, \emph{On {W}eissler's conjecture on the {H}amming
		cube {I}}, Int. Math. Res. Not. IMRN (2022), no.~9, 6991--7020. \MR{4411473}
	
	\bibitem[IRR{\etalchar{+}}21]{IRRRY2021}
	S.~Iyer, A.~Rao, V.~Reis, T.~Rothvoss, and A.~Yehudayoff, \emph{Tight bounds on
		the fourier growth of bounded functions on the hypercube}, arXiv preprint
	arXiv:2107.06309 (2021).
	
	\bibitem[IS24]{IS2024}
	P.~Ivanisvili and Y.~Stone, \emph{The {KKL} inequality and {R}ademacher type
		2}, Discrete Anal. (2024), Paper No. 2, 14. \MR{4741754}
	
	\bibitem[IvHV20]{IvHV2020}
	P.~Ivanisvili, R.~van Handel, and A.~Volberg, \emph{Rademacher type and {E}nflo
		type coincide}, Ann. of Math. (2) \textbf{192} (2020), no.~2, 665--678.
	
	\bibitem[IZ26]{IZ2026}
	P.~Ivanisvili and H.~Zhang, \emph{On the {E}ldan-{G}ross inequality}, J. Funct.
	Anal. \textbf{290} (2026), no.~4, Paper No. 111255, 16. \MR{4983377}
	
	\bibitem[JLLZ24]{JLLZ2025}
	Y.~Jiao, W.~Lin, S.~Luo, and D.~Zhou, \emph{Quantum kkl-type inequalities
		revisited}, arXiv preprint arXiv:2411.12399 (2024).
	
	\bibitem[JLZ25]{JLZ2025}
	Y.~Jiao, S.~Luo, and D.~Zhou, \emph{Functional {$L_1$}-{$L_p$} inequalities in
		the {CAR} algebra}, J. Funct. Anal. \textbf{288} (2025), no.~2, Paper No.
	110700. \MR{4812337}
	
	\bibitem[JW17]{JW2017}
	Y.~Jiao and M.~Wang, \emph{Noncommutative harmonic analysis on semigroups},
	Indiana Univ. Math. J. \textbf{66} (2017), no.~2, 401--417. \MR{3641481}
	
	\bibitem[Kat70]{Ka1970}
	T.~Kato, \emph{A characterization of holomorphic semigroups}, Proc. Amer. Math.
	Soc. \textbf{25} (1970), 495--498. \MR{264456}
	
	\bibitem[Kin14a]{Ki2014}
	C.~King, \emph{Hypercontractivity for semigroups of unital qubit channels},
	Comm. Math. Phys. \textbf{328} (2014), no.~1, 285--301. \MR{3196986}
	
	\bibitem[Kin14b]{K2014}
	\bysame, \emph{Hypercontractivity for semigroups of unital qubit channels},
	Comm. Math. Phys. \textbf{328} (2014), no.~1, 285--301. \MR{3196986}
	
	\bibitem[KK13]{KK2013}
	N.~Keller and G.~Kindler, \emph{Quantitative relation between noise sensitivity
		and influences}, Combinatorica \textbf{33} (2013), no.~1, 45--71.
	\MR{3070086}
	
	\bibitem[KKK{\etalchar{+}}21]{KKKMS2021}
	E.~Kelman, S.~Khot, G.~Kindler, D.~Minzer, and M.~Safra, \emph{Theorems of
		{KKL}, {F}riedgut, and {T}alagrand via random restrictions and
		{L}og-{S}obolev inequality}, 12th {I}nnovations in {T}heoretical {C}omputer
	{S}cience {C}onference, vol. 185, Schloss Dagstuhl. Leibniz-Zent. Inform.,
	Wadern, 2021, pp.~Art. No. 26, 17.
	
	\bibitem[KKL88]{KKL1988}
	J.~Kahn, G.~Kalai, and N.~Linial, \emph{The influence of variables on {B}oolean
		functions}, 29th {A}nnual {S}ymposium on {F}oundations of {C}omputer
	{S}cience, IEEE Comput. Soc. Press, Washington, DC, [1988] \copyright 1988,
	pp.~68--80. \MR{4732718}
	
	\bibitem[KKM16]{KKM2016}
	G.~Kalai, N.~Keller, and E.~Mossel, \emph{On the correlation of increasing
		families}, J. Combin. Theory Ser. A \textbf{144} (2016), 250--276.
	\MR{3534069}
	
	\bibitem[Kle66]{Kl1966}
	D.~J. Kleitman, \emph{Families of non-disjoint subsets}, J. Combinatorial
	Theory \textbf{1} (1966), 153--155. \MR{193020}
	
	\bibitem[KMS14]{KMS2014}
	N.~Keller, E.~Mossel, and A.~Sen, \emph{Geometric influences {II}: correlation
		inequalities and noise sensitivity}, Ann. Inst. Henri Poincar\'{e} Probab.
	Stat. \textbf{50} (2014), no.~4, 1121--1139.
	
	\bibitem[KR09]{KR2009}
	R.~Krauthgamer and Y.~Rabani, \emph{Improved lower bounds for embeddings into
		{$L_1$}}, SIAM J. Comput. \textbf{38} (2009), no.~6, 2487--2498.
	
	\bibitem[Kra06]{K2006}
	S.~G. Krantz, \emph{Geometric function theory}, Cornerstones, Birkh\"{a}user
	Boston, Inc., Boston, MA, 2006, Explorations in complex analysis.
	\MR{2167675}
	
	\bibitem[Led88]{Le1988}
	M.~Ledoux, \emph{Isop\'{e}rim\'{e}trie et in\'{e}galit\'{e}s de {S}obolev
		logarithmiques gaussiennes}, C. R. Acad. Sci. Paris S\'{e}r. I Math.
	\textbf{306} (1988), no.~2, 79--82. \MR{929114}
	
	\bibitem[LMN93a]{L-M-N1993}
	N.~Linial, Y.~Mansour, and N.~Nisan, \emph{Constant depth circuits, {F}ourier
		transform, and learnability}, J. Assoc. Comput. Mach. \textbf{40} (1993),
	no.~3, 607--620. \MR{1370363}
	
	\bibitem[LMN93b]{LMN1993}
	\bysame, \emph{Constant depth circuits, {F}ourier transform, and learnability},
	J. Assoc. Comput. Mach. \textbf{40} (1993), no.~3, 607--620. \MR{1370363}
	
	\bibitem[LP86]{Lu1986}
	F.~Lust-Piquard, \emph{In\'{e}galit\'{e}s de {K}hintchine dans
		{$C_p\;(1<p<\infty)$}}, C. R. Acad. Sci. Paris S\'{e}r. I Math. \textbf{303}
	(1986), no.~7, 289--292. \MR{859804}
	
	\bibitem[LP98]{Lu1998}
	\bysame, \emph{Riesz transforms associated with the number operator on the
		{W}alsh system and the fermions}, J. Funct. Anal. \textbf{155} (1998), no.~1,
	263--285. \MR{1623158}
	
	\bibitem[LPP91]{LuP1991}
	F.~Lust-Piquard and G.~Pisier, \emph{Noncommutative {K}hintchine and {P}aley
		inequalities}, Ark. Mat. \textbf{29} (1991), no.~2, 241--260. \MR{1150376}
	
	\bibitem[Mey84]{Me1984}
	P.-A. Meyer, \emph{Transformations de {R}iesz pour les lois gaussiennes},
	Seminar on probability, {XVIII}, Lecture Notes in Math., vol. 1059, Springer,
	Berlin, 1984, pp.~179--193. \MR{770960}
	
	\bibitem[MN14]{MN2014}
	M.~Mendel and A.~Naor, \emph{Nonlinear spectral calculus and super-expanders},
	Publ. Math. Inst. Hautes \'{E}tudes Sci. \textbf{119} (2014), 1--95.
	\MR{3210176}
	
	\bibitem[MO10]{MO2010}
	A.~Montanaro and T.~J. Osborne, \emph{Quantum {B}oolean functions}, Chic. J.
	Theoret. Comput. Sci. (2010).
	
	\bibitem[O'D14]{ODo2014}
	R.~O'Donnell, \emph{Analysis of {B}oolean functions}, Cambridge University
	Press, New York, 2014.
	
	\bibitem[Ole17]{Ol2017}
	K.~Oleszkiewicz, \emph{On mimicking {R}ademacher sums in tail spaces},
	Geometric aspects of functional analysis, Lecture Notes in Math., vol. 2169,
	Springer, Cham, 2017, pp.~331--337. \MR{3645131}
	
	\bibitem[OW13]{OW2013}
	R.~O'Donnell and K.~Wimmer, \emph{K{KL}, {K}ruskal-{K}atona, and monotone
		nets}, SIAM J. Comput. \textbf{42} (2013), no.~6, 2375--2399.
	
	\bibitem[Pis82a]{P1982}
	G.~Pisier, \emph{Holomorphic semigroups and the geometry of {B}anach spaces},
	Ann. of Math. (2) \textbf{115} (1982), no.~2, 375--392. \MR{647811}
	
	\bibitem[Pis82b]{Pi1982}
	\bysame, \emph{Holomorphic semigroups and the geometry of {B}anach spaces},
	Ann. of Math. (2) \textbf{115} (1982), no.~2, 375--392.
	
	\bibitem[Pis86]{Pi1986}
	Gilles Pisier, \emph{Probabilistic methods in the geometry of {B}anach spaces},
	Probability and analysis ({V}arenna, 1985), Lecture Notes in Math., vol.
	1206, Springer, Berlin, 1986, pp.~167--241. \MR{864714}
	
	\bibitem[Pis16]{Pi2016}
	G.~Pisier, \emph{Martingales in {B}anach spaces}, Cambridge Studies in Advanced
	Mathematics, vol. 155, Cambridge University Press, Cambridge, 2016.
	
	\bibitem[PR17]{PR2017}
	G.~Pisier and \'{E}. Ricard, \emph{The non-commutative {K}hintchine
		inequalities for {$0<p<1$}}, J. Inst. Math. Jussieu \textbf{16} (2017),
	no.~5, 1103--1123. \MR{3709005}
	
	\bibitem[PX97]{PX1997}
	G.~Pisier and Q.~Xu, \emph{Non-commutative martingale inequalities}, Comm.
	Math. Phys. \textbf{189} (1997), no.~3, 667--698.
	
	\bibitem[Ros20]{Ros2020}
	G.~Rosenthal, \emph{Ramon van {H}andel's remarks on the {D}iscrete cube}.
	
	\bibitem[RWZ24]{RWZ2024}
	C.~Rouz\'{e}, M.~Wirth, and H.~Zhang, \emph{Quantum {T}alagrand, {KKL} and
		{F}riedgut's theorems and the learnability of quantum {B}oolean functions},
	Comm. Math. Phys. \textbf{405} (2024), no.~4, Paper No. 95, 47. \MR{4729863}
	
	\bibitem[SVZ24]{S-V-Z2024}
	J.~Slote, A.~Volberg, and H.~Zhang, \emph{Bohnenblust-{H}ille inequality for
		cyclic groups}, Adv. Math. \textbf{452} (2024), Paper No. 109824, 35.
	\MR{4770907}
	
	\bibitem[Sze25]{S1925}
	G.~Szeg\"{o}, \emph{\"{U}ber einen {S}atz von {A}. {M}arkoff}, Math. Z.
	\textbf{23} (1925), no.~1, 45--61. \MR{1544730}
	
	\bibitem[Tal93]{Ta1993}
	M.~Talagrand, \emph{Isoperimetry, logarithmic {S}obolev inequalities on the
		discrete cube, and {M}argulis' graph connectivity theorem}, Geom. Funct.
	Anal. \textbf{3} (1993), no.~3, 295--314.
	
	\bibitem[Tal94]{Ta1994}
	\bysame, \emph{On {R}usso's approximate zero-one law}, Ann. Probab. \textbf{22}
	(1994), no.~3, 1576--1587.
	
	\bibitem[Tal96]{Ta1996}
	\bysame, \emph{How much are increasing sets positively correlated?},
	Combinatorica \textbf{16} (1996), no.~2, 243--258. \MR{1401897}
	
	\bibitem[Tal97]{Ta1997}
	\bysame, \emph{On boundaries and influences}, Combinatorica \textbf{17} (1997),
	no.~2, 275--285.
	
	\bibitem[Vol24]{Vo2024}
	A.~Volberg, \emph{Tail spaces estimates on {H}amming cube and
		{B}ernstein-{M}arkov inequality}, J. Math. Anal. Appl. \textbf{529} (2024),
	no.~2, Paper No. 127597, 22. \MR{4650802}
	
	\bibitem[VZ24a]{V-Z2024}
	A.~Volberg and H.~Zhang, \emph{Noncommutative {B}ohnenblust-{H}ille
		inequalities}, Math. Ann. \textbf{389} (2024), no.~2, 1657--1676.
	\MR{4745749}
	
	\bibitem[VZ24b]{VZ2024}
	\bysame, \emph{Noncommutative {B}ohnenblust-{H}ille inequalities}, Math. Ann.
	\textbf{389} (2024), no.~2, 1657--1676. \MR{4745749}
	
	\bibitem[Xio17]{Xi2017}
	X.~Xiong, \emph{Noncommutative harmonic analysis on semigroups and
		ultracontractivity}, Indiana Univ. Math. J. \textbf{66} (2017), no.~6,
	1921--1947. \MR{3744814}
	
	\bibitem[Zha23]{Zh2023}
	H.~Zhang, \emph{Heat-smoothing for holomorphic subalgebras of free group von
		{N}eumann algebras}, Proc. Amer. Math. Soc. \textbf{151} (2023), no.~9,
	3801--3815. \MR{4607625}
	
\end{thebibliography}


\end{document}
